% !TEX program = xelatex
\documentclass[12pt,a4paper,fontset=fandol]{ctexbook}

% ==================== 页面 ====================
\usepackage[top=2.5cm,bottom=2.5cm,left=2.5cm,right=2.5cm]{geometry}

% ==================== 字体 ====================
\usepackage{fontspec}
\usepackage{iffont}
% CJK 钉 fandol（两侧都取自 TeX 树里的同一份 FandolSong-Regular.otf）；
% mono 钉 DejaVu Sans Mono：CI 由 fontconfig 命中同名，本机没有该系统字体，
% 就用 MiKTeX dejavu 包里的同一版本 TTF（kpathsea 按文件名命中）。
\IfFontExistsTF{DejaVu Sans Mono}
  {\setmonofont{DejaVu Sans Mono}[Scale=MatchLowercase]}
  {\setmonofont[Extension=.ttf,
      UprightFont=DejaVuSansMono,
      BoldFont=DejaVuSansMono-Bold,
      ItalicFont=DejaVuSansMono-Oblique,
      BoldItalicFont=DejaVuSansMono-BoldOblique]
     {DejaVuSansMono}[Scale=MatchLowercase]}

% ==================== 颜色 ====================
\usepackage[dvipsnames,svgnames,x11names]{xcolor}
\definecolor{codebg}{HTML}{F5F5F5}
\definecolor{codeframe}{HTML}{E0E0E0}
\definecolor{cppkeyword}{HTML}{1A1AFF}
\definecolor{cppcomment}{HTML}{2E8B2E}
\definecolor{cppstring}{HTML}{B22222}
\definecolor{linenumgray}{HTML}{999999}
\definecolor{algheadbg}{HTML}{1A3C5E}

% ==================== listings ====================
\usepackage{listings}

% ---- 所有代码块共用的排版（先定义，各语言 style 组合它）----
\lstdefinestyle{codebase}{
  basicstyle=\small\ttfamily,
  keywordstyle=\color{cppkeyword}\bfseries,
  commentstyle=\color{cppcomment}\itshape,
  stringstyle=\color{cppstring},
  numberstyle=\tiny\color{linenumgray},
  numbers=left,
  frame=single,
  rulecolor=\color{codeframe},
  framesep=6pt,
  breaklines=true,
  tabsize=4,
  columns=flexible,
  keepspaces=true,
  backgroundcolor=\color{codebg},
  showstringspaces=false,
  escapeinside={(*@}{@*)},
}

\lstdefinestyle{cppstyle}{
  style=codebase,
  language=C++,
  morekeywords={constexpr,consteval,constinit,decltype,auto,
    concept,requires,co_await,co_yield,co_return,
    noexcept,override,final,nullptr,static_assert,
    template,typename,namespace,using,mutable,volatile,
    explicit,operator,friend,inline,virtual,
    thread_local,alignas,alignof},
  emph={size_t,int32_t,uint32_t,int64_t,uint64_t,ssize_t,
    string,vector,map,set,unique_ptr,shared_ptr,optional,variant,any,
    span,string_view,FILE,wchar_t,std::size_t},
  emphstyle=\color{teal!80!black},
}
\lstset{style=cppstyle}

% ---- 纯 C（多数数值/信号库的公共接口层就是 C）----
\lstdefinestyle{cstyle}{
  style=codebase,
  language=C,
  morekeywords={restrict,inline,typeof,offsetof,int8_t,int16_t,
    int32_t,int64_t,uint8_t,uint16_t,uint32_t,uint64_t,size_t,ssize_t,
    intptr_t,uintptr_t,wchar_t,_Complex,complex,bool,true,false},
  emph={cblas_dgemm,cblas_sgemm,cblas_ddot,lapack_int,fftw_plan,
    fftwf_plan_r2r_1d,fftwf_malloc,fftwf_free,cplx8,cplx64,snd.sf,
    PaStream,MiniaudioContext,cmplx,float32_t,float64_t},
  emphstyle=\color{teal!80!black},
}

% ---- Fortran（netlib 谱系的原文语言）----
\lstdefinestyle{fortranstyle}{
  style=codebase,
  language=[90]Fortran,
  morekeywords={subroutine,function,call,module,use,implicit,none,
    intent,in,out,inout,allocatable,deallocate,parameter,real,
    double,complex,precision,character,logical,do,enddo,then,endif},
  emph={dgemm,zgemm,dsyev,dlanv2,complex128,iso_c_binding,c_double},
  emphstyle=\color{teal!80!black},
}

% ---- Python（科学栈的门面层）----
\lstdefinestyle{pythonstyle}{
  style=codebase,
  language=Python,
  morekeywords={async,await,match,case,yield,from,import,lambda,
    with,class,def,return,None,True,False,self,raise,try,except,finally,
    pass,break,continue,global,nonlocal,assert,del,in,is,not,and,or,
    if,elif,else,for,while},
  emph={np,sc,sp,la,fft,signal,librosa,soundfile,sndformat,ctypes,cffi,
    scipy,numpy,torch,xr,pd,int,float,bool,str,bytes,list,dict,tuple,
    set,type,range,len,print},
  emphstyle=\color{teal!80!black},
}

% ---- pyproject.toml / 配置 ----
\lstdefinestyle{tomlstyle}{
  style=codebase,
  morecomment=[l]{\#},
  morestring=[b]",
  morekeywords={build-system,requires,build-backend,project,name,version,
    description,dependencies,optional-dependencies,tool,requires-python,
    dynamic,classifiers,scripts},
  keywordstyle=\color{cppkeyword}\bfseries,
}

% ---- Shell ----
\lstdefinelanguage{ShellLD}{
  morekeywords={if,then,else,elif,fi,for,in,do,done,while,case,esac,
    function,return,local,export,source,echo,cd,set,exec,trap,read,
    printf,test,exit,shift,eval,alias},
  sensitive=true,
  morecomment=[l]{\#},
  morestring=[b]",
  morestring=[d]',
}
\lstdefinestyle{shellstyle}{
  style=codebase,
  language=ShellLD,
  keywordstyle=\color{cppkeyword}\bfseries,
  emph={curl,git,docker,pip,pip3,python,python3,uv,cmake,ninja,make,
    nvcc,msbuild,dotnet,pwsh,bash,ls,cat,cp,rm,mkdir,touch,
    grep,find,awk,sed,pytest,mypy,ruff,black,auditwheel,twine,
    clang,gcc,g++,gfortran,ldd,objdump,nm,pkg-config,vcpkg,conan,
    sox,ffmpeg,flac,alsa-info,aplay,rec},
  emphstyle=\color{teal!80!black}\bfseries,
}

% ---- CMake ----
\lstdefinelanguage{CMakeLD}{
  morekeywords={cmake_minimum_required,project,add_executable,
    add_library,target_link_libraries,target_include_directories,
    target_compile_definitions,target_compile_options,set,if,elseif,
    else,endif,find_package,option,message,install,file,
    target_sources,add_subdirectory,enable_language,set_target_properties,
    FetchContent_Declare,FetchContent_MakeAvailable,list,string},
  sensitive=false,
  morecomment=[l]{\#},
  morestring=[b]",
}
\lstdefinestyle{cmakestyle}{
  style=codebase,
  language=CMakeLD,
  keywordstyle=\color{cppkeyword}\bfseries,
}

% ==================== tcolorbox ====================
\usepackage[most]{tcolorbox}

\newtcolorbox{guideline}[1][]{
  enhanced,breakable,
  colback=blue!5!white,colframe=blue!75!black,
  fonttitle=\bfseries,title=要点,#1,
  boxed title style={colback=blue!75!black},
  attach boxed title to top left={yshift=-2mm,xshift=2mm},
  arc=2mm,boxrule=0.8mm,
}
\newtcolorbox{compare}[1][]{
  enhanced,breakable,
  colback=green!3!white,colframe=green!50!black,
  fonttitle=\bfseries,title=对比,#1,
  boxed title style={colback=green!50!black},
  attach boxed title to top left={yshift=-2mm,xshift=2mm},
  arc=2mm,boxrule=0.8mm,
}
\newtcolorbox{warning}[1][]{
  enhanced,breakable,
  colback=red!5!white,colframe=red!75!black,
  fonttitle=\bfseries,title=常见陷阱,#1,
  boxed title style={colback=red!75!black},
  attach boxed title to top left={yshift=-2mm,xshift=2mm},
  arc=2mm,boxrule=0.8mm,
}
\newtcolorbox{note}[1][]{
  enhanced,breakable,
  colback=orange!5!white,colframe=orange!80!black,
  fonttitle=\bfseries,title=注意,#1,
  boxed title style={colback=orange!80!black},
  attach boxed title to top left={yshift=-2mm,xshift=2mm},
  arc=2mm,boxrule=0.8mm,
}
\newtcolorbox{bestpractice}[1][]{
  enhanced,breakable,
  colback=purple!5!white,colframe=purple!60!black,
  fonttitle=\bfseries,title=最佳实践,#1,
  boxed title style={colback=purple!60!black},
  attach boxed title to top left={yshift=-2mm,xshift=2mm},
  arc=2mm,boxrule=0.8mm,
}

% ==================== 章节标题 ====================
\usepackage{titlesec}
\usepackage{fancyhdr}
\usepackage{graphicx}
\usepackage{amssymb}
\usepackage{amsmath}
\usepackage{tabularx}
\usepackage{booktabs}
\usepackage{longtable}
\usepackage{array}
\usepackage{multirow}
\usepackage{enumitem}
\usepackage{seqsplit}
\usepackage{float}

\titleformat{\chapter}[display]
  {\normalfont\huge\bfseries\color{algheadbg}}
  {\chaptertitlename\ \thechapter}{20pt}{\Huge}
\titleformat{\section}
  {\normalfont\Large\bfseries\color{blue!70!black}}
  {\thesection}{1em}{}
\titleformat{\subsection}
  {\normalfont\large\bfseries\color{blue!60!black}}
  {\thesubsection}{1em}{}

\pagestyle{fancy}
\fancyhf{}
\fancyhead[LE,RO]{\thepage}
\fancyhead[RE]{\leftmark}
\fancyhead[LO]{\rightmark}
\setlength{\headheight}{14.5pt}

\setlist[itemize]{leftmargin=2em,itemsep=2pt}
\setlist[enumerate]{leftmargin=2em,itemsep=2pt}
\emergencystretch=3em

% ==================== 超链接 ====================
\usepackage{hyperref}
\hypersetup{
  colorlinks=true,
  linkcolor=blue!70!black,
  urlcolor=blue!60!black,
  bookmarks=true,
  pdfauthor={C++ Reference},
  pdftitle={C/C++ 数学、几何与信号处理库参考手册：BLAS、Eigen、FFTW、CGAL 与电磁仿真},
}

% ==================== 自定义命令 ====================
\newcommand{\Fn}[1]{\texttt{#1()}}
\newcommand{\Tp}[1]{\texttt{#1}}
\newcommand{\Opt}[1]{\texttt{-\/-#1}}
\newcommand{\Bk}{\allowbreak}
% 库名：全书统一用粗体，避免与函数名、类型名混淆
\newcommand{\Lib}[1]{\textbf{#1}}
% 语言层：这四个宏是全书讲"为什么是 C/C++"时的统一口径
\newcommand{\Lang}[1]{\textsf{\bfseries #1}}
\newcommand{\CL}{\Lang{C}}
\newcommand{\CPP}{\Lang{C++}}
\newcommand{\FTN}{\Lang{Fortran}}
\newcommand{\PY}{\Lang{Python}}
% 信号域标签：声音与电磁波
\newcommand{\AUD}{\Lang{Audio}}
\newcommand{\EMW}{\Lang{EM}}

% ====================================================================
\begin{document}

\frontmatter

\begin{titlepage}
\centering
\vspace*{4cm}
{\Huge\bfseries C/C++ 数学、几何与信号处理库参考手册\par}
\vspace{1cm}
{\LARGE 数值计算、几何计算、声音与电磁波信号处理\par}
\vspace{0.4cm}
{\Large 库生态、以 C/C++ 为底层的成因与使用示例\par}
\vfill
{\large \today\par}
\end{titlepage}

\tableofcontents

\mainmatter

\cleardoublepage

% ====================================================================
\part{为什么这一域以 C/C++ 为主}

\chapter{热路径的经济学：性能从哪里来}

本章要回答的是全书的问题：为什么数值计算、几何计算与信号处理这一大批库，实现语言几乎总是 \CL{}、\CPP{}（或 \FTN{} 原文之上加一层 \CL{} 接口），而 \PY{}、Java、C\# 这些``更高层''的语言只能站在门外做编排。答案不是语言偏好，也不是历史惯性可以一句话打发，而是一组可核对的物理量：每个采样点上允许花多少时间、每个字节要搬多少次、每纳秒能换回多少算力、以及回调线程被剥夺 CPU 时会发生什么。把这组量摆清楚，就会发现 C/C++ 不是``更快的语言''，而是唯一能把这几笔约束同时说清楚的语言层。

先声明边界。本章只讲性能与约束的来源，不讲语法风格；线程原语与内存序的细节在并行算法分册里，本章只在需要时引用结论。全书用 \CL、\CPP、\FTN、\PY{} 四个标签指代语言层，用 \Lib{库名} 指代具体实现，避免``某个库快''这种无法核对的说法。

\section{把``快''拆成四个可核对的量}

一段处理代码的实测耗时，几乎总能归到四个独立的方向上。把它们分开来看，才能判断换语言、换布局还是换算法。

\textbf{其一，运算密度}：每搬运一个字节要做多少次运算。它是决定``这段代码该在哪块硬件上跑''的第一判据，也是声学、电磁这类逐点算子与稀疏图搜索的分水岭。运算密度低的核，把指令换成再漂亮的向量码也无济于事，瓶颈早退到内存带宽上。

\textbf{其二，访存模式}：顺序、跨步还是随机。同一份 \Tp{gemm} 内核，交换行列主序就可能差三倍；一条光线遍历点云的成本主要由邻居访问的局部性决定，而不是由求交公式决定。

\textbf{其三，分支可预测性}：数据相关的分支会打断流水，也会让向量化直接失败。信号处理里常见的``遇到 NaN 就切换路径''、几何谓词里``退化就走高精度分支''都是典型代价点；工程做法不是消灭分支，而是把它挪到循环外面（分块判类，块内单态）。

\textbf{其四，能不能停}：这一项被大多数语言讨论忽略，但它是音频与实时控制域的决定性约束。GC 暂停、JIT 去优化、页错误与内存分配器的锁，都是``不可控停顿''；实时回调要求的是最长路径可预算，而不是平均吞吐高。

表~\ref{tab:ch01-four} 把四项与``谁能改善它''对应起来。

\begin{table}[H]
\centering
\caption{四个可核对的量，以及各自的主战场}
\label{tab:ch01-four}
\small
\begin{tabular}{@{}>{\raggedright\arraybackslash}p{0.17\textwidth}
  >{\raggedright\arraybackslash}p{0.30\textwidth}
  >{\raggedright\arraybackslash}p{0.33\textwidth}@{}}
\toprule
量 & 判据 & 主要改善手段 \\
\midrule
运算密度 & 每字节 flop 数与硬件脊点比较 &
  算法重排、融合、把中间量留在寄存器 \\
访存模式 & 顺序、跨步、随机三类访存占比 &
  布局改造（AoS 换 SoA）、分块、预取 \\
分支可预测性 & 循环体内数据相关分支的条数 &
  单态分块、无分支写法、查表 \\
能不能停 & 最长路径与周期预算 &
  无分配、无锁、无 GC、锁定页与优先级 \\
\bottomrule
\end{tabular}
\end{table}

要测出这四个量，先要有一个不会被编译器骗过去的测量骨架。测量代码本身的坑比被测代码还多：优化器会把整个循环删掉，首次调用会被冷缓存与动态链接污染，跨核时钟与频率调节会让两轮结果不可比。

\begin{lstlisting}[style=cppstyle,caption={最小可用的微基准：多次取最小值、防止死代码消除},label={lst:ch01-bench}]
// 关键约定：1) 用最小值而不是平均值，频率抖动只抬高尾部
//           2) 结果必须被"消费"，否则整个循环会被删掉
//           3) 先 warmup 再计时；warmup 与计时用不同数组
float axpy_ref(const float* __restrict__ x, float* y, size_t n,
               float a) {
  float acc = 0.f;                     // acc 与 y 的写入一起充当 sink
  for (size_t i = 0; i < n; ++i) {
    y[i] = a * x[i] + y[i];
    acc += y[i];                       // 累加，保证循环不被删
  }
  return acc;
}

static volatile float g_sink = 0.f;    // volatile：写它不可被优化掉

double time_axpy(const float* x, float* y, size_t n, float a,
                 int reps) {
  std::vector<double> t;
  t.reserve(reps);
  for (int r = 0; r < reps; ++r) {
    auto t0 = std::chrono::steady_clock::now();
    g_sink = axpy_ref(x, y, n, a);     // 每轮都消费返回值
    auto t1 = std::chrono::steady_clock::now();
    t.push_back(std::chrono::duration<double>(t1 - t0).count());
  }
  return *std::min_element(t.begin(), t.end());
}
\end{lstlisting}

\begin{warning}[title={微基准的四个坑}]
常见错误依次是：只测一次（冷缓存与首次调用会把结论抬高一个数量级）；用平均值而不是最小值（平均值把调度噪声当成性能）；没有 sink 导致内核被整体优化消失（表现为``快得离谱''）；把不同规模的数据放同一坐标轴比较（$10^3$ 点在 L1 里，$10^7$ 点在内存里，两者不是同一个算法）。这四个坑在 \CL、\PY{} 与任何语言里都存在，但只有在 C/C++ 里，你才有机会有责任把它们逐个排掉。
\end{warning}

\section{SIMD 与 FMA：宽度不是免费的}

这一域的核函数几乎都落在可向量化的形态上：逐点乘加、跨样点的 stencil、批量复数乘法、点的距离判定。语言层提供的能力差别，最终体现在能不能把源数据以连续、对齐、无别名的形式交给向量单元。

\begin{lstlisting}[style=cppstyle,caption={AVX2 下的复数乘法：8 个 float 通道装两组复数},label={lst:ch01-avx}]
// 复数以 split 形式存放：re[n], im[n] 两条连续数组（SoA）。
// 交错存放（re,im,re,im）也可向量化，但要额外一次 deinterleave。
void cmul_split_avx2(const float* __restrict__ a_re,
                     const float* __restrict__ a_im,
                     const float* __restrict__ b_re,
                     const float* __restrict__ b_im,
                     float* __restrict__ c_re,
                     float* __restrict__ c_im, size_t n) {
#ifdef __AVX2__
  for (size_t i = 0; i + 8 <= n; i += 8) {
    __m256 ar = _mm256_load_ps(a_re + i);
    __m256 ai = _mm256_load_ps(a_im + i);
    __m256 br = _mm256_load_ps(b_re + i);
    __m256 bi = _mm256_load_ps(b_im + i);
    __m256 cr = _mm256_fnmadd_ps(ai, bi, _mm256_mul_ps(ar, br));
    __m256 ci = _mm256_fmadd_ps(ar, bi, _mm256_mul_ps(ai, br));
    _mm256_store_ps(c_re + i, cr);     // 与 load 对称：需 32 字节对齐
    _mm256_store_ps(c_im + i, ci);
  }
  // 尾块（不足 8 个）必须单独处理，否则要么漏算要么越界
#endif
}
\end{lstlisting}

三个工程事实值得记住。第一，\Tp{fmadd} 与 \Tp{fnmadd} 是``乘加一次舍入''（清单里即 \Tp{\_mm256\_fnmadd\_ps}），与先乘后加的两次舍入结果不同：同一份代码在不同向量化路径下最后几位可能不一致，这正是可复现性要求（逐位对拍）会强制关闭某些优化的原因。第二，宽度越大，尾块与对齐的相对代价越高；把 $n$ 取成宽度的整数倍，收益往往大于换更宽的指令集。第三，intrinsics 是逃生舱而不是默认写法：绝大多数情况下应当先让编译器向量化，读它的失败报告，再决定要不要手写。

自动向量化最常见的三个杀手是别名、归约顺序与循环体里的函数调用。\Tp{restrict} 只是否定别名的承诺，编译器不会验证它；把它写错，得到的是能跑但结果错误的二进制。

\begin{note}[title={宽松浮点开关改变的是语义}]
\Tp{-ffast-math} 允许重排浮点加法、把 \Tp{x*0} 当作 $0$、把有限性与 NaN 判定去掉。对做参数拟合、迭代求解与能量守恒检查的代码，这会同时破坏收敛判据与对拍基准；信号处理链里一个被``优化掉''的钳位就是削波。若某个内核确实需要宽松语义，应当只在单个翻译单元或单个函数上开，并把这一点写进构建说明。
\end{note}

\section{布局、对齐与分配：数据决定速度}

几何与信号数据都天然多字段：点有 $x,y,z$ 与法线，采样帧有 $N$ 个通道，光线有原点、方向、区间。字段排布方式（AoS 与 SoA）决定了向量化的可行性，也决定了同一份内核能不能原地运算。

\begin{lstlisting}[style=cppstyle,caption={同一批点的两种布局：AoS 便于遍历物体，SoA 便于向量化},label={lst:ch01-layout}]
// AoS：一个点占一条连续记录，字段跨步为 12 float
struct PointAoS { float v[12]; };            // pos,nrm,uv 简化表示
// SoA：同一字段连成一条数组，向量单元可直接吞下
struct PointSoA {
  float* x; float* y; float* z;              // 各 n 个元素
  float* nx; float* ny; float* nz;
  size_t n;
};
// SoA 下"到某平面的带符号距离"是纯数组运算，天然可向量化；
// AoS 下同一操作要先把 x 抽出来（gather），代价常常高于计算本身。
void dist_to_plane(const PointSoA& p, float a, float b, float c,
                   float d, float* __restrict__ out) {
  for (size_t i = 0; i < p.n; ++i)
    out[i] = a * p.x[i] + b * p.y[i] + c * p.z[i] + d;
}
\end{lstlisting}

对齐与分配是同一件事的两端。需要 32 或 64 字节对齐的缓冲，不能靠普通 \Tp{malloc} 承诺得到：\Lib{FFTW} 提供专用的分配器，\Lib{Eigen} 为固定尺寸向量类型重载 \Tp{new}，\Lib{OpenBLAS} 与 \Lib{MKL} 的内部缓冲按缓存行对齐。这一类要求不会出现在语言教程里，但它们是接口文档的第一行——库要求你按它的布局交付数据，而不是它适配你的数据。

不自动分配，是这一域的第二个共性纪律。音频回调、射线遍历、迭代求解器内层循环里出现 \Tp{malloc}，代价不只是几十字节的时间，而是分配器锁、页错误与不可控停顿。于是常见形态变成：初始化阶段一次性建好 plan、工作缓冲与内存池，运行阶段只做原地运算。\CPP{} 在这里的价值是把``资源在构造期获取、在析构期归还''写进类型；\CL{} 在这里的价值是没有隐藏开销，一个 \Tp{struct} 加两个函数就是全部契约。

\section{解释器的每元素成本}

\PY{} 侧的写法与 \CL{} 侧的差距，可以用``每个元素上付多少钱''来描述。列表元素的 \Tp{for} 循环要在解释器里做对象取回、类型判定、装箱与引用计数；同一段代码落到 numpy 的 ufunc 上，这些动作被换成一次 C 循环里的指针算术。数量级的差别不是``慢一点''，而是同一算法能不能在采样率下成立。

\begin{lstlisting}[style=pythonstyle,caption={两种写法的分工：循环留在 C，编排留在 Python},label={lst:ch01-py}]
import numpy as np

def gain_iou(x: np.ndarray, g: float) -> np.ndarray:
    return (x * g).astype(np.float32, copy=False)   # 向量化：循环在 C

def gain_bad(x, g):
    out = np.empty_like(x, dtype=np.float32)
    for i in range(x.shape[0]):          # 每元素一次解释器往返
        out[i] = x[i] * g
    return out
\end{lstlisting}

这里要区分两件常被混谈的事：GIL 只影响多核并行，与单线程逐元素开销无关；而单线程开销才是``为什么内核必须是 C''的主因。把循环交给 C 之后，Python 层负责类型检查、编排、绘图与文件路径——这也正是 \Lib{numpy}、\Lib{scipy}、\Lib{librosa} 的分工形态，本册后面几章会反复看到这个模式的具体实现。

\section{四笔必须付的税}

选择 C/C++ 不是免费的，它把代价从运行期搬到了构建期与工程期。

\textbf{可移植性税}。SIMD 代码与 ISA 绑定，同一份内核要为 SSE、AVX2、NEON 各写一份，再在运行期按 CPUID 分派；\Lib{Embree} 与 \Lib{OpenBLAS} 都提供这类分派机制，代价是构建矩阵成倍增长。

\textbf{许可证税}。这一域大量基础库落在弱传染或商业授权下（参照附录的矩阵），把它们静态链进发布二进制之前要核对原文，不要引用二手结论。

\textbf{构建税}。\FTN{} 原文加 \CL{} 接口再加 \CPP{} 封装，意味着一个工程里同时有三种编译单元、两套名字修饰规则和三种符号可见性；CMake 的混合语言工程与 \Tp{FindBLAS} 这类探测模块就是为此存在的。

\textbf{漏抽象税}。薄封装会把两份生命周期、两次内存拷贝与一层 GIL 语义同时带进来。这一项在门面章里算细账。

\begin{compare}[title={同一个内核在四种语言里的契约形态}]
把``求一批点到平面的带符号距离''分别写成 \FTN{} 内核、\CL{} 接口、\CPP{} 模板与 \PY{} 编排，差别只在契约的表达方式：\FTN{} 用数组与意图声明，\CL{} 用指针加长度（谁都不拥有内存），\CPP{} 用视图类型加静态尺寸（所有权写进类型），\PY{} 用形状与 dtype 约定（正确性靠运行时检查）。四档里只有前两档能作为跨语言汇合点，因为它们的 ABI 是平台 ABI 的子集；后两档的抽象要么编译期即消失（模板），要么依赖运行时（对象头）。
\end{compare}

\section{判据：什么时候不该写 C/C++}

反向判据同样重要，写错方向的成本是几个月的编译时间。

其一，算法本身还没定形。这一域的探索阶段更适合 \PY{}：形状、单位与边界条件先用一次对话试出来，再把定下的内核搬到 \CL{} 或 \CPP{} 里实现。判据是``改动一次要重编译多少分钟''与``改动一次要重跑多少实验''的比值。

其二，瓶颈不在每元素成本。若总时间被磁盘、网络或数据库占满，换语言只是换个错误。回到第~\ref{lst:ch01-bench} 清单的测量骨架，先量再改。

其三，实时性与体积不是约束。没有回调预算、没有嵌入式目标、没有跨语言 ABI 需求时，托管语言的成本模型完全可以接受，此时选 C/C++ 只是把可移植性税与构建税白白付一遍。

\begin{guideline}[title={本章的判断口径}]
\begin{itemize}
\item 一段代码值不值得下沉到 C/C++，取决于它每秒被调用的次数与每次允许的最长时间，而不是它的算法难度。
\item SIMD、对齐与无分配是同一件事的三个面：把数据布局定下来，另外两个自己会跟上。
\item 浮点语义是接口的一部分：重排、融合乘加与宽松优化都会改变结果，可复现性要求必须显式写进构建说明。
\item 测量先于优化：取最小值、留 sink、分规模，否则你在比较噪声而不是比较代码。
\end{itemize}
\end{guideline}
% ====================================================================

% ====================================================================
\chapter{接口层：C ABI、Fortran 血统与四种绑定路线}

上一节讲的是执行层的物理约束，本章讲的是接口层的制度性约束，两者合起来才构成``为什么这一域以 C/C++ 为主''的完整答案。接口层的结论更锐利：跨语言的汇合点只有一个，就是 \CL{} 的调用约定与内存布局（俗称 C ABI）。无论内核用 \FTN{}、\CPP{} 还是别的语言写成，只要它想让别人调用，最终都要暴露一层 \CL{} 接口；反过来，只要一个库暴露的是 \CPP{} 类模板，它的可用范围就自动缩小到``同一个标准库、同一个编译器、同一份头文件''的圈子里。

\section{netlib 谱系：为什么是 Fortran 原文加一层 C}

这一域的骨干库有共同的出身。20 世纪 70 到 90 年代，数值库由 netlib 与阿贡、国家实验室一脉分发：\Lib{BLAS}、\Lib{LAPACK}、\Lib{FFTW} 的早期内核、\Lib{ARPACK}、\Lib{SPICE} 类工具，源语言都是 \FTN{}。理由很朴素：那是当时唯一能把数组、复数与向量化硬件说清楚的语言，而 \CL{} 直到 C99 才有像样的复数类型。

留下的技术后果有四项，今天仍然能从接口签名里读出来。

\textbf{其一，名字修饰的约定}。\FTN{} 编译器普遍在符号尾部加下划线、把名字转小写、并按引用传递每一个实参。\Lib{Cython} 与 \Lib{ctypes} 用户写 `\Tp{dgemm\_}` 而不是 `\Tp{dgemm}` 的困惑正来自这里；`\Tp{cblas\_}` 前缀则是后来 \CL{} 接口层为了不与 \FTN{} 符号冲突而主动加的名字空间。

\textbf{其二，按引用传递}。\FTN{} 把标量也当指针传，所以 \FTN{} 原码里的标量实参在 \CL{} 侧要写成 `\Tp{\&n}`。

\begin{lstlisting}[style=fortranstyle,caption={一段 \FTN{} 原码与它的调用形态：标量也是按引用传递},label={lst:ch02-fortran}]
      SUBROUTINE DAXPY(N, DA, DX, INCX, DY, INCY)
      DOUBLE PRECISION DA
      DOUBLE PRECISION DX(*), DY(*)
      INTEGER INCX, INCY, N
      INTEGER I
      IF (N .LE. 0) RETURN
      IF (DA .EQ. 0.0D0) RETURN
      IF (INCX .EQ. 1 .AND. INCY .EQ. 1) GO TO 20
      ...
   20 CONTINUE
      DO I = 1, N
         DY(I) = DY(I) + DA * DX(I)      ! 唯一的密集运算
      END DO
      END SUBROUTINE
\end{lstlisting}

\textbf{其三，步长参数}\Bk。\Tp{INCX}、\Tp{INCY} 与后来 \CL{} 接口里的 \Tp{ld}（leading dimension）是同一个概念的两副面孔：接口不假设你的数组连续，只承诺``相邻两个逻辑元素隔着这么多字节''。这一设计让同一份内核能服务子矩阵、转置视图与交错音频缓冲，代价是调用方要对自己算出的步长负责。

\textbf{其四，返回码与状态参数}。\FTN{} 没有异常，于是 \Lib{LAPACK} 用 \Tp{info} 整型出参表示失败原因，并把``第几个参数非法''编码进去。\Lib{LAPACKE} 做的第一件事就是把这个出参改回返回值，再把行列主序做成宏。

现代 \FTN{} 与 \CL{} 的互通不再靠猜修饰规则：\FTN{} 2003 起的 \Tp{iso\_c\_binding} 模块与 \Tp{bind(c)} 属性可以精确指定符号名与参数类型。

\begin{lstlisting}[style=fortranstyle,caption={用 \Tp{bind(c)} 把 \FTN{} 内核暴露成 \CL{} 符号},label={lst:ch02-iso}]
module dsp_kernel
  use, intrinsic :: iso_c_binding
  implicit none
contains
  subroutine fft_stage(re, im, n) bind(c, name='fft_stage')
    use, intrinsic :: iso_c_binding
    real(c_float), intent(inout) :: re(*), im(*)   ! 与 C 的 float* 等价
    integer(c_int), value        :: n              ! value：按值传，像 C
  end subroutine
end module
! 编译后要确认两件事：符号确实叫 fft_stage（无尾下划线），
! 且 n 是按值而不是按引用——ctypes 侧的签名要与后者一致。
\end{lstlisting}

\begin{warning}[title={跨语言边界的三种静默错位}]
第一，符号名对不上：链接期还能靠 \Tp{extern \textquotedbl C\textquotedbl} 修好，运行期用 \Tp{dlsym} 或 \Tp{ctypes} 找不到就是 ``\Tp{AttributeError}'' 这类异常，而 \FTN{} 加了下划线的版本又找得到，于是出现``一半能跑''的诡异状态。第二，按值与按引用混用：\FTN{} 默认按引用，\CL{} 默认按值，声明错一个标量参数，得到的是错误结果而非崩溃。第三，主序与步长算错：矩阵乘法返回``看起来合理''的数，只在特定尺寸下才暴露。三者都不会有任何诊断信息。
\end{warning}
\section{C ABI 为什么是唯一汇合点}

``C ABI''指的是平台文档里明确规定的那一套：参数放在哪些寄存器或栈槽、谁来清栈、结构体字段之间插多少填充、基本类型的宽度与对齐、返回值走哪条通道。\CL{} 之所以能成为汇合点，不是因为它的表达力强，而是因为它足够小：没有对象头、没有虚表、没有异常表、没有引用计数、没有 GC 句柄、没有泛型运行时。一条边界上不存在两种运行时会谈不拢的东西。

\CPP{} 的 \Tp{extern \textquotedbl C\textquotedbl} 只改变一件事：这个名字不做 \CPP{} 的名字修饰、不参与重载解析。它不改变调用约定（Windows 上仍可能是 \Tp{stdcall} 与 \Tp{cdecl} 的差别），不改变结构体布局规则，也不让异常跨边界安全传播。理解这一点，才能理解为什么绝大多数``\CPP{} 库的对外接口''写成句柄加平面函数：那是把 \CPP{} 的运行时留在库内部。

\section{三种接口形态与它们的可达范围}

\textbf{句柄式平面函数}（opaque handle）。\Tp{SNDFILE*}、\Tp{fftw\_plan}、\Tp{PJ*}、\Tp{CVODE\_Mem}、\Tp{btDynamicsWorld} 的 C 封装都属于此类：类型在头文件里只有前向声明，生命周期由成对的 create/destroy 管，错误用返回码。这是跨语言绑定的最短路径，因为它只需要类型宽度（全是指针）与调用约定。

\textbf{平面数值函数}（flat function）。\Lib{C ABI 的 BLAS}、\Lib{CMSIS-DSP}、\Lib{sundials} 的向量与矩阵操作以裸数组加长度暴露，没有对象也没有状态；可分离编译、可被任何语言逐点调用，代价是没有类型系统保护（尺寸、步长、主序全靠人）。

\textbf{模板头文件库}。\Lib{Eigen}、\Lib{CGAL} 的 Kernel、\Lib{Boost} 系列属于此类：接口即实现，编译期内联，零运行时开销，尺寸与对齐可以进类型系统。代价是没有 ABI 可言——每次改动都要重编译所有使用者，而且任何外部绑定工具都必须在 \CPP{} 层面实例化模板才能包起来。

\begin{table}[H]
\centering
\caption{三种接口形态的可达范围}
\label{tab:ch02-forms}
\small
\begin{tabular}{@{}>{\raggedright\arraybackslash}p{0.16\textwidth}
  >{\raggedright\arraybackslash}p{0.20\textwidth}
  >{\raggedright\arraybackslash}p{0.22\textwidth}
  >{\raggedright\arraybackslash}p{0.21\textwidth}@{}}
\toprule
形态 & 跨编译器 & 绑定难度 & 典型代表 \\
\midrule
句柄加平面函数 & 稳定（只需 C ABI） &
  低（自动生成或手写声明） &
  \Lib{libsndfile}、\Lib{PROJ}、\Lib{FFTW} \\
裸数组数值函数 & 稳定 &
  低，但易传错尺寸与步长 &
  \Lib{OpenBLAS}、\Lib{CMSIS-DSP} \\
模板头文件库 & 需同版本头与标准库 &
  高（要先实例化） &
  \Lib{Eigen}、\Lib{CGAL}、\Lib{Boost.Odeint} \\
\bottomrule
\end{tabular}
\end{table}

\section{四种（半）绑定路线}

把上面的接口交给 \PY{}，有四条主流路线加一条例外。它们的差别不在``能不能''，而在需不需要 C 编译器、能不能处理所有权、以及出错时报在哪一层。

\textbf{\Tp{ctypes}}：纯标准库、运行时按 ABI 直接调 \Tp{.so}/\Tp{.dll}，不需要编译器，也不读头文件。类型表只覆盖基本类型与指针，结构体要手抄一遍布局，回调要自己保证生命周期。优点是零构建、跨发行版；缺点是任何签名错误都表现为``要么崩溃要么结果错''。

\textbf{\Tp{cffi}}：两种模式。ABI 模式与 \Tp{ctypes} 同构；API 模式把声明交给 C 编译器生成胶水，于是能校验结构体布局、能处理宏与内联函数，产物是需要编译的扩展模块。它是 \Lib{libsndfile} 的 Python 门面 \Lib{soundfile} 曾经走过的路线之一（另一条是 \Tp{ctypes}）。

\textbf{\Tp{Cython}}：写``带类型的 \PY{}''再转译成 C，既能为已有 C 库生成绑定，也能自己写内核；与 numpy 的数组协议结合最自然，科学计算栈里大量热路径实际是 Cython 产物（\Lib{gprMax} 的 FDTD 内核即是一例）。代价是又引入一门需要学习的中间语言与一套构建。

\textbf{\Tp{pybind11}}：面向 \CPP{} 头文件的绑定，模板与 STL 都能自动映射，支持继承、智能指针持有者与异常翻译；它是 \Lib{Open3D}、\Lib{libigl} 的 Python 侧做法。代价是必须编译、必须为每个用到的模板实例化、并且要处理 GIL 与所有权语义。它的细节（\Tp{py::array\_t}、零拷贝、桩文件与 wheel 分发）在并行算法分册里，本册只在门面章引用其结论。

\textbf{（半条形）\Tp{numba} 与运行时编译}：不做绑定，而是把 \PY{} 函数 JIT 成本地码，只认 numpy 数组与少数 C 函数指针。适合临时提速一段算法实验，不适合做库的分发形态。

\begin{lstlisting}[style=cstyle,caption={同一个 C 接口的三种声明：库侧、\Tp{ctypes} 侧、Cython 侧},label={lst:ch02-three}]
// 库侧（C 头文件）：句柄 + 返回码，全是 ABI 可见的东西
typedef struct SfInfo { int samplerate; int channels; int format;
                        long frames; } SfInfo;
int sf_probe(const char* path, SfInfo* out);   // 0 表示成功

# ctypes 侧：布局要手抄，字段顺序与宽度自己保证
# import ctypes
# lib = ctypes.CDLL('./libdspker.so')
# class SfInfo(ctypes.Structure):
#     _fields_ = [('samplerate', ctypes.c_int),
#                 ('channels', ctypes.c_int),
#                 ('format', ctypes.c_int),
#                 ('frames', ctypes.c_long)]
# lib.sf_probe.argtypes = [ctypes.c_char_p, ctypes.POINTER(SfInfo)]
# lib.sf_probe.restype = ctypes.c_int

# Cython 侧：直接声明头文件里的类型，编译期由 C 校验布局
# cdef extern from "dspker.h":
#     struct SfInfo:
#         int samplerate
#         int channels
#         int format
#         long frames
#     int sf_probe(const char* path, SfInfo* out)
\end{lstlisting}

\begin{compare}[title={五条绑定路线的代价对照}]
\begin{tabular}{@{}>{\raggedright\arraybackslash}p{0.15\textwidth}
  >{\raggedright\arraybackslash}p{0.17\textwidth}
  >{\raggedright\arraybackslash}p{0.19\textwidth}
  >{\raggedright\arraybackslash}p{0.24\textwidth}@{}}
\toprule
路线 & 需要编译器 & 能处理的接口 & 典型失败模式 \\
\midrule
\Tp{ctypes} & 不需要 & 句柄、平面函数 &
  签名或布局抄错，崩溃或静默错值 \\
\Tp{cffi}（ABI） & 不需要 & 同上 &
  与 \Tp{ctypes} 相同 \\
\Tp{cffi}（API） & 需要 & 头文件、宏、内联 &
  构建环境与目标平台不一致 \\
\Tp{Cython} & 需要 & 头文件加自研内核 &
  类型标注与真实 C 类型不符 \\
\Tp{pybind11} & 需要且要 \CPP{} 头 & 模板、类、STL &
  实例化缺失、GIL 与所有权误判 \\
\bottomrule
\end{tabular}
\end{compare}

\section{给 \CPP{} 库加一层 C 门面}

如果目标是被多种语言复用，``\CPP{} 内核加 \CL{} 门面''是标准做法：门面只做参数校验、异常截断与句柄转换，不含业务逻辑。

\begin{lstlisting}[style=cppstyle,caption={把 \CPP{} 内核包成句柄式 C 接口：异常在边界内消化},label={lst:ch02-handle}]
// 内核用 C++：RAII、vector、异常都在库内部
class Spectrum {
 public:
  Spectrum(int n, int fs);
  void feed(const float* x, int n);
  double peakHz() const;
};

// 门面：不透明句柄、返回码、无模板
struct SpectrumCtx;                       // 只在 .cpp 里定义
extern "C" {

SpectrumCtx* spectrum_create(int n, int fs) {
  try {
    return reinterpret_cast<SpectrumCtx*>(new Spectrum(n, fs));
  } catch (...) {
    return nullptr;                       // 失败信息留给错误码接口
  }
}

int spectrum_feed(SpectrumCtx* h, const float* x, int n) {
  if (!h || !x || n <= 0) return -1;      // 参数校验在边界做一次
  try {
    reinterpret_cast<Spectrum*>(h)->feed(x, n);
    return 0;
  } catch (...) {
    return -2;                            // 异常绝不穿过 C 边界
  }
}

void spectrum_destroy(SpectrumCtx* h) {
  delete reinterpret_cast<Spectrum*>(h);  // 与 create 成对
}

}  // extern "C"
\end{lstlisting}

四条纪律写在这段代码里，值得单独点出：所有权在创建方与销毁方之间是显式配对的；异常一律在边界转成返回码；句柄类型对调用方只暴露前向声明；参数校验只做一次而不是每层都做。

\section{边界上的四类共同问题}

\textbf{所有权}。\CL{} 没有析构约定，所以``谁分配谁释放''必须由文档写成类型规则：\Tp{fftw\_free} 只能配 \Tp{fftw\_malloc}，\Tp{sf\_close} 会连带关掉你交给它的 \Tp{FILE*}（取决于 flags）。\PY{} 侧要把它翻译成 finalizer，且必须保证 finalizer 的调用顺序不会用到已释放的上下文。

\textbf{错误传递}。返回码、\Tp{errno} 式线程局部状态、库自己的错误消息查询函数与异常四种风格并存（前者如 \Tp{proj\_errno\_string}，后者如几何库那组同名消息接口）。工程上唯一稳妥的做法是：在门面层把错误收敛成一种，并让线程局部错误码可查询。

\textbf{线程与 GIL}。\CPP{} 侧起线程回调 \PY{}，必须先取回 GIL；\PY{} 调用长耗时 C 函数时，绑定层应当在调用前释放 GIL（\Tp{pybind11} 的 \Tp{gil\_scoped\_release}、\Tp{Cython} 的 \Tp{with nogil}），否则多线程数值库在 Python 里退化成单线程。

\textbf{可重入与全局状态}。\Lib{GSL} 的错误处理器、\Lib{PROJ} 与 \Lib{GEOS} 的上下文、\Lib{FFTW} 的 wisdom 文件与 planner 锁都是全局或半全局资源。带 `\Tp{\_r}` 后缀的可重入版本（\Tp{GEOS\_init\_r}）是线程安全的正确入口。

\section{ABI 稳定性是一种工程承诺}

选 \CL{} 接口层的另一收益是``实现可替换''。\Lib{OpenBLAS} 换成 \Lib{MKL}、\Lib{FFTW} 换成 \Lib{KissFFT}，只要符号与语义一致，调用方一行不改。反过来，\CPP{} 接口一旦暴露类布局、内联函数与标准库类型，换编译器版本就可能链接失败或运行期错乱；这也是 ``\Tp{\_GLIBCXX\_USE\_CXX11\_ABI}'' 这类宏存在的历史原因。\Lib{PIMPL} 惯用法是 \CPP{} 侧对 ABI 稳定的补救，但它的代价（多一次间接、多一次分配）恰好与本章的动机冲突。

\begin{bestpractice}[title={一份 C 门面头文件的写法}]
\begin{itemize}
\item 对外只暴露 \CL{} 可达的东西：指针大小的句柄、定宽整型、裸浮点数组加长度、返回码。
\item 类型名写定宽（\Tp{int32\_t}、\Tp{float}），不要在 ABI 上用 \Tp{int}、\Tp{long} 与 \Tp{size\_t} 之外的宽度含糊类型。
\item 生命周期成对命名（\Tp{create}/\Tp{destroy}、\Tp{plan}/\Tp{free}），并在文档里把这一对写死。
\item 异常、\Tp{std::string}、\Tp{std::vector} 与模板实例不进 ABI；需要时只出现在 \CPP{} 门面里。
\item 给符号加库前缀，并把``符号名到底是什么''用 \Tp{nm} 与 \Tp{dumpbin} 的实测结果写进测试。
\end{itemize}
\end{bestpractice}

\begin{guideline}[title={接口层的四条结论}]
\begin{itemize}
\item ``用 C/C++''在这一域的真正含义是：把 C ABI 作为对外契约，把实现语言留在契约之后。
\item \FTN{} 的历史约定（按引用、尾下划线、步长参数、info 出参）仍在接口签名里，读接口文档时要认得出它们。
\item 绑定路线的选择取决于接口形态：句柄加平面函数任选 \Tp{ctypes} 或 \Tp{cffi}，模板头文件库只能走编译式绑定。
\item 所有权、错误传递、GIL 与全局状态是四种跨越边界必然要处理的问题，早定规则比事后调崩溃便宜。
\end{itemize}
\end{guideline}
% ====================================================================

% ====================================================================
\chapter{数值接口的共同约定：布局、复数、错误与线程}

把 \Lib{OpenBLAS}、\Lib{FFTW}、\Lib{libsndfile}、\Lib{PROJ} 与 \Lib{Sundials} 的签名摆在一起看，会发现它们反复使用同一组约定：指针加长度、步长、返回码、显式的上下文对象、以及``单位写在文档里''。这些约定不是某个库的风格，而是 \CL{} 接口层在数值域的必要产物——没有析构、没有异常、没有类型级别的形状信息，就必须用别的东西把缺失的信息补上。本章把它们集中讲清，因为这一域绝大多数缺陷不是崩溃，而是这些约定被误读之后的静默错误。

\section{主序、步长与视图}

\textbf{主序}决定同一份线性内存被解释成哪种矩阵。\CL{} 的多维数组与 numpy 默认是行主序（row-major），\FTN{} 与 \Lib{BLAS}/\Lib{LAPACK} 默认是列主序（column-major）。\textbf{步长}（\Tp{ld}、\Tp{incx}、\Tp{stride}）则回答一个更基本的问题：相邻两个逻辑元素在内存里隔多少\emph{元素}。列主序下 \Tp{ld} 是列与列之间的距离，行主序接口（\Tp{cblas\_} 系列）里它是行与行之间的距离。

关键推论有三条。第一，转置不需要搬数据，只需要交换步长——``视图''的全部本钱就在这里，这也是 numpy 的 \Tp{a.T} 是常数时间、而 \Tp{a.copy()} 不是的原因。第二，任何数组切片都是带步长的视图，非连续视图交给要求连续的库时会触发隐式拷贝（或被拒绝）。第三，交织音频帧本质就是一个二维数组：帧数为外维、通道数为内维，通道交织的 stride 为 $1$，平面（planar）布局则是另一侧为 $1$。

\begin{lstlisting}[style=cstyle,caption={行主序数据交给列主序内核：靠转置恒等式避免搬数据},label={lst:ch03-order}]
// 三个矩阵都按行主序存放在内存里：A 是 m x k，B 是 k x n，C 是 m x n。
// cblas_dgemm 的 RowMajor 参数只是接口糖；若只能调列主序的 dgemm_，
// 就把"行主序的 C = A*B"读成"列主序的 C' = B'*A'"，其中 X' 指同一块
// 内存的列主序解释（内存不动，只是换个看法）。
extern void dgemm_(const char*, const char*, const int*,
                   const int*, const int*, const double*,
                   const double*, const int*, const double*,
                   const int*, const double*, double*, const int*);
void row_major_gemm(int m, int n, int k,
                    const double* A, int lda,
                    const double* B, int ldb,
                    double* C, int ldc) {
  const double alpha = 1.0, beta = 0.0;
  dgemm_("N", "N", &n, &m, &k,     // 行列尺寸对调成 n, m
         &alpha, B, &ldb,          // 第一个因子换成 B
         A, &lda,                  // 第二个因子换成 A
         &beta, C, &ldc);          // beta 为零：覆盖，不累加
}
// 自检两条：lda 至少是不连续那一维的长度；
// 结果若"整体转置了"，几乎一定是顺序读反而不是算错。
\end{lstlisting}

\begin{warning}[title={步长与尺寸的三种误用}]
把 \Tp{ld} 写成元素个数之外的东西（例如字节数）是第一种；复用同一块缓冲当不同形状的矩阵、却没有重算 \Tp{ld} 是第二种；第三种最阴：\Tp{ld} 比实际需要的间距大，运算看起来正确，直到某次分配器把两块内存紧挨着放，越界写进了邻居。前两种能在小规模测试里被发现，第三种通常只在生产环境出现。接口层能做的防护是把``每维长度加每维步长''一起校验，并在调试构建里用带边界的分配器。
\end{warning}

\section{复数的三种表示与半谱}

信号处理里复数是基带、频谱与解析信号的日常载体，但它在接口上有三种互不兼容的形态。

\textbf{交错存储}：实部与虚部相邻，即 \Tp{double z[2]} 或 \Tp{fftw\_complex}（它就是长度为 2 的数组）。\Lib{FFTW} 的接口文档明确要求这一形态。\textbf{分离存储}（split）：\Tp{re[n]} 与 \Tp{im[n]} 两条数组，向量化时两条通道各自连续，SIMD 写法更直接。\textbf{语言原生复数类型}：C99 的 \Tp{\_Complex}、\FTN{} 的 \Tp{COMPLEX}、\PY{} 的 \Tp{complex64}。前两种是\emph{内存布局}，第三种是\emph{类型}，二者不能混谈。

必须注意的一点：C99 复数类型虽然内存上等价于两个实数，但它的\emph{传值规则}由实现定义（可能走寄存器对而非指针）。跨语言 ABI 不要用原生复数类型传参，改用``指针加长度''或两个实数数组——这也是 \Tp{fftw\_complex} 选择数组形式的原因之一。

频谱侧的公共约定是共轭对称：长度为 $N$ 的实序列做 FFT，非冗余的 bin 只有 $\lfloor N/2 \rfloor + 1$ 个，因此实数变换（r2c/c2r）的接口只给半谱，输出数组的复数元素个数是 $N/2+1$ 而不是 $N$。把这一条读错，得到的是``后半段频谱全零''或越界，而不是任何错误提示。

\begin{table}[H]
\centering
\caption{三种复数形态的适用场合}
\label{tab:ch03-complex}
\small
\begin{tabular}{@{}>{\raggedright\arraybackslash}p{0.19\textwidth}
  >{\raggedright\arraybackslash}p{0.26\textwidth}
  >{\raggedright\arraybackslash}p{0.35\textwidth}@{}}
\toprule
形态 & 谁用 & 代价 \\
\midrule
交错（re,im 相邻） &
  \Lib{FFTW}、\Lib{cuFFT}、基带采样流 &
  向量化要做 deinterleave；缓存行了却友好 \\
分离（两条数组） &
  手写 SIMD 内核、几何与物理求解器 &
  两次访存流；接口要多传一个指针 \\
语言原生类型 &
  \PY{}、\FTN{}、\CL{} 内部计算 &
  ABI 传参规则不统一，不适合当对外接口 \\
\bottomrule
\end{tabular}
\end{table}

\section{单位、时间基与坐标系}

数值与信号库的接口不会检查单位，因此``单位约定''是接口文档的一部分，也是事故清单上出现频率最高的一项。本书后面几章会反复遇到这些具体形式：音频的样本值域与 dBFS、采样率与``每样本多少秒''；几何的单位（米还是毫米）与坐标 handedness；射电与光学的频率还是波长、近场还是远场、电场还是功率。工程判据只有一条：任何跨越模块边界的量，要么在类型上写死单位（强类型或封装），要么在名字与文档上写死，两者都不做就等于把错误推给下游。

\section{错误处理的四种风格}

\textbf{返回码}。\Lib{libsndfile}、\Lib{GEOS}、\Lib{PROJ}、\Lib{Sundials} 的主体接口都返回整型或指针，零/非零或负值表示失败。信息量小但跨语言无成本。

\textbf{库专属的状态查询}。\Tp{sf\_error\_str} 与 \Tp{proj\_errno\_string} 是这一类：调用只给一个码，详细描述要从库里再取一次（几何库也有对应的消息查询函数，名字同样长，此处不列）。这带来一个隐藏约束——错误字符串的所有权与有效期通常属于该线程的该上下文，复制前不要跨线程使用。

\textbf{线程局部或全局错误状态}。\Tp{errno} 式设计与 \Lib{GSL} 的错误处理器属于此类；危险点在于``上一次的错误''会伪装成``这一次的错误''，所以每次调用都要显式清零或立即读取。

\textbf{异常}。只在 \CPP{} 门面与 \PY{} 侧出现。跨边界必须翻译：\Tp{pybind11} 能把 \CPP{} 异常译成 \PY{} 异常，但反方向（从 \CL{} 回调里抛进 \CPP{}）与穿过 \Tp{extern \textquotedbl C\textquotedbl} 都不安全。

\section{线程模型与环境变量}

同一批库在线程上分三类，接口签名里看不出来，只能从文档与实现里读。第一类\emph{内部已并行}：\Lib{OpenBLAS}、\Lib{MKL}、\Lib{FFTW} 的 threads 版本，它们自己起线程，你只该在它们外面写串行循环。第二类\emph{假定调用方并行}：内核函数无共享状态，可以按数据分片从多个线程调用（多数 \Lib{Eigen} 表达式、\Lib{CMSIS-DSP} 的滤波器实例）。第三类\emph{有全局状态，不保证可重入}：老式随机数、全局上下文、wisdom 文件的读写、\Lib{GSL} 的错误处理。

\textbf{嵌套超订}是三类混用时最常见的性能事故：外层 \PY{} 用多线程库，内层 BLAS 又起满核，线程数超过核数后吞吐塌到谷底且伴随长尾。诊断顺序应当是：先看运行时的实际线程数（\Tp{threadpoolctl}、\Tp{top -H}），再看是否为同一并行域套了两层，最后才怀疑算法。环境变量矩阵见表~\ref{tab:ch03-env}，它们必须由部署方而不是库来统一设置。

\begin{table}[H]
\centering
\caption{常见线程与后端相关环境变量（名称以各库文档为准）}
\label{tab:ch03-env}
\small
\begin{tabular}{@{}>{\raggedright\arraybackslash}p{0.27\textwidth}
  >{\raggedright\arraybackslash}p{0.20\textwidth}
  >{\raggedright\arraybackslash}p{0.33\textwidth}@{}}
\toprule
变量 & 影响谁 & 误用的表现 \\
\midrule
\Tp{OMP\_NUM\_THREADS} &
  以 OpenMP 为后端的库 &
  设成核数再套一层并行，长尾暴涨 \\
\Tp{OPENBLAS\_NUM\_THREADS} &
  \Lib{OpenBLAS} &
  小矩阵上并发启动开销压过收益 \\
\Tp{MKL\_NUM\_THREADS} &
  \Lib{MKL} &
  与 \Tp{OMP\_NUM\_THREADS} 不一致 \\
\Tp{FFTW\_WISDOMPATH} 类 &
  \Lib{FFTW} 的 wisdom &
  首次调用做 planner，进程启动变慢 \\
\Tp{OMP\_PROC\_BIND} &
  线程到核的绑定 &
  与容器 CPU 配额冲突导致抖动 \\
\bottomrule
\end{tabular}
\end{table}

\section{精度、可复现与``逐位一致''的边界}

这一域的精度选择由数据规模与误差累积决定，而不是由``越高越好''决定。音频与实时处理用 \Tp{float}（单精度），因为采样率高、内存带宽与 SIMD 宽度都是约束，且 $10^6$ 量级的样本数下累积误差可用分块与双精度累加器控制；射频基带用复单精度；结构、几何谓词与参数拟合用 \Tp{double}；只有精度表与验证基准会动用扩展精度。

\textbf{可复现性有三个层次}，工程上必须说清要哪一层。第一层``统计一致''：结果在容差内，适合蒙特卡洛与机器学习式流程。第二层``算法确定''：归约顺序、分块尺寸与迭代终止条件固定，因此同一二进制同一输入逐位一致，适合回归测试。第三层``跨实现一致''：换 BLAS、换 FFT 库、换编译开关仍然逐位一致——这一层几乎不可能，因为归约顺序与融合乘加的使用是实现的自由。

\begin{bestpractice}[title={对拍测试的分工}]
\begin{itemize}
\item 用两种实现互验（自家 C 内核与 \PY{} 参考实现），而不是只和自己的旧版本比。
\item 判据分层：形状、单位、量级用精确比较；浮点结果用相对与绝对容差的合式判据，容差按数据尺度归一化，不写死 $10^{-8}$。
\item 把 golden 文件与产生它的输入、环境变量、库版本一起入库，否则``基线''会在三个月后无法解释。
\item 明确禁止宽松浮点开关进入发布构建，并在测试里加一条会因重排而失败的断言。
\end{itemize}
\end{bestpractice}

\begin{guideline}[title={共同约定五条}]
\begin{itemize}
\item 主序与步长是同一件事的两面：先确定``哪个维度 stride 为 1''，再谈转置、切片与拷贝。
\item 复数的存储形态是接口的一部分；对外 ABI 不用语言原生复数类型，交错与分离只能选一种并写进文档。
\item 单位、时间基与坐标系必须在名字或文档里出现，接口不会替你检查它们。
\item 错误处理与线程模型要在集成前问清：谁起线程、错误状态属于哪个线程、上下文能不能共享。
\item 可复现性要选层：统计一致、算法确定还是跨实现一致，三层的成本差一个数量级。
\end{itemize}
\end{guideline}
% ====================================================================

% ====================================================================
% ====================================================================
% frag_ch04.tex —— 第 4 章：BLAS 与 LAPACK
% ====================================================================
\part{数学库}

\chapter{BLAS 与 LAPACK：从 netlib 谱系到多实现}

这一章是全书主线最干净的样本：\Lib{BLAS} 与 \Lib{LAPACK} 由 \FTN{} 之人设计、以 \FTN{} 写就、数值语义至今留在 \FTN{} 里，但今天任何一种语言要想借到机器里的向量单元，都必须从 netlib 顺手留下的那层 C 接口穿过去。本章沿五节拆开几层机制：BLAS 的三层划分如何把性能全部压进 \Tp{gemm} 一个 kernel；符号修饰与传参约定如何逼出 \Tp{cblas\_} \Bk{} 前缀；\Tp{ld} 与 alpha/beta 这两个看起来``无聊''的接口细节各自承诺了什么；同一套头文件为什么能被五六家厂商各自重写、而结果只差速度。章末给出线程与环境变量的因果链、\PY{} 门面之下 \Tp{a @ b} \Bk{} 的真实去向，以及构建系统如何发现这些库。

\section{BLAS 的三层与``以 gemm 为中心''的工程理由}

\CL{} Basic Linear Algebra Subprograms 诞生于二十世纪七十年代末，初衷朴素得近乎反高潮：把线性代数里最常出现的几种小循环做成可替换的标准例程，让同一个程序在换机器时只换库、不改码。最初的接口只有向量级；八十年代末加入第三级，是为了适配新的存储层次。三层之间的差别不是``复杂度''，而是每搬一个字节能摊到多少次浮点运算：

\begin{itemize}
\item \textbf{Level 1（向量对向量）：}\Fn{ddot}、\Fn{daxpy}、\Fn{dnrm2}、\Fn{dscal}。
  运算量 $O(n)$，数据量 $O(n)$，算术强度是个常数。
  这一层在任何现代机器上都是内存带宽受限的：算力再强，数据没到就是没到。
  它存在的意义是作为其它两层的构件与语义积木，而不是性能来源。
\item \textbf{Level 2（矩阵对向量）：}\Fn{dgemv}、\Fn{dsymv}、\Fn{dtrmv}。
  $O(n^2)$ 次浮点运算对 $O(n^2)$ 次数据读取，算术强度仍然约为 $1$。
  迭代求解法的内循环住在这里，所以大规模迭代法的瓶颈通常是带宽而不是浮点单元——
  这一事实在解算器章节还会反复出现。
\item \textbf{Level 3（矩阵对矩阵）：}\Fn{dgemm}、\Fn{dsymm}、\Fn{dtrmm}、\Fn{dsyrk}。
  $O(n^3)$ 次运算对 $O(n^2)$ 次数据访问，算术强度随 $n$ 线性增长：
  只要把大矩阵切成装得进缓存的块，绝大部分运算都能在片上完成。
  这是唯一一层能让流水线、双发射 FMA 和宽向量单元同时吃饱的接口。
\end{itemize}

``以 gemm 为中心''因此不是审美，而是三层划分的推论。超级计算机排行榜上的 HPL 基准本质上就是把 \Fn{dgemm} 包了一层外壳，厂商为此对这一个 kernel 做手工指令级调优，收益外溢给所有把它当子程序的库。LAPACK 的阻塞算法在分块之后，主体工作仍然落在 Level 3 调用上；\Lib{BLIS} 干脆把``一个算子族都能化成 gemm 的形状''写成架构本身（所谓 gemm-like 与 trsm-like 两大族）。工程理由归结为一句话：\textbf{接口只有几十个例程、语义是精确的代数式、实现可以整块替换}，三条同时成立，厂商才愿意只优化一个 kernel 而让上游全部受益。参考实现与调优实现在同一签名下的差距，随 CPU 代次不同可以达到一个量级以上；但两者的数值关系是``同一个数学问题的不同舍入顺序''，这一点是换实现的安全边界。

补一个只有接口考古才给得出的细节：复数内积在 Level 1 里是\emph{两个}函数——\Fn{zdotc} 共轭、\Fn{zdotu} 不共轭。这个区分说明 BLAS 的接口纪律是把\emph{全部}数学语义写进符号名：调用者不需要、也不允许用布尔参数``顺便''改语义，实现者也因此能对两个分支分别打磨。hermitian 族的 \Fn{zherk} 同理，为复数对称性单独开例程而不是靠 flag。这套纪律正是``可替换''的前提：语义若藏在参数里，kernel 就藏不住。

\section{接口层为什么是 C：名字修饰、\Tp{cblas\_} 前缀与列主序}

把 \FTN{} 例程直接拿来给 C 调用，会撞上三堵墙，它们都属于``约定''而非``能力''问题。第一堵是符号修饰：netlib 源码里写的是 \Tp{DGEMM}，gfortran 导出的符号是 \Tp{dgemm\_} \Bk{}（全小写加尾下划线），而别的时代的编译器或全大写、或加前导下划线、甚至把参数表长度编码进符号名——C 侧没有任何可移植的办法写出``那个符号''的字面拼法。第二堵是传参约定：\FTN{} 默认按引用传递标量，\Tp{CHARACTER} 实参附带一个隐藏的长度参数，而 C 按值传标量。第三堵是数组形状：\FTN{} 的假定大小数组 \Tp{A(LDA,*)} \Bk{} 没有 C 能对等的描述符。C 语言恰好在这三点上都有全平台一致的 ABI，于是接口层落到 C 上——这层包装就是 netlib 的 CBLAS，其命名规则简单粗暴到接近透明：加 ``\Tp{cblas\_}'' 前缀，换 layout 枚举，标量按值传。

下面这段摘录自真实 \FTN{77} 签名（固定格式源码的续行行首是 \Tp{.}），读它只需要对照三件事：标量没有一个按 C 的方式传，字符带隐藏参数，数组是假定大小的。

\begin{lstlisting}[style=fortranstyle,caption={\FTN{77} 原文的 DGEMM 声明（摘录）：C 无法直接调用的三个原因},label={lst:ch04-fortran}]
      SUBROUTINE DGEMM(TRANSA,TRANSB,M,N,K,ALPHA,A,LDA,B,LDB,BETA,
     .                 C,LDC)
      DOUBLE PRECISION ALPHA,BETA
      INTEGER K,LDA,LDB,LDC,M,N
      CHARACTER TRANSA,TRANSB
      DOUBLE PRECISION A(LDA,*),B(LDB,*),C(LDC,*)
*     .. 省略：内部是朴素三重循环 ..
      END SUBROUTINE DGEMM
\end{lstlisting}

\CPP{} 与 \FTN{} 各自还有一半问题要用别的机制解决。\Tp{extern \textquotedbl C\textquotedbl} 站在 \CPP{} 一侧：它关掉 C++ 的名字改编（name mangling），让 \CPP{} 编译出的函数以裸 C 符号参与链接，从而能被 C 头文件声明、被一切认 C ABI 的运行时调用。\FTN{2003} 的 \Tp{iso\_c\_binding} 模块（\Tp{bind(C)} 子句与 \Tp{c\_double}、\Tp{c\_int} 类型）站在 \FTN{} 一侧：它让一个 \FTN{} 过程直接以 C 的名字与 C 的调用约定导出，省掉手写包装。两者解决的是桥的两端，互不替代；而且都不改变数据布局。

列主序从 \FTN{} 一路继承下来，是调用者最容易栽跟头的``布局遗产''：二维数组在内存里按列连续排布，相邻元素上下相邻而非左右相邻。于是 \Tp{ld}（leading dimension，首维跨距）的真实含义是\emph{相邻两列起点之间的距离}，它在合法调用里不小于列主序视图下的行数，但完全可以更大——而``可以更大''恰是这个参数的存在意义：把一个子矩阵、一个转置视图、或为对齐填充过的缓冲区直接交给 kernel，零拷贝；超算里所谓 padding 避缓存冲突，落点就是抬高的 \Tp{ld}。把 \Tp{ld} 当成``矩阵尺寸''来传，是语义级误读。另一处布局遗产是索引基数：\FTN{} 数组下标从 $1$ 开始，LAPACK 的置换输出（比如 \Tp{ipiv} 里的行交换记录）也延续了这个约定，接进 C 的 $0$ 基数组前要逐个减一，这一步没有编译器帮你检查。

alpha/beta 是同类接口细节的第二个标本。\Tp{gemm} 的代数式是 $C \leftarrow \alpha\,\mathrm{op}(A)\,\mathrm{op}(B) + \beta\,C$，beta 为 $0$ 等于向实现承诺``不必读取 C 的旧值''——调优实现据此省掉一整轮加载；beta 为 $1$ 则把``累加''折进同一次写回。这不是风格，是接口承诺，因为它直接改写内存流量；也正因如此，``看起来一样''的 gemm 加 axpy 组合，交给单次 kernel 永远在内存账上占优。

清单~\ref{lst:ch04-dgemm} 用 CBLAS 在 C 里做一批齐次变换：$N$ 个齐次点排成行主序的 $N \times 4$ 缓冲，整条点云就是一次 \Fn{cblas\_dgemm} \Bk{}；第二次调用以 beta 为 $1$ 演示``融合累加''的形状。

\begin{lstlisting}[style=cstyle,caption={C 里用 \Tp{cblas\_dgemm} 做批量齐次变换与加权累加},label={lst:ch04-dgemm}]
#include <cblas.h>                /* 这是 CBLAS 头，不是 Fortran 的 */

void transform_batch(int n_pts, const double* pts,
                     const double* m, double* out) {
  /* CBLAS 签名里 B 是非 const 指针：实现理论上可覆写它，
   * 只读缓冲先拷一份最稳妥（个别发行版的头已修成 const）。 */
  double mbuf[16];
  for (int i = 0; i < 16; ++i) mbuf[i] = m[i];

  /* out = 1.0 * pts * m + 0.0 * out
   * beta=0 是对实现的承诺：out 不必被读，省一轮内存流量。 */
  cblas_dgemm(CblasRowMajor, CblasNoTrans, CblasNoTrans,
              n_pts, 4, 4,
              1.0, pts, 4,       /* alpha, A, lda=4：行跨距 */
              mbuf, 4,           /* B, ldb=4 */
              0.0, out, 4);      /* beta, C, ldc=4 */

  /* beta=1：把 0.5 * pts * m 累加进 out，
   * 不再需要一次独立的 axpy 扫描。 */
  cblas_dgemm(CblasRowMajor, CblasNoTrans, CblasNoTrans,
              n_pts, 4, 4,
              0.5, pts, 4, mbuf, 4,
              1.0, out, 4);
}
\end{lstlisting}

\begin{warning}[title={主序与步长写错不会有任何报错}]
\begin{itemize}
\item \Tp{ld} 传成矩阵尺寸而非跨距、或行主序缓冲按列主序解释，\textbf{不会有任何报错}：
  kernel 忠实地按你给的步长读内存，产出``形状正确、数值错误''的结果。
\item 症状通常是可复现的错误答案而非崩溃——除非 \Tp{ld} 小到让读取越出分配边界，
  那时才表现为随机段错误，回头查已经隔了好几层调用。
\item 自保手段三件套：先用一个 \Tp{ld} 大于实际尺寸的子矩阵做小规模对照（与 \Fn{dgemv} \Bk{}
  逐行核对）；把行主序调用集中到一个薄封装里，边界只转置一次；
  用对称矩阵做主序探针——两种解释读出的结果应当完全一致。
\end{itemize}
\end{warning}

\section{LAPACK 的两层结构与 LAPACKE 的 C 封装}

\Lib{LAPACK} 在 BLAS 之上再分层，目录结构就写着这个分层。\emph{driver} 层（\Fn{dgesv} 解方阵、\Fn{dsyev} 实对称特征值、\Fn{dgesdd} 分治 SVD、\Fn{dgels} 最小二乘）做参数检查、必要的平衡缩放、调用 computational 层、再做后处理；\emph{computational} 层（\Fn{dgetrf} 一个 LU 分解、\Fn{dgetrs} 用分解回代、\Fn{dpotrf} Cholesky）只做一件事；外面还有一圈 \emph{auxiliary}（\Fn{dlacn} 估计条件数、\Fn{ilaver} 查询版本）。工程结论很直白：\textbf{绝大多数人只需要 driver}，因为平衡与后处理正是手写调用最容易漏掉的部分。值得下沉到 computational 层的场景只有三类：一次分解、多次右端回代（\Fn{dgetrf} 加多次 \Fn{dgetrs}，摊薄 $O(n^3)$ 的分解代价）；要精确控制 workspace、拒绝封装层替你 malloc；以及把分解嵌进自定义流水线、不接受 driver 替你做的缩放。

\Lib{LAPACKE} 把这套语义翻译成了 C。翻译不只是改名字，它做了三件接口层面的事。其一，标量参数改为按值传递，整型参数统一为 \Tp{lapack\_int} \Bk{}——这个 typedef 是 LP64 与 ILP64 两种索引宽度的分野，把它的实际宽度写死在调用方代码里，是``符号在、行为崩''式链接事故的经典来源。其二，Fortran 放在末位的 \Tp{info} 变成\emph{返回值}：负值是非法参数的序号（按 LAPACKE 的 C 参数序计），正值是算法级失败（比如特征值迭代不收敛），调用者从此无法``忘记''检查它——除非把返回值直接丢弃，这在评审里应当视同 bug。其三，首参增加 \Tp{lapack\_layout} \Bk{} 枚举：\Tp{LAPACK\_COL\_MAJOR} 直通底层，\Tp{LAPACK\_ROW\_MAJOR} 让封装层负责转换——多数例程用分块转置把缓冲复制成列主序再算、算完复制回来（每次调用一份额外 $O(mn)$ 的拷贝代价，但只有一份，且对调用者透明）；封装层内部对个别热点例程有更省的路径，调用者不应指望。带状矩阵等可选例程组是否编入，取决于发行时的编译开关，与调用点写法无关。

清单~\ref{lst:ch04-dsyev} 是最小可编译的 LAPACKE 用法：行主序缓冲、driver 层求特征值。需要复用分解的 computational 层写法在清单~\ref{lst:ch04-lu}，它的结构正是``driver 替你做掉的那部分被你重新接管''。

\begin{lstlisting}[style=cstyle,caption={LAPACKE 的 \Tp{dsyev}：行主序、info 为返回值、输出覆写输入},label={lst:ch04-dsyev}]
#include <lapacke.h>
#include <stdio.h>

int symmetric_eigen(void) {
  /* 对称阵按行主序存全阵即可；uplo='U' 只要求上三角有效 */
  double a[9] = {4.0, 1.0, 2.0, 1.0, 3.0, 0.0, 2.0, 0.0, 5.0};
  double w[3];

  lapack_int info = LAPACKE_dsyev(LAPACK_ROW_MAJOR,
                                  'V', 'U', 3, a, 3, w);
  if (info < 0) {          /* 返回值为负：参数检查就没过 */
    fprintf(stderr, "arg %d illegal\n", (int)(-info));
    return 1;
  }
  if (info > 0) {          /* 返回值为正：算法级失败 */
    fprintf(stderr, "eigen iter failed: %d\n", (int)info);
    return 2;
  }
  /* w 为升序特征值；jobz='V' 时 a 被覆写为特征向量。
   * 预分配 workspace 的变体：LAPACKE_dsyev_work，
   * 尺寸先问 LAPACKE_dsyev_workspace。 */
  return 0;
}
\end{lstlisting}

\begin{lstlisting}[style=cstyle,caption={computational 层：一次 \Tp{dgetrf} 分解、多次 \Tp{dgetrs} 回代},label={lst:ch04-lu}]
#include <lapacke.h>

/* a 与 ipiv 来自同一次 LAPACKE_dgetrf；每列 rhs 原地写回解 */
int solve_many(lapack_int n, const double* a, lapack_int lda,
               const lapack_int* ipiv, double** rhs, int nrhs) {
  /* ipiv 沿用 Fortran 的 1 基下标：喂给 C 数组前要减一 */
  for (int k = 0; k < nrhs; ++k) {
    lapack_int info = LAPACKE_dgetrs(LAPACK_ROW_MAJOR, 'N',
                                     n, 1, a, lda, ipiv,
                                     rhs[k], 1 /* ldb */);
    if (info != 0) return info;  /* 失败即止，不复用旧分解 */
  }
  return 0;
}
\end{lstlisting}

\section{实现版图：线程模型、SIMD 后端、许可证与环境变量}

同一份 \Tp{cblas.h} 或 LAPACKE 头背后站着多种实现，差异集中在四个维度：线程模型、SIMD 后端的组织方式、可发现性（能否被构建系统自动找到）、以及许可证。表~\ref{tab:ch04-impls} 按这四个维度摆五家。

\begin{longtable}{@{}>{\raggedright\arraybackslash}p{0.15\textwidth}
  >{\raggedright\arraybackslash}p{0.16\textwidth}
  >{\raggedright\arraybackslash}p{0.18\textwidth}
  >{\raggedright\arraybackslash}p{0.21\textwidth}
  >{\raggedright\arraybackslash}p{0.13\textwidth}@{}}
\caption{五类 BLAS/LAPACK 实现的差异维度}\label{tab:ch04-impls}\\
\toprule
实现 & 线程模型 & SIMD 后端组织 & 可发现性与配套 & 许可证基调 \\
\midrule
\endfirsthead
\toprule
实现 & 线程模型 & SIMD 后端组织 & 可发现性与配套 & 许可证基调 \\
\midrule
\endhead
\Lib{netlib 参考实现} & 单线程，无并行层 &
  纯 \FTN{} 循环，仅靠编译器自动向量化 &
  一切发行版都能编；充当数值语义的地面真值 &
  宽许可（netlib 式），以仓库 LICENSE 为准 \\
\Lib{OpenBLAS} & 自带线程池（默认不依赖 OpenMP），可由运行时接口设线程数 &
  按 CPUID 运行时分派到每架构目录的调优 kernel，源自 GotoBLAS2 谱系 &
  \Tp{pkg-config} \Bk{} 有独立 \Tp{openblas.pc}；发行版常注册进 alternatives；另发 CBLAS 与 LAPACKE 兼容层 &
  BSD 系宽许可，以仓库 LICENSE 为准 \\
\Lib{MKL} & Intel OpenMP 或 pthread 两种线程层可切，亦有运行时接口 &
  逐代次手工调优 kernel，同时发 LP64 与 ILP64 两套接口 &
  工具链变量注入；CMake 可点名 \Tp{BLA\_VENDOR} \Bk{}；附带 FFT、向量数学与稀疏解算器 &
  Intel 自有许可：可免费使用但非开源，条款以其协议为准 \\
\Lib{BLIS} & 编译期三选一：无、pthreads、OpenMP &
  显式的三级宏 kernel、面板 kernel、微 kernel 结构；微 kernel 面向寄存器组形状，可移植性是卖点 &
  纯库需自行打包；不含 LAPACK 全家桶，常配参考 LAPACK &
  BSD 系宽许可，以仓库 LICENSE 为准 \\
\Lib{ArmPL}、富士通 \Lib{libgfmv} 等厂商库 & 与自家运行时、绑核与网卡拓扑协同（NUMA 策略内置） &
  针对自家指令集调优（SVE、A64FX 一类） &
  只随厂商工具链或 HPC 平台分发；对外保持 BLAS 兼容符号 &
  专有，随平台协议 \\
\bottomrule
\end{longtable}

从表里读出一条规律：\textbf{越靠近硬件的厂商，越把 BLAS 兼容接口当广告位}——kernel 必须是自己的，接口必须是公共的，否则没人换用。反过来，LAPACK 自带的对照测试套件之所以重要，正因为它是区分``实现差异''与``接口误解''的唯一手段：换一家实现，残差与正交性检验仍应过。

还有一个常被问反的问题值得记下：既然接口层是 C，为什么实现层至今仍有大量 \FTN{}？
因为 \FTN{} 的数组整体运算语法、对无别名的默示假设，
以及与之配套了几十年的厂商编译器优化，恰好是写 kernel 最省力的一族工具；
OpenBLAS 与 BLIS 的核心虽已 \CPP{} 化或 C 化，LAPACK 全家桶本体仍是 \FTN{}。
C 是``见面''的层，不是``干活''的层——这正是全书主线的措辞学：
ABI 汇合点与实现语言本来就可以分离。
想知道当前链接的是哪一版 LAPACK，调 \Fn{ilaver} \Bk{}（LAPACKE 侧是 \Fn{LAPACKE\_ilaver}）
把三个版本整数打出来即可；排查``到底链进了谁''时它比 \Tp{ldd} \Bk{} 更直接，
因为静态链接时 \Tp{ldd} 什么也看不见。

线程是这张版图里最容易把性能搞崩的一维。默认情况下，被调库按环境变量决定自己的并行度：专用变量（\Tp{MKL\_NUM\_THREADS} \Bk{}、\Tp{OPENBLAS\_NUM\_THREADS} \Bk{}）优先，通用的 \Tp{OMP\_NUM\_THREADS} \Bk{} 是各家共同的退路。灾难的标准配方是嵌套并行：外层一个 OpenMP 区域起 $T$ 个线程，每个线程里再调一次 gemm，被调库默认又各起 $T$ 个线程——瞬时 $T \times T$ 个 runnable 线程挤在 $T$ 个硬件线程上，上下文切换与 L2、L3 颠簸同时发生，吞吐可以掉一个台阶，而 CPU 利用率看着仍是满的。诊断顺序：第一，数线程（\Tp{ps -eLf} \Bk{} 或任务管理器，线程数远超核数即有嫌疑）；第二，把内层压到单线程再测一次（专用变量优先设置，例如 \Tp{MKL\_NUM\_THREADS=1}）；第三，若外层并行度本来就不固定，改用运行时接口（\Fn{mkl\_set\_num\_threads}、\Fn{openblas\_set\_num\_threads}）做逐调用控制，而不是环境变量。原则一句话：\textbf{并行只在一层展开}；选哪层，看每次调用的问题规模够不够喂饱一个多线程 kernel。

\begin{note}[title={接口、线程与 kernel 三层各有归属}]
三层各有各的归属：调用约定与符号名由 C ABI 决定，任何实现都必须服从，否则链接期就炸；并行度由环境变量与运行时 API 决定，是部署期决策，把线程数写死进代码几乎总是错的；而真正跑多快由 kernel 的分块与指令调度决定，藏在实现内部，调用者唯一能做的是\emph{给够可以分块的规模}：少碎调用、多批量 gemm。把这三层混在一起谈，是这一族库最常见的讨论事故。顺带一提线程安全的边界：对 disjoint 缓冲的并发 \Fn{cblas\_dgemm} \Bk{} 调用一般安全，因为这些例程是对各自缓冲的纯操作；真正带``状态对象''的坑（plan、handle 一类）要到 FFT 与解算器 API 里才会再遇到。
\end{note}

\section{进入路径：CMake、pkg-config 与 \PY{} 门面}

构建系统怎么``找到''这些库，本身就是一章的考古。CMake 的 \Tp{FindBLAS} 与 \Tp{FindLAPACK} 是\emph{探测}模块而非配置包：它编译试探程序，在系统库、MKL、OpenBLAS、macOS 的 Accelerate 之间挨个碰运气，结果以 \Tp{BLAS\_LIBRARIES} 一族变量给出；较新的 CMake（3.22 起）补上了导入目标 \Tp{BLAS::BLAS} 与 \Tp{LAPACK::LAPACK} \Bk{}，优先用它们。选定实现靠 \Tp{BLA\_VENDOR} \Bk{}，静态链接靠 \Tp{BLA\_STATIC} \Bk{}，两个开关都必须在 \Fn{find\_package} 之前设置。

\begin{lstlisting}[style=cmakestyle,caption={CMake 里点名实现并静态链接 BLAS 与 LAPACK}]
set(BLA_VENDOR OpenBLAS)      # 必须赶在 find_package 之前
set(BLA_STATIC ON)
find_package(BLAS  REQUIRED)
find_package(LAPACK REQUIRED) # 会连带探测它依赖的 BLAS

# 优先导入目标（CMake 3.22+）；老版本退回结果变量兜底
target_link_libraries(solver PRIVATE
  LAPACK::LAPACK BLAS::BLAS)
# 兜底写法：${LAPACK_LIBRARIES} ${BLAS_LIBRARIES}
\end{lstlisting}

静态与动态的选择有真实代价差：动态链接让``换实现''退化成改搜索路径或换 alternatives 链接（Linux 发行版常把 \Tp{libblas.so.3} \Bk{} 注册成可切换的替代项），代价是部署时要管 RPATH；静态链接把 kernel 焊死进产物，利于复现与打包，代价是对 MKL 这类``接口宽度有两种''的实现必须在链接期就决定 LP64 还是 ILP64——选错是``符号在、行为崩''。不走 CMake 的场合，\Tp{pkg-config --libs openblas} 给出即用的标志串；MKL 则更常以模块文件或工具链环境变量进入构建。

\PY{} 侧的叙事最短：\Tp{a @ b} \Bk{} 在 NumPy 里最终变成\emph{构建期}就选定好的那家实现的 \Fn{gemm}。调度链短到不需要图：\Tp{@} 进 \Fn{matmul} 的 C 循环，C 循环进 \Fn{gemm}。实现是哪一家、LP64 还是 ILP64、用哪个线程层，全部记录在 wheels 打包时的元数据里，用 \Fn{numpy.show\_config} \Bk{} 可以读出来；\Fn{scipy.linalg.get\_blas\_funcs} 则能把 \Tp{dgemm} 当作显式函数取出，绕过一切抽象直接喂缓冲。线程纪律与上一节完全同源：外层已经用多进程或多线程铺开的，用 \Tp{threadpoolctl} \Bk{} 在临界区里把被调库的线程限到 1，这是事实标准。整条链上没有任何一次``重新选择实现''的机会——选择在打包 wheels 时就已冻结。

\begin{lstlisting}[style=pythonstyle,caption={\PY{}：读出门面下接的是哪家实现，并限制其线程}]
import numpy as np
import threadpoolctl

print(np.show_config())          # 列出链接的 BLAS/LAPACK 与线程层
a = np.random.rand(4096, 4096)   # 两个 float64 矩阵：
b = np.random.rand(4096, 4096)   # 元素类型决定调 dgemm 还是 sgemm

# a @ b 的分派终点是所选实现的 cblas_dgemm；
# 小矩阵时这一次调用的固定开销可能比运算本身还大。
with threadpoolctl.threadpool_limits(limits=1):
    # 外层已有并行时把 BLAS 内部线程临时压到 1，
    # 避免 T x T 超订导致的吞吐塌陷
    c = a @ b
\end{lstlisting}

收束本章的图景：\FTN{} 负责``值不值得信任''，C 负责``能不能链接''，\CPP{} 负责``抽象得爽不爽''（下一章的主题），\PY{} 负责``用的人多不多''。接口之所以是 C，不是因为它最擅长数值，而是因为它是这四方唯一都承认的汇合点；LAPACK 谱系能活三十多年，靠的正是这个汇合点之上``实现可替换、kernel 藏性能''的结构。下一章的 Eigen 走的是光谱另一端：没有 ABI、没有可替换，一切在编译期实例化——但那是另一种自由。

% ====================================================================
% ====================================================================

% ====================================================================
% ====================================================================
% frag_ch05.tex —— 第 5 章：Eigen
% ====================================================================
\chapter{Eigen：\CPP{} 侧的线性代数抽象}

第~4 章的 BLAS 讲的是``运行期可替换''：符号是 C 的，实现是厂商的，谁快换谁。
Eigen 站在光谱的另一端：没有 .so 让你换，没有符号让你链接，
一切在调用者的编译单元里现场实例化。
这不``退步''，而是另一种工程取舍——它换来的是 BLAS 那三十几个例程永远表达不出来的东西：
编译期已知尺寸的三阶小矩阵、跨多个算子融合的单遍求值、
以及零拷贝包住外部缓冲的视图。
本章按四个机制展开：表达式模板与求值时机、对齐、\Tp{Map} \Bk{}，
以及解算器谱系与它和 BLAS 的混合点；最后回到全书主线，说明一个没有 ABI 的库如何与世界握手。

\section{为什么 BLAS 之上还需要一个 \CPP{} 层}

先算一笔最具体的账：姿态解算、图形学、结构力学里最常见的乘法是 $3 \times 3$ 或 $4 \times 4$。
走 \Fn{cblas\_dgemm} \Bk{} 这条路，你要摆正 14 个实参、过一次 layout 与参数校验、
进实现后再走一遍``小规模退化为朴素三重循环''的分支；
而走 Eigen 的 \Tp{Matrix3d}，尺寸在类型里，
编译器把循环整个展开成寄存器上的乘加，栈上甚至不需要数组落地。
函数调用边界在这个尺度上就是主要成本，
任何``把小矩阵攒批再调 gemm''的建议都意味着改写调用方结构。
第二笔账是组合表达：\Tp{R.transpose() * R + mu * I} \Bk{} 这类式子，
BLAS 词汇表里没有对应例程，C 手写则至少两次遍历、一份临时；
Eigen 用表达式模板把它融成一趟写。
第三笔是零拷贝：几何管线、信号帧缓冲的内存往往由 C API 或 \PY{} 分配，
Eigen 的 \Tp{Map} \Bk{} 可以直接把那块地址\emph{当作}矩阵使用而不拥有它（见 \S\,5.4），
这在任何``以缓冲为中心''的接口族旁边都是刚需。
把这些托底的同一件事是：尺寸与标量类型进了类型系统，
越界检查、动态分派与堆分配才可能被编译器整体消掉。

这类 \CPP{} 库不止 Eigen 一家，抽象路数各有不同，表~\ref{tab:ch05-cpplibs} 摆四家对照。

\begin{longtable}{@{}>{\raggedright\arraybackslash}p{0.13\textwidth}
  >{\raggedright\arraybackslash}p{0.17\textwidth}
  >{\raggedright\arraybackslash}p{0.14\textwidth}
  >{\raggedright\arraybackslash}p{0.13\textwidth}
  >{\raggedright\arraybackslash}p{0.11\textwidth}
  >{\raggedright\arraybackslash}p{0.13\textwidth}@{}}
\caption{四个 \CPP{} 线性代数库：抽象方式、布局、GPU、许可证与到 \PY{} 的距离}\label{tab:ch05-cpplibs}\\
\toprule
库 & 抽象方式 & 内存布局 & GPU 现实 & 许可证基调 & 与 \PY{} 的距离 \\
\midrule
\endfirsthead
\toprule
库 & 抽象方式 & 内存布局 & GPU 现实 & 许可证基调 & 与 \PY{} 的距离 \\
\midrule
\endhead
\Lib{Eigen} & 表达式模板；尺寸与标量类型都在类型参数里，固定尺寸在栈上展开 &
  列主序缺省，逐矩阵以模板参数切换行主序 &
  Tensor 模块有 CUDA 后端，算子覆盖与工具链支持随版本变化，需逐版核对 &
  MPL2 系宽许可，以仓库 LICENSE 为准 &
  最近：pybind11 的 eigen 插件把 ndarray 缓冲零拷贝映射成 \Tp{Map} \Bk{}；本身也常被用作别人后端的后端 \\
\Lib{Armadillo} & 运算符重载配模板表达式，接口刻意贴近 MATLAB &
  列主序固定 &
  较新版本提供经 cuBLAS 的实验路径，覆盖面有限 &
  Apache-2.0 基调，以仓库 LICENSE 为准 &
  中间：官方不做绑定，常经依赖它的项目（如 mlpack）间接到达 \\
\Lib{xtensor} & xexpression 惰性表达式，广播语义是卖点 &
  行主序缺省，layout 可作模板参数替换 &
  本体无 GPU 路径，卸载靠自建 &
  BSD 系，以仓库 LICENSE 为准 &
  近：xtensor-python 专为 pybind11 设计零拷贝映射 \\
\Lib{Blaze} & 表达式模板加自研 SIMD kernel，偏 HPC 口味 &
  行、列主序均可，编译期选定 &
  CPU 向量化为主；GPU 经 HIP 与 SYCL 后端，成熟度以版本文档为准 &
  BSD 系，以仓库 LICENSE 为准 &
  远：无官方 \PY{} 绑定 \\
\bottomrule
\end{longtable}

\section{表达式模板：求值时机、\Tp{.eval()} 与别名}

机制一句话：运算符不立即算，先\emph{构造一个描述算式的模板对象}。
\Tp{A + B + D} 返回的是嵌套类型 \Tp{add<A, add<B, D>>} \Bk{}，
一个没跑过任何浮点运算的语法树；
真正的循环只在赋值给矩阵时生成，一趟把三个源读进、把结果写出，中间矩阵不存在。
这就是``融合''的全部秘密：代价换在编译期（类型更深、错误信息更长、目标文件膨胀），
收益在内存流量。
逐元素世界是惰性的，乘法是例外——乘法的求值策略是及早的（先算进临时再合并），
这个不对称正是别名安全的来源，也解释了 \Tp{.noalias()} \Bk{} 的存在：
当你\emph{保证}左侧与乘积的输入不重叠（比如独立的 \Tp{C = A * B}），
挂上 \Tp{noalias()} 允许实现跳过临时直接写进 C，省一次 $O(n^2)$ 的分配与写出；
反之，把 noalias 用在重叠缓冲上是\emph{未定义行为}，而且往往错得很安静。

强制求值用 \Tp{.eval()} \Bk{}：把一段惰性子式冻结成独立临时。
两个典型动机：调试时想看``这个表达式此刻的值''，以及切断长表达式对局部变量的引用。
后者藏着表达式模板最经典的坑：表达式对象\emph{按引用}持有操作数，
用 \Tp{auto} \Bk{} 接住它（推导出的正是那个表达式类型）再让操作数离开作用域，
之后求值读到的就是悬垂内存。清单~\ref{lst:ch05-eval} 把三种时机摆在一起。

\begin{lstlisting}[style=cppstyle,caption={惰性、求值与生命周期：\Tp{auto} 接表达式是个真问题},label={lst:ch05-eval}]
void lifetime_trap(const MatrixXd& a) {
  MatrixXd b(3, 3);
  b.setOnes();
  auto e = a + b;         // e 的类型是 add：只攥着 a、b 的引用
  b.setZero();            // 改 b，e "跟着变"——它引用的就是 b
  // 若把 e 带出 b 的作用域再求值，读到的便是悬垂内存
}

MatrixXd fixed_way(const MatrixXd& a, const MatrixXd& b) {
  auto e = (a + b).eval();  // 冻结成真的 MatrixXd
  return e;                 // 此后 a、b 怎样都无关
}

void alias_zone(MatrixXd& c, const MatrixXd& a,
                const MatrixXd& b) {
  c = a * c;                // 安全：乘积先落临时再写 c
  c.noalias() = a * b;      // 前提：c 与 a、b 都无重叠
}
\end{lstlisting}

清单~\ref{lst:ch05-rigid} 是两段的合成例：刚体变换的复合与批量应用、平面的最小二乘拟合，
全程没有手写临时，尺寸检查全部前移到编译期。

\begin{lstlisting}[style=cppstyle,caption={刚体变换与平面最小二乘：表达式直接写代数式},label={lst:ch05-rigid}]
#include <Eigen/Dense>
#include <vector>
using namespace Eigen;

// 齐次变换：绕单位轴 yaw 转角，再平移 t
Matrix4d make_transform(const Vector3d& axis, double yaw,
                        const Vector3d& t) {
  Matrix3d r = AngleAxisd(yaw, axis.normalized())
                   .toRotationMatrix();
  Matrix4d m = Matrix4d::Identity();
  m.block<3,3>(0,0) = r;            // 块视图：写穿，不拷贝
  m.block<3,1>(0,3) = t;
  return m;                         // 固定尺寸，RVO 后零开销
}

void apply_batch(const Matrix4d& m, std::vector<Vector3d>& pts) {
  const Matrix3d r = m.block<3,3>(0,0);
  const Vector3d t = m.block<3,1>(0,3);
  for (auto& p : pts)
    p = r * p + t;  // 乘积先落临时：p 同时是输入与输出也安全
}

// 拟合 z = c0 + c1*x + c2*y：超定方程 Ax=b 的最小二乘
Vector3d fit_plane(const std::vector<Vector3d>& pts) {
  auto n = (Index)pts.size();
  MatrixXd a(n, 3);                 // 设计矩阵 [1, x, y]
  VectorXd b(n);                    // 观测 z
  for (Index i = 0; i < n; ++i) {
    a.row(i) << 1.0, pts[i].x(), pts[i].y();
    b(i) = pts[i].z();
  }
  // 列主元 QR：稳、可暴露秩；代价约为 LU 的两倍
  return a.colPivHouseholderQr().solve(b);
}
\end{lstlisting}

\section{对齐：无条件向量化把代价压在哪里}

Eigen 对固定尺寸向量化的类型（\Tp{Vector4d}、\Tp{Matrix3d} 这类大小恰为向量寄存器整数倍的）
默认使用\emph{对齐}的 SIMD 装载指令：
SSE 路径要求 16 字节、AVX 路径 32 字节、AVX-512 路径 64 字节，
上界由 \Tp{EIGEN\_MAX\_ALIGN\_BYTES} \Bk{} 这个编译期宏控制
（现代版本默认 32；设 16 可退回 SSE 时代的布局，设 0 则全线禁用对齐优化，用性能换混链安全）。
栈对象与全局对象由编译器保证对齐，问题集中在堆上：
Eigen 3.x 系列的契约按 C++98 时代的现实订立，
那里的 \Tp{operator new} \Bk{} 只承诺 8 或 16 字节，
于是结构体里只要嵌一个固定尺寸向量化成员，``new 出来''就可能踩中静态断言——
这是 Eigen 最有名的编译错误，而它防的其实是运行期在 \Tp{movaps} \Bk{} 一类指令上的随机段错误。
出路按年代分三段：
老代码用 \Tp{EIGEN\_MAKE\_ALIGNED\_OPERATOR\_NEW} \Bk{} 给类补一个对齐版 \Tp{operator new} \Bk{}；
动态尺寸类型（\Tp{Matrix<double,Dynamic,1>}）不做静态对齐，用运行期检查换来免打扰；
而 C++17 起语言标准已要求 over-aligned new \Bk{}，
新工具链下多数场景此问题自然消失——
但在自定义分配器、placement new、或跨 DLL 传对象的场合，它仍会回来。

\begin{warning}[title={对齐崩溃的两种炸法与统一修法}]
\begin{itemize}
\item \textbf{编译期}炸：错误是 Eigen 的``全大写加下划线''静态断言，读法见本章末的 note 框；
  常见根因是含固定尺寸向量化成员的结构体被 \Tp{new} \Bk{} 分配、或被塞进不保证对齐的容器。
\item \textbf{运行期}炸（SIGSEGV，回溯里有 \Tp{movaps} 或 \Tp{vmovaps}）：
  几乎总是把 malloc 来的、或栈上未对齐的裸指针交给\emph{默认按对齐处理}的 \Tp{Map} \Bk{}。
  外部缓冲不要假设它 32 字节对齐——C 分配器、图像行的 padding、旧 numpy 小数组都不保证。
\item 修法只有一个方向：显式声明非对齐，\Tp{Map<MatrixXf, Unaligned>} \Bk{}；
  全局层面可把 \Tp{EIGEN\_MAX\_ALIGN\_BYTES} \Bk{} 设为 0 或 16 与同链对象统一，
  但这会削掉向量化上限，等于把性能决策变成了链接期政治。
\end{itemize}
\end{warning}

\section{\Tp{Map}：与 C 缓冲和 \PY{} 握手的零拷贝接口}

\S\,5.1 说的第三笔账在这里兑现。
\Tp{Map} \Bk{} 是一个\emph{不拥有}内存的矩阵视图：
构造时只记下指针、尺寸、跨距，之后一切读取与写回都直接落在那块缓冲上，
没有一次分配、没有一次拷贝。
它有三个模板自由度，恰好对应与外部世界握手的三个维度。
其一，布局参数：外部 C API 给行主序缓冲，就写
\Tp{Map<Matrix<float,Dynamic,Dynamic,RowMajor>>} \Bk{}，
让 Eigen 迁就缓冲而不是反过来搬数据。
其二，对齐 Options：外部缓冲默认给 \Tp{Unaligned} \Bk{}，这是与 malloc 共存的安全阀（\S\,5.3）。
其三，\Tp{Stride} \Bk{}：真实缓冲里``下一行''往往不是``上一行加宽度''——
图像有行填充、多视图共享同一大分配、交织格式里列步长大于 1——
把内、外跨距交给 \Tp{Map} \Bk{}，子块与隔行采样就都不需要拷贝。

\begin{lstlisting}[style=cppstyle,caption={用 \Tp{Map} 接住 C 侧行主序缓冲、按跨距解释、原地写回},label={lst:ch05-map}]
#include <Eigen/Dense>

// C 侧给出 n 个行主序 xyz 点，存于 float* buf；
// 相邻两点的起始距离是 pitch，可能大于 3（行填充）。
void shade_points(float* buf, long n, long pitch,
                  const float* m44 /* 行主序 4x4 */) {
  // 模板三要素一次看全：布局 RowMajor、Unaligned、
  // 外跨距动态、内跨距编译期为 1
  using RM = Eigen::Map<Eigen::Matrix<float,
      Eigen::Dynamic, Eigen::Dynamic, Eigen::RowMajor>,
      Eigen::Unaligned,
      Eigen::Stride<Eigen::Dynamic, 1>>;
  RM pts(buf, n, 3, Eigen::Stride<Eigen::Dynamic, 1>(pitch));

  Eigen::Map<const Eigen::Matrix<float,4,4,Eigen::RowMajor>,
             Eigen::Unaligned> m(m44);
  const Eigen::Matrix3f rt = m.block<3,3>(0,0).transpose();
  const Eigen::RowVector3f tv = m.block<3,1>(0,3).transpose();

  for (long i = 0; i < n; ++i)
    pts.row(i) = pts.row(i) * rt + tv;   // 直接写穿 buf
}
\end{lstlisting}

清单~\ref{lst:ch05-map} 的左段值得多看一眼：\Tp{tv} \Bk{} 必须先转置成与行向量匹配的形状再相加，
因为 Eigen 在类型层拒绝任何维度戏法——这正是``尺寸进类型''的售后体验：
写的时候多打两个字符，换来编译期兜底。
与 \PY{} 的接缝同理且更省事：
pybind11 的 eigen 插件在函数签名出现 \Tp{Eigen::Map<...>} \Bk{} 时
直接把 ndarray 的数据指针与步长填进视图，
numpy 行主序对 Eigen 列主序的错位由布局模板参数解决，而不是由一次全量拷贝解决；
\Tp{ctypes} \Bk{} 递进来的裸缓冲则回到清单同款的手工 \Tp{Map} \Bk{}。
一条性能旁注：Eigen 的内核传统上按列主序打磨，行主序操作数在部分乘法里要走转置适配路径；
布局是 ``尊重外部现实''的工具，不是免费的——大尺度上值得把两种布局各测一遍。
换句话说，第~4 章``C 是汇合点''的论断在这里以镜像形式成立：
Eigen 没有 ABI，所以它靠\emph{接受 C 的布局}与世界握手。

\section{解算器谱系、头文件-only 的账、与 BLAS 的混合}

Eigen 的 \Tp{Solvers} \Bk{} 模块把第~4 章 LAPACK 的选择判据按\emph{结构}重新组织了一遍：
每个解算器是一个类，持有分解结果，可复用、可反复 solve。

\begin{itemize}
\item \Fn{LLT}（Cholesky，仅对称正定：约为 LU 一半的代价且无选主元；矩阵并非真正定时它不报错、直接给垃圾）
  与 \Fn{LDLT}（对称但不定）。
\item \Fn{PartialPivLU}：方阵、够非奇异、不需要秩信息时的快路径；
  \Fn{FullPivLU}：可暴露秩与核，代价高。
\item \Fn{ColPivHouseholderQR}：最小二乘与秩揭示的主力（\S\,5.2 清单用的就是它）；
  rank 亏缺的最小范数解交给 \Fn{CompleteOrthogonalDecomposition}。
\item \Fn{BDCSVD} 与 \Fn{JacobiSVD}：伪逆、条件数、一切``我要知道这个矩阵怎么样''的问题；
  最贵，信息也最全。
\item 稀疏一族：\Fn{SimplicialCholesky} \Bk{}（稀疏对称正定的直接法，
  符号分解与数值分解分离，同一稀疏模式反复求解可复用符号分解——有限元里的主场）、
  \Fn{SparseLU}、\Fn{SparseQR}。
\item 迭代一族：\Fn{ConjugateGradient} \Bk{}（要求 SPD，只需 matvec 可用；
  收敛速度随条件数恶化，配 \Fn{IncompleteCholesky} 预条件）、\Fn{BiCGSTAB}（非对称）。
\end{itemize}

\begin{guideline}[title={解算器选择判据}]
\begin{itemize}
\item 结构：对称正定走 \Fn{LLT} \Bk{}（或稀疏版 \Fn{SimplicialCholesky}），最便宜；
  对称不定用 \Fn{LDLT}；一般方阵用带部分主元的 LU。
\item 规模与形态：稠密且装得下，直接法；
  大到装不下、或 matrix-free 只能黑盒相乘，才谈迭代法与预条件。
\item 要的信息：只要解，选最便宜能满足结构假设的那个；
  要秩、核、条件数、最小范数，付出 SVD 或列主元 QR 的价钱。
\item 复用：同一矩阵多个右端，持有分解对象反复 \Fn{solve}；
  同一稀疏模式多个数值，复用符号分解。
\end{itemize}
\end{guideline}

头文件-only 是 Eigen 与第~4 章的分水岭：没有库可链，就没有 ABI、没有可发现性
（构建系统里不需要 \Tp{FindBLAS} \Bk{} 那样的探测，只需要 include 路径）、
也没有实现替换的余地——``换实现''意味着重新编译。
三笔账都记在调用方：编译时间随表达式模板的实例化深度增长；
优化等级与 \Tp{-march} \Bk{} 一类目标开关由调用方的编译命令决定，库作者无法预编译成调优产物；
发布方也无法通过换 .so 悄悄修 bug，版本升级以重编译为代价换来逐位可复现的构建。
与 BLAS 的混合点是显式的：定义 \Tp{EIGEN\_USE\_BLAS} \Bk{} 后大尺度动态乘积转交 \Tp{cblas\_*} \Bk{}，
\Tp{EIGEN\_USE\_LAPACKE} \Bk{} 再把部分分解换成 LAPACK 例程。
划不划算没有先验答案：
厂商微 kernel 在极大矩阵上靠专用分块与寄存器分配通常领先，
Eigen 自身的打包向量化在中小尺度常有更好的综合表现，
正确姿势是在自己的规模、自己的实现上各测一遍，
再决定要不要为了这一次领先放弃 header-only 的可移植性。
\Tp{Tensor} \Bk{} 模块与它的 CUDA 后端则要保持现实预期：
它面向多维数组与自动微分风格的表达式，
GPU 路径的算子覆盖、对新 CUDA 工具链与新架构的支持历来比专用 GPGPU 库慢半拍，
把它评估为``实验后端、逐版核对支持矩阵''，比当作现成的设备加速器稳妥。

\begin{note}[title={模板报错的三步读法}]
Eigen 的编译错误动辄上百行，因为整条表达式链是一次实例化，
一个维度错误会把 \Tp{add<A, product<B, C>>} \Bk{} 这类类型名一路传染到栈顶。
可操作的读法有三步。
第一，只看\emph{最内层}的实例化与第一处静态断言：
Eigen 的断言文本本身就是说明书，是全大写加下划线的短句
（比如提示你把标量当矩阵用、或尺寸对不上的那种），它比后面的类型海有价值得多。
第二，把出错表达式拆成单算子、逐个赋给具名变量，第一次失败的算子即元凶——
维度戏法在类型层就报错，正是这件事的售后成本与收益。
第三，工具侧：\Tp{g++} 的 \Tp{-fmax-errors=1} \Bk{} 能掐断错误洪水，
让最内层的第一条诊断浮到眼前；
工程上则把重型实例化集中到少数 .cpp 里，别让每个翻译单元都重新推导一遍全项目。
\end{note}

回扣全书主线：Eigen 证明了``\CPP{} 层可以有多薄''——薄到没有运行时存在；
同时又证明了``\CPP{} 抽象可以有多厚''——厚到把尺寸、布局、对齐全部搬进类型系统。
它与世界握手的两只手，一只是 \Tp{Map} \Bk{}（接受 C 的内存现实），
一只是编译期实例化（把性能责任转给调用方的编译器）。
后面进入信号域的章节里，你会再见到这对组合的另一端：接口层是 C 的 FFT 库。

% ====================================================================
% ====================================================================

% ====================================================================
% ====================================================================
\chapter{GSL、Boost.Math 与特殊函数、随机数}

前面讲 BLAS 与 LAPACK 的那一章处理的是``有结构的大数组''：一次调用吃进成千上万个元素，吞吐由内存带宽决定，接口开销可以忽略。本章换到另一类完全不同的数值代码：一次调用吃进一个数，吐出一个数。特殊函数、数值积分、求根、随机数都属于这一类。因为单次调用的工作量极小，接口的每一层开销（一次间接调用、一次装箱、一次状态拷贝）都会被放大；而它们的正确性也不来自线性代数定理，而来自逼近论加上浮点边界上的手工处理。

本章要回答两个问题。第一，为什么同一件事在 \CL{} 里长成\emph{句柄加错误码}的形状（\Lib{GSL}），在 \CPP{} 里长成\emph{模板加 policy}的形状（\Lib{Boost.Math}）：这不是口味差异，而是两者的开销与可复现性预算落在不同的地方。第二，为什么``自己写一个贝塞尔函数''与``自己挑一个随机数发生器''是数值工作里最常见的两个事故源头：前者的错误藏在边界条件与精度预算里，后者的错误藏在\emph{全局状态}里。

一条主线贯穿始终：\CL{} 提供可稳定链接的汇合点，\CPP{} 与 \FTN{} 提供实现层，\PY{} 提供门面。\PY{} 侧的 \Tp{scipy.special} 与 \Tp{numpy.random} 之所以能与 C 侧给出几乎一致的结果，正因为两边读的是同一族算法；第 6.5 节末尾的对拍清单就是在演示这件事。

\section{\Lib{GSL}：一门 \CL{} 接口的语法课}

\Lib{GSL}（GNU Scientific Library）是这一域里最典型的纯 \CL{} 库。它的接口风格不是历史包袱，而是一组有意识的选择，每一条都能对应到嵌入宿主时的具体收益。

\textbf{前缀即命名空间。} \CL{} 没有命名空间，\Tp{gsl\_} 前缀加上严格的``名词在前、动词在后''构词法（\Tp{gsl\_vector\Bk\_alloc}、\Tp{gsl\_vector\Bk\_free}、\Tp{gsl\_vector\Bk\_get}、\Tp{gsl\_vector\Bk\_set}、\Tp{gsl\_vector\Bk\_view\Bk\_array}）承担了命名空间的作用。可预测的构词法让调用者只读一份头文件就能猜出其余 API 的形状；对工具链而言，统一前缀让 \Fn{nm} 与动态加载器能一次筛出全部符号，也让自动生成的绑定（ctypes、cffi、\Tp{extern \textquotedbl C\textquotedbl} 声明块）可以按前缀批量产出。

\textbf{句柄是显式的、布局可见的。} \Tp{gsl\_vector} 与 \Tp{gsl\_matrix} 的结构写在公共头文件里，而不是不透明指针：尺寸、步长、数据指针、所属的 \Tp{gsl\_block}。这带来一个关键能力——\emph{视图与外部缓冲}：
\begin{itemize}
\item \Fn{gsl\_vector\Bk\_view\Bk\_array} 与 \Fn{gsl\_matrix\Bk\_view\Bk\_array} 把一段别人分配的连续内存包成向量或矩阵，\emph{不复制、不接管所有权}；带步长的 \Fn{gsl\_matrix\Bk\_view\Bk\_stride\Bk\_array} 能把列主序的 \FTN{} 数组、或 BLAS 输出缓冲里的一段直接当矩阵用。
\item \Fn{gsl\_vector\Bk\_subvector}、\Fn{gsl\_matrix\Bk\_submatrix} 与 \Fn{gsl\_matrix\_row} 产出的视图共享同一个 \Tp{block}，只做偏移与步长运算。分块算法因此不必把数据搬来搬去。
\item 反面是生命周期全归调用方：\Fn{gsl\_vector\Bk\_free} 之后，指向它的所有视图都成了悬垂指针，而且\emph{不会有任何提示}。\CPP{} 侧的 RAII 包装（本册讲 Eigen 的那一章会再遇到）解决的就是这一条。
\item ``复制''与``借用''是两个不同的函数：\Fn{gsl\_vector\Bk\_alloc\Bk\_from\Bk\_array} 复制一份，\Fn{gsl\_vector\Bk\_view\Bk\_array} 不复制。混用要么得到一个仍指向调用方缓冲的句柄，要么得到一次意料之外的 $n$ 元素拷贝。
\end{itemize}

视图换来零拷贝，代价是\textbf{步长不再是 1}。这直接影响与 \FTN{} 谱系接口的互操作：\Fn{gsl\_blas\Bk\_dgemv} 一类的 BLAS 封装要求矩阵步长能退化成一个合法的 LDA，否则返回 \Tp{GSL\_EBADSTRIDE}。把 GSL 当作 BLAS 的 \CL{} 门面时，这是最常撞的一次报错，也是``视图是廉价的、但不是免费的''这条规律的体现。同理，需要 SIMD 对齐的热缓冲不能指望公共句柄替你保证：\Tp{gsl\_block\Bk\_alloc} 给到的对齐随分配器与平台不同，向量化路径要按本书讲 SIMD 那一章的口径自己核对。

\textbf{错误码而不是异常。} GSL 的函数返回一个 \Tp{int} 错误码（\Tp{GSL\_SUCCESS}、\Tp{GSL\_EINVAL}、\Tp{GSL\_ERANGE}、\Tp{GSL\_EDOM}、\Tp{GSL\_MAXITER}、\Tp{GSL\_EBADTOL} 等），数值函数还有 ``\Tp{\_e} 后缀''的第二形态：接受一个 \Tp{gsl\_sf\Bk\_result} 结构指针，里面除了 \Tp{val} 还有 \Tp{err}。这一点值得停下来看：\Tp{err} 把``算出来的值可信到几位''写进了返回类型，正是本章精度预算话题在接口层的落点。\CPP{} 侧没有这个位置（返回值只有一个），误差信息只能靠 policy 间接控制。

问题出在 GSL 的另一条约定：内部宏 \Tp{GSL\_ERROR} 触发时会调用一个\emph{进程级}的错误处理器函数指针，而默认处理器打印诊断后调用 \Fn{abort}。

\begin{warning}[title={默认错误处理器会把宿主打死}]
一个嵌在宿主进程里的 GSL 调用，最坏后果不是算错，而是\emph{把宿主打死}。三条独立的机制叠在一起：
\begin{enumerate}
\item \textbf{默认处理器会 abort。} 参数越界（\Tp{GSL\_EDOM}）、迭代上限（\Tp{GSL\_MAXITER}）、容差不合法（\Tp{GSL\_EBADTOL}）这类``数据依赖的正常事件''在生产代码里迟早发生；默认行为是直接结束进程。宿主是解释器、数据库服务或某个图形界面主循环时，用户看到的是一次崩溃而不是一个错误码。
\item \textbf{处理器指针是进程级全局量。} \Fn{gsl\_set\Bk\_error\Bk\_handler\Bk\_off} 与 \Fn{gsl\_set\Bk\_error\Bk\_handler} 改的是同一个全局函数指针。线程 A 关掉它，线程 B 正在把处理器换成别的实现，两者互相覆盖；后果是同一个程序在不同线程调度下行为不同，属于最难复现的一类缺陷。
\item \textbf{处理器里不能抛 \CPP{} 异常。} 在处理器里 \Tp{throw} 会让异常穿过 GSL 自己的 C 栈帧；GSL 的内部宏依赖顺序执行与局部清理，跨 C 帧展开不保证资源被释放。要在 C 与 \CPP{} 边界上把错误变成对象，正确的位置是\emph{你自己的包装层}，不是 GSL 的处理器。
\end{enumerate}
可行的纪律：\Fn{gsl\_set\Bk\_error\Bk\_handler\Bk\_off} 只在单线程的初始化阶段调用一次（谁拥有 GSL 谁负责），它返回旧处理器指针，测试里可以先换回默认处理器、用完再恢复；此后逐函数检查返回值，用 \Fn{gsl\_strerror} 把码转成宿主的错误类型。许可证方面 GSL 属于 GPL 谱系，能否静态进产品要按你所用版本的 LICENSE 判断，通常更稳妥的做法是把它放在独立进程或可单独分发的组件里。
\end{warning}

\begin{lstlisting}[style=cstyle,caption={\Lib{GSL} 求柱贝塞尔函数：关掉默认 abort，连误差上界一起取},label={lst:ch06-gsl}]
/* 编译：cc -O2 -std=c11 本文件 -lgsl -lgslcblas -lm */
#include <stdio.h>
#include <gsl/gsl_errno.h>
#include <gsl/gsl_sf_bessel.h>
#include <gsl/gsl_sf_gamma.h>
#include <gsl/gsl_sf_erf.h>

static int check(int st, const char *what)
{
  if (st != GSL_SUCCESS) {
    fprintf(stderr, "%s 失败: %s\n", what, gsl_strerror(st));
    return 0;                    /* 只报告，不终止进程 */
  }
  return 1;
}

int main(void)
{
  /* 默认处理器会 abort，嵌入宿主前必须先关掉（见上方陷阱框） */
  gsl_set_error_handler_off();

  gsl_sf_result r;               /* val + err：误差上界随结果一起返回 */
  for (int n = 0; n <= 4; ++n) {
    int st = gsl_sf_bessel_Jn_e(n, 12.5, &r);
    if (!check(st, "bessel_Jn_e")) return 1;
    printf("J%d(12.5) = %+.15e   err < %.1e\n", n, r.val, r.err);
  }

  /* 大自变量的 Gamma 用对数形式：Gamma(171.7) 已超出 double 上限 */
  if (check(gsl_sf_lngamma_e(400.0, &r), "lngamma_e"))
    printf("ln Gamma(400) = %+.15e   err < %.1e\n", r.val, r.err);

  /* Gamma 在非正整数处是极点：越界变成错误码，而不是 abort */
  if (gsl_sf_lngamma_e(-2.0, &r) != GSL_SUCCESS)
    puts("负整数极点返回 GSL_EDOM，进程仍然活着");

  if (check(gsl_sf_erf_e(0.75, &r), "erf_e"))
    printf("erf(0.75) = %+.15e   err < %.1e\n", r.val, r.err);
  return 0;
}
\end{lstlisting}

清单里那套 \Tp{int} 错误码加 \Tp{gsl\_function} 回调的约定不是 GSL 的孤例：数值库的 \CL{} 接口几乎都收敛到同一形状，因为它是唯一能被 \CPP{}、\FTN{}、\PY{} 三方都廉价接住的形状。本书讲微分方程与优化的那一章里，\Lib{Sundials} 用的是同一套约定（回调返回状态码、求解器是长生命周期句柄），这里的纪律可以直接搬过去。

\section{\Lib{Boost.Math}：把开销从运行时挪到编译期}

\Lib{Boost.Math} 走另一条路：基本是头文件-only 的 \CPP{} 模板库（Boost 软件许可证，商用友好；个别组件依赖 Boost.Multiprecision 等外围库），调用形状像下面这样。

\begin{lstlisting}[style=cppstyle,caption={\Lib{Boost.Math} 的 policy：精度、性能、错误处理都是编译期参数}]
#include <boost/math/policies/policy.hpp>
#include <boost/math/special_functions/bessel.hpp>
#include <boost/math/special_functions/erf.hpp>

namespace bm  = boost::math;
namespace pol = boost::math::policies;

/* 默认：float 实参先在更宽的类型里算完再舍入，末位更准 */
double exact = bm::erf(0.5F);

/* 不提升：省掉一次类型往返，换来末位误差变大 */
using no_promote = pol::policy<
    pol::promote_float<false>,
    pol::promote_double<false>>;
double cheap = bm::erf(0.5F, no_promote());

/* 越界不抛异常：批量求值的热循环里没有异常边沿，
   编译器更愿意内联与向量化，代价是 NaN 自己查 */
using quiet = pol::policy<
    pol::error_handling<pol::ignore_error>>;
double bad = bm::cyl_bessel_j(0.5, -1.0, quiet());
\end{lstlisting}

关键机制是 \Tp{policies::policy<...>} 与查询它的元函数（\Tp{precision\Bk\_policy}、\Tp{promotion\Bk\_policy}、\Tp{error\Bk\_policy}、\Tp{evaluation\Bk\_policy}）：精度目标、内部用什么类型累加、越界时抛异常还是安静返回一个可预测值，三件事各自是一个类型参数，于是调用点开销为零（代码路径在编译期就定死），代价是每个（类型，policy）组合各自实例化一份，编译时间与目标码体积同时上升，而且\textbf{没有可共享的 ABI}——你无法把 Boost.Math 编译进一个 \Tp{.so} 再宣称它对其他 Boost 版本二进制兼容。

两处细节解释了这条路线的价值。其一是 \Tp{promote\_float}：\Tp{float} 的 \Tp{erf} 若直接用 24 位有效数字做极小极大逼近，末位误差会明显大于正确舍入；\Lib{Boost.Math} 的做法是内部在 \Tp{double} 或 \Tp{long double} 里算完再舍回，用一次类型往返买回几位余量，而 GSL 侧根本没有这个旋钮（接口只有 \Tp{double}，加宽由调用方做）。取向差别不是谁更准，而是\emph{精度预算写在哪一层}：模板参数，还是调用方代码。

其二是 \Tp{error\_handling<pol::ignore\_error>}：去掉异常边沿之后，函数体更容易被内联进循环、也更容易被向量化；批量求值特殊函数（本册后面做窗函数修正或谱权系数时会遇到）常常靠这一条就拿到可观收益。反面是错误处理责任转给了调用方：越界变成 \Tp{NaN} 顺着数据流往下走，很远之后才以``结果不对''的形式暴露。配套纪律是每段计算末尾做一次 \Fn{std::isnan}、\Fn{std::isfinite} 抽查，或者把抽查塞进调试构建里常开的断言。

\section{特殊函数为什么必须由底层实现}

现在回答本章最硬的问题：贝塞尔函数、Gamma 函数、误差函数这些``定义就在教科书里''的东西，为什么不能自己写？答案分三层。

\textbf{其一，自变量域要求分段策略。} 以 $J_n(x)$ 的幂级数为例：
\[
J_{n}(x)=\left(\frac{x}{2}\right)^{n}\sum_{k=0}^{\infty}
\frac{(-1)^{k}}{k!\,\Gamma(n+k+1)}\left(\frac{x^{2}}{4}\right)^{k},
\]
级数在整个实轴收敛，但 $x \gg n$ 时项要先涨到 $\sim e^{x}$ 量级再衰减，中间是剧烈的逐项抵消：在 \Tp{double} 里有效位数会被白白吃掉。所以实现必须按 $\nu$ 与 $x$ 的比值切域：小 $x$ 用级数、大 $x$ 用渐近展开或 Hankel 表示、转折点附近换另一套。GSL 的 \Fn{gsl\_sf\Bk\_bessel\Bk\_Jnu} 与 \Lib{Boost.Math} 的 \Fn{cyl\_bessel\_j} 内部都是这样的多段调度，公开接口只有一个函数名，复杂度全藏在里面。

\textbf{其二，逼近是拟合出来的，不是推出来的。} 主流实现用有理分式加多项式逼近：以极小化最大相对误差（Remez 一类的极小极大拟合）离线算出系数，再用 Clenshaw 或 Estrin 格式求值。Cephes 谱系（历史上 \FTN{} 的 PORT 数学库、\Lib{Boost.Math} 里 TR1 谱系的组件、\PY{} 侧 \Tp{scipy.special} 沿用的那一批）都遵循这个套路。系数是数值实验的产物：一段手写代码\emph{没有}这些系数，也没有它们被验证过的工作域。这里就有``精度预算''的概念：为了在 \Tp{double} 里给出末位一到两个 ulp 的结果，内部要多带几位余量；同一组逼近拿去做 \Tp{long double} 就不够了，做任意精度更是不够。这解释了 \Tp{gsl\_sf\Bk\_result} 为什么带 \Tp{err}，也解释了为什么 Boost 的精度旋钮必须是\emph{模板}参数。

\textbf{其三，边界要手写处理。} 三个真实的坑：
\begin{itemize}
\item \Tp{erfc} 与 \Tp{1 - erf}。$x=3$ 时 $\mathrm{erf}(x)$ 约为 $0.99998$，相减之后只剩约 4 位有效数字。所以库必须给 \Fn{gsl\_sf\Bk\_erfc}，还要给尾部专用的 \Fn{gsl\_sf\Bk\_erf\Bk\_Q}（标准正态上尾概率），以及缩放的 \Tp{erfcx}（乘上 $e^{x^{2}}$，避开下溢；\PY{} 的 \Tp{scipy.special} 与 \Lib{mpmath} 提供这一族）。\Lib{Boost.Math} 的离散分布补集 \Fn{complement} 是同一动机的 \CPP{} 版。
\item \Tp{tgamma} 与 \Tp{lgamma}。$x=171.7$ 附近 $\Gamma(x)$ 溢出，对数形式在 $730$ 上下才溢出；非正整数是极点，$x<0$ 且接近整数时用反射公式会把误差连同符号一起放大。库里有的是带符号、取整或倒数形式，让调用方不必自己推反射关系——自己推的那一次，通常就是把极点附近算成 $0$ 的那一次。
\item 递推方向。Bessel 函数对阶 $n$ 的三项递推在一个方向上数值不稳定（要用 Miller 反向递推再加归一化）；球贝塞尔函数在 $n$ 接近 $x$ 时尤其明显。写成 \Tp{for} 循环的朴素递推进会给出``前几项对、后面全飞''的结果，而这正是最难查的一类错。
\end{itemize}

三条合起来看：库外重写丢掉的不只是那几百行逼近代码，还有一整套``边界处理加测试基准''。公开记录里最贵的一例是电子表格软件的正态分布函数（\Tp{NORM.S.DIST} 一族），厂商的修正说明里提到的误差就在末几位量级；那是一段``看起来显然''的公式，配上极小极大逼近之后仍然出错。同一件事发生在你的代码里时通常没有基准可比，因此不会被发现——它会一路走进统计量、走进结论。

覆盖面也值得记成一张表，因为选型时真正的问题是``这个库里有没有我要的那个函数''。表~\ref{tab:ch06-coverage} 里``有''指存在可直接调用的接口，``无''指没有或要靠外部库。

\begin{longtable}{@{}>{\raggedright\arraybackslash}p{0.14\textwidth}
  >{\raggedright\arraybackslash}p{0.10\textwidth}
  >{\raggedright\arraybackslash}p{0.23\textwidth}
  >{\raggedright\arraybackslash}p{0.19\textwidth}
  >{\raggedright\arraybackslash}p{0.17\textwidth}@{}}
\caption{特殊函数相关库的覆盖面与语言层对照}\label{tab:ch06-coverage}\\
\toprule
库 & 语言层 & 覆盖 & 精度与误差控制 & 接入成本 \\
\midrule
\endhead
\Lib{GSL} & \CL{} &
  特殊函数一族（Bessel、Gamma、Airy、不完全 Gamma、椭圆积分、勒让德、Hermite、Laguerre、超几何）、积分与求根、随机数与分布、统计、插值、直方图、特征分解 &
  \Tp{double} 为主，\Tp{\_e} 系列随值返回误差上界 &
  共享库 \Tp{libgsl} 加 \Tp{libgslcblas}；GPL 谱系许可需按版本评估 \\
\Lib{Boost\Bk.Math} & \CPP{} 头文件 &
  覆盖面在 \CPP{} 生态里最广，含分布函数与反函数、多项式求根、数值积分与求根、常微分方法 &
  policy 控制提升类型、精度目标、是否抛异常 &
  零链接成本，代价是编译期与目标码体积，无稳定 ABI \\
\CPP{}17 标准 & \CPP{} &
  \Tp{tgamma}、\Tp{lgamma}、\Tp{beta}、\Tp{erf}、\Tp{expint}、\Tp{riemann\_zeta}、柱与球贝塞尔、修正贝塞尔、勒让德、Hermite、Laguerre &
  无误差返回，精度是实现自定义质量 &
  零依赖；不完全 Gamma、Airy、椭圆积分、分布函数都不在标准里 \\
Cephes 谱系 & \CL{} &
  约几十个小而快的 \Tp{double} 特殊函数，历史上是许多语言数学库的底子 &
  面向 \Tp{double}，没有系统的误差估计 &
  体积极小、许可宽松（以源码分发处声明为准），但没有整合的求解器层 \\
\Lib{mpmath} & \PY{} &
  任意精度特殊函数、积分、求根与常微分，函数族相当全 &
  十进制任意精度，直接调 \Tp{mp.dps} &
  纯 \PY{} 实现，比 \CL{} 路径慢一到两个数量级；适合当验证基准而非热路径 \\
\bottomrule
\end{longtable}

顺带一条常被误用的观察：\CPP{}17 的标准特殊函数（\Tp{std::cyl\_bessel} 与 \Tp{std::legendre} 一族）在各家标准库里并非同一份代码，实现来源与达到的精度都随标准库与版本不同，有的直接从 \Lib{Boost.Math} 的算法谱系改写而来；而\emph{标准只规定接口、不规定精度}，所以同一调用在不同实现上的末几位可以不同，也没有误差返回。把它们当``少一个依赖''的便利是合理的，当``可逐位复现的基准''是不合理的。

\section{数值积分与求根：自适应求积的接口形状}

自适应求积的标准做法是 Gauss-Kronrod。Gauss 的 $m$ 点公式与 Kronrod 的 $2m+1$ 点公式是\emph{嵌套}的：Kronrod 的节点包含全部 Gauss 节点，只多出 $m+1$ 个。于是同一次划分里两个和都能算出来，两者之差就是现成的误差估计，不必额外求值。$m=10$（即 GK15，一次划分取 21 个节点）是事实上的默认档。\FTN{} 的 \Lib{QUADPACK} 里的 \Tp{dqags} 与 GSL 的 \Fn{gsl\_integration\Bk\_qags} 用的都是这套 21 点规则：GSL 的实现是把 QUADPACK 的思路用 \CL{} 重写了一遍，这是本册反复看到的``\FTN{} 提供算法、\CL{} 提供接口''的又一个样本。

\textbf{接口的形状就是开销的形状。} GSL 的求积入口是
\[{\small
\Fn{gsl\_integration\_qag}(f,\ \texttt{params},\ a,\ b,\ \texttt{epsabs},\ \texttt{epsrel},\ \texttt{limit},\ w,\ \texttt{result},\ \texttt{abserr})
}\]
其中 $f$ 是 \Tp{double (*f)(double, void *)} 形式，挂在 \Tp{gsl\_function} 结构里；$w$ 是 \Fn{gsl\_integration\Bk\_workspace\Bk\_alloc} 分配的长生命周期工作区，存划分栈与中间结果，可复用。三条推论都能落到代价上：
\begin{itemize}
\item 工作区要显式分配与释放，因此``同一线程串行复用''合法，``跨线程共享一个工作区''不合法。想在积分这一层榨并行度，就得按线程开工作区，并把区间划分固定下来（否则合并顺序一变，结果就变）。
\item 回调是间接调用，每次划分至少 21 次。若回调跨语言（\PY{} 闭包），装箱成本按求值次数成倍放大。工程上的处理是把被积函数改写成一次吃整个节点数组的形式（向量化回调），代价是接口要能表达批量求值：\PY{} 侧的积分接口提供了这一类开关，动机就在这里。
\item 停止判据是 $|E|\le\max(\texttt{epsabs},\ \texttt{epsrel}\,|I|)$。把 \Tp{epsabs} 设为 $0$ 而被积函数真值接近零时判据永远不满足，最后以 \Tp{GSL\_MAXITER} 收场。正确做法是给一个绝对容差当保底，哪怕它远小于你需要的精度。
\end{itemize}

端点奇异是自适应求积最常见的失败原因：$1/\sqrt{x}$ 在 $0$ 处会让算法不停在同一端附近二分。GSL 的对策是把``哪一端奇异''写进接口——带区间类型的变体（\Fn{gsl\_integration\Bk\_qag\Bk\_int}）接受 \Tp{regular}、\Tp{singular\_lower}、\Tp{singular\_upper}、\Tp{singular\_both} 四选一，二分时强制包含奇异端点；代数与对数奇异有专门的 \Fn{gsl\_integration\Bk\_qaws}（权重 $(x-a)^{\alpha}(b-x)^{\beta}\ln^{\mu}(x)$）；无穷区间有 \Fn{qagi}、\Fn{qagiu}、\Fn{qagid}；带 $\sin$ 或 $\cos$ 因子的振荡积分有 \Fn{qawsin}、\Fn{qawcos}，它们走 Clenshaw-Curtis 型路线而不是继续二分。\AUD{} 与 \EMW{} 侧的谱积分以后会回到这一条。

\CPP{} 侧 \Lib{Boost.Math} 把同一族方法做成模板：\Tp{boost::math::quadrature\Bk::quadgk<T, Policy>} 是自适应 Gauss-Kronrod，\Tp{tanh\_sinh<T>} 与 \Tp{exp\_sinh<T>}、\Tp{sinh\_sinh<T>} 是双指数变换（端点奇异与无穷区间上常常收敛最快的一档），\Tp{fejer\_2} 处理振荡积分，\Tp{gauss\_kronrod<T, N>} 与 \Tp{gauss\Bk\_legendre<T, N>} 是固定阶的快速内核。因为类型与 policy 都在模板参数上，同一个 \Tp{quadgk} 能在 \Tp{double}、\Tp{long double} 或 Boost.Multiprecision 后端上工作，这是头文件路线相对 GSL 的真实优势。

求根这一侧，接口差异体现在``要不要区间''。有括区的方法（二分、试位、Illinois、TOMS 748）只要函数值、保证收敛；无括区的方法（Newton、Halley、Schroeder）更快但会跳出定义域，还得给导数。\Lib{Boost.Math} 的默认选择是 \Fn{toms748\_solve}（命名空间在 \Tp{boost::math::tools} 下）：有括区、只用函数值、阶约 $2.7$，机制与代价都落在中间最有利的位置；\Fn{bracket\Bk\_and\Bk\_solve\Bk\_root} 负责自动扩区间；重根与``极小平顶''这类病态情形有专门的 \Tp{rho\_solver}（ILLOIS 一族）。GSL 侧的对应物是 \Tp{gsl\_root\Bk\_fsolver} 族（\Tp{bisection}、\Tp{falsepos}、\Tp{ridger}）与带导数信息的变体，多元方程组走 \Tp{gsl\_multiroot\Bk\_fsolver}（Powell 混合、\Tp{dlevmar} 即 Levenberg-Marquardt），类型名以官方头文件为准。

值得单列的 GSL 设计：收敛判据被抽成 \Tp{gsl\_fsolver\Bk\_fmonitor}，你可以注册一个返回状态码的回调，让求解器在``调用方认为够好了''时停，而不是只看区间长度。这是一次解耦，代价是回调又进了一次间接调用的热路径。本书讲优化的那一章里，优化器的提前停止、以及 \Tp{-ffast-math} 为什么会破坏这类判据，都回到``monitor 是个回调''这个形状上。

\section{随机数：可复现性是一条接口约束}

随机数是本章里唯一一个``选错不会立刻报错''的部件：结果照样跑出来，只是再也复现不了，或者两条流之间有说不清的统计相关性。

\textbf{\Fn{rand()} 为什么不能用。} 四条独立的理由：
\begin{enumerate}
\item 状态是\emph{进程级全局量}。\Fn{srand()} 只有一个种子槽，进程里任何一处调用（第三方库、宿主、甚至一段无关的老代码）都会重置或推进它。``两条独立的流''在这个接口下无从表达。
\item 线程不安全。不加锁的实现里，两个线程会读到同一个值或整段跳过；后果是同样的种子、同样的线程数，结果随调度而变，既不可复现也难解释。
\item 序列质量与位宽。若干经典实现里低位的有效周期远小于全长，\Tp{RAND\_MAX} 可能只有 32767；于是 \Tp{rand() \% K} 带明显的模偏置，缩放到 $[0,1)$ 也留下格点结构。蒙特卡洛积分对高维格点结构尤其敏感（线性同余类发生器的超平面结构是有文献记载的失效模式），这不是风格洁癖。
\item 序列不是任何标准规定的。同一份代码在两种 libc 上给出完全不同的数字，跨平台对拍失去意义。
\end{enumerate}
如果只想洗牌或抽样而不关心复现，显式状态版本（\Tp{drand48}、\Fn{nrand48}、\Fn{jrand48}，把 \Tp{unsigned short xi[3]} 交给调用方保管）比 \Fn{rand()} 体面，但它是 48 位线性同余，别拿去做蒙特卡洛。\Fn{random} 与 \Fn{initstate} 一类的显式缓冲接口可复现性更好，但序列同样不在任何标准里。

\textbf{\Lib{GSL} 的取法：引擎是表加句柄。} 状态类型 \Tp{gsl\_rng\_state} 由 \Fn{gsl\_rng\Bk\_malloc} 分配、\Fn{gsl\_rng\_set} 播种、\Fn{gsl\_rng\_free} 释放；引擎本身是全局符号表里的条目，例如 \Tp{gsl\_rng\Bk\_mt19937} 与更宽的 \Tp{gsl\_rng\Bk\_mt19937\_64}，\Tp{gsl\_rng\_taus} 与 \Tp{gsl\_rng\_gfsr2} 这类位发生器。取值入口统一成三个：整数用 \Fn{gsl\_rng\_get}，$[0,1)$ 用 \Fn{gsl\_rng\Bk\_uniform}，有界整数用 \Fn{gsl\_rng\Bk\_uniform\Bk\_int}。抽样分布则被拆成近乎纯函数的一族，比如 \Fn{gsl\_ran\Bk\_poisson}、\Fn{gsl\_ran\Bk\_gaussian\Bk\_boxmuller} 与需要 state 以复用查表的 \Fn{gsl\_ran\Bk\_gaussian\Bk\_ziggurat}。两个值得注意的接口：
\begin{itemize}
\item \Fn{gsl\_rng\_fwrite} 与 \Fn{gsl\_rng\_fopen}：把内部状态整体落盘再读回，长模拟的断点续算靠它；代价是格式与字长、字节序绑定，跨平台续算要自己校验一次。
\item 较新版本提供的 \Tp{gsl\_rng\_qrng} 类型（配合 \Fn{gsl\_rng\Bk\_qrng\Bk\_fopen} 一类入口，具体名字以所用版本为准）：从一个文件读入外部真实随机源。它的用途不是``更随机''，而是\emph{不参与可复现的循环}——一旦把它当外部输入，模拟结果就不能指望逐位复现，这一点必须写进实验记录。
\end{itemize}
环境变量 \Tp{GSL\_RNG\_TYPE} 允许运行时换引擎，方便但不自证：日志里要同时记录引擎名与种子，否则一次``悄悄改了环境变量''就足以让两批数字对不上。

\textbf{\CPP{}11 \Tp{<random>} 的分工。} 引擎一侧有 \Tp{mt19937}、\Tp{mt19937\_64}、\Tp{minstd\_rand}、\Tp{ranlux24}、\Tp{ranlux48} 与打乱用的 \Tp{knuth\_b}；分布一侧有 \Tp{uniform\_real\Bk\_distribution} 与 \Tp{normal\_distribution}、\Tp{poisson\_distribution}。两族彻底解耦：一个引擎可以喂任意分布，引擎与分布各自可拷贝、可独立播种。可复现性上有一条关键的非对称：\textbf{标准精确规定了引擎的算法与常数}，所以 \Tp{mt19937\_64} 在任何合规实现上产出同样的序列；而\emph{分布的抽样算法是实现自定义的}。同一个引擎配同一个种子，在 libstdc++、libc++ 与 MSVC STL 上 \Tp{normal\_distribution} 的取值可以不同。要逐位对拍就得把分布也固定下来：自己实现抽样，或选一个把分布也规定清楚的库。

两条与代价直接相关的性质：\Tp{mt19937} 的状态是 624 个 32 位字（约 2.5 KB），\Tp{mt19937\_64} 是 312 个 64 位字，所以``每线程一个引擎''在内存上没问题，但检查点与状态落盘要按这个尺寸预算；\Tp{minstd\_rand} 只有一个字的状态，落盘最便宜、派生最干净，代价是线性同余的格点结构与短周期。还要注意 \Tp{std::random\_device}：它的语义是``非确定性随机数''，用于播种而不是用于结果复现，某些实现会在熵源不足时阻塞，把它塞进热循环是典型的性能事故。

\textbf{\Lib{Random123}：把状态换成索引。} 这一族计数式发生器的代表是 \Tp{philox4x32-10} 与 \Tp{threefry4x64-20}，以及同类的 ARS。它们不维护状态：一个输出等于把``流号与位置''喂进哈希，即 $\mathrm{hash}(\text{key},\text{counter})$，而 counter 里依次放子区间 id、块 id、线程 id。这带来三个别的方案给不了的能力：
\begin{itemize}
\item \emph{位置可寻址}：任何线程、任何 \Lang{GPU} 线程块、任何像素都能直接算出自己那一段的第一个随机数，不需要切分、不需要 \Tp{discard} 上亿步、也不需要共享计数器。
\item \emph{逐位可复现且跨执行域一致}：没有内部状态，就意味着主机侧、\Lang{GPU} 侧、不同编译器给出同一份数字。这也是为什么 \PY{} 侧 \Tp{numpy.random} 里能看见 \Tp{Philox}、\Tp{ThreeFry} 这类位发生器，而设备侧随机库（例如 \Tp{curand}）也把 Philox 列为可选发生器之一；具体支持范围随版本不同。
\item \emph{断点续算是免费的}：``状态''就是一个整数索引，检查点等于记一个下标，重启顺序无关。
\end{itemize}
代价同样具体：计数器式设计没有``自动跳过一段序列''的语义，流号必须自己设计成互不重叠（同一位置被两条流取到就出现重复样本）；单个输出是一次整数哈希，统计质量依赖轮数与参数选择，不像 MT19937 那样用长达 623 字的状态换来已被广泛验证的性质。蒙特卡洛积分与 \Lang{GPU} 上的逐元素抽样用它通常更合适；需要严格独立性的抽样检验场景两者都合格，但要按该库文档给出的分层或逆映射用法来。

\textbf{高精度这一侧的现实取舍。} \Tp{long double} 是\emph{平台上的一种类型}，不是一个精度承诺：x86-64 的 System V ABI 下它常是 80 位扩展精度（占 16 字节，只有约 64 位有效数字），AArch64 与部分平台上它可能是四倍精度，也可能干脆等于 \Tp{double}；Windows 上的 MSVC 就是后者。跨平台数值代码要写 \Tp{sizeof(long double)} 与 \Tp{numeric\_limits<long\Bk{} double>::digits} 的运行时断言，别把假设写进注释。想要确定的四倍精度，\CL{} 侧有 GCC/Clang 的 \Tp{\_\_float128}（IEEE binary128）加 \Tp{libquadmath} 的 \Fn{sqrtq}、\Fn{sinq}、\Fn{logq}：软件实现，比 \Tp{double} 慢一到两个数量级，但序列语义稳定，链接时多一个 \Tp{-lquadmath}（\FTN{} 侧的对应物是 \FTN{}2008 以来的四倍精度实数类型）。再往上就是软件后端：Boost.Multiprecision 的 \Tp{cpp\_bin\_float}、\Tp{cpp\_dec\_float<50>}（自带实现，头文件可用）与 \Tp{mpfr\_float}、\Tp{gmp\_float}（依赖外部 MPFR/GMP，是 C 库，链接与许可证要一起算）。工程上的常见分工是：主循环用 \Tp{double}，累加器用 \Tp{long double}（注意 x87 时代留下的每结果一次存储依赖，也注意在 MSVC 上它什么也没买到），验证基准用 \Lib{mpmath} 或 MPFR 算 30 位真值。

\begin{bestpractice}[title={可复现蒙特卡洛的四条纪律}]
\begin{enumerate}
\item \textbf{每个工作单元一条独立的流，种子由索引派生。} 用 \Tp{seed\_seq} 把``全局种子''与``单元号''混合，或用 Random123 把单元号放进计数器。不要图省事用``种子加线程号''这种算术拼接：相邻种子在 MT 一族里会产出高度相关的序列，这是该族发生器被文档明确警告过的用法。
\item \textbf{归约顺序固定。} 浮点加法不可结合，所以``每块各自求和、再按索引相加''是可复现的，而并行归约（\Tp{std::reduce} 的执行策略、OpenMP 的 \Tp{reduction(+:...)}）都可能随分块数改变顺序。要逐位复现就得把合并树写死，本册讲并行归约的地方专门讨论过这条约束。
\item \textbf{分块按索引，不按``下一个''。} 单元 $k$ 拿到固定区间 $[kB,(k+1)B)$，于是增删线程数不改变任何样本的归属，扩展性测试才有意义。用动态负载窃取来消除不均会立刻毁掉这一点——这是可复现性与负载不均之间一次真实的取舍。
\item \textbf{把发生器的身份写进产物。} 引擎名、种子、派生规则、块大小四样进元数据；缺任一项，半年后就无法解释这批数字是怎么来的。
\end{enumerate}
\end{bestpractice}

\begin{lstlisting}[style=cppstyle,caption={\CPP{}11 随机数做分块蒙特卡洛：每块一条派生流，合并顺序固定},label={lst:ch06-mc}]
// 编译：g++ -O2 -std=c++17 -pthread 本文件 -lm -o mc
#include <random>
#include <vector>
#include <cstdint>
#include <iomanip>
#include <iostream>

struct BlockStat {                // 每块一个结果，不共享任何状态
  std::uint64_t id;
  long double sum;                // 累加器用 long double，见本节末
  std::uint64_t hits;             // 同一条流上的第二个估计量
};

// 块 k 的样本区间是 [k*B, (k+1)*B)，与线程数无关
static BlockStat run_block(std::uint64_t k, std::uint64_t B,
                           std::uint32_t master) {
  std::seed_seq seq{master, static_cast<std::uint32_t>(k),
                    static_cast<std::uint32_t>(k >> 32)};
  std::mt19937_64 eng(seq);       // 引擎序列由标准规定，跨实现一致
  std::uniform_real_distribution<double> u(0.0, 1.0);

  BlockStat out{k, 0.0L, 0};
  for (std::uint64_t i = 0; i < B; ++i) {
    double x = u(eng), y = u(eng);
    // 平滑被积函数：真值是 Catalan 常数 G
    out.sum += 1.0L / (1.0L + (x * y) * (x * y));
    out.hits += (x * x + y * y <= 1.0);   // 同一个流上的 pi 估计
  }
  return out;
}

int main() {
  const std::uint64_t B = 1u << 20, nb = 64;
  std::vector<BlockStat> parts;
  parts.reserve(nb);
  // 这里顺序推进；改成多线程仍按 id 收集，结果不受影响
  for (std::uint64_t k = 0; k < nb; ++k)
    parts.push_back(run_block(k, B, 0xC0FFEEu));

  long double total = 0.0L, circle = 0.0L;  // 合并顺序固定 => 可复现
  for (const auto& p : parts) {
    total += p.sum / static_cast<long double>(B);
    circle += static_cast<long double>(p.hits);
  }
  const long double n = static_cast<long double>(nb);
  std::cout << std::setprecision(17)
            << "积分估计 " << double(total / n) << "\n"
            << "pi 估计 "
            << double(4.0 * circle / double(B * nb)) << "\n";
}
\end{lstlisting}

\begin{lstlisting}[style=pythonstyle,caption={\Tp{scipy.special} 与 \Tp{scipy.integrate} 的对拍：同一族算法，看的是差多少},label={lst:ch06-scipy}]
# pip install numpy scipy ；与前面的 C 清单输出逐项对比
import math
import numpy as np
from scipy import special, integrate

x = np.array([12.5])
print("J0", special.j0(x), "J1", special.j1(x))
print("Jn", [special.jv(n, 12.5) for n in range(5)])
print("Y0", special.y0(x))
print("lngamma", special.gammaln(400.0))
print("erf erfc", special.erf(0.75), special.erfc(3.0))
print("erfcx", special.erfcx(25.0))          # 缩放形式，避开下溢

# 端点奇异：把奇异点写进 points，或改用加权接口
val, err = integrate.quad(lambda t: 1.0 / np.sqrt(t),
                          0.0, 1.0, points=[0.0], limit=200)
print("quad", val, err)
print("quad_vec", integrate.quad_vec(
    lambda t: np.exp(-t * t), np.array([0.0, 1.0]))[0])
\end{lstlisting}

\begin{note}[title={门面与 C 侧对拍的判据}]
清单~\ref{lst:ch06-scipy} 与清单~\ref{lst:ch06-gsl} 的关系是这一章主线的落点：\PY{} 门面上的函数与 C 侧读的是同一族算法，所以应当\emph{要求}差异在 $1$ 到 $2$ 个 ulp 之内，而不是指望逐位相同。真差了五个数量级时，先查参数约定（$J_\nu$ 与球贝塞尔、\Tp{I} 与 \Tp{K} 对 \Tp{J} 与 \Tp{Y}、\Tp{erfc} 对 $1-\mathrm{erf}$ 是最常见的混法），再查尺度（GSL 的 \Tp{scaled} 系列返回的是乘过指数的值），最后才怀疑实现。这套顺序在本册后面所有``门面层与 C 层对拍''的场合都适用。另外注意：\PY{} 侧的一次调用看着便宜，实际要为装箱与数组检查付固定开销；把标量特殊函数放进百万次循环，正确做法是改用批量数组接口，或者干脆回到 C 侧循环——这正是``\PY{} 是门面''的具体含义。
\end{note}

\begin{guideline}[title={本章要点}]
\begin{enumerate}
\item \Lib{GSL} 的 \Tp{gsl\_} 前缀、布局可见的句柄、不复制的视图、错误码加 \Tp{\_e} 误差返回，合起来是 \CL{} 作为汇合点的标准语法：可预测、可 FFI、可跨语言，代价是生命周期与错误检查全归调用方。默认错误处理器的 \Fn{abort} 在宿主进程与多线程环境里必须一次性关掉（见本章的陷阱框）。
\item \Lib{Boost.Math} 用头文件-only 模板把开销挪到编译期，用 \Tp{policies::policy} 把精度（\Tp{promote\_float}、\Tp{promote\_double}）、性能与错误处理（\Tp{raise\_error} 与 \Tp{ignore\_error}）做成类型参数；换来零链接成本与可向量化的热循环，代价是编译期、目标码膨胀、没有稳定 ABI。
\item 特殊函数必须由底层实现：分段调度、极小极大有理逼近的系数、极点与溢出与递推方向的手写处理，三样都在库里面。库外重写丢掉的不是代码量，而是边界处理与测试基准。
\item 自适应求积的接口形状决定开销：\Fn{gsl\_integration\Bk\_qag} 是回调、工作区与 \Tp{epsabs}、\Tp{epsrel} 判据；端点奇异要用区间类型或 \Fn{qaws} 权重显式表达，无穷与振荡各有专门入口。\CPP{} 侧 \Tp{quadgk} 与双指数变换把同一族方法做成模板。有括区求根的稳健默认是 \Fn{toms748\_solve}。
\item 随机数的第一位属性是可复现性：\Fn{rand()} 与 \Fn{srand()} 因为进程级全局状态而不该出现在任何计算结果里；\CPP{}11 规定了引擎序列、没规定分布算法；\Lib{Random123} 用索引代替状态，因此适合 \Lang{GPU}、逐位复现与廉价检查点。
\item 蒙特卡洛的可复现分块有四条纪律（按索引派生流、固定合并顺序、样本归属与线程数无关、把发生器身份写进产物）；高精度按``主循环 \Tp{double}、累加 \Tp{long double}、验证用 \Tp{\_\_float128} 或 \Lib{mpmath}''分工，并把 \Tp{long double} 的宽度当运行时事实而不是假设。
\end{enumerate}
\end{guideline}
% ====================================================================

% ====================================================================
\chapter{数值软件栈：ODE、优化、最小二乘与自动微分}

前两章把线性代数内核讲完了：BLAS/LAPACK 管稠密，Eigen 管编译期抽象。但真实应用要解的不是矩阵乘法，而是``把这条微分方程推到 $t=2$\,s''、``让这批观测的残差平方和最小''。这一层——ODE 解算器、优化器、最小二乘与自动微分——有个反直觉的分布：接口语言一半是原生 \CL{}（\Lib{SUNDIALS}、\Lib{NLopt}、\Lib{SuiteSparse}），一半是纯 \CPP{} 模板（\Lib{Boost.Odeint}、\Lib{Ceres}），几乎没有 \FTN{} 当代入口。原因不复杂：这一层的``热路径''是每步几十到几百次浮点运算加上一次稀疏解，瓶颈在\emph{状态管理与调度}而不是 kernel 本身，于是决定语言选择的变量换成了生命周期、句柄与编译期配置。表~\ref{tab:ch07-stack} 先把六个栈摆齐，本章按行展开。

\begin{longtable}{@{}>{\raggedright\arraybackslash}p{0.14\textwidth}
  >{\raggedright\arraybackslash}p{0.08\textwidth}
  >{\raggedright\arraybackslash}p{0.14\textwidth}
  >{\raggedright\arraybackslash}p{0.12\textwidth}
  >{\raggedright\arraybackslash}p{0.12\textwidth}
  >{\raggedright\arraybackslash}p{0.11\textwidth}
  >{\raggedright\arraybackslash}p{0.08\textwidth}@{}}
\caption{六个数值软件栈的接口、求解器族、并行、许可基调与实时可行性}\label{tab:ch07-stack}\\
\toprule
栈 & 接口语言 & 求解器族 & 并行方式 & 许可证基调 & 适合的问题规模 & 进实时环 \\
\midrule
\endfirsthead
\toprule
栈 & 接口语言 & 求解器族 & 并行方式 & 许可证基调 & 适合的问题规模 & 进实时环 \\
\midrule
\endhead
\Lib{CVODE} & \CL{} 原生 & 可变阶 BDF 与 Adams 二族可切，Newton 配可插拔线性求解器；同族 IDA 解微分代数方程，CVODES 带灵敏度 &
  本体串行；向量、矩阵、线性求解三层抽象可换后端 &
  BSD-3 类，以仓库 LICENSE 为准 &
  状态维几千内单机从容，更大靠换后端 &
  否（自适应步长） \\
\Lib{ARKODE} & \CL{} 原生 & 加性隐显 Runge-Kutta，快慢项在方程右端就拆开；与 CVODE 同一套可插拔后端 &
  同上，可挂稀疏求解器与预处理 &
  同上 &
  对流占优的 PDE 线、刚性比已知的系统 &
  否 \\
\Lib{Odeint} & \CPP{} 头文件 & 显式 RK 一族（欧拉到 dopri5）、辛算法、广义多步；stepper 可与误差控制器组合 &
  无内部并行；可经容器 traits 接 GPU 后端卸载 &
  BSL 类，以仓库 LICENSE 为准 &
  步进级；状态容器由类型参数决定 &
  定步长路径可（见 \S\ref{sec:ch07-rt}） \\
\Lib{NLopt} & \CL{} 核心加各语言 binding & 局部无导数（BOBYQA、Nelder-Mead）、局部梯度（LBFGS、MMA）、全局（DIRECT、CRS、多起点） &
  无内部并行；句柄非线程安全，一一线程一句柄 &
  LGPL 类，以仓库 LICENSE 为准 &
  几十到几百维、带一般约束的目标 &
  否（迭代次数依数据） \\
\Lib{Ceres} & \CPP{} 模板加少量 \CL{} 式接口 & LM、Dogleg、TRF；稠密、稀疏正规方程、稀疏消元与迭代解 &
  自带线程池，按分片并行评估残差与偏导 &
  Apache-2.0 类，以仓库 LICENSE 为准 &
  残差十万到百万级，参数千到万 &
  否（离线批处理） \\
\Lib{Suite\Bk Sparse} & \CL{} 句柄 & UMFPACK 稀疏 LU、CHOLMOD 稀疏 Cholesky、LDL 参考实现、SPQR 稀疏 QR，外加置换与消元树工具族 &
  个别包有 OpenMP 编译开关，默认串行 &
  GPL 系为主（个别包 LGPL 类），以仓库 LICENSE 为准 &
  稀疏线性系统，亿级非零元有实证 &
  边界：符号相可预计算 \\
\bottomrule
\end{longtable}

\section{常微分方程求解：SUNDIALS 与 Boost.Odeint}\label{sec:ch07-ode}

\Lib{SUNDIALS}（LLNL）是 \CL{} 侧数值栈的标本，四个套件各管一类问题：\Fn{CVode} 解显式 ODE，\Fn{IDA} 解微分代数方程组（DAE，带代数约束与微分指数），\Fn{ARKODE} 解加性（隐显分裂）ODE，\Fn{CVODES} 在 CVODE 之上做灵敏度与积分优化。它值得细看的不是算法，而是\emph{接口工程}。

第一，\CL{} 是原生实现语言，\FTN{} 接口是\emph{生成}的产物。SUNDIALS 提供 \FTN{77} 与 \FTN{2003} 两套接口，官方文档写明它们由 C 源码自动包装生成、不另立语义——这与第 4 章 netlib 谱系``\FTN{} 原文加 C 包装''的方向恰好相反。第二，三层可插拔抽象全是 C 级多态：\Tp{N\_Vector}（向量容器）、\Tp{SUNMatrix}（矩阵存储）、\Tp{SUNLinSol}（线性求解器）各自是一个函数指针表加句柄的 struct——serial、parallel NVector、dense、band、sparse（KLU、SuperLU）互为替身。这正是第 2 章 C ABI 纪律在解算器层的复用：稳定符号、可见布局、可跨语言 FFI。第三，读 API 要认新旧两代：v7 起所有构造函数带 \Tp{SUNContext} \Bk{}（为 MPI 等并行环境提供资源句柄），v6 同名函数少这一个参数；按清单~\ref{lst:ch07-cvode} 的读法，两代都能一分钟对上号。

CVODE 的行为开关值得逐个点名。\emph{BDF 与 Adams 的切换}：创建时选 \Tp{CV\_BDF}（后退微分公式，1 到 2 阶可变阶，刚性）或 \Tp{CV\_ADAMS}（前向积分同族，非刚性），新接口允许以自动模式让求解器按刚性估计在两族间换；换族意味着重开阶历史，代价体现在步数统计里。\emph{稠密输出}：每个被接受的步都留插值多项式，\Fn{CVodeGetDense} 可以在步内任意时刻取值——绘图、控制采样、与另一个定步长系统对表都靠它，而不是把输出步长强加给积分器。\emph{根查找}：\Fn{CVodeRootInit} 注册若干个$g_k(t,y)$，求解器在每个步区间内变号定位、二分收窄、以 \Tp{CV\_ROOT\_RETURN} \Bk{} 把事件交给调用方——碰撞、阈值、状态机切换是事件驱动仿真的骨架（与实时仿真那一章的接口问题同源）。\emph{Newton 线性代数}：隐式法每步解一个非线性方程，Jacobian 可解析提供、可用近似差分、可矩阵自由（只给 Jacobian-向量积配 Krylov 预处理）。

清单~\ref{lst:ch07-cvode} 是最小可编译的 CVODE 用法：阻尼摆，BDF 加 Newton，稠密矩阵加稠密 LU，手给解析 Jacobian。

\begin{lstlisting}[style=cstyle,caption={\CL{} 用 CVODE（按 SUNDIALS v7 新 API 读；v6 少 \Tp{SUNContext}）},label={lst:ch07-cvode}]
#include <cvode/cvode.h>          /* CVODE 主头 */
#include <nvector/nvector_serial.h>
#include <sunmatrix/sun_densematrix.h>
#include <sunlinsol/sunlinsol_dense.h>
#include <math.h>
/* 右端函数。v7 原型统一带 user_data；v6 没有末参 */
static int PendRhs(sunrealtype t, N_Vector y,
                   N_Vector ydot, void* user) {
  realtype* s = (realtype*) NV_DATA(y);
  realtype* f = (realtype*) NV_DATA(ydot);
  f[0] = s[1];
  f[1] = -0.2 * s[1] - sin(s[0]);   /* 阻尼摆 */
  return CV_SUCCESS;
}

/* 解析 Jacobian：Newton 每步迭代都会要它 */
static int PendJac(sunrealtype t, N_Vector y, N_Vector fy,
                   SUNMatrix J, void* user,
                   N_Vector t1, N_Vector t2, N_Vector t3) {
  realtype* s = (realtype*) NV_DATA(y);
  SM_ELM(J, 0, 0) = 0.0;        SM_ELM(J, 0, 1) = 1.0;
  SM_ELM(J, 1, 0) = -cos(s[0]); SM_ELM(J, 1, 1) = -0.2;
  return CV_SUCCESS;
}

int simulate(double t_end) {
  SUNContext ctx; SUNContext_Create(SUN_COMM_NULL, &ctx);
  N_Vector y = N_VNew_Serial(2, ctx);
  NV_DATA(y)[0] = 3.0;  NV_DATA(y)[1] = 0.0;

  void* cv = CVodeCreate(CV_BDF, CV_NEWTON);
  CVodeInit(cv, PendRhs, 0.0, y);
  CVodeSStolerances(cv, 1e-8, 1e-10); /* 相对/绝对容差 */
  CVodeSetFunctions(cv, &PendRhs, NULL, NULL, NULL,
                    &PendJac);        /* 旧代码：SetJacFn */
  SUNMatrix A = SUNDenseMatrix(2, 2, ctx);
  SUNLinSol S = SUNLinSol_Dense(y, A, ctx);
  CVodeSetLinearSolver(cv, S, A);
  realtype tret;
  for (int i = 1; i <= 20; ++i) {
    int flag = CVode(cv, t_end * i / 20.0, y, &tret,
                     CV_NORMAL);
    if (flag < 0) return 1;   /* 负值为错，正值为根事件 */
  }
  CVodeFree(&cv);  SUNLinSolFree(S); SUNMatDestroy(A);
  N_VDestroy(y);   SUNContext_Free(&ctx);
  return 0;
}
\end{lstlisting}

\CPP{} 侧的对照面是 \Lib{Boost.Odeint}：头文件-only，没有句柄、没有链接、没有运行时。它的组织方式是把数值方法拆成正交的两半——stepper（单步推进器：定步长 RK、Dormand-Prince、辛算法、多步法）与 policy（误差控制器、迭代修正器、步长序列），用模板参数\emph{组装}而不是用参数\emph{配置}；同一个 Dopri5 stepper，接上 \Tp{make\_controlled} \Bk{} 就是自适应积分器，接 \Tp{integrate\_const} 就是定步长积分器。状态容器经 traits 解耦：\Tp{std::vector}、裸数组、\Tp{std::array}、Eigen 的 \Tp{VectorXd} 都能当状态，容器操作（缩放、加法、范数）由对应的 traits 特化提供——这是第 5 章 Eigen 表达式模板能被喂进解算器的原因，也是``模板即配置''在本章收束节的具体所指。GPU 后端（Thrust 一族）走的同样是 traits 这条路。清单~\ref{lst:ch07-odeint} 给出定步长与自适应两种写法，外加编译期步进器配置。选型时真正要核对的是``这一版 Boost 里它还在不在''：\Lib{Boost.Odeint} 长期以头文件-only 形态随 Boost 分发，近年正在拆为独立仓库维护，发行形态确实在变。

\begin{lstlisting}[style=cppstyle,caption={Boost.Odeint：\Tp{integrate\_const} 配 \Tp{runge\_kutta\_dopri5}，同一段右端函数接两种驱动},label={lst:ch07-odeint}]
#include <boost/numeric/odeint.hpp>
#include <array>
#include <cmath>
#include <vector>
using namespace boost::numeric::odeint;
using state_t = std::vector<double>;

void pendulum(const state_t& s, state_t& ds, double) {
  ds[0] = s[1];
  ds[1] = -0.2 * s[1] - std::sin(s[0]);
}

int main() {
  state_t y = {3.0, 0.0};
  // 定步长：Dopri5 在这里只当五阶单步法用
  runge_kutta_dopri5<state_t> step;
  integrate_const(step, pendulum, y, 0.0, 20.0, 0.01,
    [&y](double t, const state_t& s) {
      (void) t; (void) s;   // 观测器：落盘、喂控制、进队列
    });

  // 自适应：同一个算法包上误差控制器，容差进类型
  integrate_adaptive(
    make_controlled<runge_kutta_dopri5<state_t> >(1e-9, 1e-6),
    pendulum, y, 0.0, 20.0, 0.05);
  // 编译期状态：array 作容器，尺寸进类型，无堆分配
  using fixed_t = std::array<double, 2>;
  fixed_t z = {3.0, 0.0};
  runge_kutta4<fixed_t> rk4;
  for (int i = 0; i < 2000; ++i) rk4.do_step(z, 0.01);
  return 0;
}
\end{lstlisting}

刚性判定与选 solver 的实操口径，按成本从低到高排：（1）物理判据——系统里最快与最慢时间尺度之比超过约三个数量级，显式法的步长会被\emph{稳定性}而不是精度卡死，这是刚性的定义性症状；（2）事后统计——先用非刚性路径（CVODE 的 Adams 档或 odeint 自适应）跑一遍，比较 \Fn{CVodeGetNumSteps} 与输出点需求：每单位时间步数远超精度所需、且步长反复贴地，就该换 BDF；（3）谱判据——对稳态点取显式 Jacobian 算特征值，负实部跨度大即刚性，代价是你要真的能写出 Jacobian；（4）结构判据——快慢项在方程里本来就可分离（对流项慢、扩散项快）时，不要二选一，直接上 ARKODE 的隐显分裂，让快项进隐式、慢项留显式。经验法则的最后一条兜底：拿不准时，BDF 解非刚性问题只是慢一点，显式法碰刚性问题是\emph{静默地}花掉你一百倍时间。

\section{优化与最小二乘：NLopt、Ceres 与因子图一笔}\label{sec:ch07-opt}

\Lib{NLopt} 是``一个句柄''设计的典范：\Tp{nlopt\_opt} 是一个不透明的 \CL{} 结构指针，算法、上下界、停止判据、约束、目标函数全部以 setter 挂上去；求解就是 \Fn{nlopt\_optimize} 一句。各语言 binding（\PY{}、\Lang{Julia}、Lua 等）包装的是\emph{同一个}指针语义，算法常量在全语言共用一张编号表——第 2 章讲的``C 核心加薄门面''在这里是字面意义上的实现。导数可选机制是它最有教学价值的一处：目标函数的原型是 \Tp{double f(unsigned n, const double\* x, double\* grad, void\* data)}，\Tp{grad} 参数\emph{可能为空}——无导数算法传 \Tp{NULL}，梯度算法要你在非空时填好。同一个 \CPP{} lambda 因此可以喂给 BOBYQA（无导数）也可以喂给 LBFGS（带导数），换档不改目标函数代码。算法按两轴分档：局部对全局（梯度对无导数），带一般约束的只有少数档（MMA、CCSAQ 一族）；句柄非线程安全，多起点并行的正确姿势是一线程一句柄。

\Lib{Ceres}（Google）解的是另一类问题：大规模非线性最小二乘 $\min_x \sum_i \rho(\|f_i(x)\|^2)$，残差块天生稀疏、天生可并行。它的形状由三件事决定。\emph{其一，自动微分}：用户交上来的是一个模板化的仿函数（functor），\Tp{AutoDiffCostFunction} \Bk{} 用 \Tp{Jet} 类型把它实例化，编译期得到解析到机器精度的 Jacobian；手写导数走 \Tp{SizedCostFunction} 家族，差分兜底走 \Tp{NumericDiffCostFunction}（见 \S\ref{sec:ch07-ad}）。\emph{其二，问题图}：残差块与参数块的关联就是一张二部图，Ceres 在它上面做列分块、Schur 消元与稀疏正规方程 $J^{\top}J\,\Delta = -J^{\top}r$，线性代数全走稀疏路径——这正是它能吃满多线程的原因：残差评估彼此独立，库自带线程池按分片（shard）铺 job，Jacobian 各写各的格子，无锁可写。\emph{其三，依赖按档位点亮}：只依赖 Eigen 即可跑稠密档；打开 SuiteSparse（CHOLMOD、COLAMD）或自带的 CXSparse 才有稀疏正规方程与稀疏消元档；LAPACK 与 BLAS 被这些稀疏内核在超节点内部回调——第 4 章的 kernel 与本章的解算器在这里咬合。

清单~\ref{lst:ch07-ceres} 给最小闭环：同一个残差仿函数，自动微分与中心差分两档，稀疏正规方程档求解。

\begin{lstlisting}[style=cppstyle,caption={Ceres：模板仿函数、\Tp{AutoDiffCostFunction} 与稀疏正规方程},label={lst:ch07-ceres}]
#include <ceres/ceres.h>
#include <memory>
#include <vector>

// 重力拟合：y_i = 0.5 * g * t_i^2，残差只写一遍
struct DropResidual {
  DropResidual(double y, double t) : y_(y), t_(t) {}
  template <typename T>
  bool operator()(const T* const g, T* res) const {
    res[0] = T(y_) - T(0.5) * (*g) * T(t_) * T(t_);
    return true;                 // false：本次评估作废
  }
  double y_, t_;
};

int fit_gravity(double* g, const std::vector<double>& ys,
                const std::vector<double>& ts) {
  using Auto = ceres::AutoDiffCostFunction<DropResidual,
                                           1, 1>;
  using Num  = ceres::NumericDiffCostFunction<
      DropResidual, ceres::CENTRAL, 1, 1>;
  ceres::Problem problem;
  for (size_t i = 0; i < ys.size(); ++i) {
    // 默认 loss：平方；换 ceres::HuberLoss(1.0) 即鲁棒档
    // 偶数行走自动微分，奇数行走差分档：同一仿函数
    if (i % 2 == 0)
      problem.AddResidualBlock(
          Auto::Create(new DropResidual(ys[i], ts[i])),
          nullptr, g);
    else
      problem.AddResidualBlock(
          Num::Create(new DropResidual(ys[i], ts[i])),
          nullptr, g);
  }
  ceres::Solver::Options opt;
  opt.linear_solver_type = ceres::SPARSE_NORMAL_EQUATIONS;
  opt.num_threads = 1;      // 钉死线程数才可逐位复现
  ceres::Solver::Summary sum;
  ceres::Solve(opt, &problem, &sum);
  return sum.termination_type == ceres::CONVERGENCE ? 0 : 1;
}
\end{lstlisting}

\begin{compare}[title={NLopt 与 Ceres：两条不相交的供给线}]
\begin{itemize}
\item \textbf{问题形状：}NLopt 吃``一个黑箱标量函数加盒式或一般约束''，维度到几百；Ceres 吃``显式拆开的成百上千个残差块''，维度到十万级。
\item \textbf{导数来源：}NLopt 要么你自己给梯度，要么用根本不问导数的算法；Ceres 默认在编译期\emph{生成}梯度，差分只是兜底。
\item \textbf{线程归属：}NLopt 无内部并行，并行是调用方多句柄的事；Ceres 自带线程池，确定性反而是调用方钉线程数的事。
\item \textbf{接口形态：}NLopt 是 C 句柄加 setter，天然适合做各语言 binding；Ceres 的接口主体是模板，binding 只能包住最外层的 Problem 语义。
\end{itemize}
\end{compare}

因子图优化（机器人 SLAM 与视觉里程计）是 Ceres 的领域特化近亲：\Lib{g2o} 与 \Lib{GTSAM} 把``残差块—参数块''图换成``因子—位姿与观测点''图，求解器族同源（LM 加稀疏正规方程或 Cholesky），卖点在流形上的重参数化与增量式重线性化；它们的存在说明本章的问题图不是 Ceres 的私设，而是这一类问题的公共形状。

\section{自动微分：三种形态如何改变库的形状}\label{sec:ch07-ad}

先给判据：自动微分不是符号微分，也不是有限差分，而是\emph{把程序的每个初等运算按链式法则接成一张计算图再求导}。按图的携带方式分三种形态，表~\ref{tab:ch07-ad} 摆开。第一个形态是运算符重载编译期展开：Ceres 的 \Tp{Jet<T,N>} \Bk{} 是一个``值加 $N$ 个偏导''的结构体，加减乘除与基本函数都被重载为同时传播偏导；\Tp{N} 在\emph{模板参数}上，偏导数组是栈上定长数组，链式法则的循环在编译期展开成直线代码。Eigen 侧的对应物是 \Tp{AutoDiffScalar} \Bk{}（梯度挂动态向量，代价每处高一点，换来类型简单）。第二个形态是记录式反向（tape）：前向扫时把中间值记进带，反向扫按带累积伴随；成本随\emph{输出数}线性，适合``多输入单输出''（训练里一个标量损失对一万参数），代价是带的存储——它必须活到反向扫结束。第三个形态是源到源改写与代码生成：库在表达式层求导，吐出的却是 \CL{} 文件——\Lib{CasADi} \Bk{} 的 \Fn{codegen} \Bk{} 即为此设，生成物只依赖 C ABI，任何宿主都能链接。

\begin{longtable}{@{}>{\raggedright\arraybackslash}p{0.12\textwidth}
  >{\raggedright\arraybackslash}p{0.21\textwidth}
  >{\raggedright\arraybackslash}p{0.17\textwidth}
  >{\raggedright\arraybackslash}p{0.14\textwidth}
  >{\raggedright\arraybackslash}p{0.16\textwidth}@{}}
\caption{求导的四档形态：机制、代价、语言要求与代表实现}\label{tab:ch07-ad}\\
\toprule
形态 & 机制 & 代价与精度 & 语言要求 & 代表实现 \\
\midrule
\endfirsthead
\toprule
形态 & 机制 & 代价与精度 & 语言要求 & 代表实现 \\
\midrule
\endhead
有限差分（退化档） & 扰动输入重评函数，商近似偏导 &
  每列至少一次重评；最优步长约根号 eps，此处截断与消去误差打平，精度丢半数有效位 &
  无；任何语言任何黑箱 &
  Ceres \Tp{NumericDiff\Bk Cost\Bk Function} 与各库差分兜底 \\
前向模式，运算符重载 & 双数型携带值与偏导数组，链式法则随每个运算传播 &
  一次求值同时得函数值与全部偏导；成本随\emph{输入数}线性；机器精度 &
  模板与运算符重载（或等价的类型参数机制） &
  Ceres 的 \Tp{Jet}、Eigen 的 \Tp{AutoDiff\Bk Scalar} \\
反向模式，记录式 & 前向扫记带存中间值，反向扫累积伴随 &
  成本随\emph{输出数}线性；换来一次分配整幅梯度，带本身占内存 &
  重载加动态数组，或编译器插桩 &
  ADOL-C；训练框架的梯度带 \\
源到源、代码生成 & 在表达式层求导，生成可读的 \CL{} 源再编译 &
  运行期零解释开销；代价挪进工具链与构建步骤 &
  产出物只要求宿主能吃 C &
  CasADi 的 \Fn{codegen}、符号库的导出器 \\
\bottomrule
\end{longtable}

``为什么 \CPP{} 模板能做到而 \CL{} 不能''是可以逐条核对的，不玄：（1）\emph{编译期展开}——同一个仿函数被实例化两次（\Tp{T} 取 \Tp{double} 得残差，取 \Tp{Jet} 得 Jacobian），没有虚分派、没有运行期类型判断；\CL{} 没有函数模板，做不到``一份源码两种类型''，除非手写两遍或宏拼接（宏丢掉类型检查）。（2）\emph{类型即图}——\Tp{Jet<double,\,7>\Bk} 的偏导宽度写在类型里，循环展开与栈分配由类型决定；\CL{} 的类型系统里不存在这个自由度。（3）\emph{运算符重载}——链式法则藏在 \Tp{operator*} \Bk{} 里，用户写的是代数式不是求导式；\CL{} 的运算符不可重载。所以 \CL{} 侧自动微分只能取两个退化形状：把导数当\emph{外部函数}（用户自备 Jacobian 函数指针，NLopt 与 SUNDIALS 的解析 Jacobian 槽都是这个形状），或者把\emph{生成}当外挂（代码生成器产出 C，回到第 2 章``C 是汇合点''的论点：不是 C 会微分，是生成的产物只认 C）。

\section{稀疏专用内核：SuiteSparse 一族与``为什么不叫 BLAS''}\label{sec:ch07-sparse}

最小二乘、Newton 步、隐式 ODE 每走一步都要解 $Ax=b$，而这些 $A$ 几乎总是稀疏的。SuiteSparse 是这一族的中心仓库：\Lib{UMFPACK} \Bk{}（非对称稀疏 LU，左视超节点加 Markowitz 型置换）、\Lib{CHOLMOD} \Bk{}（对称正定，稀疏 Cholesky， simplicial 与超节点两路，AMD/CAMD 后置置换）、\Lib{LDL} \Bk{}（纯 LDL 转置的教学级参考实现，小而全）、\Lib{SPQR} \Bk{}（多重前沿稀疏 QR，给欠定与秩亏问题）、\Lib{KLU} \Bk{}（为电路仿真调校的小矩阵 LU，不做填元优化但启动极快，SUNDIALS 的稀疏槽默认伙伴），外加 COLAMD、AMD 这些置换工具。

为什么这一族不在 BLAS/LAPACK 名下？三条机制上的分岔。（1）\emph{填元是数据依赖的}：稀疏 LU 的中间非零元集合要做了符号分析（消元树加置换）才知道，$A$ 换一组数值，结构就可能变——BLAS 的接口前提是``问题形状给定即可分块''，这里根本不成立。（2）\emph{访存模式反吞吐}：稀疏三角回代是指针追逐，算术强度随矩阵掉到远小于 1，Level 3 那种规整流式访问的假设整体失效；但\emph{超节点内部}又是局部稠密的，所以 CHOLMOD 恰恰在超节点上回调 BLAS——稀疏库不是 BLAS 的替代，是 BLAS 的调度器。（3）\emph{接口形态因此必须长得不一样}：CCS 三元组句柄、符号相与数值相分离（\Fn{cholmod\_analyze} 对 \Fn{cholmod\_factorize}）、内存钩子（\Fn{cholmod\_set\_memory\_func}，把 malloc 交还宿主）——这三样在稠密世界都不需要。清单~\ref{lst:ch07-cholmod} 就是这套句柄纪律的速写。

\begin{lstlisting}[style=cstyle,caption={CHOLMOD（长索引接口）：符号相、数值相、回代三段},label={lst:ch07-cholmod}]
#include <cholmod.h>

/* a：CCS 下三角；b：右端。解 x 经 out 交还调用方 */
int solve_spd(cholmod_sparse* a, cholmod_dense* b,
              cholmod_dense** out) {
  cholmod_common cm;
  if (cholmod_l_start(&cm) < 0) return 1;

  cholmod_factor* L = cholmod_l_analyze(a, &cm);
  if (!L) { cholmod_l_finish(&cm); return 2; } /* 符号相 */

  int ok = cholmod_l_factorize(a, L, &cm);
  if (ok && cm.status == CHOLMOD_OK) {         /* 数值相 */
    cholmod_dense* x = cholmod_l_solve
        (CHOLMOD_A, L, b, &cm);                /* 回代 */
    if (!x) ok = 0; else { *out = x; }
  }
  cholmod_l_free_factor(&L, &cm);
  cholmod_l_finish(&cm);
  return ok ? 0 : 3;
}

/* 复用形态：矩阵只换数值不换结构时，analyze 一次、
 * 每轮只 factorize 加 solve；内存钩子可绑宿主池。 */
\end{lstlisting}

\begin{warning}[title={``稀疏一词的两处歧义}]
\begin{itemize}
\item 稀疏\emph{存储}不等于稀疏\emph{求解}：Eigen 的 Sparse 与 \Lib{MKL} 的稀疏求解器是内核，本章的 SuiteSparse 是带置换与相位的完整解算器，二者在依赖图的不同层，比较``谁快''没有意义。
\item ``Ceres 依赖 SuiteSparse''是分档的：默认构建只吃 Eigen，稀疏档要显式打开开关；部署时看到\Tp{libsuitesparseconfig} 的链接与否，才能确定你拿到的是哪一档。
\item CHOLMOD 只对对称正定可靠；数值不定（Newton 步的 $J^{\top}J$ 秩亏常见）时，正确档位是 LDL 转置带扰动、SPQR，或 Ceres 的迭代加 Truncated 预处理档。
\end{itemize}
\end{warning}

\section{实时与确定性：哪些能进环}\label{sec:ch07-rt}

硬实时环（音频回调、控制周期）对这一层的判据只有三条：单步成本是否\emph{有界}、内存是否\emph{不增长}、错误路径是否\emph{不崩溃宿主}。按这三条过筛：

\begin{itemize}
\item \textbf{可进环：}odeint 的定步长显式路径（\Tp{std::array} 一类零分配容器加 \Tp{do\_step}，成本就是固定次数的运算，清单~\ref{lst:ch07-odeint} 的最后一段就是它的形状）；预计算增益后的一次矩阵乘（线性 MPC 的显式解）；固定迭代步数的 Krylov 法。
\item \textbf{软实时边缘：}\Fn{cholmod\_factorize} 加回代——符号相预计算、内存钩子绑固定池之后，单轮时延可界定，但仍是数据相关；\Tp{CV\_NORMAL} 模式下一次 \Fn{CVode} 调用内部步数不定，只能在\emph{界内退化}（限 \Fn{CVodeSetMaxNumSteps} 为 1 逐帧推进）时用作观察，不是用作实时求解。\item \textbf{不进环：}一切自适应步长解算器（CVODE、ARKODE、IDA 与 odeint 的 \Tp{integrate\_adaptive}）；NLopt（迭代次数依数据）；Ceres（秒到分级，内部线程池加动态分配）；UMFPACK 与 SPQR 的数值相（填元随数据）。
\end{itemize}

\begin{bestpractice}[title={进环前的四项体检}]
\begin{enumerate}
\item \textbf{测最坏而不是平均：}用 WCET 的思路跑边界输入（刚性爆发、Jacobian 病态、约束退化），看单回调耗时上界，不看均值。
\item \textbf{看分配行为：}对回调做分配器插桩（或链接计数的 \Fn{malloc} 钩子），一帧内出现任何一次堆增长即不合格；\Lib{SuiteSparse} 一族与 SUNDIALS 都提供把内存钩子交给宿主的入口，用好它。
\item \textbf{查错误路径：}这些库的默认错误处理（SUNDIALS 的 err handler、GSL 时代的 \Fn{abort}）会杀掉宿主；环内必须换成``返回错误码加保持旧状态''。
\item \textbf{钉确定性：}Ceres 的结果会随 \Tp{num\_threads} 改变浮点归约顺序，逐位复现要求钉线程数（清单~\ref{lst:ch07-ceres} 里 \Tp{num\_threads=1} 的原因）； odeint 定步长路径天然逐位可复现，前提是别在中途换容器。
\end{enumerate}
\end{bestpractice}

\section{小结：为什么这一层仍是 \CL{} 与 \CPP{}}\label{sec:ch07-close}

\begin{guideline}[title={本章要点}]
\begin{itemize}
\item \textbf{模板即配置：}odeint 用类型参数组装 stepper 与 controller，Ceres 用模板仿函数在编译期生成 Jacobian——``配置''这件事被搬到编译期，运行期只剩裸循环；这是 \CPP{} 在这一层的核心竞争力，与第 5 章 Eigen 的表达式模板是同一招。
\item \textbf{零运行时与可静态链接：}这一层没有 GC、没有解释器、没有抖动分配；头文件-only 的 odeint 连链接对象都没有，SUNDIALS 可以整库静态焊进宿主。音频插件、固件与宿主进程因此能把它``吃进去''而不是``调出去''——热路径经济学的老论点（第 1 章）在这里的形态是：每帧的固定开销预算表上，``库的存在感''本身是一项要归零的成本。
\item \textbf{句柄加函数指针是 \CL{} 的多态答案：}SUNDIALS 的三层可插拔抽象、NLopt 的一个句柄、CHOLMOD 的内存钩子，全是第 2 章 C ABI 纪律与第 3 章共同约定（布局可见、错误码返回、一线程一句柄）的复用；\FTN{} 在这一层只剩生成出来的兼容接口，方向与 netlib 时代相反。
\item \textbf{选栈的顺序：}先判刚性（时间尺度分离、步数统计、谱跨度），再定导数来源（自动微分、手导、差分兜底），再定线性代数档位（稠密、稀疏正规方程、消元），最后按实时判据过筛——四步都问的是机制，不是口碑。
\end{itemize}
\end{guideline}

下一章进入几何库：当问题从``把 y 推到 t''变成``这四个点谁在谁里面''，数值问题会换上另一副面孔——那里的敌人不是刚性，而是舍入符号。
% ====================================================================
% ====================================================================

% ====================================================================
% ====================================================================
\part{几何库}

\chapter{几何谓词、鲁棒计算与坐标投影}

数值计算一章关心的是``算得准''：残差、条件数、收敛阶，误差再大也只是让结果差几位。几何这一域的第一性问题却形如一个布尔判定：点在定向直线的哪一侧？点在三角形的外接圆内还是圆外？这两比特的判定一旦出错，上层算法不是精度下降，而是直接产出拓扑上非法的对象——自交的三角剖分、内外翻转的多边形、永远走不出的翻转循环。本章先把问题定位到谓词上，再给出三条鲁棒化路线（自适应精确谓词、区间算术、任意精度），然后依次进入这一域的三类代表：把谓词实现做成模板参数的 \Lib{CGAL}、把 \CL{} ABI 用到极致的 \Lib{GEOS} 与 \Lib{PROJ}，并在末尾回答为什么这一域几乎只能由 \CL/\CPP{} 承担、\PY{} 只配当门面。

\section{公共难题是谓词，不是算法}

平面几何算法——凸包、约束 Delaunay 三角剖分、点定位、多边形叠加——拆开到底，反复调用的是同一个核心判定：点 $c$ 在有向直线 $ab$ 的左侧、右侧还是线上。记
\[
  D=(a_x-c_x)(b_y-c_y)-(a_y-c_y)(b_x-c_x),
\]
它恰是以 $(a_x-c_x,\ a_y-c_y)$ 与 $(b_x-c_x,\ b_y-c_y)$ 为两行的二阶行列式的展开，而 \Tp{orient2d}(a,b,c) 就是 $D$ 的符号：正、负、零分别对应左、右、共线。Delaunay 的空圆条件再加一个四元组判定 \Tp{incircle}（点 $d$ 在 $a,b,c$ 外接圆内与否），它是三阶量的符号。这两个谓词撑起几乎整个平面剖分算法族；三维情形还有 \Tp{orient3d}（点与三角形面片的方位）与 \Tp{insphere}，四面体化与体素化输出的合法性同样系于它们的符号。

问题在于，用 \Tp{double} 直接算 $D$ 时，它是两个乘积之差：当真实值接近零——三点几乎共线、四点几乎共圆——两个乘积的有效数字在相减时被大量抵消（catastrophic cancellation），留下的低位噪声完全由舍入决定，符号成了掷硬币。更糟的是输入往往\emph{天生}接近退化：坐标是 GPS 收来的六位小数、是从别的投影逆变换回来再截断的、是两条线段求交的浮点结果。在这类数据上，``一般位置''是奢侈品而非常态。

符号判错触发的是\emph{自相矛盾}而不是``差一点''：\Tp{orient2d}(a,b,c) 与 \Tp{orient2d}(c,b,a) 本应反号，浮点版可能同号；Lawson 局部优化依据 \Tp{incircle} 决定要不要翻一条边，与边定向判定一旦打架，同一条边会被反复翻转，剖分产物出现自交三角形，导出的网格在下一个读入它的软件手里炸开。这不是某个实现的 bug，而是\textbf{用浮点等式冒充符号判定}这一做法的必然：谓词输出的是拓扑决策，拓扑没有``接近正确''。

还有一个放大机制让谓词问题比单个 bug 更严重：谓词位于调用栈的最底层，被凸包、剖分、叠加、可见性等多个算法族\emph{共用}。一个把 \Tp{orient2d} 写成三点叉积减法的``小工具函数''一旦被复用进五六个算法，它的失败模式也被复制五六个份；而上层算法无法察觉自己拿到的符号是噪声，因为它们之间的契约本来就只有一比特。所以这一域的库把``谁提供谓词''当成架构问题：要么全栈共享一个经过滤波的谓词内核，要么就别声称自己支持退化输入。

\section{鲁棒化的三条路线}

\subsection{Shewchuk 的自适应精确谓词}

这一域的标杆是 Jonathan Shewchuk 在二十世纪九十年代中后期给出的自适应精确谓词（adaptive exact predicates），至今仍是 \Lib{CGAL} 之外多数精确几何内核的参考实现。它的结构是三层滤网：

\begin{itemize}
\item \textbf{第一层：纯浮点估计。} \Fn{orient2dfast} 用普通 \Tp{double} 行列式给出估计值 $po$。绝大多数远离退化的调用在这层就返回，成本与裸实现相同。
\item \textbf{第二层：误差界。} Shewchuk 推导出一个只依赖坐标绝对值与舍入单位 $u=2^{-52}$ 的多项式上界 $errbound$；若 $|po|>errbound$，真实值不可能过零，浮点符号可信，直接返回。这一层只在第一层结果可疑时才计算，且本身仍是廉价浮点运算。
\item \textbf{第三层：精确展开。} 还不决断，就进入用\emph{未求和的两数组合}（unevaluated sum）搭建的展开算术：\Fn{TwoSum}、\Fn{TwoProduct} 把每一步乘加精确拆成高低两段，误差项被逐项界住、逐级补算（\Fn{orient2dadapt}），必要时落到全精度展开 \Fn{orient2dexact}，给出与任意精度实数算术逐位一致的符号。
\end{itemize}

它的天才之处在于\textbf{只把精确算术花在布尔判定上}：谓词的输出是一比特的符号，可以承受最昂贵的算法；而构造类运算（外心坐标、交点坐标）留给浮点或较低档的精确性。随机均匀数据里落在第三层的调用占比极低，退化数据里代价上升但结果依然正确。这一``先滤后精''的结构后来成为整个精确几何计算的通用骨架。

值得再说清一层代价结构：展开算术里没有除法，每一步都是若干个 \Tp{double} 的精确组合，第三层的开销是\emph{可数}的——大约几十到上百次浮点加乘；真正不可控的从来不是谓词，而是把精确算术用到构造上的那些方案。所以工程上的推论很直接：凡是``只在符号处较真''的实现都能进热点循环，凡是``每一步都较真''的实现只能进离线流程。

\subsection{区间算术与任意精度：两条兜底}

区间算术给每个运算配一对向外舍入的端点：$x+y\mapsto[\underline{x}+\underline{y},\ \overline{x}+\overline{y}]$（代表库：\Lib{Boost.Interval}、\Lib{FIL}、\Lib{MPFI}）。判定规则只有一条：结果区间不含零则符号确定，含零则报告``无法判定''并请求更细的输入划分。它与自适应谓词的差别是结构性的：区间宽度在依赖复用的表达式里会放大（同一变量出现两次被当成两个独立区间，即所谓依赖问题），几层运算后区间可能宽到永远含零，所以它适合当构造运算的粗滤器，而不是符号判定的终点。

任意精度路线不做估计也不做界，直接把坐标装进精确数域再判：整数与有理数用 \Lib{GMP}（\Tp{mpz}、\Tp{CGAL::Gmpq}，后者的分子分母各是一个大整数），超越函数用带舍入方向控制的 \Lib{MPFR}（\Tp{mpfr\_float}）。一次 \Tp{incircle} 判定变成一次完整的有理数行列式，符号永远正确；代价是位数与内存随运算链增长，且增长量\emph{不可预知}——这正是它只能当兜底、不能当默认的原因。三条路线的分工可以背下来：自适应谓词管符号，区间管构造的粗筛，任意精度管最后的死角。

\begin{guideline}[title={选型时要问的三个精确性问题}]
现代几何库的谓词调用几乎都长成同一个形状：浮点快路径 $\to$ 误差界证明 $\to$ 精化或精确兜底，三层共享同一份输入。选型时真正要问的是：（1）谓词是否保证与实数算术一致（只信前两层不算，界必须严格到可证明）；（2）构造是否也精确——若构造只用 \Tp{double}，拓扑对了坐标仍可能微偏；（3）精确算术是否只花在符号上——若每次构造都展开成任意精度，规模上千就撑不住。三条都满足的 kernel 很少，\Lib{CGAL} 的分层正是把这三个问题拆开卖给你选。
\end{guideline}

三条路线不是竞争关系而是分层关系，表~\ref{tab:ch08-routes} 按机制与代价形态把它们对齐一次，后文 \Lib{CGAL} 的 kernel 档位就是这张表的具象。

\begin{table}[H]
\centering
\caption{三条鲁棒化路线的机制与代价对照}
\label{tab:ch08-routes}
\small
\begin{tabular}{@{}>{\raggedright\arraybackslash}p{0.14\textwidth}
  >{\raggedright\arraybackslash}p{0.24\textwidth}
  >{\raggedright\arraybackslash}p{0.22\textwidth}
  >{\raggedright\arraybackslash}p{0.20\textwidth}@{}}
\toprule
路线 & 机制 & 代价形态 & 代表实现 \\
\midrule
自适应精确谓词 & 浮点估计加可证明误差界，可疑时才逐级精化 &
  快路径与裸浮点同价，慢路径是常数次浮点加乘 &
  Shewchuk 原始代码、\Lib{CGAL} EPICK 谓词、\Lib{GEOS} 内核 \\
区间算术 & 每步运算传播一对向外舍入的端点 &
  依赖问题使宽度膨胀，判定常退化为``无法确定'' &
  \Lib{Boost.Interval}、\Lib{FIL}、\Lib{MPFI} \\
任意精度 & 坐标装入大整数或有理数域后完整求值 &
  位数与内存随运算链增长，事前不可预知 &
  \Lib{GMP}、\Lib{MPFR}、\Lib{CGAL} EPECK \\
\bottomrule
\end{tabular}
\end{table}

\section{\Lib{CGAL}：把谓词实现做成模板参数}

\Lib{CGAL} 给出的答案是把``用哪种算术''从算法里抽走。它的 Kernel 概念打包了一组谓词（\Tp{orient\_2\_kernel}、\Tp{circle\allowbreak\_side\allowbreak\_of\allowbreak\_circle\allowbreak\_2\allowbreak\_kernel} 等）、一组构造（外心、交点）与配套的数类型；\Tp{Triangulation\_2}、\Tp{Alpha\allowbreak\_shape\_2}、约束插点等算法只依赖这个概念，不知道底下是浮点还是有理数。换 kernel 就是改一个模板实参：

\begin{itemize}
\item \Tp{Simple\_cartesian}：谓词也用 \Tp{double} 行列式，最快，但上一节讲的浮点谓词风险原样继承，只配得上教学与玩具数据。
\item \Tp{Exact\_predicates\allowbreak\_inexact\allowbreak\_constructions\allowbreak\_kernel}（惯写 EPICK）：谓词走滤波加精确兜底、结果与实数算术一致；构造（外心、交点坐标）仍是 \Tp{double}。拓扑决策正确、几何坐标允许微偏，是三角剖分与网格生成的常用档。
\item \Tp{Exact\_predicates\allowbreak\_exact\allowbreak\_constructions\allowbreak\_kernel}（EPECK）：谓词与构造全精确。交点坐标要精确\emph{当谓词用}（决定第三条线从它的哪一边穿过）时这是刚需，否则叠加运算里的自相交判断还会翻车。
\end{itemize}

EPECK 的精确坐标靠表达式模板（lazy kernel）实现：构造不求值，只把运算组装成一棵表达式树，直到某次调用需要符号或需要输出时才展开求值。这避免了中间结果的组合爆炸——代价是表达式树深度随叠加次数增长，最终求值的位数可能急剧膨胀；工程惯例是用 EPECK 保证计算过程正确、输出前用 \Tp{to\_double} 快照成浮点。

Kernel 全部以模板交付、头文件为主，这带来两件事：谓词在剖分内循环里被完全内联，没有虚表跳转、没有跨库调用税；同时绝大多数 \Lib{CGAL} 组件不需要链接任何预编译二进制，升级是换头文件而非换 \Tp{.so}，代价是全额度的模板编译时间。这套命名在三维与任意维照抄一遍：\Tp{Exact\_predicates\_inexact\_constructions\_3\_kernel} 服务 \Tp{Triangulation\_3} 与四面体剖分，带 \Tp{\_d} 后缀的族把同一取舍推广到任意维度——记住``同一族 kernel 换个维度后缀''这一条，文档里的一半名字就不用再查。许可证是要认真核对的一档：\Lib{CGAL} 本体长期以 GPL 系与商业授权双许可分发；EPICK 通常不需要外部大数库，EPECK 一旦接入 \Lib{GMP} 与 \Lib{MPFR}，就引入它们各自的条款（\Lib{GMP} 走 GPL 与 LGPL 双许可谱系，\Lib{MPFR} 属 LGPL 谱系，具体版本与例外条款以你钉住的发布版 LICENSE 文件为准）。闭源产品静态链 EPECK 之前，先确认自己拿到了哪一档授权，而不是先确认编译器版本。

\section{\Lib{GEOS}：叠加运算的鲁棒性史与可重入 \CL{} 上下文}

\Lib{GEOS} 是\emph{纯 \CL{} 库}：算法血统来自 \Lib{JTS}（Java Topology Suite，一门 \Lib{Java} 面向对象库）的逐类移植，对外符号却是 \CL{} ABI——\Tp{GEOSBuffer\_r}、\Tp{GEOSUnion\_r}、\Tp{GEOSSimplify\_r}（后者是 Douglas-Peucker 抽稀），几何对象以不透明指针 \Tp{GEOSGeometry} 传递。序列化走 OGC 那一套：\Tp{GEOSWKTReader}/\Tp{WKTWriter} 读写人类可读的文本 Well-Known Text，\Tp{GEOSWKBReader}/\Tp{WKBWriter} 读写带字节序标记的二进制 WKB——数据库列里存的就是它。谓词族覆盖 OGC 简单要素模型的拓扑关系：\Tp{GEOSIntersects\_r}、\Tp{GEOSContains\_r}、\Tp{GEOSWithin\_r}，以及输出九维交点矩阵的一般式 \Tp{GEOSRelate\_r}。对一棵大空间对象做海量小查询，用 prepared geometry：\Tp{GEOSPrepare} 先把多边形装进内部树索引，之后 \Tp{GEOSPrepared\allowbreak\_contains\_r}、\Tp{GEOSPrepared\allowbreak\_intersects\_r} 每次查询都省掉重新建树，这是官方文档反复强调的复用点；反过来做``一大对多''的批量筛选则用 \Tp{GEOSSTRtree\_r} 系列建 Sort-Tile-Recursive 树，插入、查询、回调回收三步走。

叠加运算（union、intersection）是这一域鲁棒性问题的集中爆发地：两条折线近乎重合地相交，浮点求交给出两个坐标差 $10^{-13}$ 量级的``不同''交点，noding 结果里同一条边出现重复顶点，后续拓扑遍历读到零长度线段，JTS 与早期 \Lib{GEOS} 的经典报错 ``no next edge'' 型异常就是这么来的。修复思路不是把交点算得更准，而是\textbf{先统一精度模型、再重连拓扑}：snap-and-relax 路线把输入顶点吸附到粗栅格（grid snap），把吸附产生的微缝当作可松弛的容差重新连接边；JTS 1.14 与 \Lib{GEOS} 3.5 一带的版本逐步换成这套做法，近年又引入自 JTS 移植的新一代 overlay 内核（默认开关与行为随版本变化，上生产前查对应 release notes）。配套的显式工具是 \Tp{GEOSGeom\_setPrecision\_r}：按给定 grid size 做 snap-rounding 型化简，把``差一点的重复顶点''在进叠加之前先合并掉。

\Lib{GEOS} 的另一半工程价值在 API 形态：旧接口是一组操作全局状态的无后缀函数；线程安全的新接口全部带 \Tp{\_r} 后缀、首参是 \Tp{GEOSContextHandle\_t}（由 \Fn{GEOS\_init\_r} 创建；更早的发行里对应 \Fn{initGEOS\_r} 一代写法），错误消息经 \Tp{GEOSContext\allowbreak\_setErrorMessageHandler\_r} 挂回调。清单~\ref{lst:ch08-geos} 给出完整的最短闭环。

\begin{lstlisting}[style=cstyle,caption={\Lib{GEOS}：可重入上下文、WKT 读入、缓冲区与面积},label={lst:ch08-geos}]
#include <geos_c.h>
#include <stdio.h>

/* 错误回调：几何非法时 GEOS 只经这里报告 */
static void on_error(const char *msg, void *ud) {
  (void)ud;
  fprintf(stderr, "GEOS error: %s\n", msg);
}

/* 带圆角段数的缓冲：GEOSBufferParams 一族自 3.3 起提供 */
static GEOSGeometry *buffer_round(GEOSContextHandle_t h,
    const GEOSGeometry *g, double w) {
  GEOSBufferParams *bp = GEOSBufferParams_create_r(h);
  GEOSBufferParams_setQuadrantSegments_r(bp, 8);
  GEOSGeometry *out = GEOSBufferWithParams_r(h, g, bp, w);
  GEOSBufferParams_destroy_r(bp);
  return out;
}

int main(void) {
  /* 可重入上下文：每线程各持一个，不碰旧全局单例 */
  GEOSContextHandle_t h = GEOS_init_r();
  if (!h) { fprintf(stderr, "init failed\n"); return 1; }
  GEOSContext_setErrorMessageHandler_r(h, on_error, NULL);

  GEOSWKTReader_t *rd = GEOSWKTReader_create_r(h);
  GEOSGeometry *poly = GEOSWKTReader_read_r(h, rd,
      "POLYGON ((0 0, 4 0, 4 3, 0 3, 0 0))");
  if (!poly) {              /* 非法 WKT：返回 NULL 加回调 */
    fprintf(stderr, "WKT 解析失败\n");
    return 1;
  }

  GEOSGeometry *buf = buffer_round(h, poly, 0.5);
  double area = 0.0;
  if (buf) {
    if (GEOSArea_r(h, buf, &area) == 0) area = -1.0;
    printf("缓冲面积 = %.6f\n", area);
    GEOSGeom_destroy_r(h, buf);
  }
  GEOSGeom_destroy_r(h, poly);
  GEOSWKTReader_destroy_r(h, rd);
  GEOS_finish_r(h);         /* 上下文最后回收全部状态 */
  return 0;
}
\end{lstlisting}

注意清单~\ref{lst:ch08-geos} 里的资源纪律：几何对象、reader、缓冲参数各有自己的 \Tp{\_destroy\_r}，句柄本身最后 \Fn{GEOS\_finish\_r}；所有 \Tp{\_r} 函数返回 \Tp{int} 时 1 成 0 败、返回指针时 \Tp{NULL} 即失败，错误正文只能从回调里拿。这套``句柄进、整数出''的接口正是 \Lib{PostGIS} 能在 \CL{} 侧直接转发的原因——\Tp{ST\_Buffer}、\Tp{ST\_Intersection} 基本是对同名 \Lib{GEOS} 函数的包装，数据库进程与 GEOS 之间没有第二层语言。

\section{\Lib{PROJ}：从全局函数到上下文式坐标转换}

\Lib{PROJ} 的 \CL{} API 演进是一部并发与数据化的历史。4.x 时代只有一个投影概念：\Fn{pj\_init} 按字符串建投影、\Fn{pj\_fwd} 做正变换，错误走进程级 \Tp{errno}，参数表编译进库，全局状态让多线程直接出局。6.0 起了三件大事：第一，\emph{坐标操作}（coordinate operation）成为一等对象，一次转换是``基准面变换 $+$ 投影''的可分解链条而不是单个投影函数；第二，权威定义进了基于 SQLite 的坐标系数据库（EPSG、ESRI、NKG 等代码空间，随发行携带 \Tp{proj.db} 与基准面网格文件），网络模式下缺网格可按需下载（离线环境要谨慎处理这一步的失败）；第三，一切收进显式上下文 \Tp{PJ\_CONTEXT}，\Fn{proj\_create} 一族都吃句柄，线程各自建上下文即可安全并行。基准面变换是链条里最容易被低估的一环：WGS84 与地方基准（如日本 JTSMUPLAND、北美 NAD83 相关体系）之间的偏移靠 NTv2 一类网格文件插值，缺网格时 PROJ 会按参数退化处理——要么用无网格的近似偏移、要么报错，行为取决于版本与精度容忍设置。跨基准转换永远要显式验证拿到的是哪条操作：\Fn{proj\_get\_area\_of\_use} 查操作的适用范围，操作链本身可以用 \Fn{proj\_pjav} 之类接口拆出来检查（具体函数名与可用性以你钉住的版本文档为准）。

轴序是 6.0 最出名的一处坑：EPSG 对 4326 的权威定义是\textbf{纬度在前}，而可视化与 GeoJSON 世界的惯例是经度在前；把 $(116.4,\,39.9)$ 按权威轴序原样送进转换链，北京会被投到赤道以南的海面上——``坐标换了轴、地图上看还是合理的点位''是 GIS 事故统计里最难排查的一类，因为它不报错。补救有两个入口：建完操作立刻用 \Fn{proj\_normalize\_for\_visualization} 重排成视觉轴序（lon, lat），或在 CRS 串里显式写 \Tp{EPSG:4326!LONGLAT} 后缀换轴；两条都要在建操作之后马上做，因为轴序藏在链条的第一环。清单~\ref{lst:ch08-proj} 走完整流程并演示错误检查。

\begin{lstlisting}[style=cstyle,caption={\Lib{PROJ}：上下文创建、轴序归一化与转换错误处理},label={lst:ch08-proj}]
#include <proj.h>
#include <stdio.h>
#include <stddef.h>

int main(void) {
  /* 显式上下文取代 4.x 的全局状态，线程各自持有 */
  PJ_CONTEXT *cc = proj_context_create();

  /* CRS 到 CRS 的操作：基准面变换加投影的完整链条 */
  PJ *op = proj_create_crs_to_crs(cc, "EPSG:4326",
                                  "EPSG:3857", NULL);
  if (!op) {
    fprintf(stderr, "建操作失败: %s\n",
            proj_context_get_string(cc));
    return 1;
  }
  /* EPSG:4326 权威轴序 lat,lon；重排为视觉轴序 lon,lat */
  PJ *vis = proj_normalize_for_visualization(cc, op);
  if (!vis) vis = op;          /* 失败则退回原始操作 */

  /* 单点正变换：返回 PJ_COORD union，失败置 HUGE_VAL */
  PJ_COORD q = proj_trans(vis, PJ_FWD,
                          proj_coord(116.4, 39.9, 0, 0));
  if (proj_errno(vis) != 0) {  /* 错误码要紧接着查 */
    fprintf(stderr, "转换失败: %s\n",
            proj_errno_string(proj_context_get_error(cc)));
    return 1;
  }
  printf("北京 3857: %.2f %.2f\n", q.xy.x, q.xy.y);

  /* 批量转换：按 stride 吃 SoA 数组，z、t 两路留空 */
  double lon[2] = {116.4, 121.5}, lat[2] = {39.9, 31.2};
  double x[2], y[2];
  size_t nn = proj_trans_generic(vis, PJ_FWD,
      lon, sizeof(double), 2, lat, sizeof(double), 2,
      NULL, 0, 0, NULL, 0, 0,
      x, sizeof(double), 2, y, sizeof(double), 2,
      NULL, 0, 0, NULL, 0, 0);
  if (nn == 2)
    printf("上海 3857: %.2f %.2f\n", x[1], y[1]);

  proj_destroy(vis); proj_destroy(op);
  proj_context_destroy(cc);
  return 0;
}
\end{lstlisting}

单个点走 \Fn{proj\_trans}：它按 \Tp{PJ\_DIRECTION}（\Tp{PJ\_FWD} 正变换、\Tp{PJ\_INVERSE} 反变换）返回 \Tp{PJ\_COORD}——x、y、z、t 四路 union；失败时各分量被置成 \Tp{HUGE\_VAL} 一类哨兵值，随后用 \Fn{proj\_errno}（上下文级错误另走 \Fn{proj\_context\_get\_error}）区分``转换失败''与``结果恰好是无穷''。批量与多维靠 \Fn{proj\_trans\_generic}：输入给 x、y、z、t 四路各一块数组、一个字节 stride 与元素数，输出同样四路、可空，一次调用吃掉整块坐标流——这就是 \CL{} 版 API 对 SoA 布局的原生支持；\PY{} 门面 \Lib{pyproj} 一次转十万个点走的也是它。\Lib{GDAL} 在 \CPP{} 侧用它，\Lib{PostGIS} 的 \Tp{ST\_Transform} 在 \CL{} 侧用它，QGIS 与 \Lib{GRASS} 经 \Lib{GDAL} 间接用它：一个 \Tp{.so} 挂住整条 GIS 工具链。

\begin{warning}[title={谓词布尔化与共享上下文}]
第一，把浮点算出来的 \Fn{orient2d} 当布尔用是这一域最贵的习惯：退化构型下它的符号是舍入噪声，而下游拿它做拓扑决策，一次翻号换来一份自交网格。要么选带滤波的谓词（\Lib{CGAL} EPICK 起步、\Lib{GEOS} 用 3.5 后的内核），要么自己给行列式加可证明的误差界（Shewchuk 论文给了可照抄的常数），没有``先上线再修''这一档。第二，\Lib{PROJ} 与 \Lib{GEOS} 的新 API 都是``上下文进句柄''：\Tp{PJ\_CONTEXT} 与 \Tp{GEOSContextHandle\_t} 都\textbf{不是}可以随便全局共享的对象——同一上下文跨线程并发使用不受支持，旧式无 \Tp{\_r} 后缀的 \Lib{GEOS} 函数还额外共享进程级错误状态。正确姿势是一线程一上下文或一任务一上下文；\Lib{PROJ} 里 \Tp{PJ} 操作对象被多线程\emph{只读}复用的边界，以当前版本文档为准再依赖。
\end{warning}

\section{为什么这一域必须是 \CL/\CPP{}}

把本章的论证收拢成三条机制。其一，\textbf{谓词是逐点调用的}：一次十万点的三角剖分要做百万级 \Tp{orient2d}，解释型语言每次调用的分派税与装箱税直接乘在谓词头上；只有模板与内联能把整套展开算术压进接近一次普通函数调用的成本，这是精确内核清一色 \CPP{} 头部实现（\Lib{CGAL} kernel、\Lib{GEOS} 背后 \Lib{JTS} 的移植谱系）的共同理由。其二，\textbf{精确算术的复杂度必须编译期摊销}：自适应谓词的三层滤网、lazy kernel 的表达式树，本质都是把``花多少算术''的决定搬到类型与特化里，运行时只走已经定好的路径；而任意精度底座（\Lib{GMP}、\Lib{MPFR}）以 \CL{} 符号暴露句柄（\Tp{mpz\_t}、\Tp{mpfr\_t}），谁都能链。其三，\textbf{GIS 生态的扩展点全是 C 符号}：PostgreSQL 扩展必须提供 \Tp{PG\_FUNCTION\_INFO()} 约定的 \CL{} 入口（\Tp{Datum} 与内存分配走服务端 API），PostGIS 挂上去的同时就把 \Lib{GEOS} 与 \Lib{PROJ} 钉死在 \CL{} ABI 上；QGIS 的 provider、\Lib{GRASS} 的模块同理。\PY{} 在这一域的真实角色是\emph{门面}：\Lib{shapely} 经 Cython 调 \Lib{GEOS}，\Lib{pyproj} 经 C 扩展调 \Lib{PROJ}，\Lib{rasterio} 调 \Lib{GDAL}——门面越薄越好，谓词循环从不跨过这门面。表~\ref{tab:ch08-libs} 把本章五个库与下一登场的 \Lib{libigl} 放在一起对照。

\begin{table}[H]
\centering
\caption{几何域代表库的形态对照（许可证类别以各仓库当前 LICENSE 为准）}
\label{tab:ch08-libs}
\small
\begin{tabular}{@{}>{\raggedright\arraybackslash}p{0.12\textwidth}
  >{\raggedright\arraybackslash}p{0.13\textwidth}
  >{\raggedright\arraybackslash}p{0.24\textwidth}
  >{\raggedright\arraybackslash}p{0.14\textwidth}
  >{\raggedright\arraybackslash}p{0.18\textwidth}@{}}
\toprule
库 & 语言层 & 接口形态 & 许可证类别 & 可否嵌入进程 \\
\midrule
\Lib{CGAL} & \CPP{} 模板 & 头文件为主，kernel 编译期内联 &
  GPL 系加商业双许可 & 源码级嵌入，无运行时分发 \\
\Lib{GEOS} & \CL{} 符号 & \Tp{\_r} 可重入 API、WKT 与 WKB 序列化 &
  LGPL 系 & 动态库直链，PostGIS 即如此 \\
\Lib{JTS} & \Lang{Java} & jar、面向对象、异常报错 &
  LGPL 系 & 不进 \CL{} 进程，靠逐语言移植 \\
\Lib{PostGIS} & \CL{} 扩展 & PostgreSQL 服务端函数、SQL 门面 &
  GPL 系 & 只在数据库进程内运行 \\
\Lib{PROJ} & \CL{} 为主 & 上下文式 API、SQLite 坐标系数据库 &
  MIT 风格宽松 & 动态库直链，GDAL 与 pyproj 共用 \\
\Lib{libigl} & \CPP{} 头文件 & Eigen 矩阵进出的函数集 &
  MPL 系 & 编译期嵌入，随项目分发源码 \\
\bottomrule
\end{tabular}
\end{table}

\begin{bestpractice}[title={鲁棒性的三档投入}]
（1）谓词档：产品路径上的任何剖分、叠加、包含判定，起步就是精确谓词（EPICK、\Lib{GEOS} 3.5 以后的内核），不接受裸浮点行列式。（2）构造档：只有当输出坐标还要\emph{参与后续几何判定}（例如叠加里的交点再喂给 noding）才升级 EPECK 或 snap-rounding，否则 EPICK 加最后快照更快。（3）集成档：把几何当进程间数据就用序列化（WKB、GeoJSON），把几何当函数参数就用句柄 API（\Tp{\_r} 后缀、\Tp{PJ\_CONTEXT}），并且一线程一句柄；跨语言门面只包最外一层，谓词循环永远留在 \CL/\CPP{} 侧。
\end{bestpractice}

下一章进入几何对象的本体——网格与点云怎么摆进内存、射线怎么打上去，那把本章的谓词问题换成了布局问题。
% ====================================================================

% ====================================================================
\chapter{网格、点云与光线求交}

上一章管``判定怎么不坏''，本章管对象本身：网格与点云怎样摆进内存，射线怎样打上去。两个问题共享同一个答案——布局即语义：同一只模型兔的顶点存在哪、按什么顺序存，决定了邻接查询、拉普拉斯算子、点云配准与光线求交各自要花什么代价。本章以三个库为骨架：\Lib{libigl} 把几何化成 \Lib{Eigen} 矩阵上的线性代数，\Lib{Open3D} 用同一份 \CPP{} 内核对外摆出两套门面，\Lib{Embree} 把 CPU 光线求交推到工程极限；末尾谈内核之外的格式约定——OBJ、PLY、glTF 的解析层与单位、手性。三者都是 \CPP{} 库，但形态差别极大，这个差别正是全书主线在这一域的落点：为什么求交这种``分枝加访存''的活儿在解释型语言里根本不成立。

\section{几何数据的四种摆法与它们的内存代价}

\textbf{其一：顶点-面索引}，即所谓 triangle soup 加索引：$n\times3$ 的浮点顶点数组配 $m\times3$ 的无符号整数面数组，OBJ 与 PLY 文件在磁盘上就是这个形状。它最紧凑（每顶点 12 字节、每面 12 字节量级）、按面顺序处理时对缓存最友好，\Lib{libigl} 与 \Lib{Embree} 的入口数据都是它。代价是邻接信息缺失：问``顶点 $i$ 的 1-ring 有哪些面''要先遍历全部面，所以一切像样的算法都先用它现算一张边表或邻接表（\Tp{igl::adjacency}、\Tp{igl::adjacency\_list}）。

\textbf{其二：半边结构}。每条无向边拆成两条有向半边，各存对偶（twin）、前驱、后继、起点与所属面；\Tp{CGAL::Polyhedron\_3}、\Tp{CGAL::Surface\_mesh} 与 \Lib{Open3D} 的 \Tp{HalfEdgeTriangleMesh} 都以此为内核。邻接与环遍历变成严格的 $O(\deg)$，欧拉操作（挤边、翻边、劈分）写起来指哪打哪——网格简化和特征检测离不开它。代价同样明确：每个半边三四个指针或索引，总内存常比索引网格高出数倍，而且遍历 1-ring 是指针追逐，预取器帮不上忙。索引结构把代价付在预处理，半边把代价付在内存与局部性——选谁取决于你要按面扫还是按邻接走。

\textbf{其三：点云加 KDTree}。$n\times3$ 浮点坐标（外加法向、颜色、强度等每点属性），配一棵按坐标轴循环取中位数切分的二叉空间划分树：\Lib{Open3D} 的 \Tp{KDTreeFlann}、\Lib{CGAL} 的 \Tp{Search\_tree}、以及单头文件的 nanoflann。近邻与半径查询期望 $O(\log n)$，建树 $O(n\log n)$；维度升高后剪枝失效（高维退化），三维扫描数据不受此累。一条工程纪律常被违反：KDTree 缓存的是坐标的划分，点集或刚体变换一变就得重建或重问——迭代最近点这类算法每步换姿态，树与姿态的一致性要自己管。每点属性建议按通道拆成并列数组：查询命中的是索引，回取的属性越紧凑，随机回表越便宜。

\textbf{其四：体素与八叉树}。当数据本质是空间场（占据、密度、符号距离）而不是表面，均匀体素栅格可按三指数直接寻址、天然向量化，但内存随边长立方增长；八叉树用层次划分换稀疏性，代价又回到指针追逐与缓存缺失。\Lib{OpenVDB} 一系的稀疏层次结构、\Lib{Embree} 的体素遮挡查询、以及下一节要讲的截断符号距离体积，都是这一族的工程化。

四种摆法的账目集中列在表~\ref{tab:ch09-layouts}。还要补一句它们\emph{不是互斥}的：真实管线里同一份几何常同时挂两三种视图——索引数组当存储主体、半边当处理阶段的临时结构、BVH 当查询加速器，三者间的转换层（重建邻接、重建树）就是流水线的节拍器。转换层也是精度层：从 \Tp{double} 索引网格生成 \Tp{float} BVH 的那一步，正是上一章谓词最容易翻车的交接点——盒测试用浮点无所谓，判定用的坐标却值得留双份。

\begin{table}[H]
\centering
\caption{四种几何数据摆法的代价对照}
\label{tab:ch09-layouts}
\small
\begin{tabular}{@{}>{\raggedright\arraybackslash}p{0.14\textwidth}
  >{\raggedright\arraybackslash}p{0.22\textwidth}
  >{\raggedright\arraybackslash}p{0.24\textwidth}
  >{\raggedright\arraybackslash}p{0.20\textwidth}@{}}
\toprule
摆法 & 邻接与查询 & 内存与局部性 & 适用场景 \\
\midrule
顶点-面索引 & 无邻接，需现算边表或邻接表 &
  最紧凑，按面顺序扫描缓存友好 &
  文件交换、着色器上传、光线求交主体 \\
半边结构 & 环遍历严格 $O(\deg)$，欧拉操作齐备 &
  每半边多个指针或索引，指针追逐伤预取 &
  网格简化、特征检测、拓扑修改 \\
点云加 KDTree & 近邻与半径查询期望 $O(\log n)$ &
  树与坐标两份存储，几何一变树即失效 &
  扫描配准、统计滤波、离群点剔除 \\
体素与八叉树 & 场值按坐标直接寻址或按层下降 &
  均匀体素内存立方增长，八叉树换稀疏 &
  TSDF 重建、占据栅格、大场景遮挡 \\
\bottomrule
\end{tabular}
\end{table}

还要补一句\textbf{数值精度}的选择：顶点数组用 \Tp{float32} 还是 \Tp{double64}，是布局之外的第一个架构决定。渲染与光线求交通常用 \Tp{float}：内存减半、SIMD 宽度翻倍，求交误差对像素级结果不可见；但同一份坐标若还要进上一章的精确谓词或地理配准，\Tp{float} 的大场景坐标会先丢几位有效数字再进入判定，退化构型更密、翻车更多。工程惯例是``双精度存家底、单精度打热点''：源数据与变换保持 \Tp{double}，进 BVH 与着色循环前一次性降到 \Tp{float}，并且把场景平移到原点附近再降——平移量不先消掉，远处的细节会整片糊在同一格浮点上。

\subsection{AoS 与 SoA：布局决定光线能不能成批打}

同一批 $4$ 条光线，两种摆法。数组的结构（AoS）把每条光线的 $(o_x,o_y,o_z)$、$(d_x,d_y,d_z)$ 与 $t_{far}$ 存成连续结构体：一次 SIMD 寄存器加载只拿到一条光线的半个分量集，四条光线的同一分量散在四个 cacheline 里，向量化只能靠低效的横向收集。结构的数组（SoA）把同一分量摆成 4 元连续数组：一条 SIMD 寄存器正好装下四条光线的 $o_x$，三角形面积、距离检验全部按通道并行。\Lib{Embree} 的 \Tp{RTCRay} 对用户是 AoS 结构体，但 \Fn{rtcIntersect4} 在核内部先做 4$\times$7 转置再进 SoA 遍历——转置成本付一次，SIMD 收益吃整棵树。\Lib{libigl} 用 \Lib{Eigen} 稠密矩阵按行存顶点，相当于把网格本身摆成对批量数值友好的布局；而半边结构的每个半边各带若干指针，是对照组：那是为\emph{遍历图}优化的布局，不是为\emph{批量数值}优化的。

半边结构把每条边显性化，索引结构把每面隐式化——这个取舍会一路传导到算法可写性：能做``翻一条边''的地方必须有半边或等价的临时邻接，能做``扫一遍求和''的地方索引数组就够。写网格算法时先问自己在动图还是在动数组，答案决定选哪层库。

\section{\Lib{libigl}：以 \Lib{Eigen} 为数据层的研究原型底座}

\Lib{libigl} 是一坨头文件模板：没有要链的 \Tp{.so}，每个算法一个头（\Tp{igl/cotmatrix.h} 一类），数据的通用货币是 \Lib{Eigen} 的 \Tp{MatrixXd}（double 稠密矩阵）、\Tp{MatrixXi} 与 \Tp{SparseMatrix<double>}。它的前身是苏黎世联邦工大几何处理组（Alexa、Sorkine 一脉）的论文配套代码，所以``函数名等于论文记号''的基因一直保留着。它的接口哲学可以概括成一条：\textbf{一个数学公式对应一个函数}。网格读入是 \Tp{igl::readOBJ(path, V, F)}——返回 \Tp{bool} 而不是抛异常，因为文件是外部世界；同族的小工具按头文件铺开——\Fn{igl::per\_vertex\_normals} 出单位法向场、\Fn{igl::doublearea} 逐面给出两倍面积（顺手就能把退化三角形筛出来）、\Fn{igl::boundary\_loop} 拉出边界环，函数名就是论文里的记号。离散拉普拉斯是 \Fn{igl::cotmatrix}：输出稀疏矩阵 $L$，非对角元 $L_{ij}=-\tfrac12(\cot\alpha_{ij}+\cot\beta_{ij})$，对角元取负行和，每行和恰为零——这正是``余切权重''离散 Laplace-Beltrami 的矩阵形态，\Fn{igl::cotin} 再配质量矩阵给出 $M^{-1}L$ 形式的算子，供曲率流与形变直接迭代。

余切权重的几何来历也顺带把上一章的阴影带进来：$\cot$ 在退化三角形上爆炸，在钝角三角形上取负——前者是谓词与数值精度的问题，后者是离散算子的合法性质（负权是保角逼近的组成部分，不是 bug）。用 \Tp{double} 坐标算出的 $L$ 对病态网格可能给出不对称的行和，这在谱几何处理（\Fn{igl::eig} 一类的谱分解）里会放大成可见的伪影。\Lib{libigl} 的定位因此很清楚：假设输入网格``基本健康''，把算力全留给下游方程。

形变是招牌应用：\Fn{igl::arap} 实现 Sorkine-Alexa 一脉的局部/全局旋转不变坐标形变——\Fn{igl::arap\_precompute} 把稀疏系统因子化一次、缓存预计算量，之后 \Fn{igl::arap\_solve} 每步只解各顶点的局部刚体旋转再整体回代，交互拖拽控制点时每次迭代都复用那次因子化。射线查询也有一席：\Tp{igl::point\_mesh\_intersection} 给一批\emph{无限}射线与三角网格求首个交，底层是 \Tp{igl::BVH}（中位数划分的 AABB 层次树）——功能对标 \Lib{Embree}，工程（ISA 分派、SoA 批量、无分配遍历）差一个量级，下一节正面讲这套工程。许可证上主干是 MPL 系（\Lib{Eigen} 同为 MPL 系），\Tp{igl::copyleft::} 命名空间下的部件按更严格的 GPL 系分发、依赖 GPL 的可选后端也用这层隔离——商用前按目录核对，不要按``header-only 就无所谓''处理。

\begin{lstlisting}[style=cppstyle,caption={\Lib{libigl}：读 OBJ、组装余切拉普拉斯并检查负权与面积},label={lst:ch09-igl}]
#include <igl/readOBJ.h>
#include <igl/cotmatrix.h>
#include <Eigen/Sparse>
#include <iostream>

int main() {
  Eigen::MatrixXd V;          // 顶点 n x 3，按行存
  Eigen::MatrixXi F;          // 面   m x 3，零基索引
  if (!igl::readOBJ("bunny.obj", V, F)) {
    std::cerr << "读入 bunny.obj 失败\n";
    return 1;
  }

  Eigen::SparseMatrix<double> L;   // 余切拉普拉斯
  igl::cotmatrix(V, F, L);         // 行和恒为 0
  long tot = 0, neg = 0;           // 统计钝角三角形产生的负余切
  for (int k = 0; k < L.outerSize(); ++k)
    for (Eigen::SparseMatrix<double>::InnerIterator
           it(L, k); it; ++it)
      if (it.row() != it.col()) {  // 非对角元 = -w_ij
        ++tot;
        if (it.value() > 0.0) ++neg;   // w_ij 为负即钝角
      }

  double a2 = 0.0;                 // 面积：叉积模长之和的一半
  for (int i = 0; i < F.rows(); ++i) {
    Eigen::RowVector3d p = V.row(F(i, 0));
    Eigen::RowVector3d b = V.row(F(i, 1)) - p;
    Eigen::RowVector3d c = V.row(F(i, 2)) - p;
    a2 += b.cross(c).norm();
  }

  std::cout << "顶点 " << V.rows() << " 面 " << F.rows()
            << " 面积 " << 0.5 * a2 << "\n";
  std::cout << "非对角 " << tot << "，其中负权 " << neg << "\n";
  return 0;
}
\end{lstlisting}

清单~\ref{lst:ch09-igl} 这四十行能当半个体检报告：面积用来卡单位（下一节的坑），负权比例用来预警下游谱方法的质量。研究原型阶段选 \Lib{libigl} 的理由到此可以说尽：依赖图小、跨平台一条编译命令、数据是随手可拼的 \Lib{Eigen} 矩阵；弱点也照实说：默认单线程、没有场景与资源管理、性能上限受遍历工程拖累。

把它和 \Lib{CGAL} 的分工再钉一颗钉：同一个``给我这个域的拉普拉斯''的问题，\Lib{CGAL} 从半边内核出发、保证拓扑与谓词正确，\Lib{libigl} 从索引数组出发、保证你二十分钟能跑通公式。原型验证用后者，产品里怕退化网格咬人时前者的档位救场——两把梯子常常同架在一个项目上。

\section{\Lib{Open3D}：同一份 \CPP{} 内核，两套门面}

\Lib{Open3D} 是另一个极端：完整的三维数据处理管线——输入输出、空间索引、配准、重建、渲染、传感器仿真——编译成一份 \CPP{} 共享库，再经 pybind11 把同名对象树递到 \PY{} 手里，官方直接发 pip 轮子。两扇门后是同一内核：\PY{} 里 \Tp{o3d.geometry.TriangleMesh} 的方法调用穿 pybind11 落到 \CPP{} 侧 \Tp{open3d::geometry::TriangleMesh} 的成员函数，数据在边界上按引用或零拷贝语义交接，转换税只付在调用点、不付在循环里。``同一份实现、两张脸''是全书讲门面模式的样本案例——连第三方引擎都是它的原料：\Tp{RaycastingScene} 内部就嵌着 \Lib{Embree}（Apache 系），MIT 系的 \Lib{Open3D} 把它们组装成自家类，许可证在组装链上的传递关系要在发布物料里逐层写清。

分发形态也是两套门面的一部分：\CPP{} 侧走 CMake（\Tp{find\_package(Open3D)} 或包管理器的端口），\PY{} 侧是预编译 wheel——同一份内核的两种交付物出现版本错位时，报错方式完全不同（前者是符号找不到，后者是运行时异常），排障第一步是确认你握的是哪一层的对象。连可视化与离屏渲染也按同一双门面发：\PY{} 几行画出点云与网格、\CPP{} 侧同套场景对象接进自家应用，渲染后端换实现不换接口——这进一步坐实了``两扇门后是同一内核''的说法。

API 有新旧两代并存。传统 geometry 模块（\Tp{PointCloud}、\Tp{TriangleMesh}、\Tp{KDTreeFlann}）驻留 CPU 内存，存量扫描数据管线大多长在上面；新一代带 T 前缀的 API 把一切装进 \Tp{core::Tensor}，用设备字符串 \Tp{CPU:0}、\Tp{CUDA:0} 决定数据与计算的家，GPU 加速对上层只是换设备前缀。两代覆盖面此消彼长、共存是近年文档的常态，选型时按\emph{功能}而不是按\emph{命名风格}核对可用性。

点云配准是它的看家管线：迭代最近点（ICP）。循环体三段——用 \Tp{KDTreeFlann::\allowbreak SearchKNN} 找源点到目标点云的最近邻作对应，按 \Tp{Transformation\allowbreak Estimation\allowbreak PointToPlane}（点对面，收敛快且压平切向漂移）或点对点建最小二乘目标，解一个 4$\times$4 刚体位姿。对应质量决定一切，所以粗配准先用特征：FPFH 描述子配 \Tp{Global\allowbreak Registration\allowbreak ICP} 或 RANSAC，把姿态先送进 ICP 的收敛半径。多站扫描的闭合误差则进位姿图（\Tp{PoseGraph}）做全局优化——这些都是 \CPP{} 内核方法，\PY{} 侧同名调用。

体素表示的活样本是 RGBD 融合：\Tp{RGBDImage} 配 \Tp{PinholeCamera} 从深度图反投三维点，\Tp{ScalableTSDFVolume} 把每帧沿视线方向在截断带宽内写截断符号距离场——体素存``离表面多远''，多帧观测做加权平均，表面就在零等值面处；最后 \Fn{ExtractTriangleMesh} 出封闭网格，把散点变水密模型的正是这一族（\Lib{PCL} 与 \Lib{OpenVDB} 里是同路的不同实现）。光线查询在这里预演：\Tp{RaycastingScene} 直接内置 \Lib{Embree}，\Tp{rays\_o}、\Tp{rays\_d} 两批矩阵进去、首个命中出来——它把本章下半部的引擎做成了自家类库的一个成员。许可证 MIT 系，两套餐具都不贵。

\section{\Lib{Embree}：CPU 光线求交的极限工程}

\Lib{Embree}（Intel 出品，Apache 系许可）是把射线求交当成内核库来雕的：主接口是 \CPP{} 头文件（\Tp{embree4/rtcore.h}，概念只有 Device、Scene、Geometry、Buffer、Ray），面向原生 \CPP{} 宿主直接链接；3.x 时代另备有一组句柄式的 C 绑定给纯 \CL{} 宿主与 FFI 封装用，4.x 是否保留该接口面以其发布说明为准。包名与 CMake 目标名随发行渠道而异（源树里的实现目标常见 \Tp{embree\_cpp} 一类命名，打包发行则多叫 \Tp{embree4}），接口面以你链上的版本为准。场景是显式的层次对象：三角形、四边形、曲线、细分面、几何实例（引用另一棵 BVH）各是一种 \Tp{RTCGeometry}，自定义求交的回调挂接点让``几何''与``怎么求交''解耦——遍历器归库，求交测试归你。

三个工程决定值得单独说。\textbf{其一，ISA 运行时分派}：编译产物里同时躺着 SSE2、AVX、AVX2、AVX-512 几套遍历内核，\Fn{rtcNewDevice} 时探测 CPU 特性选最快的一套；\Tp{isa=sse2} 可以强制保底档用于调试与老机器。一份预编译二进制走到哪都跑，这是``库作为发行物''的基本要求，解释型语言扩展模块反而难做得这么干净。\textbf{其二，BVH 构建档位}：\Fn{rtcSetSceneBuildQuality} 与逐几何的 \Fn{rtcSetGeometryBuildQuality} 在低档（快、宽）与高档（慢、紧）之间调，布局按 ISA 取 4 路或 8 路分支节点；动态场景可以只 refit 变化的几何，其余树保持不动。\textbf{其三，求交接口的光线形态}：\Tp{RTCRayHit} 的 ray 段是 $o,d$ 加分量各一的 \Tp{tnear}、\Tp{tfar}、掩码与标志，hit 段带 $t$、几何与图元 ID、重心坐标 $(u,v)$ 与几何法向；单条走 \Fn{rtcIntersect1}，批量走 \Fn{rtcIntersect4}——四条光线一组、一个 32 位有效掩码，核内转置成 SoA 后按通道并行，同屏同帧的相干光线收益最大；命中后把 \Tp{tfar} 置 $-\infty$ 是``交完退场''的惯用法，未命中保留原 \Tp{tfar}，阴影射线因此可以整批打。只要存在性不要交点时改走 \Fn{rtcOccluded1}，遍历到首个命中即终止；几何实例（instance）让同一棵子树挂进多处、每实例带自己的变换矩阵，``复用 $+$ 每实例一行''的场景组织靠它。数据侧的纪律也值得注意：每个几何缓冲的描述符显式给出格式、基址、stride 与元素数——这份元数据就是 \CL{} 式 SoA 合同的实体，库按它把任意宿主内存接进 SIMD 内核，不需要你先拷成它的私有布局。

场景的更新路径同样被参数化了：静态大场景一次 \Fn{rtcCommitScene} 建高档 BVH、之后只读；动得频繁的物体单独挂低档构建质量或 \Tp{RTC\allowbreak\_BUILD\allowbreak\_QUALITY\allowbreak\_NONE} 的直接图元列表，提交时只重建该几何的叶层；再进一步，光线可以带命中与遮挡两个回调上下文（\Tp{RTCRayQueryContext}），把``这条射线属于哪一帧、哪个通道''传给自己的回调。把这些档位用对，\Lib{Embree} 在动态场景里也能维持增量提交，而不必整树重建。

并发这一头也有明确合同：已提交的场景是只读快照，多线程同时对其求交安全；改几何必须等遍历排空或开新场景。建 BVH 自身是多线程的，设备配置里还能把宿主应用的任务系统注册进去，让库的建树任务跑在你的线程池上而不是另起炉灶（这组回调的形态在 4.x 有调整，以当前版本文档为准）。

\subsection{为什么求交只在 \CL/\CPP{} 里成立}

射线遍历的热点循环是\textbf{分枝与访存的混合}：每层先做分支无关（branchless）的射线-盒测试、压栈靠哪边更近先走哪边，随后从池里随机取子节点、再随取三角形数据。三件事在别的语言层几乎不可获得：其一，热点循环里\emph{零分配}——遍历栈预分配、命中记录按固定布局写，解释器与 GC 语言连``这一轮不分配内存''都承诺不了；其二，4$\times$7 的 SoA 转置、通道掩码与分支无关栈要\emph{直接}对着寄存器写，\Lib{Embree} 的求交核用 ISPC 生成，接口层仍是 \CPP{}；其三，每射线每节点的决定以纳秒计，任何动态分派（装箱参数、虚调用、跨 FFI 边界）乘上树深就吞掉全部优化。于是``把求交做成库''只剩一条路：内核 \CPP{}/ISPC、接口 \CL{} 符号，谁都能嵌——\Lib{Open3D} 的 \Tp{RaycastingScene} 与渲染器的 CPU 可见性通道走的都是这条路。

\begin{lstlisting}[style=cppstyle,caption={\Lib{Embree}：建三角网格场景、提交 BVH 并批量求交},label={lst:ch09-embree}]
#include <embree4/rtcore.h>
#include <cstdio>

int main() {
  /* 保底 ISA 为 SSE2；新 CPU 上运行时仍自动升档 */
  RTCDevice dev = rtcNewDevice("isa=sse2");
  RTCScene sc = rtcNewScene(dev);

  RTCGeometry g = rtcNewGeometry(dev,
      RTC_GEOMETRY_TYPE_TRIANGLE);
  float *vb = (float*)rtcSetNewGeometryBuffer(g,
      RTC_BUFFER_TYPE_VERTEX, 0, RTC_FORMAT_FLOAT3,
      3 * sizeof(float), 3);        // 3 个顶点
  float *nb = (float*)rtcSetNewGeometryBuffer(g,
      RTC_BUFFER_TYPE_NORMAL, 0, RTC_FORMAT_FLOAT3,
      3 * sizeof(float), 3);
  unsigned *ib = (unsigned*)rtcSetNewGeometryBuffer(g,
      RTC_BUFFER_TYPE_INDEX, 0, RTC_FORMAT_UINT32,
      sizeof(unsigned), 1);         // 1 个三角形
  const float V[9] = {0,0,0, 2,0,0, 0,2,0};
  const float N[9] = {0,0,1, 0,0,1, 0,0,1};
  const unsigned I[3] = {0, 1, 2};
  for (int i = 0; i < 9; ++i) { vb[i] = V[i]; nb[i] = N[i]; }
  for (int i = 0; i < 3; ++i) ib[i] = I[i];

  rtcSetGeometryBuildQuality(g, RTC_BUILD_QUALITY_HIGH);
  rtcCommitGeometry(g);
  rtcAttachGeometry(sc, g);
  rtcReleaseGeometry(g);            // 场景接管所有权
  rtcCommitScene(sc);               // BVH 在这里生成

  const int n = 64;                 // 垂直上射的批量光线
  RTCRayHit hits[64];
  for (int i = 0; i < n; ++i) {
    hits[i].ray.org_x = 1.0f - 0.01f * i;  // 起点正对三角形下方
    hits[i].ray.org_y = 0.5f;
    hits[i].ray.org_z = -1.0f;
    hits[i].ray.dir_x = hits[i].ray.dir_y = 0.0f;
    hits[i].ray.dir_z = 1.0f;
    hits[i].ray.tnear = 0.0f;
    hits[i].ray.tfar = 1e30f;
    hits[i].ray.mask = 0xFF;
    hits[i].ray.flags = 0;
  }
  for (int i = 0; i < n; ++i)
    rtcIntersect1(sc, &hits[i], nullptr);
  int nhit = 0;
  for (int i = 0; i < n; ++i)
    if (hits[i].hit.geomID != RTC_INVALID_GEOMETRY_ID)
      ++nhit;
  std::printf("命中 %d / %d，primID %u，u %.3f v %.3f\n",
              nhit, n, hits[0].hit.primID,
              hits[0].hit.u, hits[0].hit.v);

  rtcReleaseScene(sc);
  rtcReleaseDevice(dev);
  return 0;
}
\end{lstlisting}

清单~\ref{lst:ch09-embree} 里 \Tp{u}、\Tp{v} 重心坐标连同 \Tp{primID} 一起回来，插值属性（法向、纹理、材质 ID）因此不需要二次查询。回调这条缝再展开一点：用户几何要求你提供 bounds 与 intersect 两个函数，遍历器算到哪张盒子、就把光线连同一份你自定的上下文指针递回来——球、圆柱、隐式曲面都能这样进同一棵树，代价是你的求交函数必须扛住 SIMD 批入口（一次进来的是 4 条或 8 条的分量数组，不是单条结构体）；三角形回调则允许在不改几何的前提下改判定，比如按材质掩码否决一次命中。遍历归库、求交归你的分工就体现在这里。与 GPU 的分工也清楚了：光栅化把``画三角形''外包给固定吞吐的管线，硬件光追再把 BVH 遍历与求交外包给专门的 RT 单元——吞吐更高，但受显存与生态约束；\Lib{Embree} 的生态位是 CPU 侧与混合渲染：可见性、阴影、环境光遮蔽、降噪这类``射线量中等、场景超显存、要精确 ID 语义''的通道，以及把场景序列化成紧凑二进制、供快速加载与向 GPU 侧加速结构搬运的动机。批渲染全交给 GPU 时 \Lib{Embree} 退场，混合管线里它是 CPU 那一半的主角。

位图再大也只是一次上传，网格再复杂也只是两段数组——把``数据''与``对数据的决定''分开存放，是本章所有布局讨论的尽头。

\section{数据交换约定：解析层与内核无关}

格式层与内核层是两桩生意。\textbf{OBJ} 是行式文本：\Tp{v}、\Tp{vn}、\Tp{vt} 三族记录加 \Tp{f} 面记录，索引从 1 起（还允许负数相对索引），材质走另一个 \Tp{.mtl} 文件——任何库都能写个解析器，\Tp{igl::readOBJ}、\Lib{Open3D} 的 io、\Lib{tinyobjloader}、\Lib{Assimp} 都停在``读进哪个容器''为止。\textbf{PLY} 是扫描数据的事实标准：ASCII 或二进制头声明每元素及其 property 列表，顶点可以带强度、颜色、法向等自定义通道，二进制负载按声明顺序紧凑排列、可直接读进每属性一条的连续数组。\textbf{glTF} 是为传输设计的现代格式：JSON 清单加二进制块（GLB 单文件打包），accessor 显式给出字节偏移、分量类型与数量，天然是``GPU 友好的 SoA 合同''；FBX 是其专有对位，单位与轴向约定由导出器写入并在导入端被解释。

约定层真正的两个暗礁是单位与手性，而它们\emph{不报错}。OBJ 与 PLY 没有强制单位字段：glTF 约定一米一单位，工业网格常按毫米建模，DCC 软件导出时又常缩放——同一文件两个宿主差三个数量级是合法解读的产物。手性更隐蔽：左手系与右手系换一次，网格镜像、有向面积与法向集体反向，但三角形仍然闭合、BVH 照样能建——只有当你算通量、包围体积或做布尔时，``体积怎么是负的''才把问题炸出来。导出侧的自律同样简单：把自己信的单位与轴向写进交换合同（文件名、元数据或文档三选一），解析层做完转换就交给内核，绝不让``大概是一米''这种假设流进算法。

\begin{note}[title={单位与手性：量级错却看起来对}]
一个按毫米建模的模型被当成厘米读入：线性尺寸差 $10$ 倍、面积差 $10^2$ 倍、体积差 $10^3$ 倍——质量、惯量、包围盒表面全都``合理''地错了几个数量级，渲染出来却和原图一模一样，这是三维管线最难查的一类 bug。镜像手性则让有向面积符号翻转：逐面看仍然正常，求和时``内表面''贡献负体积，封闭网格算出的体积接近零或为负。防线只有两条：入口处把单位写进项目合同并断言一个已知尺寸（比如人体模型高在 $1.5$ 到 $2.2$ 之间、单位应为米），处理带朝向语义的量之前先检查封闭网格的带符号体积为正——两条都能在导入层自动跑。
\end{note}

选型口诀可以这样背：给人看、给 CAD 互导用 OBJ；给扫描数据带属性用 PLY；给运行时与网络传输用 glTF/GLB。三者都只是``容器''，容器里没有语义：单位、手性、索引基、面朝向全靠交换双方口头以外的合同钉死。

\begin{table}[H]
\centering
\caption{本章四个库与两种接入形态的对照（许可证类别以各仓库当前 LICENSE 为准）}
\label{tab:ch09-libs}
\small
\begin{tabular}{@{}>{\raggedright\arraybackslash}p{0.10\textwidth}
  >{\raggedright\arraybackslash}p{0.22\textwidth}
  >{\raggedright\arraybackslash}p{0.14\textwidth}
  >{\raggedright\arraybackslash}p{0.17\textwidth}
  >{\raggedright\arraybackslash}p{0.17\textwidth}@{}}
\toprule
库 & 抽象层 & 许可证类别 & \PY{} 门面 & 是否实时 \\
\midrule
\Lib{libigl} & 头文件函数集，\Lib{Eigen} 矩阵为数据货币 &
  MPL 系（copyleft 子树另有 GPL 系） &
  无官方包，多从 \CPP{} 侧嵌入 & 离线原型定位 \\
\Lib{Open3D} & 编译类库，传统 API 与 T 系并存 &
  MIT 系 & 官方 pybind11 加 pip 轮子 &
  中等规模近交互，管线偏离线 \\
\Lib{Embree} & 场景与求交内核库（Device、Scene 句柄） &
  Apache 系 & 无官方门面，宿主各自绑定 &
  为实时与交互光线追踪设计 \\
\Lib{CGAL} & 模板概念库（Kernel、Polyhedron 等） &
  GPL 系加商业双许可 & 无官方门面 &
  重几何计算离线，不追帧率 \\
\bottomrule
\end{tabular}
\end{table}

\begin{bestpractice}[title={用三个提问定位一个几何库}]
问``对象怎么摆''：要邻接遍历就半边（\Lib{CGAL}、\Lib{Open3D}），要批量数值就索引加 \Lib{Eigen}（\Lib{libigl}），要空间查询就建 KDTree 或 BVH，要场就 TSDF 或体素。问``算力压在哪''：热点是求交与遍历选 \Lib{Embree}（或经 \Lib{Open3D} \Tp{RaycastingScene} 间接用它）；热点是方程与谱分解选 \Lib{libigl} 加 \Lib{Eigen}；热点是拓扑正确性回到 \Lib{CGAL} 的 kernel 档位。问``跨语言怎么办''：内核留在 \CPP{}，\PY{} 门面经 pybind11 只包调用面；给纯 \CL{} 宿主当插件时，接口层自己写一组 C 链接符号（用 \Tp{extern} 声明把 \CPP{} 名字按 \CL{} ABI 导出），别让模板泄漏到边界上——这与上一章 \Lib{GEOS}、\Lib{PROJ} 的形态是同一个决定的两次落地。
\end{bestpractice}

最后收拢本章对全书主线的回答：求交不是能整体外包给一次调用的``矩阵乘''，而是决定与取数交织的紧耦合循环——性能只有当数据布局、指令档位、分配行为都攥在代码作者手里时才存在，这三样恰好是 \CPP{} 给得了、门面语言给不了的。于是分工再次成形：\PY{} 把成批的光线与矩阵整块递过边界，\CPP{} 侧的树在边界另一头逐层下降；门面是门面，内核是内核。
% ====================================================================

% ====================================================================
% ====================================================================
\chapter{碰撞检测、刚体动力学与仿真的实时约束}

几何一旦进入物理引擎，它的职责就从``描述形状''变成``在一帧的固定预算内回答两两关系''。这个转变决定了本域全部工程细节：查询必须分批做（先粗后细）、结果必须可缓存（上一帧的邻接关系是这一帧的猜测）、误差必须有方向（穿透要能被推开，而不是报一个中性的布尔值）。更重要的是，它给实现语言设了三条硬约束：步长固定、回调域内不得分配与阻塞、状态必须可复现。本章先把几何在引擎里的三类调用摊开，再逐一检视 \Lib{Bullet}、\Lib{ODE}、\Lib{PhysX}、\Lib{MuJoCo} 与 \Lib{FCL} 的语言层与接口形态，最后把确定性与浮点求和顺序这条最容易踩的线讲清楚。

边界声明：本章只谈``几何查询怎么被调用、为什么必须用这些语言写''，不推导刚体运动学与约束方程的数学。涉及连续碰撞检测（CCD）、柔性体与流体动画的段落只做接口层面的说明。

\section{三类调用：宽相、窄相与约束求解}

物理引擎对几何库的调用可以严格分成三段，每一段的输入、输出与代价结构都不一样。把这三段混起来谈，是很多性能问题找不到根因的原因。

\subsection{宽相：把 $O(n^{2})$ 压回近线性}

宽相（broadphase）只知道每个物体的包围盒（AABB）与位姿，它的任务是把所有可能相交的\emph{候选对}列出来，并且每帧只更新真正动过的那些。三条主流路线：

\textbf{扫描-裁剪（Sweep and Prune，SAP）。} 把每个 AABB 在某条轴上的投影区间端点放进一张有序表，只有区间重叠的物体才可能是候选。物体每帧位移很小，所以有序表只需做相邻交换式的增量修正，代价接近 $O(n)$。\Lib{Bullet} 的 \Tp{btAxis\Bk Sweep3}（以及把三条轴都排起来的 SAP 变体）就是这条路：它假设物体数量固定、位移小，一旦场景里出现长距离抛体或传送，排序要大幅回移，收益就没了；而它按单轴排出来的重叠对还要在另外两轴上复查，三轴联动的实现复杂度不低。

\textbf{层次包围盒（BVH）。} 把物体组织成一棵（或两棵）包围盒树，自顶向下剪枝。\Lib{Bullet} 的默认宽相 \Tp{btDbvt\Bk Broadphase} 用的是动态层次包围盒：一棵活动树装动体、一棵预测树装静体，帧末合并。它对\emph{任意量级}的位移都稳，代价是插入与删除要沿树重算包围盒，且树的形状依赖插入顺序——这直接是后面确定性一节的一个隐患。

\textbf{体素与网格空间。} 把世界切成均匀格子或四叉/八叉树，物体按格子落桶，只对同桶者报候选。\Lib{ODE} 把这一层抽象成 \Tp{dSpaceID} 句柄并用多个实现供选择（哈希空间 \Tp{dHash\Bk Space\Bk Create}、四叉树 \Tp{dQuad\Bk Tree\Bk Create}，构造函数清一色是 \Tp{d...Create} 的形状），\Lib{PhysX} 则在场景内部维护自己的层次结构与静态网格加速结构。这类结构在大而稀疏的场景里内存换时间，收益明显；痛点是物体跨越桶边界时的迁移，以及超大 AABB（比如地形）会污染一路格子。

三者的输出形态是一致的：一份候选对表，加一个``代理（proxy）''对象用来接收位姿更新。凡是把这一层暴露成 C API 的实现，几乎都会给出 \Tp{dSpaceID} 或某个 \Tp{...BroadphaseInterface} 这样的句柄式抽象——它是唯一一个能被外部语言稳定持有的东西。

\subsection{窄相：凸体是通用货币}

窄相（narrowphase）拿到候选对之后才做精确几何。这一层的共识是：\emph{任意形状都先用凸包近似}，因为凸体之间的距离与穿透有可靠的、只依赖支持函数（support function）的算法，而凹形可以拆成若干凸片。\Lib{Bullet} 里的 \Tp{bt\Bk Convex\Bk Polyhedron} 与 \Tp{bt\Bk Shape} 家族、\Lib{FCL} 里的 \Tp{fcl::\Bk Convex} 都以此为中心。

\textbf{GJK。} 给定两个凸体，只要求它们的支持函数（在方向 $d$ 上使 $d \cdot x$ 最大的那个顶点 $x$），就能在 Minkowski 差 $A\ominus B$ 上迭代逼近原点：维护一个至多四个顶点的单纯形，每轮把最接近原点的区域缩进去。返回的是距离、以及两个凸体上的\emph{见证点}（witness points），后者可以直接给出接触法向。它的优点是解析几何之外什么都不需要，圆柱、胶囊、凸包、Minkowski 和共享同一套代码；代价是迭代次数不可预算，而且在近平行的贴面接触上会退化（详见清单~\ref{lst:ch10-gjk} 与其后的讨论）。

\textbf{EPA。} 一旦 GJK 报告距离为零（相交），改用扩张多面体算法：在包含原点的单纯形凸包上，反复朝\emph{最小}外法向距离的方向加新支持点，得到穿透深度与穿透法向。它的收敛是几何级数式的，但精度受容差与面数限制，深度很浅时迭代次数会明显上升——这正是``堆叠箱子发抖''的隐藏成因之一。

\textbf{SAT。} 分离轴定理走另一条路：枚举两个多面体的候选分离轴（各自的面对法向加上边-边叉积），找到穿透最小的那根轴。\Lib{Bullet} 的盒-盒、盒-凸包路径（AXYZ、SAT 加 Z-Axis Fall-Back 一类实现）就属于这一族，代价是轴数随面与边的数量增长，复杂度约 $O(n_{A}+n_{B})$，但每一轮只是点积比较，没有迭代收敛问题。

\textbf{解析特例。} 球-球、球-盒、胶囊-胶囊有闭式解，任何严肃实现都会为它们写专门分支：\Lib{Bullet} 的形状对分派表（\Tp{bt\Bk Collision\Bk Dispatcher} 按形状类型对查算法）存在的理由就是把特例当一等公民。三角形网格对凸体则是``GJK 定粗位置 + 局部三角形遍历''，纯逐个三角形跑 GJK 是常见的性能事故。

窄相的输出是\emph{接触流形}：一组接触点，各带穿透深度、法向、以及由法向张出的两个摩擦方向基。一个盒-盒相交最多能给出四个点，这就是``一对候选不等于一个约束''，也是求解器规模常常被低估的原因。

\subsection{约束求解：把几何结果换成力}

接触点本身不是物理量，它要立刻被翻译成互补约束。\Lib{Bullet} 把上帧的流形缓存在 \Tp{bt\Bk Persistent\Bk Manifold} 里（同时充当宽相与窄相之间的热启动载体），求解器则有三条技术路线：

\textbf{LCP。} 法向冲量 $\lambda$ 满足 $\lambda \ge 0$、剩余间隙 $f \ge 0$、互补 $\lambda \cdot f = 0$ 的线性互补问题，一次精确解需要构造并分解稠密的接触质量矩阵（$K$ 矩阵），代价随接触数立方增长。\Lib{Bullet} 曾提供 \Tp{bt\Bk MLCPSolver} 走 Dantzig 主轴法，实践中因为规模不划算而少用。

\textbf{PGS 与序贯冲量。} 投影梯度式（Projected Gauss-Seidel）逐条约束做近似并立即夹紧到可行域，几轮迭代就能给出视觉上稳定的结果；每轮代价约 $O(m)$（$m$ 为约束数）。这是实时预算下真正胜出的算法，代价是收敛质量：多层堆叠、闭合链的耦合会明显收敛慢，因此\emph{热启动}（用上一帧的冲量当初始猜测）几乎是必需的。同理，穿透修正要么走 Baumgarte 式的速度偏置，要么走 split impulse 把位置修正与速度修正分开，避免修正量污染动量守恒。

\textbf{ABA。} Featherstone 的艺术化体算法（Articulated Body Algorithm）用 $O(n)$ 的正向递推与反向传播处理铰接链，把``通用接触``降级为局部冲量。它擅长机械臂、机器人这类深链，不擅长大规模自由体堆叠。\Lib{Bullet} 的多体路径（\Tp{bt\Bk Multi\Bk Body\Bk Constraint\Bk Solver}）与 \Lib{MuJoCo} 的 CRB/ABA 都在这条线上。

三段的量级关系值得记住：$n$ 个物体堆叠，宽相候选可以远超 $O(n)$，窄相每对最多产四个点，求解器每轮迭代都要扫一遍全部约束。于是帧预算被写成 $\text{cost} \approx C_{\text{bp}} + n_{\text{pair}} \cdot C_{\text{np}} + \text{iter} \cdot n_{\text{constr}} \cdot C_{\text{row}}$，而 \Tp{iterations} 是调用方手里最便宜的那个旋钮。

\textbf{还有第四段：连续碰撞检测（CCD），它不是窄相的附注。} 一颗高速抛体穿过薄墙，在 $t$ 与 $t+\Delta t$ 两个采样时刻都``不相交''，离散步进的窄相根本看不见这件事。答案是扫掠体：把凸体沿相对位移拉伸成 Minkowski 和（凸体加一条线段），或者在位移插值上反复调用 GJK 求最早接触时刻（time of impact）。\Lib{Bullet} 的 CCD 路径与 \Lib{FCL} 的连续碰撞接口都只提供这个查询的一个近似档位，代价是每对要多跑几轮迭代，因此两边都要求调用方显式标记哪些物体参与 CCD。把 CCD 打开给全场，帧预算立刻爆掉。

\subsection{GJK 的循环与它的两种退化}

\begin{lstlisting}[style=cppstyle,label={lst:ch10-gjk},caption={GJK 距离迭代的核心循环：单纯形、收敛判据与退化处理}]
// 骨架示意：真实的区间管理在 Bullet 的 btGjkPairDetector、
// btGjkEpaSolver2 与 FCL 的 GJK 求解器里都各有专门实现
struct Sup {
  virtual Vec3 support(const Vec3& d) const = 0;  // 方向 d 上最远的点
};

// Minkowski 差 A(-)B 的支持函数：只用两次单形状查询拼出来
Vec3 supportDiff(const Sup& a, const Sup& b, const Vec3& d) {
  return a.support(d) - b.support(-d);
}

// 返回两凸体距离；返回 0 表示相交，此时该由 EPA 接手求深度
float gjk(const Sup& a, const Sup& b, Vec3& wA, Vec3& wB,
          int maxIter, float tol) {
  Simplex s;
  Vec3 d(1.f, 0.f, 0.f);          // 初始方向只要非零
  growSimplex(a, b, d, s, wA, wB); // 取第一个见证点，方向翻向原点
  for (int it = 0; it < maxIter; ++it) {
    float v = s.closestDist(wA, wB);      // 单纯形到原点的最近距离
    if (v <= tol * (1.f + v)) return v;   // 判据 1：相对改进量耗尽
    Vec3 nd = s.closestDir();             // 判据 2：新的搜索方向
    Vec3 p = supportDiff(a, b, nd);       // 朝原点一侧再取一个点
    if (p.dot(nd) <= 0.f) return v;       // 新点在背面：不可能更近
    if (s.almostCoplanar(p, tol)) {       // 退化 1：新点几乎落在旧面上
      jitter(nd);                         // 扰动方向，否则增量趋于 0
      continue;                           // 贴面接触就是从这里退化的
    }
    s.add(p);
    if (s.containsOrigin(wA, wB)) return 0.f;  // 退化 2：原点进入，改走 EPA
  }
  return s.closestDist(wA, wB);   // 迭代耗尽：返回当前估计值而非真值
}
\end{lstlisting}

清单~\ref{lst:ch10-gjk} 的两处退化决定了它在引擎里能不能用。第一种是\emph{共面退化}：两个多面体的面几乎贴合时，新支持点落在已有单纯形的平面上，方向投影长度趋零，迭代再多也不改进——所有实现都要在这里扰动方向或改用专门的面-面路径（SAT 一族正是为这种情况准备的）。第二种是\emph{返回上界而非真值}：迭代耗尽时返回的是当前单纯形给出的距离估计，它偏大；调用方若把它当成精确距离去做``是否接触''的布尔判断，会漏检。工程上的结论很直白：宽相容差、迭代上限与穿透阈值必须一起调，而不是只调其中一个。

\section{主流实现的语言层与接口形态}

\textbf{\Lib{Bullet}}。接口是纯粹的 \CPP{} 类层次：\Tp{bt\Bk Collision\Bk World} 提供只做检测不做动力学的世界（适合编辑器、遮挡与射线查询），\Tp{bt\Bk Discrete\Bk Dynamics\Bk World} 在其上加刚体与求解器，多体与软体各自再有世界类型。形状对象（\Tp{bt\Bk Capsule\Bk Shape}、\Tp{bt\Bk BvhTriangle\Bk Mesh\Bk Shape}）是与位姿无关的共享几何，\Tp{bt\Bk Collision\Bk Object} 才持有变换，这个分工在清单~\ref{lst:ch10-bullet} 里看得到。它按 zlib 系条款发行，代价是接口没有稳定的 C ABI：\CPP{} 的虚表布局、名字修饰与异常跨越编译单元时都要一致，于是各语言的绑定长期分叉，绑定层与主干版本对不上是常态。

\textbf{\Lib{ODE}}。反过来的历史样本：公共接口是纯 \CL{}，一切以句柄传递（\Tp{dWorldID}、\Tp{dBodyID}、\Tp{dSpaceID}、\Tp{dJoint\Bk GroupID}），步进是 \Tp{dWorld\Bk Quick\Bk Step}（它的岛分割加序贯冲量正是 PGS 家族）。这让它在 \CL{}、\PY{}、\FTN{} 侧都极易绑定，也让它成为很多教学与老项目的默认选择。代价同样真实：句柄必须成对手动释放、错误靠返回值与全局状态、命名不完全一致、内存管理要靠联合组（joint group）自己回收。历史上的 \CPP{} 封装从未成为主流，因为真正在用的引擎大多直接吃 C 层。

\textbf{\Lib{PhysX}}。接口是 \CPP{} 抽象类加描述符结构体的两段式（先填 \Tp{Px\Bk Rigid\Bk Actor\Bk Desc} 再建对象），这种写法的一个动机就是跨动态库边界的 ABI 稳定：不把布局敏感的对象暴露出去，只传描述符与不透明指针。它起家于专有授权，后来的版本以源码形式公开、条款与早期发行不同，具体以官方仓库为准，本章不做条款判定；GPU 加速的刚体与粒子在部分平台与版本上可用，接口与 CPU 路径不同。对本书主线而言它的价值在于示范了``商用闭源核心 + 稳定接口''如何被多语言驱动。

\textbf{\Lib{MuJoCo}}。这是本域最干净的 \CL{} 优先设计：\Tp{mj\Bk Model} 装常量（质量、惯性、关节拓扑、碰撞基元），\Tp{mj\Bk Data} 装状态与结果，两者分离使同一模型可以有任意多份独立状态，天然适合并行采样与批量强化学习。步进是显式的 \Fn{mj\_step}（内部可拆成 \Fn{mj\_step1} 与 \Fn{mj\_step2} 以便在中间注入控制）。它按 Apache 2.0 条款发行，并有官方 \PY{} 门面——门面薄到只是把结构体字段暴露成数组视图，因为 C ABI 已经把可绑定问题解决了。

\textbf{\Lib{FCL}}（\Tp{libfcl}）。它的位置在机器人运动规划而不是游戏物理：运动规划器要的是\emph{距离}与\emph{连续性检查}（这条线段两端之间有没有扫过障碍），而不是接触力。它把距离查询当一等公民（\Fn{fcl::\Bk distance} 与 \Fn{fcl::\Bk collide} 并列），支持层次结构与不同后端求解器（\Tp{GJK\Bk Solver\Bk Type} 可选择内置或外接 \Lib{libccd}），并提供连续碰撞检测接口。接口是模板化的 \CPP{}，被 MoveIt 一系的规划栈直接嵌入。

\begin{lstlisting}[style=cppstyle,label={lst:ch10-bullet},caption={Bullet：只用碰撞世界做检测与流形读取，不引入动力学}]
// 依赖：Bullet 3.x；链接 btBulletCollision 与 btLinearMath
#include "btBulletCollisionCommon.h"
#include <cstdio>

int main() {
  btDefaultCollisionConfiguration conf;     // 窄相算法的配置表
  btCollisionDispatcher disp(&conf);        // 按形状对分派窄相算法
  btDbvtBroadphase broadphase;              // 宽相：动态层次包围盒
  btCollisionWorld world(&disp, &broadphase, &conf);

  btCapsuleShape capsule(0.25f, 1.0f);      // 半径与柱段长度
  btBoxShape box(btScalar(0.5), btScalar(0.5), btScalar(0.5));
  btTransform tfA = btTransform::getIdentity();
  btTransform tfB = btTransform::getIdentity();
  tfA.setOrigin(btVector3(0.f, 0.f, 0.f));  // 形状是共享几何，
  tfB.setOrigin(btVector3(0.6f, 0.f, 0.f)); // 位姿挂在 object 上

  btCollisionObject objA; objA.setCollisionShape(&capsule);
  objA.setWorldTransform(tfA);
  btCollisionObject objB; objB.setCollisionShape(&box);
  objB.setWorldTransform(tfB);
  world.addCollisionObject(&objA);
  world.addCollisionObject(&objB);

  world.performCollisionDetection();        // 宽相候选 + 窄相流形

  const int nm = world.getDispatcher()->getNumManifolds();
  for (int i = 0; i < nm; ++i) {
    const btPersistentManifold* m =
        world.getDispatcher()->getManifoldByIndexInternal(i);
    for (int k = 0; k < m->getNumContacts(); ++k) {
      const btManifoldPoint& p = m->getContactPoint(k);
      const btVector3& n = p.m_normalWorldOnB;   // 法向由流形给出
      // dist <= 0 表示已互相穿入，其绝对值就是穿透深度
      std::printf("dist=%8.4f n=(%.3f,%.3f,%.3f)\n",
                  p.getDistance(), n.getX(), n.getY(), n.getZ());
    }
  }
  // 单个物体的即时查询还有：world.contactTest(obj, cb)
  // 射线与形状查询：world.rayTest(o, d, cb) 系列
  return 0;
}
\end{lstlisting}

\begin{lstlisting}[style=cstyle,label={lst:ch10-mujoco},caption={MuJoCo 的 C API：模型与数据分离，固定步长显式推进}]
// 依赖：mujoco；链接 -lmujoco
#include <mujoco/mujoco.h>
#include <stdio.h>

int main(void) {
  char err[1024] = {0};
  mjModel* m = mj_loadXML("arm.xml", NULL, err, sizeof err);
  if (m == NULL) { fprintf(stderr, "load: %s\n", err); return 1; }
  mjData* d = mj_makeData(m);            // 状态容器，与模型解耦

  m->opt.timestep = 0.001;               // 固定步长：节拍与墙钟无关
  m->opt.iterations = 60;                // 迭代上限决定最坏帧耗时
  m->opt.tolerance = 1e-6;               // 只作早退，不作唯一判据

  for (int tick = 0; tick < 3000; ++tick) {
    mj_step(m, d);                       // 一次 = 前向 + 约束求解
    // 读结果：d->qpos、d->qvel、d->ncon、d->contact[k].dist
    // 需要中途注入控制时，用 mj_step1 与 mj_step2 拆开
  }
  mj_resetData(m, d);                    // 复位到初始状态，可整帧重放
  mj_deleteData(d);                      // 顺序很重要：先 data 后 model
  mj_deleteModel(m);
  return 0;
}
\end{lstlisting}

\begin{table}[H]
\centering
\caption{五个实现的接口语言、许可类别、GPU 能力、确定性与 \PY{} 门面}
\label{tab:ch10-libs}
\small
\begin{tabular}{@{}>{\raggedright\arraybackslash}p{0.11\textwidth}
  >{\raggedright\arraybackslash}p{0.12\textwidth}
  >{\raggedright\arraybackslash}p{0.15\textwidth}
  >{\raggedright\arraybackslash}p{0.12\textwidth}
  >{\raggedright\arraybackslash}p{0.17\textwidth}
  >{\raggedright\arraybackslash}p{0.13\textwidth}@{}}
\toprule
库 & 接口语言 & 许可类别 & GPU & 确定性 & \PY{} 门面 \\
\midrule
\Lib{Bullet} & \CPP{} 类 & zlib 系 & 无系统级 GPU 刚体 &
  同平台可复现，跨平台不保证 & 第三方绑定，需与版本对齐 \\
\Lib{ODE} & 纯 \CL{} 句柄 & LGPL 系 & 无 &
  与调用顺序绑定，可复现 & 第三方绑定，薄 \\
\Lib{PhysX} & \CPP{} 描述符加接口 & 见官方条款 &
  部分版本与平台提供 & 官方提供确定性模式（程度见文档） &
  第三方包装，非官方主推 \\
\Lib{MuJoCo} & 纯 \CL{} 结构体 & Apache 2.0 &
  有官方渲染后端，非物理 GPU & 官方声明同平台确定 &
  官方门面，字段级直连 \\
\Lib{FCL} & \CPP{} 模板 & BSD 系 & 无 &
  查询是纯函数，可复现 & 机器人栈内绑定 \\
\bottomrule
\end{tabular}
\end{table}

\section{为什么这一域必须由 C/C++ 承担}

表~\ref{tab:ch10-libs} 里的五个实现没有一个把脚本语言放在实现层。原因不是习惯，而是四条可核对的工程约束。

\textbf{其一，固定步长与指令周期预算。} 物理世界的时间轴由 \Tp{timestep} 决定，与墙钟无关：仿真要的是每 tick 做同样的事。在一帧 $16.7$ ms、主频 $3$ GHz 量级的核上，预算大约是 $5\times 10^{7}$ 个周期，要装下一次宽相增量更新、几千对窄相查询和几十轮约束迭代。这个数量级的预算只有在能直接控制内存布局、SIMD 与循环展开的语言里才可估算，也才可\emph{守住}：\CPP{} 允许把包围盒写成连续数组、按缓存行对齐、用无分支的区间测试，并把\emph{最坏}路径（退化接触、迭代耗尽）算进预算而不是期望值。

\textbf{其二，回调域内禁止分配与阻塞。} 引擎在主循环或独立的高优先级线程里调步进函数，那个域里有三个共同禁忌：堆分配（\Tp{malloc}/\Tp{operator new} 拿全局堆锁，可能触发页错误与内存清零）、日志（格式化、加锁、写设备）、锁（优先级反转与无界等待）。这三条在下一节的警告框里合成一个事故清单。

\textbf{其三，要能被别的栈以稳定 ABI 嵌入。} \Lib{ODE} 与 \Lib{MuJoCo} 选纯 \CL{} 接口，正是为了这个游戏引擎、机器人中间件与科学计算栈通吃的汇合点：C ABI 只暴露句柄或 POD 结构体，跨编译器、跨语言、跨动态库都稳。\CPP{} 接口要嵌入就得自己克制（描述符两段式、不透明指针、虚表不跨边界），\Lib{PhysX} 就是这个路数的示范。反过来，脚本侧的需求全部被推到绑定层：\PY{} 门面之所以能做，是因为底下的步进函数是\emph{无隐藏状态}的一次调用。

\textbf{其四，可枚举的三处并行都落在这些语言里。} 宽相的排序与树遍历（SAP 的三条轴可以分别求，BVH 的 refit 天然可分片）、窄相的候选对表（每对独立，是标准的并行按块处理）、求解器要先做\emph{岛分割}（把互相耦合的约束集拆开，通常用并查集，岛与岛之间才能并行）——\Lib{ODE} 的 \Tp{dWorld\Bk Quick\Bk Step} 里就有这一步。数据布局要为这三处服务：包围盒存成连续数组而不是对象、接触点按帧复用同一块 arena，让 \Lib{Bullet} 的形状与物体由调用方持有（清单~\ref{lst:ch10-bullet} 里两个 \Tp{btCollisionObject} 都是栈对象），库不拥有状态，于是步进路径上没有隐藏分配，也没有析构顺序问题。\Lib{MuJoCo} 把 model 与 data 拆成两个结构体是同一动机的另一种写法。要注意并行度与逐位复现之间没有免费的调和：这三处每加一次并行，都会改变遍历与归约顺序，直接撞上下一节的确定性要求，只能显式选择站在哪一边。

\begin{warning}[title={回调域里的分配、日志与锁}]
\begin{itemize}
\item \textbf{分配。} 在步进回调里 \Fn{new} 一个临时向量，平均耗时看起来微不足道，但它的最坏路径由堆锁、页错误与释放时的合并决定；一旦超过帧预算，后果不是``慢一点''而是\emph{掉帧与物理抖动}——下一帧的 $\Delta t$ 变了，接触解的初值也变了，堆叠开始晃。做法是预分配池、arena 分配器与固定容量数组，把分配挪到世界初始化时。
\item \textbf{日志。} 在窄相循环里打印一对接触，会把一次微秒级工作变成一次控制台同步写。若日志是异步的，还要看它的队列是否有界；无界队列在内存压力下自己变成阻塞源。正确做法是帧末计数、按统计频率输出。
\item \textbf{锁。} 回调里为读共享状态拿互斥锁，与更高或更低优先级的线程形成竞争，最坏是无界等待。实时域里应换成无锁环形缓冲、双缓冲或把状态复制进回调自己的 arena。
\end{itemize}
\end{warning}

\section{确定性、锁步仿真与可复现的代价}

\textbf{为什么要可复现。} 三个场景对``逐位相同''有刚需：锁步（lockstep）同步的网络游戏中，各端只传输入与 tick 号、各自重放同一世界，一旦浮点轨迹分叉就会\emph{desync}且极难定位；bug 复现靠录制输入序列回放现场；强化学习批量的对照组必须只有随机种子在变。这三件事共同的隐含要求是：仿真函数在相同状态与相同输入下返回逐位相同的结果。

\textbf{分叉从哪里来。} 第一个源头是\emph{求和顺序}：浮点加法不满足结合律，于是任何并行归约、任务窃取、线程数变化、\Tp{reduction} 式的分块都会改变最后几位。第二个源头是\emph{迭代式求解器对顺序的敏感}：序贯冲量逐条约束投影，约束遍历顺序换了，即使每轮总工作量不变，结果也不同——而遍历顺序常常来自宽相容器的内部结构（BVH 的插入顺序、SAP 表里相同键值的相对次序、哈希空间的桶顺序），这些都不是被保证稳定的。第三个源头是\emph{工具链与硬件}：超越函数在数学库实现间的差异、是否把 $a\cdot b+c$ 收缩成 FMA、x87 的扩展精度与 SSE 的分歧、\Tp{long double} 的宽度、SIMD 尾块归约的顺序。

\textbf{\Tp{-ffast-math} 为什么直接毁掉回放。} 它把 IEEE 语义当成可选优化：允许重排与重新结合加法（于是归约顺序不再是代码里写的那个）、把 $\sqrt{x}$ 换成倒数平方根的近似迭代、取消对 NaN 与 $\pm\infty$ 的特殊路径判断、把非规格化数当零（flush-to-zero）。这些改动每一项都合法地\emph{改变了结果}，而且是在硬件与编译器版本间不一致地改变。用它换来的一点提速，代价是录制-回放体系整体失效。真要做性能优化，应该走明确语义的路子：\Tp{-fno-finite-math-only} 之类的部分开关、把热循环改成双精度累加或 Kahan 求和、或者干脆固定迭代次数而不是让容差驱动早退。

\textbf{工程上的组合拳。} 把宽相容器与接触对表做成\emph{顺序确定}的数据结构（显式索引、固定分桶），把求解器迭代数固定、把容差只当早退，把关键状态（tick 号、物体句柄）留在整数域，每帧对状态算一个哈希用于回放比对，跨端一致性优先于跨端性能。最后：如果确实需要跨平台逐位一致，就得把浮点降到受控的定点或整数路径，而不是祈祷三个数学库行为相同。

\textbf{几何数据结构是跨域复用的。} 本章宽相与窄相用的层次包围盒、体素树与区间测试，与后面``网格、点云与光线求交''一章为射频传播规划准备的可见性、遮挡查询，是同一套数据结构在两个领域的两次调用：碰撞问``这两个立体是否重叠''，传播问``两点之间的线段是否被遮挡''，加速结构、增量更新与遍历顺序约束完全同构，差别只在查询形状（OBB 换到线段）与容差语义（穿透深度换到路径损耗的判定阈值）。这意味着一处实现可以服务两个域，也意味着上面那整节确定性约束在射线求交上同样成立。

\begin{guideline}[title={本章要点}]
\begin{itemize}
\item 几何进物理引擎要分三段看：宽相买\emph{候选}（SAP、BVH、体素树各有失效场景），窄相买\emph{距离与深度}（凸体是通用货币，GJK、EPA 与 SAT 各有退化），求解器买\emph{可执行的冲量}（LCP 精确但贵，PGS 便宜但要热启动，ABA 擅长铰接链）。
\item 接口语言的选择是这个域的功能需求而非口味：\Lib{ODE} 与 \Lib{MuJoCo} 用纯 C 句柄换可绑定性，\Lib{Bullet} 与 \Lib{FCL} 用 \CPP{} 模板换内联与泛型，\Lib{PhysX} 用描述符两段式换跨库 ABI 稳定。
\item 帧预算、回调域禁忌、可枚举的并行点与 C ABI 汇合点，是``为什么必须是 C/C++''在本域的四条硬理由；\PY{} 只能站在门面，因为一次步进都在实时域里。
\item 可复现性是加法顺序、容器遍历顺序与工具链浮点语义的合取性质。\Tp{-ffast-math} 一开，锁步与回放就不成立，这不是调优而是换语义。
\end{itemize}
\end{guideline}
% ====================================================================

% ====================================================================
% ====================================================================
\part{信号处理库：声音}

\chapter{采样、音频格式与实时 IO}

声音这一域和前面几章的数值计算有一个根本差别：它的\emph{消费者是实时设备}。矩阵乘法的输出晚到十微秒毫无后果，而音频回调晚到十微秒就是一声爆音。因此本章的所有约定都可以归成两件事：第一，把``声音''这团连续信号压成一套可跨语言、跨文件、跨设备传递的离散表示（采样率、样本类型、布局、通道顺序、时间基）；第二，在操作系统给的硬预算内把这块表示搬进搬出（回调、缓冲块、XRUN）。看一遍这一章的接口就会明白为什么实现层是 \CL{} 与 \CPP{}：\Lib{libsndfile}、\Lib{PortAudio}、\Lib{miniaudio}、\Lib{Cubeb} 与 \Lib{SoX} 全部把 \CL{} 接口当第一等公民，\PY{} 那一侧的 \Lib{soundfile} 与 \Lib{sounddevice} 只是同一批符号的搬运工。

边界声明：滤波器、重采样算法与响度度量属于本章之后的 FIR 与 IIR 一章；这里只把\emph{表示层与搬运层}讲透，遇到需要 DSP 的地方明确指向那一章。

\section{表示层的六条约定}

\textbf{采样率与带宽。} 采样定理给出可表示的最高频率 $f_s/2$：$48$ kHz 的名义上限是 $24$ kHz，$44.1$ kHz 则是 $22.05$ kHz（这个档位的来历与 CD 时代用模拟录像机存 PCM 母带的历史约束有关，并不是技术上的最优选）。工程上有两个后果：其一，采集链路的抗混叠滤波器必须在 $f_s/2$ 之前把能量压下去，这是硬件的事；其二，$44.1$ kHz 与 $48$ kHz 之间的转换\emph{不是}整数比，任何在这两档之间搬数据的流程都要经过一次分数重采样（比例 $160/147$），这是素材管线里噪声最常见的产地。

\textbf{位深度与动态范围。} $b$ 位线性定点的量化的信噪比约为 $6.02b + 1.76$ dB，于是 16 位约 $98$ dB、24 位约 $146$ dB。参考点是 \textbf{dBFS}（decibel full scale）：满刻度为 $0$，所有实际信号都是负值。要注意 dBFS 有两种口径（以单样本满刻度为 $0$ dB，或以满刻度正弦峰值为 $0$ dB，两者相差约 $3$ dB 的峰值因子），跨工具比较数值时先确认口径。

\textbf{三种样本类型。} \Tp{int16\_t}（16 位有符号，最省，仍是文件与蓝牙链路的主流）、24 位样本、\Tp{float32}（\CL{} 里 \Tp{float}，约定的满刻度区间是 $[-1,1]$，允许中间结果超过 $\pm 1$、到定点化时才饱和）。24 位在容器里有两个样子：真正的 $3$ 字节打包（\Tp{SF\_\Bk FORMAT\_\Bk PCM\_\Bk 24}，一些 WAV 与 FLAC 就是这个），以及装在 32 位字里的高位有效、低位补零的形式（ALSA 的 \Tp{SND\_\Bk PCM\_\Bk FORMAT\_\Bk S32\_\Bk LE} 与 FFmpeg 的 \Tp{AV\_\Bk SAMPLE\_\Bk FMT\_\Bk S32} 常承载这种 24-in-32）——两者的字节宽度与位移不同，直接按字节 memcpy 会整体静音或满屏噪声。

\textbf{交织与非交织。} 交织（interleaved）按帧排列：L R L R L R，它是文件与设备驱动的自然布局，因为一次 DMA 搬运整帧。非交织（planar）按通道排列：先一整段 L 再一整段 R，它是 DSP 的自然布局，因为对某一路做滤波时寻址是连续的、可向量化的。FFmpeg 用 \Tp{AV\_\Bk SAMPLE\_\Bk FMT\_\Bk FLT}（交织）与 \Tp{AV\_\Bk SAMPLE\_\Bk FMT\_\Bk FLTP}（非交织）区分这两件事，而 \Lib{libsndfile} 只暴露交织语义，非交织要调用方自己做置换（清单~\ref{lst:ch11-sndfile}）。

\textbf{通道数与 speaker mask。} ``两个通道''不``等于''一个明确定义：是 L R、还是 Ls Rs、还是 Ambisonics 的第一阶四路？WAV 的扩展头 \Tp{WAVE\_\Bk FORMAT\_\Bk EXTENSIBLE} 用一个 32 位 \Tp{dwChannel\Bk Mask} 声明每一路挂在哪个扬声器位置，FFmpeg 把它读成 \Tp{channel\_layout} 位掩码，未声明时不同实现会各自猜默认顺序。多声道素材一经过 ``声道数''这个整数就再也恢复不了原始语义，所以任何跨工具的音频管线都应把掩码一起带走。

\textbf{时间基与块大小。} 音频里有三种时钟：样本索引（离散的、可精确计数的）、秒（浮点的、会漂移）、设备帧计数（由声卡驱动，独立于 CPU）。任何一次跨域换算都要指定基准。搬运用的最小单位是\emph{块}（一个回调处理的帧数，PortAudio 里的 \Tp{framesPerBuffer}，ALSA 里的 period），它和缓冲块数一起给出延迟下界：约 $n_{\text{buf}} \times \text{frames} / f_s$ 秒。$48$ kHz 下取 $256$ 帧一块、两块缓冲，往返延迟大约在十毫秒量级；把块调大能降低 CPU 占用与丢块概率，代价是延迟线性上升。录音监听、乐器建模这类交互任务就卡在这个折中上。

表~\ref{tab:ch11-sample} 把四种样本约定按调用方关心的四个维度排开：字节宽度、取值语义、常见来源与最容易踩的坑。

\begin{table}[H]
\centering
\caption{四种样本约定：宽度、取值语义、来源与陷阱}
\label{tab:ch11-sample}
\small
\begin{tabular}{@{}>{\raggedright\arraybackslash}p{0.14\textwidth}
  >{\raggedright\arraybackslash}p{0.10\textwidth}
  >{\raggedright\arraybackslash}p{0.19\textwidth}
  >{\raggedright\arraybackslash}p{0.19\textwidth}
  >{\raggedright\arraybackslash}p{0.18\textwidth}@{}}
\toprule
样本类型 & 每样本字节 & 取值语义 & 常见来源 & 典型事故 \\
\midrule
\Tp{int16\_t} & 2 &
  $-32768$ 到 $32767$，正负不对称 &
  文件交付的默认档、蓝牙链路 &
  缩放分母取 $32767$ 会让一端越界 \\
24-in-32 & 4（有效 3） &
  高位有效，低位未定义 &
  声卡的 32 位帧、\Tp{AV\_\Bk SAMPLE\_\Bk FMT\_\Bk S32} &
  按 3 字节解读会整段错位成噪声 \\
\Tp{float32} & 4 &
  约定 $[-1,1]$，实际可越界 &
  内部 DSP、设备共享缓冲、\Tp{SF\_\Bk FORMAT\_\Bk FLOAT} &
  越界只在定点化那一步才暴露 \\
\Tp{float64} & 8 &
  同上，余量与精度更大 &
  测量与母带链路、\Tp{SF\_\Bk FORMAT\_\Bk DOUBLE} &
  内存带宽翻倍，SIMD 通道数减半 \\
\bottomrule
\end{tabular}
\end{table}

\section{\Lib{libsndfile}：把封装与编码解耦的 C API}

\Lib{libsndfile} 是文件层的\emph{事实标准}，接口是纯 \CL{}。它最值得学的不是函数名，而是把``容器''与``编码''拆成两个正交参数的设计：

\begin{itemize}
\item \Tp{SNDFILE} 是不透明句柄，\Tp{SF\_INFO} 是可由调用方构造的纯数据描述（\Tp{samplerate}、\Tp{channels}、\Tp{frames}、\Tp{format}、\Tp{sections}）。读之前 \Fn{sf\_open} 把文件里的信息填进它；写之前\emph{由你}把它填好——这就是为什么它能写裸格式，也解释了为什么 \Tp{frames} 在写入时是输出而不是输入。
\item \Tp{format} 字段是\emph{容器位}与\emph{子格式位}的按位或：\Tp{SF\_\Bk FORMAT\_\Bk WAV} 配 \Tp{SF\_\Bk FORMAT\_\Bk PCM\_\Bk 16}、\Tp{SF\_\Bk FORMAT\_\Bk FLAC} 配 \Tp{SF\_\Bk FORMAT\_\Bk PCM\_\Bk 16}、\Tp{SF\_\Bk FORMAT\_\Bk OGG} 配 \Tp{SF\_\Bk FORMAT\_\Bk VORBIS}、\Tp{SF\_\Bk FORMAT\_\Bk AIFF} 配 \Tp{SF\_\Bk FORMAT\_\Bk PCM\_\Bk 24}。同一份 \CPP{} 读写代码因此能跨 WAV、AIFF、FLAC、Ogg、CAF、WAV64 与较新版本加入的其它容器工作；编码是否可用取决于编译时链了哪些后端（例如 Ogg 支不支持取决于是否带 libvorbis 或 libopus），运行前用 \Fn{sf\_format\_check} 问一句比假设便宜。
\item 读写按``帧''计数而不是字节：\Fn{sf\_readf\_float} 一次取 $N$ 帧并给出 $N \times \text{channels}$ 个 \Tp{float}，\Fn{sf\_readf\_short} 取整型，另有 \Fn{sf\_read\_raw} 绕过一切转换。类型转换、字节序、饱和与（写入定点时的）抖动都由库在这一步做掉——这是把音频转换逻辑放 C 层的一个典型收益：调用方不需要知道任何定点缩放细节。
\item 浮点子格式写作 \Tp{SF\_\Bk FORMAT\_\Bk FLOAT}（32 位）与 \Tp{SF\_\Bk FORMAT\_\Bk DOUBLE}；以 \Tp{SFC\_\Bk} 前缀命名的一批命令通过 \Fn{sf\_command} 调节行为，例如 \Tp{SFC\_\Bk SET\_\Bk CLIPPING} 控制写入定点时是否饱和。子格式与开关都随版本增加，命名以所装版本的头文件为准。
\end{itemize}

\begin{lstlisting}[style=cstyle,label={lst:ch11-sndfile},caption={libsndfile：读 WAV 转非交织 float32，处理后写回同格式}]
// 依赖：libsndfile；链接 -lsndfile。库只给交织语义，
// 非交织布局要在调用方做一次显式置换。
#include <sndfile.h>
#include <stdio.h>
#include <stdlib.h>

#define CAP 4096                       // 每次搬运的帧数：块越大 IO 越少

// 占位的通道内处理：真正的 DSP 在后面的滤波器一章
static void apply_gain(float* p, size_t n) {
  for (size_t i = 0; i < n; ++i) p[i] *= 0.5f;   // 换算法只改这一处
}

int main(void) {
  SF_INFO inf;                          // 采样率、帧数、通道、格式
  SNDFILE* in = sf_open("in.wav", SFM_READ, &inf);
  if (!in) { puts(sf_strerror(in)); return 1; }
  const int ch = inf.channels;
  size_t words = (size_t)CAP * (size_t)ch;
  float* pl = malloc(sizeof(float) * words);   // planar，非交织
  float* il = malloc(sizeof(float) * words);   // interleaved，交织
  SNDFILE* out = sf_open("out.wav", SFM_WRITE, &inf);
  // 写回沿用同一个 SF_INFO：容器与编码一个都没换
  sf_count_t got;
  while ((got = sf_readf_float(in, il, CAP)) > 0) {
    for (sf_count_t i = 0; i < got * ch; ++i)   // 交织 -> 平面
      pl[i % ch * CAP + i / ch] = il[i];
    for (int c = 0; c < ch; ++c)                // 各通道独立处理
      apply_gain(pl + (size_t)c * CAP, (size_t)got);
    for (sf_count_t i = 0; i < got * ch; ++i)   // 平面 -> 交织
      il[i] = pl[i % ch * CAP + i / ch];
    if (sf_writef_float(out, il, got) != got) return 1;
  }                       // 定点化与抖动由库在最后一步做掉
  sf_close(in); sf_close(out);
  free(il); free(pl);
  return 0;
}
\end{lstlisting}

清单~\ref{lst:ch11-sndfile} 里最值得注意的一行是 \Fn{sf\_readf\_float} 的返回值：它是\emph{实际读到的帧数}，短读是正常语义而不是错误，末尾不足一块时它小于请求值。把循环写成 ``一次读到底'' 的实现，会在非整除的文件长度上产生一段假静音。

\section{\Lib{PortAudio} 与 \Lib{Cubeb}：回调模型的代价}

文件层可以阻塞，设备层不行。\Lib{PortAudio} 的模型是把音频线程的控制权交给库：调用方注册一个函数指针 \Tp{Pa\Bk Stream\Bk Callback}，库在自己的高优先级线程上按块调用它，参数是输入缓冲指针、输出缓冲指针、\Tp{frames} 与一组状态标志。设备层的全部麻烦都压缩在这个签名的约束里：

\textbf{延迟模式与块选择。} 打开流时给出 \Tp{framesPerBuffer}（$0$ 表示``由驱动决定，每块大小可变''）与期望延迟；\Fn{Pa\_\Bk Get\_\Bk Stream\_\Bk Latency} 报告实际达成值。同一份代码在不同操作系统上会落到不同的后端（Windows 上的 DirectSound 与 WASAPI、macOS 的 CoreAudio、Linux 的 ALSA 与 JACK），块大小与延迟能力差异很大，因此``能跑 128 帧''是不可移植的假设。

\textbf{XRUN 的成因与后果。} 状态标志里的 \Tp{paInputUnderflow}、\Tp{paInputOverflow}、\Tp{paOutputUnderflow}、\Tp{paOutputOverflow} 就是俗称的 XRUN：库需要在这一拍取（或放）数据时，缓冲里的内容却不可用。成因只有三类——回调超支（你的函数比块周期长）、系统抢占（别的实时线程或电源管理插手）、上层消费者把数据搬走后没补上。后果是\emph{一段被替换掉的样本}，听起来是爆音；更要紧的是它让``样本索引''这个时间轴出现无声明的洞，任何依赖帧计数做同步的算法（时间码、锁相、延迟补偿）都会静默错位。所以正确做法是：在回调里把 XRUN 记成计数，块结束后再统计与告警，绝不在回调里重播、补零或试图``修好''。

\textbf{设备枚举与热插拔。} 设备列表是可变的：\Fn{Pa\_\Bk Get\_\Bk Device\_\Bk Count} 与 \Fn{Pa\_\Bk Get\_\Bk Device\_\Bk Info} 给出每个设备支持的采样率、通道数与默认延迟，而 USB 接口与蓝牙耳机的插入拔出会改变它。稳健的写法是缓存 \Tp{PaDeviceID} 与设备名，靠变化通知（或轮询比较）发现设备消失，然后\emph{重开}流而不是在原流上调整——绝大多数实现做不到``无损换设备''。多后端时代，JACK 与 CoreAudio 把重路由做在系统层，于是应用层看到的行为又不同。

\Lib{Cubeb} 是这条路上的另一份答案：它的动机是``浏览器的音频后端''，因此把接口收窄成一条流加两个回调（数据回调与状态回调），设备枚举给出统一的 \Tp{cubeb\_\Bk device\_\Bk collection}（含输入、输出与默认设备，以及变化通知），差异留给了 WASAPI、CoreAudio、ALSA 与 PulseAudio 各自的适配层。对分发者来说它比 \Lib{PortAudio} 更轻，代价是可表达的能力更少（没有全套的延迟与共享模式选项）。

全双工（同一块里既读输入又写输出）还有两个专门的坑。\Lib{PortAudio} 的 \Tp{pa\Bk Prime\Bk Output\Bk Buffers\Bk Using\Bk Stream\Bk Buffer} 决定输出缓冲初值用流缓冲还是清零，不写的话开场可能播出一段未初始化内存；\Tp{pa\Bk Never\Bk Drop\Bk Input} 决定输入快于输出时是否丢输入块，录音与播放速率不匹配时它的默认值会带来看不见的样本缺失。到 ALSA 这一层，XRUN 不是一个标志而是\emph{状态机}：设备进入 \Tp{SND\_\Bk PCM\_\Bk STATE\_\Bk XRUN} 后必须先 \Fn{snd\_\Bk pcm\_\Bk prepare}（或 \Fn{snd\_\Bk pcm\_\Bk recover}）才能继续写，漏掉这一步的话后续每次 \Fn{snd\_\Bk pcm\_\Bk writei} 都返回 \Tp{-EPIPE}，表现是``声音停了但程序还活着''。跨平台库把这些状态差异吸收成标志位，所以只用上层 API 的人常常不知道自己的应用为什么静默失声。

\begin{lstlisting}[style=cstyle,label={lst:ch11-portaudio},caption={PortAudio 回调：增益加 tanh 软限幅，并把 XRUN 记成计数}]
// 依赖：portaudio；链接 -lportaudio
// 回调运行在音频库的实时线程上：不分配、不加锁、不打日志
#include <portaudio.h>
#include <math.h>
#include <stdio.h>
#include <stdint.h>
typedef struct { float gain, prev; uint32_t xruns; } Ctx;

static int onAudio(const void* in, void* out, unsigned long frames,
                   const PaStreamCallbackTimeInfo* ti,
                   PaStreamCallbackFlags flags, void* user) {
  (void)ti;                              // 时间戳只在诊断里用
  Ctx* c = (Ctx*)user;
  const float* src = (const float*)in;
  float* dst = (float*)out;
  if (flags & (paInputUnderflow | paInputOverflow)) ++c->xruns;
  if (flags & (paOutputUnderflow | paOutputOverflow)) ++c->xruns;
  for (unsigned long i = 0; i < frames; ++i) {
    float x = src[i] * c->gain;          // 线性增益，结果可能越过满刻度
    x = tanhf(x);                        // 软限幅：单调，无硬削的破音
    dst[i] = 0.5f * x + 0.5f * c->prev;  // 一阶平滑，去掉拐点突变
    c->prev = x;
  }
  return paContinue;                     // paComplete 会提前收尾
}

int main(void) {
  Pa_Initialize();
  Ctx c = { 0.5f, 0.f, 0u };
  PaStream* st = NULL;
  // 256 帧一块：块大小与采样率一起决定延迟的下界
  PaError e = Pa_OpenDefaultStream(&st, 1, 1, paFloat32, 48000.0,
                                   256, onAudio, &c);
  if (e != paNoError) { puts(Pa_GetErrorText(e)); return 1; }
  Pa_StartStream(st);
  getchar();                             // 主线程只做退出判定
  Pa_StopStream(st); Pa_CloseStream(st); Pa_Terminate();
  return 0;
}
\end{lstlisting}

\section{单头文件与分发，以及离线批处理的分工}

\Lib{miniaudio} 把上面两层合成一个头文件：\Tp{miniaudio.h} 是\emph{单头文件}的 \CL{} 实现（定义宏展开成实现体），内部自带后端分派（WASAPI、CoreAudio、ALSA、OSS 等），不需要外部依赖，发行条款极宽松（以官方文件的声明为准）。它的接口比 \Lib{PortAudio} 高一格：\Tp{ma\_device} 管设备与回调，\Tp{ma\_engine} 在回调之上再给一张节点图（\Tp{ma\_node}），把多路音频的混音、声道转换与时钟适配做在里面。

``无第三方依赖''对分发为什么重要，是一件很实际的事：音频栈一旦要跨平台发布，代价主要不在算法而在\emph{链接}——静态还是动态、符号可见性、每个平台的系统框架怎么声明、安装包要不要额外带运行时。单头文件把这整层消掉了：拷贝一个头、可选地编译一个源文件，没有版本矩阵要对。代价是所有能力都得自带（后端选择、重采样、混音算法都在同一个文件里），升级粒度很粗。

\begin{lstlisting}[style=cstyle,label={lst:ch11-miniaudio},caption={miniaudio：设备与引擎两层，混音与声道转换都在库内}]
// 依赖：把 miniaudio.h 拷进项目；实现宏只在一个 TU 里定义
#define MINIAUDIO_IMPLEMENTATION         // 其余 TU 只要头文件声明
#include "miniaudio.h"
#include <stdio.h>                       // 只有 getchar 用得到

int main(void) {
  ma_engine_config ecfg = ma_engine_config_init();
  ecfg.sampleRate = 48000;               // 引擎时钟，可与设备实值不同
  ecfg.channels = 2;                     // 输出声道数与掩码由后端定
  ma_engine eng;
  if (ma_engine_init(&ecfg, &eng) != MA_SUCCESS) return 1;

  ma_sound snd;                          // 一个声音节点：解码加混音
  ma_result r = ma_sound_init_from_file(&eng, "pad.wav", 0,
                                        NULL, NULL, &snd);
  if (r != MA_SUCCESS) return 1;         // 解码后端缺失时在这里报错
  ma_sound_set_gain(&snd, 0.5f);         // 增益在混音图里，不用自己乘
  ma_sound_start(&snd);                  // 回调与设备线程都由引擎托管
  getchar();                             // 应用侧只剩退出判定
  ma_sound_stop(&snd);
  ma_sound_uninit(&snd);                 // 顺序与 libsndfile 一样：先子后父
  ma_engine_uninit(&eng);
  return 0;
}
\end{lstlisting}

这份代码里没有一个线程是被应用创建的：\Tp{ma\_engine} 自己持有混音线程，把多个 \Tp{ma\_sound} 的缓冲按块累加、做声道数转换与采样率适配，再交给 \Tp{ma\_device} 的回调。这正是``高一格接口''的含义，也是它的风险——时钟适配算法在你看不见的地方，需要严格延迟特性时要换回 \Lib{PortAudio} 那种只给回调的层。

离线批处理里分工不同：

\begin{lstlisting}[style=shellstyle,caption={探测与批量转换：ffmpeg 管容器与编解码，SoX 管音频专用 DSP}]
# 探测：容器与编码归 ffprobe，采样率位深通道归 sox
ffprobe -hide_banner -i input.mkv
sox --i input.wav
sox input.wav -n stats -            # 峰值 dBFS 与 RMS，结果写 stderr

# 解复用与解码：从视频里取音轨，落到 24 位线性 PCM
ffmpeg -i in.mkv -vn -acodec pcm_s24le -f wav tmp.wav

# 重采样到 48 kHz、转 float32、改通道数：SoX 的 rate 是看家能力
sox tmp.wav -r 48000 -e float -b 32 out.wav rate -v -L
# 响度归一到 EBU R128 口径：ffmpeg 有专门的 loudnorm 滤镜
ffmpeg -i out.wav -af loudnorm -f wav prog.wav
\end{lstlisting}

分工的成因是关注点：\Lib{ffmpeg} 的核心资产是解复用器与编解码器（压缩格式、视频轨、时间戳与容器怪癖），它的音频 DSP（\Tp{aresample} 一类的滤镜）够用但不是主场；\Lib{SoX} 的核心资产恰恰是音频专用的变换质量（多档重采样器、噪声整形、\Fn{gain} 的峰值与响度口径、\Fn{stat} 与 \Fn{stats} 的度量、各种特效），但它不碰压缩视频。一条常见的高质量链路是：用 \Lib{ffmpeg} 把任何输入变成干净的高位深线性 PCM，再用 \Lib{SoX} 做重采样与电平处理，最后用 \Lib{ffmpeg} 压回交付格式。批处理脚本用这两个命令行工具而不是写 \CPP{} 程序，是因为它们的表达能力都集中在参数上。

\section{为什么音频域必须由 C/C++ 写，以及 \PY{} 门面的底层对应}

\textbf{理由一：实时线程的三条禁令。} 设备回调所在线程的截止期由声卡时钟决定，与你的调度无关，所以那一段代码里不允许堆分配（拿全局堆锁、可能页错误）、不允许阻塞式 IO 与日志、不允许无界等待的锁。要守住这三条，需要\emph{预分配}缓冲、\emph{有界}队列、显式的无锁双缓冲，还要能把回调线程设成实时优先级（Linux 的 \Tp{SCHED\_FIFO} 与 \Tp{mlockall}，Windows 上的多媒体线程优先级，macOS 上的实时线程属性），并且让状态与设备共享缓冲保持调用方可控的布局。这些能力是语言与运行时层面的，脚本语言给不了。

\textbf{理由二：驱动层只给 C ABI。} 每个平台的音频出口都是 C 形状接口：ALSA 的 \Fn{snd\_\Bk pcm\_\Bk writei}（配合 \Fn{snd\_\Bk pcm\_\Bk prepare} 与恢复路径）、CoreAudio 的一整套 C 函数与属性枚举（\Fn{Audio\Bk Object\Bk Get\Bk Property\Bk Data} 加渲染回调）、WASAPI 的 COM 接口（\Tp{IAudioClient}、\Tp{IAudioRenderClient}，本质是 C 兼容的虚表 ABI）。跨平台的设备层要做的事就是\emph{把这些差异藏进一个 C 接口}，因此任何上层语言想用低延迟音频，都必须先有一个 C/C++ 的实现把三份 ABI 缝起来。\Lib{PortAudio}、\Lib{Cubeb} 与 \Lib{miniaudio} 存在的理由就是这件事。

\textbf{理由三：跨时钟域的连续搬运。} 声卡时钟与 CPU 时钟有实际偏差（消费级器件在几十 ppm 量级），跑上一小时就会积累出可听的漂移或缓冲耗尽。系统的应对是周期性微调：重采样比率、跳样本、或者由 JACK 一类的服务器统一发时钟。这是一段必须每块都执行、且要\emph{无咔哒声}的代码——它的正确性判据是心理声学的，不是单元测试的。分数延迟与异步重采样的算法本身放在本章之后的 FIR 与 IIR 一章，此处只需要认清：这个机制存在于每一层设备代码里，而它跑在实时域。

\textbf{\PY{} 侧门面的底层对照。} \Lib{soundfile} 包装 \Lib{libsndfile}（经由 C 接口层），给出 \Fn{sf.read} 与 \Fn{sf.write} 直接返回 \PY{} 侧的数组，非实时、文件级；\Lib{sounddevice} 包装 \Lib{PortAudio}，把回调包成 \PY{} 可调用对象但只在阻塞式读写里安全好用；\Lib{pyaudio} 是更早的一代 \Lib{PortAudio} 绑定，接口风格更贴近 C。为什么\emph{实时回调}不能交给 \PY{}：解释器回调要从 C 线程拿到全局解释器锁，而锁的等待时间是\emph{无界}的（任何一段正在跑的 \PY{} 字节码、任何一次 GC 都能插在前面）；同时每次调用都要装箱、构造对象，这是被禁的那一类分配。结论是接口形状的正确分工——\PY{} 拥有\emph{块与块之间}的逻辑（打开流、提交缓冲、处理统计），C/C++ 拥有\emph{块之内}的逻辑：

\begin{lstlisting}[style=pythonstyle,caption={两种门面的对照：文件级可以交给 Python，回调级不行}]
import soundfile as sf        # 包装 libsndfile：文件 IO，非实时
import sounddevice as sd     # 包装 PortAudio：实时部分仍在 C 侧
data, sr = sf.read("in.wav", dtype="float32")   # 直接拿到 float32
sd.play(data, samplerate=sr)                    # 拷进 PortAudio 的环形缓冲
sd.wait()                                       # 主线程等待，GIL 无所谓

# 反过来，把块回调交给 Python 就不成立：每块都要拿一次 GIL
def cb(indata, outdata, frames, time_info, flags):
    outdata[:] = indata * 0.5    # 装箱与分配发生在截止期之内，不可控
    if flags.output_underflow:   # 只能在块结束后统计，不能在这里补救
        print("xrun")            # 演示可以，产品里这一句就是故障源

# 正确的分工：C/C++ 跑块回调，Python 只在块之间读它的状态与统计
\end{lstlisting}

\begin{note}[title={电平、削波与 int16 缩放}]
\begin{itemize}
\item \textbf{float32 不等于``天然安全''。} 约定 $[-1,1]$ 只是约定：任何一次求和、增益、滤波器的中间结果都可能越界。定点化（写文件、送设备）时才见真章——越界会被饱和成满刻度（\Lib{libsndfile} 一类的转换默认如此），于是硬削波产生刺耳的高次谐波。内部链路一律用浮点，只在\emph{边界}做限幅，并且优先用软限幅而不是硬截。
\item \textbf{0 dBFS 不等于 ``没有声音到最大''。} dBFS 是以满刻度为参考的对数单位，所有正常信号都是负值；连续电平、采样峰值电平与真峰值（true peak，评估重建后的模拟波形，可能高于采样峰值）是三个不同指标，跨工具比对前先确认是哪一个。
\item \textbf{缩放错误：以为 int16 的范围就是 $-1$ 到 $1$。} 16 位整数的取值范围是 $-32768$ 到 $32767$，\emph{不对称}。写成 \Tp{v / 32767.0f} 时正满刻度得 $1.0$ 而负满刻度得 $1.00003$，越界一次；写成 \Tp{v * 32767.0f} 的反向缩放会把 $-1.0$ 映射到 $-32767$ 而不是 $-32768$，最负的一档永远取不到。可用做法是用 \Tp{v / 32768.0f} 并让正半轴留出一点余量，或者干脆全程浮点。这类错误是听不出来的（差 $3\times 10^{-5}$），但会毁掉逐位比对与回归测试。
\end{itemize}
\end{note}

\begin{table}[H]
\centering
\caption{五个音频库的接口语言、许可类别、实时能力与格式覆盖}
\label{tab:ch11-audio}
\small
\begin{tabular}{@{}>{\raggedright\arraybackslash}p{0.14\textwidth}
  >{\raggedright\arraybackslash}p{0.12\textwidth}
  >{\raggedright\arraybackslash}p{0.15\textwidth}
  >{\raggedright\arraybackslash}p{0.19\textwidth}
  >{\raggedright\arraybackslash}p{0.20\textwidth}@{}}
\toprule
库 & 接口语言 & 许可类别 & 实时能力 & 格式覆盖 \\
\midrule
\Lib{libsndfile} & 纯 \CL{} & LGPL 系 &
  无设备层，只做文件与转换 &
  WAV、AIFF、CAF、FLAC、Ogg（Vorbis 与 Opus）、WAV64 \\
\Lib{miniaudio} & \CL{} 单头文件 & 极宽松（见官方文件） &
  设备回调加引擎混音图 & 自带解码后端，容器支持面较窄 \\
\Lib{PortAudio} & 纯 \CL{} & MIT 系条款 &
  回调模型，跨后端低延迟 & 与格式无关，只处理样本 \\
\Lib{Cubeb} & 纯 \CL{} & ISC 系条款 &
  一条流两个回调，面向浏览器 & 与格式无关，只处理样本 \\
\Lib{SoX} & 命令行加 \CL{} 库 & GPL 系条款 &
  离线批处理，不做实时回调 & 音频容器与编码很广，特效齐全 \\
\bottomrule
\end{tabular}
\end{table}

\begin{bestpractice}[title={表示层与回调的四条纪律}]
\begin{itemize}
\item 内部一律 \Tp{float32}、非交织、每通道独立缓冲；只在读写文件与设备驱动边界做一次交织置换，让置换只出现在一个地方。
\item 容器与编码分开考虑：先问``这个字节流是什么容器''，再问``里面的样本是什么编码''；写文件时显式给子格式而不是依赖默认值，读文件时校验 \Tp{SF\_INFO} 的通道与采样率而不是假设。
\item 回调里的代码要能被三句话审查：分配了几次字节（应为零）、可能阻塞在哪里（应为无）、最坏路径多少周期（应小于块周期）。用计数器和事后统计代替即时打印。
\item 延迟、丢块与时钟漂移要显式建模成运行指标：每流一个 XRUN 计数、一个实测块周期直方图，异常时降级（增大块）而不是硬扛。
\end{itemize}
\end{bestpractice}
% ====================================================================

% ====================================================================
% ====================================================================
\chapter{FFT 与谱分析：从 FFTW 到 cuFFT}

FFT 是信号处理里少数``算法早已定型、接口仍各自为政''的一域。自 Cooley--Tukey 之后没有人再发明更快的通用变换，但每次真正调用它，你仍然要做一串工程决定：缩放归前向还是归后向、复数在内存里怎么摆、变换长度与批量用什么参数表达、plan 在哪个线程上构造。这些决定没有一条来自数学，全部来自 ABI、缓存与许可证。于是这一域的库——\Lib{FFTW}、\Lib{KissFFT}、\Lib{pffft}、\Lib{PocketFFT}、\Lib{oneMKL} DFT、\Lib{cuFFT}、\Lib{rocFFT}、\Lib{VkFFT}——几乎都把 \CL{} 接口当第一等公民：变换的内核是逐点的密集算术，性能全部藏在布局与分块里，而布局与分块要求一套能精确控制内存语义的实现语言；\PY{} 侧的 \Tp{numpy.fft} 与 \Tp{scipy.fft} 则站在这批 C 实现之上做门面。

本章把``调一次 FFT''拆成三层来读：DFT 的公共约定（第 12.1 节）、库的接口形态（第 12.2 与 12.3 节）、谱分析的单位约定（第 12.4 节），最后看一眼门面层（第 12.5 节）。复数数组在数值接口里的一般约定（字节布局、别名规则、线程契约）在前面``数值接口的共同约定''一章已经铺开，本章直接引用那些结论，只谈 FFT 特有的部分。

\section{DFT 的三条公共约定：缩放、半谱与复数布局}

\textbf{缩放放在哪。} $N$ 点 DFT 的定义只有一行，$X(k) = \sum_{n=0}^{N-1} x(n)\, e^{-2\pi i k n / N}$，但把它变成 API 时第一个分歧就是逆变换的 $1/N$ 归谁管。实务上有三种约定：两向各除 $\sqrt{N}$ 的酉形式（分析文献偏爱，Parseval 恒等式干净）；前向不缩放、逆变换除 $N$（\Tp{numpy.fft} 的默认，另有 \Tp{norm} 参数可换档）；两向都不缩放（\Lib{FFTW} 与 \Lib{cuFFT}）。最后一种是 \CL{} 接口的现实选择：逆向那次除法在 SIMD 下是货真价实的 $N$ 次乘法，把它从通用实现里剥出来、让调用方合并进后续增益，比在每一层库都除一遍划算。后果必须写进代码注释：在 \Lib{FFTW} 里正变换接逆变换跑回来，得到的是原信号的 $N$ 倍——这不是 bug，是缩放约定；跨库搬数据（比如把 \Lib{FFTW} 的谱喂给别的逆变换器）时，$N$ 倍的静默能量差就是这么来的。

\textbf{半谱约定。} 实输入满足 Hermitian 对称 $X(N-k) = \overline{X(k)}$（下标按 $N$ 取模；$k = 0$ 与 Nyquist 档自共轭、必为实数），独立信息只有 $\lfloor N/2 \rfloor + 1$ 个复数档。r2c 接口返回这个长度不是省内存的偏方，而是信息论下界；对应的 c2r 接口也只消费这么多输入，另一半由实现按对称性重建——这就是``半谱重建''的全部含义。注意两个方向对缓冲长度的要求不对称：偶数长度时两边库通常都按 $N+2$ 个实数分配输出（多出的填充保证逆变换能原地跑、并保持对齐，具体以文档为准），逻辑档数是 $\lfloor N/2 \rfloor + 1$，物理字节数比它多一格。半谱约定埋的坑在能量侧：写功率谱时，除 DC 与 Nyquist 外的中间档把负频率的能量藏在共轭镜像里，要么乘 $2$ 声明自己是单边谱，要么明确声明自己是双边密度的一半——两种写法在真实项目里都常见，混用的直接后果是所有 dB 数值差 $6$ dB 整。

\textbf{还有半复数（halfcomplex）打包。} 老一代数值库（\FTN{} 血统尤其明显）里存在第三种约定：把实序列的半谱压回 $N$ 个\emph{实数}——先 DC、再 Nyquist、然后逐档 re 与 im 交替。它不新增信息，只是把 $\lfloor N/2 \rfloor + 1$ 个复数恰好摊平成 $N$ 个实数的紧凑排布。\Lib{FFTW} 的通用 r2r 接口能直接产出与消费这个格式（\Tp{FFTW\_R2HC} 与 \Tp{FFTW\_HC2R} 两个 kind），同一接口下还有离散 Hartley 变换与一族 DCT、DST——它们全是``只用实档''的变换，省掉一半档位与一半乘法。在旧代码或别的库传来的数据里遇到``长度没变、谱密度看着翻倍''的结果，先问一句它给的是半复数还是第 12.1 节的交错复数：两者是同一件事的两种观点，差一步重排，不少一个系数。

\textbf{两种复数布局。} 交错（AoS：$[\mathrm{re},\mathrm{im}]$ 相邻，\Tp{fftwf\_complex}、\Tp{cufftComplex} 与 C 的 \Tp{double\_Complex} 都是它）与分离（SoA：实部一整段、虚部一整段）。交错贴 \CL{} 与 \FTN{} 的复数数组传统，跨语言搬运只动一个指针，一个元素内部局部性极好；代价在 SIMD 上——复数乘加要把实虚部分到不同通道，一条 FMA 前面往往垫着反交错的 shuffle，或者改用``两条 lane 各算一半''的横向方案。分离布局正相反：实部与虚部各自是连续浮点流，复数乘法是两组纯向量 FMA，向量化几乎免费；代价是两条独立访存流（大数组下两倍的预取与 TLB 压力），以及没有标准的 C 类型能直接表达它。所以 \Lib{FFTW}、\Lib{pffft}、\Lib{cuFFT} 对外只有交错，而 \Lib{Kokkos} 系的 FFT 把两种布局都做进接口（复数元素相邻的 layout\_left 与实虚分离的 layout\_right 二选一，且不支持原地变换）——这个差异不是口味，是``谁的热点在乘法、谁的热点在搬运''的立场差异。

\textbf{原地与别名。} 交错布局还带来一个 FFT 特有的问题：正逆变换能不能在同一块缓冲上跑。\Lib{FFTW} 的 c2c 支持原地模式（输入指针即输出指针，内部自己安排乒乓），代价是它替你多算一份缓冲；\Lib{pffft} 明确不保证原地；\Lib{cuFFT} 的多数变换就地可行，但 r2c 与 c2r 的原地对偶要求缓冲按 $N+2$ 个实数解释而不是 $\lfloor N/2\rfloor+1$ 个复数——两个解释在偶数长度下差一个元素的间距。``原地''从来不是一个布尔参数，而是一组关于读写冲突与缓冲形状的合同，读任何 FFT 库文档时这是值得先看的一节。四条约定的对照速查见表~\ref{tab:ch12-conventions}。

\begin{table}[H]
\centering
\caption{缩放与半谱打包约定速查（物理尺寸按偶数长度，以所装版本文档为准）}
\label{tab:ch12-conventions}
\small
\begin{tabular}{@{}>{\raggedright\arraybackslash}p{0.16\textwidth}
  >{\raggedright\arraybackslash}p{0.15\textwidth}
  >{\raggedright\arraybackslash}p{0.17\textwidth}
  >{\raggedright\arraybackslash}p{0.24\textwidth}@{}}
\toprule
实现 & 前向缩放 & 逆向缩放 & r2c 输出（逻辑档数、物理实数） \\
\midrule
\Lib{FFTW} & 无 & 无（$N$ 倍自负） &
  $\lfloor N/2 \rfloor + 1$ 档、$N+2$ 实数 \\
\Lib{cuFFT} & 无 & 无（$N$ 倍自负） &
  $\lfloor N/2 \rfloor + 1$ 档、$N+2$ 实数 \\
\Tp{numpy.fft} & 无 & 除 $N$ &
  $\lfloor N/2 \rfloor + 1$ 档、门面直接给复数数组 \\
\Tp{scipy.fft} 酉档（norm=ortho） & 各除 $\sqrt{N}$ & 各除 $\sqrt{N}$ &
  同上，能量恒等式干净 \\
\bottomrule
\end{tabular}
\end{table}

\section{\Lib{FFTW}：为什么 plan 本身就是接口}

\Lib{FFTW} 的接口形态可以一句话概括：把``知道你要做什么''从每次执行提前到一次性建划。一个最小可用流程是：用 \Tp{fftwf\Bk\_alloc\Bk\_real} 与 \Tp{fftwf\Bk\_alloc\Bk\_complex} 分配缓冲（不是 \Fn{malloc}，见清单~\ref{lst:ch12-fftw} 的讨论），调 \Tp{fftwf\Bk\_plan\Bk\_dft\Bk\_r2c\Bk\_1d} 得到 plan，反复 \Fn{fftwf\_execute}，最后 \Fn{fftwf\_destroy\_plan}。plan 里编译进去的东西包括：这个长度的素因子分解该用哪套 radix 组合、缓存分块的尺寸、要不要多线程、SIMD 路径选哪条。换长度、换缓冲地址、换对齐状态，都要换 plan。

\begin{lstlisting}[style=cstyle,label={lst:ch12-fftw},caption={FFTW：r2c 单边幅度谱，plan 与缓冲的同生关系}]
// 依赖：fftw3f；链接 -lfftw3f。
// 规则一：in 与 sp 用 fftwf_alloc_* 分配——plan 按缓冲
//          的对齐做优化，malloc 的内存能跑但路径退化。
// 规则二：同一 plan 的缓冲不得并发复用。
#include <fftw3.h>
#include <math.h>
#include <stdio.h>

#define N  1024                  // 实输入长度
#define HB (N / 2 + 1)           // 半谱档数：floor(N/2)+1
#define FS 48000.0               // 采样率，Hz

int main(void) {
  float* in = fftwf_alloc_real(N);
  fftwf_complex* sp = fftwf_alloc_complex(HB);
  for (int n = 0; n < N; ++n)    // 测试信号：1 kHz 正弦加直流
    in[n] = 1.0f + 0.5f*sinf(2.f*(float)M_PI*1000.f*n/FS);

  fftwf_plan p = fftwf_plan_dft_r2c_1d(N, in, sp, FFTW_ESTIMATE);
  fftwf_execute(p);              // 正变换：两端都不缩放

  double cg = 1.0;               // 相干增益：无窗取 1，hann 是 0.5
  for (int k = 0; k < HB; ++k) {
    double a = hypot(sp[k][0], sp[k][1]);
    double w = (k == 0 || 2 * k == N) ? 1.0 : 2.0; // 中间档乘二
    double amp = w * a / (N * cg);  // 单边幅度谱：除 N 与窗增益
    if (amp < 0.01) continue;       // 只打印显著谱峰
    printf("%8.1f Hz  %.4f\n", k * FS / N, amp);
  }
  fftwf_destroy_plan(p);            // plan 是堆对象：忘了就是泄漏
  fftwf_free(in); fftwf_free(sp);   // 与 fftwf_alloc_* 配对释放
  return 0;
}
\end{lstlisting}

\textbf{建划的时间代价与 wisdom。} plan 不是数据结构壳，是一次真的求解器运行。旗标决定它花多少时间：\Tp{FFTW\_ESTIMATE} 只做解析估算，微秒级，但选出的计划通常慢于实测最优；\Tp{FFTW\_MEASURE} 在调用方的缓冲上实际跑候选内核计时，第一次执行会污染缓冲，所以建划前要把输入当真数据处理；\Tp{FFTW\_PATIENT} 与 \Tp{FFTW\_EXHAUSTIVE} 继续扩大候选的 radix 与分块组合，代价可以到数百毫秒。为了把这笔一次性开销从产品启动路径里挪走，FFTW 提供 wisdom 机制：\Fn{fftwf\_export\_wisdom} 把``某类问题对应某套计划''的知识序列化成文本，随发行物落盘，下次启动 \Fn{fftwf\_import\_wisdom} 之后，配合 \Tp{FFTW\_WISDOM\_ONLY} 旗标让建划退化成查表——查不到直接失败而不是现场重算。wisdom 按精度（单、双、长双）分池，缓存的是执行机的实测结果，跨 CPU 代次复用是收益递减而不是错误；进程退出前 \Fn{fftwf\_cleanup} 释放它的全局资源，对库的使用方是礼貌，对短命进程是省一次退出时的遍历。

\textbf{planner 的线程安全。} wisdom 与 planner 内部状态是进程级全局的：在两个线程上并发构造 plan 是未定义行为。出路有三条，按侵入性递增：configure 时用 \Tp{FFTW\_LOCK} 挂一把 pthread 互斥量（锁住所有建划入口，代价是建划串行化）；较新版本提供的 planner 线程安全开关（名字形如 \Tp{fftwf\Bk\_make\Bk\_planner\Bk\_thread\Bk\_safe}，以所装版本的头文件为准）；或者干脆按最佳实践组织代码——启动期单线程把所有 plan 建完，运行期只调 \Fn{fftwf\_execute}。最后这条不只是避险：plan 的执行允许多线程并发，前提是各线程持有自己的一套输入输出缓冲。执行线程池是另一件事：\Fn{fftwf\_init\_threads} 加 \Fn{fftwf\_plan\_with\_nthreads} 让单个 plan 内部并行，线程数的约定与 OpenMP 一类运行时可能冲突（双重并行是最常见的自伤），批处理场景下``一个 plan 一条线程''往往优于``一个 plan 八条线程''。

\textbf{批处理、stride 与``离散频率''。} 真实负载很少是一次单条变换：多通道音频要同时变换 $C$ 条等长序列，矩阵要对每一列做变换，张量要对某一模做变换。\Tp{fftwf\Bk\_plan\Bk\_many\Bk\_dft\_r2c} 用一个调用把这些全表达出来：\Tp{howmany} 是变换条数，\Tp{idist} 与 \Tp{odist} 是相邻两条之间的元素距离，\Tp{istride} 与 \Tp{ostride} 声明单条变换内部元素的跨步——对列变换这就是 istride 等于行长的转置访问，省掉一次物化拷贝，planner 会按这个访问模式挑内核。换句话说，stride 不是给调用方的``便利参数''，而是把访存模式交给优化器的契约。至于只要部分频率档（非连续的离散频率样本，例如任意频点的 chirp 或缩放傅里叶），FFTW 不提供按档抽取的接口，那是 chirp-z 机制的事（见第 12.4 节末），不要用``全量变换再筛''假装它是批处理。

\textbf{长度的选择。} 变换长度直接决定成本形状：素因子全落在 $2$ 到 $7$ 这一档的长度（``光滑数''）走 radix 分解的主干道，$2$ 的幂因为递归均匀、分块对齐而最顺；大素数长度 FFTW 会落到 Bluestein 一类的卷积化路径，仍是 $O(N \log N)$ 量级但常数明显更大。音频帧长于是有了工程惯例：$1024$、$2048$ 是安全的，$1000$ 这种``整数好看''的长度其实更贵（\Lib{KissFFT} 一类的混合 radix 对它友好，见第 12.3 节）。需要任意长 FFT 时的标准动作是补零到下一个快长度，或者用 \Tp{scipy.fft\Bk.next\Bk\_fast\Bk\_len} 一类的助手函数问一嘴——注意补零改变的是频域采样密度，不是分辨率。

\textbf{许可证与替代者的由来。} \Lib{FFTW} 的主体是 GPL 系条款、由版权方同时出售商业授权的双许可模式。这意味着任何链接它的程序（无论静态还是动态）都落在传染性条款的射程内：要么采用兼容的开源许可并履行源代码义务，要么购买商业授权。静态链接的发布尤其别扭——产物里直接含库的目标代码，义务无法用``用户自行安装''绕开。具体义务以随仓库分发的 LICENSE 与授权说明为准，本书不替代法务。正因为这笔成本，``功能等价、许可宽松''的替代实现是生态刚需，下一节那张表就是替代品地图。

\section{生态位：可嵌入、可向量化、可异构}

替代品各自回答一个不同的问题。\Lib{KissFFT} 回答``能不能只有两个 .c 文件''：任意长度（含素数，内部做混合 radix 分解），配置结构 \Tp{kiss\_fftr\_cfg} 承担 plan 的角色但没有 wisdom 与全局状态，分配器可以在编译期整体换掉（嵌入式常用自带内存池接进去），另有定点与整型构建——标量类型随编译开关改变，接口形状不变，具体名称以头文件为准。BSD 类条款，塞进任何许可的产物都没有心理负担。代价是速度：它的计划质量与 \Lib{FFTW} 的 MEASURE 档有差距，适合``够用且必须小''的场合。

\Lib{pffft} 回答``没有 planner 能不能贴着 SIMD 写''：一份为 SSE、AVX 与 NEON 手写内核的库，setup 是一次廉价的运行时查表，以实测著称（对 \Lib{FFTW} 的领先幅度随长度与机器变化，这里不给倍数），接口是裸指针加交错布局、不保证原地。\Lib{PocketFFT} 回答``头文件 \CPP{} 行不行''：纯头、BSD 类条款，接口小到一个函数，它也是 \PY{} 世界的事实底座之一（见第 12.5 节）。

批处理与异构这一侧：\Lib{oneMKL} 的 DFT 部分用描述符模式（创建、commit 成计划、execute 执行三步，入口统一是 \Tp{Dfti} 前缀），线程域独立可调，覆盖 CPU 上的单双精度与批处理。\Lib{cuFFT} 把 FFTW 的 plan 概念原样搬上显卡——\Fn{cufftPlan1d} 与 \Fn{cufftPlanMany} 建划、工作缓冲大小先查询再分配、按精度分执行入口（实入复出的单精度正向即 \Fn{cufftExecR2C}，各精度入口命名以 \Tp{cufft.h} 为准），计划与 stream 及设备指针语义绑死；它的回调机制允许在进出变换时融合一次逐点乘法，省一趟全数组往返。\Lib{rocFFT} 是 AMD 侧的镜像设计，把 plan 描述与执行拆开、精度作为参数。再往外，\Lib{VkFFT} 把变换整个写进 Vulkan compute shader，换来跨厂商显卡与``输出是着色器源码''的可移植性，代价是工作缓冲与数值细节都要调用方自己盯；\Lib{Kokkos} 系与 SYCL 一类的性能可移植层则回答另一个问题——同一份调用源码在 CPU、CUDA、SYCL 后端各自落到本地最好的实现上。它们的分工不是``谁更快''，而是``你的访存模型与分发约束允许你付哪笔钱''。

\textbf{精度与误差增长。} \Lib{FFTW} 一族按标量精度发三份接口（\Tp{fftw\_} 双精度、\Tp{fftwf\_} 单精度、\Tp{fftwl\_} 长双精度），wisdom 也按池隔离；\Lib{cuFFT} 同样以双单精度分入口，更低的位宽只在回调或部分路径上出现（细节以文档为准）。变换的相对误差随长度近似对数增长：悲观界按 $\epsilon \log N$ 的量级走，实测更接近 $\epsilon \sqrt{\log N}$（数值分析的经典结论，不是 FFT 特有的）——也就是说 $65536$ 点相对 $256$ 点的误差只涨几倍，float 的七位有效数字对谱显示、对多数特征提取绰绰有余。真正需要双精度的从来不是``谱的精度''而是``差的精度''：谱减、维纳增益、级联除法与长窗归一化这些``两个相近数相减''的路径，才是精度预算该花的地方。

\textbf{为什么 FFT 必须是 \CL{}/\CPP{}。} 把上面整张候选集写成 \PY{}，会得到数学相同、慢一到两个数量级的东西。四条机制，每条都能落到一个具体调用上：

\begin{itemize}
\item \textbf{逐点算术密度。} 一个蝶形的有效信息约十个浮点操作，任何``每元素装箱或解释''的实现都被元操作吃掉。只有 C 的内核函数或内联模板能把成本压到访存下界——这也是 \Tp{numpy.fft} 宁可维护一份 C++ 的 \Lib{PocketFFT} 也不在数组对象上写循环的原因（第 12.5 节）。
\item \textbf{缓存分块与原地性。} ``把问题切成块''在 FFT 里不是语义等价就能完工的事：块尺寸由 L1 容量决定，乒乓缓冲管理读写冲突，原地变换的调度要精确到``哪个子块先动''。这是内存语义问题，\Tp{restrict}、显式对齐与手工生命周期就是为此存在。
\item \textbf{不规则步长的访存模式。} Bluestein chirp-z 把任意长度折成一次卷积，中间要按位反转序与 chirp 调制做非连续读写；\Lib{Kokkos} 的 layout 选择同样在决定这类跨步模式。Python 里的花式索引每次物化一份拷贝，C 里它就是一个下标表达式。
\item \textbf{SIMD 复数乘加与指令调度。} 带 FMA 的复数乘法、横向归约与 shuffle 调度，要么用 intrinsics 写（\Lib{pffft}），要么交给可验证的自动向量化并盯着报告。带 GC 与装箱的运行时里没有这件事的位置。
\end{itemize}

最后才是那张老面孔：以上所有实现互相之间、以及和上层所有语言之间，靠 \CL{} ABI 汇合。\FTN{} 在批量作业与向量硬件的传统实现里仍占一席（\Lib{oneMKL} 的部分内核历史上就是 \FTN{}），但它同样要先穿上 C 的符号外衣才能被别人调到。

\begin{longtable}{@{}>{\raggedright\arraybackslash}p{0.15\textwidth}
  >{\raggedright\arraybackslash}p{0.11\textwidth}
  >{\raggedright\arraybackslash}p{0.12\textwidth}
  >{\raggedright\arraybackslash}p{0.12\textwidth}
  >{\raggedright\arraybackslash}p{0.16\textwidth}
  >{\raggedright\arraybackslash}p{0.11\textwidth}@{}}
\caption{五个 FFT 库：接口语言、许可类别、plan 语义、多维与批处理、GPU}\label{tab:ch12-libs}\\
\toprule
库 & 接口语言 & 许可类别 & plan 语义 & 多维、批处理与 stride & GPU \\
\midrule
\endfirsthead
\toprule
库 & 接口语言 & 许可类别 & plan 语义 & 多维、批处理与 stride & GPU \\
\midrule
\endhead
\bottomrule
\endfoot
\Lib{FFTW} & 纯 \CL{} & GPL 系，双许可 &
  显式 plan，wisdom 缓存，planner 需锁 &
  rank 到三维，plan\_many 批处理，stride 语义最全 & 无 \\
\Lib{KissFFT} & \CL{} & BSD 类 &
  cfg 结构，无全局 wisdom &
  一维任意长，多维靠调用方套循环 & 无 \\
\Lib{pffft} & \CL{} & BSD 类 &
  无 plan，setup 廉价 &
  一维为主，批处理靠循环调用 & 无 \\
\Lib{PocketFFT} & \CPP{} 头文件 & BSD 类 &
  无显式 plan，构造时决定 &
  一维与少量多维入口，接口极简 & 无 \\
\Lib{cuFFT} & \CL{} API & 专有，随 CUDA 分发 &
  显式 plan，工作缓冲先查再配 &
  一维至三维，PlanMany 加 stride，回调融合 & NVIDIA \\
\end{longtable}

\section{谱分析的公共约定：窗、STFT 与单位}

\textbf{频率轴在哪。} 库还回去的数组顺序是 $0$ 到 $N-1$ 的档：前一半是 $0$ 到 Nyquist 的正频率，后一半是负频率的共轭镜像。绘图工具链里那个把零频搬到中间的 \Tp{fftshift} 是门面层的便利函数，C 接口没有它——把``第 $N-1$ 档''当 $+\Delta f$ 读，是把频谱图上下颠倒的经典成因。r2c 半谱之所以没有这个问题，正是因为它只保留了不需要搬的那一半。

\textbf{窗函数与相干增益。} 对有限段做变换等于对真信号乘一个窗：矩形窗旁瓣高（约 $-13$ dB 的第一旁瓣）、泄漏重，但有最窄的主瓣；hann 与 hamming 把旁瓣压低、用主瓣变宽来付账，blackman 再把旁瓣压到约 $-58$ dB，过渡带进一步变宽。量化窗对幅度的影响用\emph{相干增益}表达：$\mathrm{CG} = \frac{1}{N}\sum_n w_n$，hann 是 $0.5$、hamming 是 $0.54$、blackman 约 $0.42$——同一正弦在加窗后的谱峰被压到 $\mathrm{CG}$ 倍，幅度谱要除以它才还原真实电平。另一笔账在噪声上：等效噪声带宽 $\mathrm{ENBW} = N\sum_n w_n^2 / (\sum_n w_n)^2$（单位是 bin，hann 为 $1.5$），窗越``瘦''、每 bin 收进的噪声功率越多。两笔账都精确到 $\mathcal{O}(1/N)$，短窗上会有百分之几的偏差。

\textbf{幅度谱还是密度。} 同一个 $\lvert X(k)\rvert$ 至少服务三种消费者：\emph{幅度谱}（找基频、找杂散），缩放取 $\frac{2\,\lvert X_k\rvert}{N\,\mathrm{CG}}$，中间档乘 $2$ 补偿镜像，纵轴直接读成物理幅度；\emph{功率谱密度}（比较噪声、验证滤波器），纵轴必须是``每 Hz''，周期图形式为 $S_k = \frac{2\,\lvert X_k\rvert^2}{f_s \sum_n w_n^2}$，其中除以 $f_s$ 这一笔就是``每 bin 还是每 Hz''的分水岭——有的工具纵轴写 $\mathrm{V^2/Hz}$，有的写 $\mathrm{V^2/bin}$，相差一个因子 $f_s/N$；\emph{频谱图}（给眼睛看的 dB 颜色），通常再叠一次对数与颜色映射的动态范围裁剪。跨工具比对前先问一句纵轴单位，是这一节的全部方法论。

\textbf{STFT 的帧移与重叠。} 连续短时傅里叶变换是``窗滑动加 FFT''，帧移 $R$ 决定两件事：时间分辨率 $R/f_s$ 秒，与频响网格 $f_s/N$ Hz 之外的插值密度。$R$ 与窗长 $N$ 的关系用重叠率表达：hann 窗取 $50\%$ 重叠时 $w^2$ 的滑动和恒为常数（加窗再平方窗的 WOLA 重构成立），$75\%$ 是相位精确重建（OLA）的常用值；矩形窗则 $0\%$ 就够——``重叠是补窗的漏''这句话对 hann 成立、对矩形不成立。$50\%$ 与 $75\%$ 之间没有免费的中间档：跳帧做特征提取（Mel 谱一类）时十毫秒级的移是常态，那是统计意义而不是重构意义，别拿它去做逆变换。参数对照举一例：$f_s = 48$ kHz、$N = 1024$、$R = 256$（$75\%$ 重叠）时，窗长 $21.3$ ms、频响网格 $46.9$ Hz、帧网格 $5.3$ ms——三个数写进设计文档，谱图的纵横比例就已经定了。

\textbf{长信号：overlap-add 与 overlap-save。} 把长滤波做成分块 FFT 卷积有两条等价的路线。\Tp{overlap-add}：每块 $M$ 样本补零到 $N \ge M + L - 1$（$L$ 为冲激长度），变换、乘频响、逆变换，输出长 $N$，相邻块重叠的 $L-1$ 个点\emph{相加}。\Tp{overlap-save}：不补零、直接对含 $L-1$ 前导的 $N$ 样本做循环卷积，丢掉前 $L-1$ 个混叠点、保留后 $M$ 个。两者计算量相同、同一实现下结果逐位同，选哪个常取决于缓冲管理。工程账：块长 $M$ 每加一倍，每样本算力降（FFT 超线性），但延迟至少多 $M/f_s$ 秒；前几个块的边界还会因补零产生非循环卷积边界——处理正确时它是伪影的唯一来源，处理不正确（忘加尾、忘丢头）时它就是咔哒声。非线性谱处理（谱减、维纳增益、掩蔽）根本不满足卷积定理，分块边界必须靠分析合成侧的交叉重叠糊过去，这是``分块 FFT 伪影''的另一个产地。

\textbf{两个实序列塞一次复 FFT。} 给两条实数据 $a$、$b$（双通道、或一帧的前后两半），令 $z = a + i\,b$ 做一次复变换 $Z$，再由对称性拆出两条各自的谱：$A_k = \frac{Z_k + \overline{Z_{N-k}}}{2}$、$B_k = \frac{Z_k - \overline{Z_{N-k}}}{2i}$。一次 $N$ 点复变换顶两次实变换（复变换的成本里虚部那一半本来就接近零头），这是 r2c 接口内部实际在做的事的显式版。反过来，拿到半谱要恢复全谱，就按 $X_{N-k} = \overline{X_k}$ 把镜像补回去再喂 c2r。

\begin{note}[title={单位金字塔：从 bin 到每 Hz 再到 dB}]
\begin{itemize}
\item \textbf{先定分母再画曲线。} ``除以 $N$''（幅度谱）、``除以 $f_s$ 乘窗能量''（每 Hz 密度）、``什么都不除''（每 bin）是三套互相差 $f_s$、$N$ 或 $\mathrm{CG}$ 因子的约定。一条 FFT 曲线被两条产品线复用而没对单位，最常见的症状就是整体上下平移一个与采样率有关的常数。
\item \textbf{dB 参考要跟着单位走。} $\mathrm{dBFS}$ 对幅度谱、$\mathrm{dB/Hz}$ 对密度谱，两者差 $10\lg(f_s/N)$ 一档。拿宽带噪声的 $\mathrm{dBFS}$ 值去对标另一个采样率下的工具，平移是必然的。
\item \textbf{峰值读数带窗修正。} 用谱峰读正弦幅度时，只有信号精确落在 bin 中心才读得准；否则即使除以 $\mathrm{CG}$ 仍有扇贝损失（hann 上界约 $1.4$ dB）。精读幅度要插值三条谱线，或者接受误差并写明窗型。
\end{itemize}
\end{note}

\begin{warning}[title={对齐、所有权与缩放：FFT 最贵的三个坑}]
\begin{itemize}
\item \textbf{plan 与缓冲的生命周期错配。} plan 持有的是\emph{构造时}那对指针：先释放缓冲再 \Fn{fftwf\_execute} 是悬垂读写；换了一块内存却不重建 plan，则是把优化建立在错误假设上。\Lib{cuFFT} 还要再错配一层——工作区是 plan 之外的第二份资源，忘了查 \Fn{cufftGetSize} 一类入口给出的大小并配足，得到的是含义模糊的库内错误码。正确顺序固定为：分配、建划、执行、销毁、释放。
\item \textbf{对齐错配是静默的。} 拿 \Fn{malloc} 的缓冲去建 plan，FFTW 不会拒绝，只会悄悄放弃对齐依赖的 SIMD 路径；性能回退没有任何日志。反过来把``只有 $\mathrm{HB}$ 个复数''的缓冲交给偶数长度 r2c 当输出（物理上需要 $N+2$ 个实数），越界写就会以``隔了两天别处崩''的形式出现。
\item \textbf{缩放错是能量错，不是形状错。} 两端都不缩放的库（\Lib{FFTW}、\Lib{cuFFT}）里漏除 $N$，波形看起来完全正常，只有能量差 $N^2$：正逆各漏一次，白噪声底抬升 $20\lg N$ dB。回归手段是 Parseval 恒等式：$\sum \lvert x \rvert^2$ 与 $\sum \lvert X \rvert^2 / N$ 应当对得上，这条断言写进测试比任何目测都便宜。
\end{itemize}
\end{warning}

\section{\PY{} 侧对照：\Tp{numpy.fft} 与 \Tp{scipy.fft} 的底层}

\Tp{numpy.fft} 的模块文档里``多算法混合实现''已经让位于一个 C++ 内核：自 1.17 系列起逐步切换、并在 1.20 起以 \Lib{PocketFFT}（头文件名 \Tp{pocketfft} 系）为默认后端，任意长度与素数走混合 radix，缩放约定是``前向不缩放、逆向除 $N$''（\Tp{norm} 参数可切酉与反向）。也就是说，第 12.1 节那三条约定在 \PY{} 侧同样成立，只是门面替你选了默认档。\Tp{scipy.fft} 重写了同一批函数为 ufunc 友好的形态，加了线程参数与\emph{可插拔后端}机制：注册之后，调用可以路由到 \Lib{FFTW}、\Lib{oneMKL} 或 \Lib{cuFFT} 的包装（具体支持矩阵与注册 API 随版本演进，以 SciPy 文档为准）。门面与引擎的分工在第 12.2 节的约束下是必然的：plan 构造与执行绑定的是裸缓冲与线程契约，而解释器侧每次调用都可能重新装箱缓冲——这正是``C 做 ABI 汇合点、上层做门面''模式在 FFT 上的体现。

\begin{lstlisting}[style=pythonstyle,caption={numpy.fft 与 scipy.fft：同一半谱，不同缩放档与后端}]
import numpy as np
import scipy.fft as sf

x = np.sin(2 * np.pi * 21 * np.arange(1024) / 1024)
a = np.fft.rfft(x)              # 底层是 C++ 的 pocketfft 内核
b = sf.rfft(x)                  # 同一数学，默认缩放也相同
c = sf.rfft(x, norm="ortho")    # 两向各除 sqrt(N)：酉约定
d = sf.rfft(x, workers=4)       # scipy.fft 自己的并行开关；若注册
                                # 了外部后端，这一句可路由到别的 C 库
print(np.abs(a - b).max())      # 两门面一致：机器精度级
print(np.abs(a / np.sqrt(1024) - c).max())   # 差的全局因子 sqrt(N)
\end{lstlisting}

\begin{lstlisting}[style=cstyle,label={lst:ch12-kiss},caption={KissFFT：无 wisdom、无对齐魔法的嵌入式路线（浮点版）}]
// 依赖：kissfft（BSD 类条款，实现就两个 .c 文件）。
// 对照 FFTW：cfg 是唯一的"计划"对象；没有全局 wisdom、
// 没有对齐要求；alloc 支持两趟式调用，自带内存可接入。
#include <kiss_fftr.h>
#include <math.h>
#include <stdio.h>

#define N  1000               // 任意长度：1000 = 2^3 * 5^3 不补零
#define HB (N / 2 + 1)        // 半谱约定与 FFTW 完全同形

static float in[N];
static kiss_fft_cpx sp[HB];   // re/im 交错；定点构建同形换标量

int main(void) {
  for (int n = 0; n < N; ++n)         // 激励精确落在第 21 档
    in[n] = sinf(2.f * (float)M_PI * 21.f * n / N);

  kiss_fftr_cfg cfg = kiss_fftr_alloc(N, 0, NULL, NULL);
  if (!cfg) return 1;                 // 分配失败（或长度非法）
  kiss_fftr(cfg, in, sp);             // 正变换：不缩放，同 FFTW
  kiss_fftri(cfg, sp, in);            // 逆变换：结果含 N 倍因子

  double e = 0;                       // 往返残差：验证半谱重建闭合
  for (int n = 0; n < N; ++n) {
    float back = in[n] / N;           // 自己除回去：缩放约定的课
    e += (back - sinf(2.f*(float)M_PI*21.f*n/N))
        * (back - sinf(2.f*(float)M_PI*21.f*n/N));
  }
  printf("roundtrip %.3e\n", e);      // 单精度下约 1e-5 量级
  kiss_fftr_free(cfg);
  return 0;
}
\end{lstlisting}

清单~\ref{lst:ch12-kiss} 与清单~\ref{lst:ch12-fftw} 放在一起读，能把本章收拢成一句话：两个实现共享同一套数学约定（半谱、缩放、交错布局），差异全部在``计划''的重量与``内存''的所有权——FFTW 把优化买断在 plan 里、代价是全局状态与对齐契约；KissFFT 把这些全部退还给调用方、代价是速度。

\begin{guideline}[title={本章要点}]
\begin{itemize}
\item 缩放约定先问库、再写注释；跨库链路上，$N$ 倍能量差是最安静的错误。
\item 半谱是 $\lfloor N/2 \rfloor + 1$ 个复数档、$N+2$ 个实数字节；单边谱乘 $2$ 只乘中间档。
\item plan 是带前提的合同：缓冲地址、对齐、stride、线程都算前提，变了就重建。
\item 选实现先选许可证与部署形态，再谈速度：FFTW 买计划质量，KissFFT 买体积，cuFFT 买吞吐。
\item 每个谱数值都要带着单位出生：除以 $N \times \mathrm{CG}$ 的是幅度，除以 $f_s \sum w^2$ 的是密度。
\item 批处理优先于线程池：先用 plan\_many 把访存合同立对，再考虑加线程；双重并行是性能报表骗人的头号来源。
\end{itemize}
\end{guideline}

而在 \EMW{} 那一部分里，同一批库会换一种身份出现：谱方法求解器每个时间步调用的是同一类 plan 与同样的交错布局，只是半谱约定让位于全复数的多维变换——主线不变，变的只是维度。
% ====================================================================

% ====================================================================
% ====================================================================
\chapter{滤波器设计与嵌入式 DSP 库}

如果说 FFT 是信号处理的``观测工具''，滤波器就是它的``肌肉''：音频链路里每一次均衡、分频、重采样、降噪，最后都落到一个每样本几十到几百次乘加的状态更新上。这个域有一条容易被忽略的分界线：\emph{设计}与\emph{执行}是两种工作负载。设计是离线的、迭代式的、可以容忍秒级响应——它产出一小摞系数；执行是线上的、每秒几万拍的、预算以指令周期计——它只消费系数并推进状态。分界线两侧的语言选择完全不同：设计侧 \PY{}（\Tp{scipy.signal}）已经是事实标准，执行侧则从桌面 \CPP{} 一路下沉到没有操作系统的单片机上的 \CL{}。

本章先把执行结构（第 13.1 节，FIR 与 IIR 的几种形态及其数值病）和设计方法（第 13.2 节，窗法、Remez、双线性变换）摊开，再走一遍实现库地图（第 13.3 节），停在嵌入式侧的硬约束上（第 13.4 节，Q 格式、实时预算与中断纪律），最后给出横跨分界线的工作流（第 13.5 节）：\PY{} 设计、\CL{} 执行、同激励对拍。

\section{实现结构：卷积、直接形式与 biquad 级联}

\textbf{FIR：卷积即结构。} $y[n] = \sum_{k=0}^{L-1} h[k]\, x[n-k]$ 本身就是实现，没有反馈，因此永远稳定、可以做成严格线性相位（系数对称时相位是频率的线性函数，群延迟恒为 $(L-1)/2$ 个样本）、并且任意有限的字长误差都只会衰减不会振荡。代价以样本计：每输出一个样本要 $L$ 次乘加，线性相位的对称性可以把乘法数砍半（先把 $x[n-k]$ 与 $x[n-L+1+k]$ 相加再乘公共系数），但上千抽头的低截止窄过渡带滤波器依然是每样本千次级的算术。FIR 的实现要点全部在数据通路：一条环形（绕圈）缓冲存最近 $L$ 个输入，热循环是对``缓冲片段与系数数组''的点积（清单~\ref{lst:ch13-cmsisfir}）。

\textbf{线性相位的四种对称。} ``对称''其实有四个象限，各有一个坑：偶数抽头配偶对称（第 I 型）最通用；奇数抽头配偶对称（第 II 型）在 Nyquist 处\emph{恒为零}——做不了高通与带阻的上边；奇对称的两种（第 III、IV 型）在零频恒为零，是微分器、Hilbert 变换器与交错延迟线的地盘。设计工具偶尔会在``对称即可''的默认下把高通交给第 II 型，交付时才对 Nyquist 处的深谷追悔——对称类型与目标频响要一起检查。

\textbf{IIR：同一传递函数，多种算术。} IIR 的传递函数是有理分式 $H(z) = B(z)/A(z)$，把它变成代码至少有四种结构，而且它们的\emph{数值行为}互不相同——这是 IIR 与 FIR 最大的区别：FIR 的结构差异只影响速度，IIR 的结构差异影响结果。直接 I 型把全零与全极点两条链分开，各存一组状态（$2N-1$ 个延迟单元），舍入噪声只进一次、不经过分母放大，代价是状态多。直接 II 型（典范型）状态最少（$N$ 个），但中间信号流经极点链再流进零点链，零极点接近时中间动态范围可以远大于输入输出——浮点上表现为``不溢出但丢精度''，定点上直接饱和。直接 II 型转置（DF2T）把延迟链合并成一条、输入到输出路径最短，对流水线与乱序执行最友好，所以成为库默认：\Lib{CMSIS-DSP} 把结构名直接写进函数名，\Tp{arm\Bk\_biquad\Bk\_cascade\Bk\_df1\Bk\_f32} 与 \Tp{arm\Bk\_biquad\Bk\_cascade\Bk\_df2T\Bk\_f32} 是同一族传递函数的两种内存布局与舍入路径。

\textbf{为什么级联二阶节。} 把 $2N$ 阶系统写成 $N$ 个 biquad 的乘积，动机是系数量化灵敏度：$8$ 阶多项式的系数是根的对称函数，动态范围横跨几个数量级（想想 $\binom{8}{4} = 70$ 这一档的中心项），单精度浮点连``表示''一对靠近单位圆的极点都勉强，更不用说保持位置；而一个二阶节只有两个极点，float32 足够把半径 $0.9999$ 的极点写对。级联还改变了舍入噪声的传播图：每节的误差被下游节衰减或放大，于是产生两个工程决策——\emph{配对}（把零点配给距离最近的极点，既压中间动态范围又做噪声整形）与\emph{排序}（定点惯例把最容易饱和的节放后面，因为前面的节已把能量滤掉；浮点惯例把高 Q 节放前面以获得最佳噪声整形——两条规则方向相反，移植代码时必查）。并行形式（部分分式展开成一阶节并联）有完全不同的性格：每节独立，天然可多发射、多核，极点靠近单位圆时的灵敏度也低，但它的分子来自留数计算、设计工具链支持薄，主要用于谐振器组与合成器，通用滤波仍以级联为主。

\textbf{舍入、极限环与溢出振荡。} 定点 IIR 有两类非线性伪影：\emph{溢出振荡}——反馈环里的加法一旦回绕，误差被环路喂回去，输出变成大幅值方波且不会自己停下，解法是饱和算术加足够的保护间隔；\emph{小信号极限环}——舍入把一个本该衰减到零的状态卡在死区里，输出残留一个低频振荡，解法是舍入（而不是截断）、状态抖动（往乘法积的低位掺一点噪声），或者干脆在浮点域跑。\Lib{CMSIS-DSP} 的 \Tp{q31} 系列函数在这条线上做了大量脏活：累加用 64 位、饱和用 \Tp{\_\_SSAT} 指令，这正是定点库值得用现成实现而不是手写的理由。

\section{设计方法：窗法、Remez 交换与双线性变换}

\textbf{FIR 窗法。} 把理想频率响应（矩形通带乘线性相位）作逆 DFT 得到无限长、非因果的冲激，截断到 $L$ 点再乘一个窗——窗的主瓣宽度决定过渡带、旁瓣电平决定阻纹。教科书的换算关系（数量级正确、精确值以工具输出为准）：矩形窗过渡带最窄但阻纹只有约 $-21$ dB，hann 与 hamming 换到约 $-44$ dB 与 $-53$ dB、过渡带宽翻倍，blackman 到约 $-74$ dB、带宽三倍。每多要 $20$ dB 阻纹，窗法长度大约翻一倍——这是``低截止窄过渡带 FIR 动辄上千抽头''的根源。

\textbf{Remez 交换（Parks--McClellan）。} 给定通带边缘、阻带边缘与两侧纹波，窗法只能给出``某族窗里最好的''，而 PM 算法用交替定理直接找\emph{等纹波}最优解：误差在极值点间以正负交替出现，交换步收敛即最优。同样指标下它比窗法显著更短（节省幅度随纹波不对称程度增大）。配套事实：等纹波只最小化最大误差，不保证能量误差，心理声学场景（比如听感优先的母带均衡）有时故意要不等纹波，那是最小二乘（\Tp{firls}）与频率加权的地盘。\PY{} 侧的入口分别是 \Tp{scipy.signal.remez}（或 \Tp{firwin} 的 win 参数族）与 \Tp{firls}，设计在桌面完成、系数序列化后交给执行层（第 13.5 节）。

\textbf{IIR：归一化原型加双线性变换。} 经典 IIR 的设计是一张两步表。第一步查归一化模拟原型（截止折叠到 $\Omega_c = 1$ rad/s）：Butterworth 通带最平（极点均匀分布在左半单位圆），Chebyshev I 用通带纹波换更陡的边缘，Chebyshev II 把纹波挪到阻带、通带保持平坦，椭圆（Cauer）两侧都有纹波、同阶数下过渡最短但相位最乱，Bessel（Thomson）放弃幅度换群延迟平坦——它的存在理由是暂态波形而不是滤波幅度。选型判据其实很短：要相位保真（乐器、测量、任何在意``起音''的）先排除椭圆、考虑 Bessel 或干脆 FIR；要最省阶数看椭圆；不知道选谁就 Butterworth。第二步用双线性变换 $s \to 2 f_s (1 - z^{-1})/(1 + z^{-1})$ 把模拟原型映到数字域：它把整条虚轴压进单位圆一一对应的位置，因此\emph{稳定性自动保持}，代价是频率轴的非线性扭曲——数字频率 $\omega$ 对应模拟频率 $\Omega = 2 f_s \tan(\omega/2)$。所以设计目标频率必须先做\emph{预畸}（prewarp）再进原型表，否则变换后的 3 dB 点会向低频漂。现代工具链把这一步藏起来了：\Tp{scipy.signal.butter} 给定 \Tp{fs} 就自动预畸；手工套老公式时它正是最常见的错误源。冲激响应不变与零极点匹配两条老路线仍在脉冲合成与仿真里出现，但它们对高频混叠或带限有各自的适配条件，通用低通均衡链的默认已是双线性。

\textbf{指标先于方法。} 不论用哪条设计路线，输入都是同一张四元组：通带边缘、阻带边缘、通带纹波、阻带衰减——而过渡带宽度是第一笔要谈的预算：同长度下过渡带减半，约等于阶数翻倍或纹波付账。所以滤波器设计的第一步不是打开工具，是把``滤掉 20 kHz 以上''翻译成四个数；写不出这四个数的需求，任何工具都帮不了它。

\textbf{半带与抽取链。} 采样率换算的滤波器有一个特化的家族：半带滤波器的冲激在偶数索引（中心除外）恒为零，通带与阻带关于 $f_s/4$ 镜像对称，单级只能把速率减半并换来稀疏乘法；把几级半带串成抽取链（$2$、$2$、$2$ 逐级降），每一级的抗混叠要求随速率下降而放宽，总抽头数远低于单级大过渡带——CD 到电话带宽的 $44.1 \to 8$ kHz 常走三到四级。插值侧是全对称的对偶故事。这条链的每一级都是第 13.1 节的 FIR 结构，只是``零抽头不乘''让算术成本近乎减半；多相分解再把``先滤波后降采样''重排成``先分路后短滤波''，就是第 13.3 节 \Lib{r8b} 们的主场。

\textbf{最小相位与 Hilbert 变换器。} 两个跨结构的特殊设计。其一，同一个幅度响应可以配无数种相位；把零反射回单位圆内（全通分解）得到\emph{最小相位}版本——能量最早到达、群延迟平方和最小，用于监听链路与``不想听到预 echo''的场合。数值做法是经典的同谱法：对数幅度经 Hilbert 变换得到最小相位，再截断成 FIR（\Tp{scipy.signal} 有 hilbert 与 mfdt 系入口，名字以文档为准）。反向的工具也存在：IIR 的群延迟纹波可以用一串全通节做相位均衡——幅度不动、只修相位，这是``线性化''一个既有滤波器的廉价路线。其二，Hilbert 变换器本身：理想频响是 $\pm i$（正频率移 $-90^\circ$、负频率移 $+90^\circ$），全通幅度、奇对称冲激，只能用奇数抽头、中心系数恒为零的反对称 FIR 逼近（第 13.1 节的第 III、IV 型），通带内相位误差随长度缓慢收敛——单边带调制、复包络解析、响度与相位检测都站在它上面。

\begin{note}[title={设计工具的输出不是实现系数}]
\begin{itemize}
\item \textbf{采样率被烤进系数里。} 双线性变换的输出系数是 $f_c / f_s$ 的函数，``2 kHz 低通''在 48 kHz 下设计出来，到 44.1 kHz 下中心就是 $2000 \times 44100/48000 \approx 1838$ Hz 的物理音。系数文件必须把 $f_s$ 写进头部元数据，固件加载时校验，不匹配就拒绝而不是``凑合用''。
\item \textbf{实现形式被假设了。} 同一行 \Tp{sos} 数组，按 DF1 还是 DF2T 的状态布局去跑，传递函数一致但中间精度不同；某些工具输出的 biquad 约定是 $a_0$ 不归一、或者把增益单独放一列（SciPy 的 SOS 每行第七列就是节增益），漏乘增益或重复归一都会让响应整体平移或缩放。对接两套工具时，先各自跑同一段冲激激励，比对前 16 个输出样本，比读文档快。
\item \textbf{定点量化没人管。} float64 设计的窄陷波在 Q15 系数量化后极点半径动了、等效 $Q$ 从 30 掉到 12，频响上表现为``陷不下去''。定点化是设计的一部分：量化系数之后必须重新画频响并检查极点半径，这一步 scipy 不替你做完。
\end{itemize}
\end{note}

\section{桌面与插件侧的实现库}

\Lib{iir1} 是``设计与执行同栈''的 \CPP{} 头文件代表：设计侧覆盖常见原型族（Butterworth、Chebyshev 等，以文档为准），把模拟原型、预畸与级联分解收在初始化路径里；执行侧每个 biquad 是一个标量类型与\emph{状态类型}互相解耦的小对象（模板参数区分输入、系数、状态的字长，状态可以比输入宽，这就是它的防溢出策略），运行期只有第 13.1 节那种每样本状态更新。它适合把滤波器嵌进任何不能带依赖的 \CPP{} 工程，MIT 类条款。\Lib{DSP-Filters-IIR}（IAPWSC 系）是同方向的另一份 \CPP{} 库，历史上带 Qt 图形背景，设计入口更贴近``教科书公式直译''，适合学习与交叉验证。

\Lib{r8b} 一族站在另一条轴上：重采样与分数延迟——多项式插值、多相结构与为采样率转换特化的 IIR/FIR 混合级，DAW 与播放器里``任意速率播放''的那一层就是它这类库；接口的共同点是输入输出块加内部相位状态，分数延迟则把``采样点之间''的读数做成可扫参的连续函数（磁带机的变调、延迟补偿的精细对齐都是它）。\Lib{SoX} 展示了另一种组织方式：设计、执行与调度全在一条命令行里——均衡与 Shelving 类特效在内部分解成 biquad 级联，\Tp{rate} 与 \Tp{sinc} 是两条品质/速度档的重采样路线，离线批处理时可以随意排序串联。它的价值不在实时（它不做实时），在于给你一个``效果链即数据''的参照系。

到了插件世界，这个参照系变成对象：\Lib{JUCE} 的 \Tp{dsp::ProcessorChain} 把若干处理器按编译期类型列表串成一元组，逐样本与逐块两种 process 签名都提供，双线性滤波器的状态与系数切换由它统一管——系数切换的正确姿势（写目标值、在下个块边界生效）正是为了避开``运行中换系数''的爆音。再往上到格式层，VST3 与 AAX 对宿主与插件的分工做了硬规定：音频线程内不得文件 IO、不得锁竞争、块大小与采样率随宿主会话变（约束细节以各 SDK 文档为准）。采样率会变这件事对滤波器库是一等需求：要么收到变更通知后重跑设计，要么预置多套系数按表切换。这三条恰好就是下一节嵌入式实时约束的桌面版镜像——插件作者踩过的坑，嵌入式作者要再踩一遍，只是预算更小。

\begin{table}[H]
\centering
\caption{五套滤波器相关库：层次、语言层、许可类别与实时能力}
\label{tab:ch13-filter}
\small
\begin{tabular}{@{}>{\raggedright\arraybackslash}p{0.17\textwidth}
  >{\raggedright\arraybackslash}p{0.14\textwidth}
  >{\raggedright\arraybackslash}p{0.13\textwidth}
  >{\raggedright\arraybackslash}p{0.15\textwidth}
  >{\raggedright\arraybackslash}p{0.18\textwidth}@{}}
\toprule
库 & 层次 & 语言层 & 许可类别 & 实时能力 \\
\midrule
\Lib{iir1} & 执行为主、附带设计 & \CPP{} 头文件 & MIT 类 &
  是：无分配的状态更新，可进音频回调 \\
\Lib{CMSIS-DSP} & 纯执行 & \CL{} 加 intrinsics & 宽松（以仓库 LICENSE 为准） &
  是：为中断与固定块深做的设计 \\
\Lib{r8b} & 执行（重采样与分数延迟） & \CPP{} 头文件 & 极宽松（以源码头为准） &
  是：块式接口，插件与播放器常用 \\
\Lib{SoX} & 设计加执行加调度 & \CL{} 内核、命令行 & GPL 系条款 &
  否：离线批处理定位 \\
\Tp{scipy.signal} & 纯设计 & \PY{} 门面、C 后端 & BSD 系条款 &
  否：解释器与分配不满足硬预算 \\
\bottomrule
\end{tabular}
\end{table}

\section{嵌入式侧：CMSIS-DSP、Q 格式与实时预算}

\Lib{CMSIS-DSP} 的接口形态本身就是教材：每个功能是一族``C 函数加实例结构体''——\Tp{arm\Bk\_fir\Bk\_instance\Bk\_f32} 存抽头数、状态缓冲指针与环形索引，\Tp{arm\_fir\_f32} 在 init 之后每块调用一次，热路径里不分配、不取模、不查表跳转；类型后缀 \_f32、\_q31、\_q15 是同一算法的三种字长分身，定点分身里还有用精度换速度的 fast 变体。三角函数走查表加插值（\Tp{arm\_sin\_q15} 一类的固定长度正弦表），因为 M 核上带库调用与浮点的通用实现是奢侈品；指令集侧，M 核的 Helium（MVE）向量扩展与部分手写汇编覆盖热循环。表驱动、绕圈缓冲、状态外置，这三件套就是嵌入式 DSP 库的全部接口美学。

\textbf{浮点还是定点。} 带 FPU 的 Cortex-M4F 与 M7 改变了这笔账：单精度 float 的乘加在 M4F 上是几条流水线内的常数周期，Q31 反而要为饱和与移位多付指令——``有 FPU 就用 \_f32 族''成为默认。定点的主场剩下三处：没有 FPU 的 M0 系列、同频下吞吐翻倍的批量作业（q15 一次装两个进 32 位寄存器）、以及与 ADC、编解码器位流严格对齐的链路。清单~\ref{lst:ch13-cmsisfir} 给的就是这场``格式战争''中立于两边的执行层骨架。

硬件侧的教科书样本是 \Lib{libdaisy}（Electro-Smith 的 Daisy 系列硬件抽象）：STM32H7 级 Cortex-M7、SAI 或 I2S 外设加 DMA 双缓冲，初始化后注册一个音频回调，每块默认 $48$ 帧——$48$ kHz 下正好 $1$ ms 一次，回调的返回直接衔接下一块的 DMA。没有操作系统，回调运行在最高优先级上下文里，这就是``实时预算''以最纯粹的样子出现的地方。

\textbf{预算的算术。} 一个块的周期预算是 $f_{\text{clk}} \times B / f_s$（$B$ 为块长，单位样本）：$480$ MHz、$48$ 帧、$48$ kHz 得 $480\,000$ 周期。对照成本：立体声 $256$ 抽头 FIR，若每周期一次乘加，每块耗 $48 \times 2 \times 256 \approx 24\,600$ 周期，约 $5\%$ 预算；换成 $2048$ 抽头就逼近 $40\%$——还没算均衡链、电平表与 MIDI。块大小与延迟是一对可调旋钮：交互性要求往返延迟压在十毫秒内，块就不能超过约 $480$ 帧，而块越小、每块的固定开销（进出回调、DMA 中断）占比越高。工程做法是把预算表贴在代码旁边：每功能标``周期每样本''，加总乘块长，与预算比——超预算的后果不是慢，是爆音，是用户可听见的缺陷。

\textbf{中断上下文里禁止什么。} 回调即中断语境：不得堆分配（堆锁与页错误或有或无，但都没有有界性可言）、不得阻塞式 IO 与日志（printf 到串口就是几十微秒起的无界等待）、不得拿可能优先级反转的锁、不得跑时长不可预测的算法路径（变长查表搜索、未设上限的迭代、首次触页的访存）。可以做的：预分配的静态或池内存、查表、定点或 FPU 算术、以及把一切``慢的''（参数平滑、UI、文件）挪到低优先级上下文，两边用双缓冲加原子交换交接。缓存与 DMA 是最后一对暗雷：DMA 直接改内存而内核读的是缓存里的旧值，或者写回脏行覆盖了 DMA 的新数据——解法要么把共享缓冲放进 MPU 标记的非缓存区，要么在 DMA 前后显式 clean 与 invalidate，并让缓冲按缓存行（$32$ 字节档）对齐，否则你得到的是一段``偶尔''用旧采样的滤波器。

\textbf{Q 格式：把小数点亮在注释里。} 定点没有二进制小数点，只有约定：Q15 用 \Tp{int16\_t} 表示 $[-1,1)$、步长 $2^{-15}$，Q31 同法用 \Tp{int32\_t}。乘 Q15 是 $(a \times b) \gg 15$ 进 \Tp{int32\_t}；累加 $L$ 项最多再长 $\lceil \log_2 L \rceil$ 位——$2048$ 抽头的 FIR 在 16 位乘 16 位下会顶破 32 位，所以必须 64 位累加或者边累加右移，右移几位是增益账，必须显式写在接口上。这就是为什么定点库的签名要把格式说尽：\Tp{arm\_conv\_q15} 带 \Tp{postShift} 参数、\_q31 系列靠 64 位中间量与饱和宏；``\Tp{int16\_t}'' 这个类型本身不能告诉你它是 Q0 的音频样本还是 Q15 的系数，接口文本上省下的字，都会以增益错误或溢出的形式还回来。饱和与舍入也不是可选项：截断引入与信号相关的系统性偏差（听感上是失真），饱和防止回绕导致的溢出振荡，两者在 \Tp{q31} 函数族里是指令级的默认行为。

\begin{table}[H]
\centering
\caption{三种执行层数字格式速查（动态范围为量级估计）}
\label{tab:ch13-formats}
\small
\begin{tabular}{@{}>{\raggedright\arraybackslash}p{0.13\textwidth}
  >{\raggedright\arraybackslash}p{0.17\textwidth}
  >{\raggedright\arraybackslash}p{0.17\textwidth}
  >{\raggedright\arraybackslash}p{0.26\textwidth}@{}}
\toprule
格式 & 存储与小数点位置 & 动态范围量级 & 典型归属 \\
\midrule
Q15 & \Tp{int16\_t}，$15$ 位小数 &
  约 $90$ dB 满摆幅 &
  系数与状态的字长受限档、老式编解码位流 \\
Q31 & \Tp{int32\_t}，$31$ 位小数 &
  约 $188$ dB（累加另账） &
  现代定点 DSP，配 $64$ 位累加器 \\
float32 & IEEE 754 单精度 &
  对数网格，上千 dB &
  带 FPU 的 M 核与桌面链路的默认 \\
\bottomrule
\end{tabular}
\end{table}

\begin{lstlisting}[style=cstyle,label={lst:ch13-cmsisfir},caption={CMSIS-DSP 风格的 FIR 加绕圈缓冲：init 一次，热循环零分配}]
// 执行层骨架：接口形态照抄 arm_fir_*——实例结构体承载状态，
// coeffs 只读、state 是绕圈缓冲；热路径无 malloc、无除法取模。
#include <stdint.h>
#include <string.h>

typedef struct {
  const float* coeffs;     // h[0] 与最新样本相乘
  float* state;            // 环形缓冲，容量 cap >= ntaps + blk
  uint32_t ntaps, blk, cap;
  uint32_t idx;            // 指向"最新样本"所在槽
} FirF32;                  // 与 arm_fir_instance_f32 同形

void fir_init(FirF32* f, const float* h, uint32_t n,
              float* buf, uint32_t cap, uint32_t blk) {
  f->coeffs = h; f->ntaps = n; f->cap = cap; f->blk = blk;
  f->state = buf; f->idx = 0;
  memset(buf, 0, sizeof(float) * cap);  // 清零只发生在 init
}

void fir_run(FirF32* f, const float* src, float* dst) {
  for (uint32_t i = 0; i < f->blk; ++i) {
    if (++f->idx >= f->cap) f->idx = 0; // 每样本一次回绕判定
    f->state[f->idx] = src[i];          // 新样本入圈
    float acc = 0.f;
    uint32_t p = f->idx + 1u;           // 从新往旧遍历历史
    if (p >= f->cap) p = 0;
    for (uint32_t k = 0; k < f->ntaps; ++k) {
      acc += f->state[p] * f->coeffs[k]; // 教学版逐点回绕；
      if (++p >= f->cap) p = 0;          // 真实库把内层拆成两段
    }                                    // 无回绕的连续内积，
    dst[i] = acc;                        // 才可能向量化（见正文）
  }
}
\end{lstlisting}

\section{\PY{} 设计、\CL{} 执行的协作范式}

把第 13.1 节的执行结构与第 13.2 节的设计方法合起来，就得到这一域的标准分工：\PY{} 拥有设计，\CL{}/\CPP{} 拥有执行，中间隔着一次序列化。完整回路是四步。第一步，桌面设计：\Tp{scipy.signal.butter} 一类的入口给出 \Tp{output} 取 ``sos'' 的二阶节数组（不要 \Tp{ba}：高于四阶的传递函数系数在 float32 下已经保不住根的相对位置，这是第 13.1 节灵敏度分析的结论）。第二步，桌面验证：\Tp{sosfreqz} 画幅频、\Tp{sosstep} 看暂态，\Tp{sosfilt} 跑参考输出——桌面侧的 ``filtfilt''（零相位）只属于离线，不能移植进实时链。第三步，序列化：把 SOS 写成 C 数组（生成头文件，或裸二进制配 $f_s$、$f_c$、阶数、节数、格式与字节序的头部）——元数据不是礼貌，是第 13.2 节那个 note 讲的三漂移的解药。第四步，板端执行与对拍：系数进常量区（\Tp{const}、ROM 段，链接脚本让它落在不映射数据的 flash 区），状态结构体静态分配，运行期只有状态更新；然后用同一激励在两侧逐点比对。

\textbf{版本纪律。} 系数生成脚本要与固件同仓同标签入库，构建时把系数数组的哈希印进版本字符串——``这版声音不对''的报障，第一步是对哈希而不是重听。回归则做成机器规则：桌面用固定激励数组生成期望输出，板端（或 C++ 测试程序）跑同一数组回传比对，让差值序列进版本库，任何系数改动都会立刻在测试里显形。

\textbf{验收里的听感一栏。} 对拍通过不等于耳朵点头：线性相位 FIR 的预振铃（pre-ringing）在打击乐上有人听得见，IIR 的群延迟纹波有人描述成``低频散''。自动对拍之外留一步人工 A/B：同一素材在新旧实现间切给不参与开发的同事听一次——便宜，但是这一域特有的最后一道防线。

\begin{lstlisting}[style=pythonstyle,caption={设计回路的最小脚本：生成 SOS、自查频响与阶跃、量化核对、导出}]
import numpy as np
from scipy import signal

fs, fc, order = 48000.0, 2000.0, 4
sos = signal.butter(order, fc, "low", fs=fs, output="sos")
w, h = signal.sosfreqz(sos, worN=2048, fs=fs)   # 幅频自查
_, s1 = signal.sosstep(sos, 400)                # 阶跃：过冲与建立
q = np.round(sos * 2**15).astype(np.int16)      # Q15 量化版
wq, hq = signal.sosfreqz(q / 2**15, 2048, fs)   # 量化后再画
print("量化后的最大幅频偏移:", np.abs(h - hq).max())
np.savetxt("lp4.sos.txt", sos, delimiter=",")   # 序列化：把 fs、
                                                # fc、order 写进文件
                                                # 头部随数组一起交付
\end{lstlisting}

对拍``通过''的标准要提前写在验收里，事后补的是狡辩：

\begin{itemize}
\item \textbf{冲激看结构。} 板端对单位冲激的响应应与桌面 \Tp{sosfilt} 输出逐样本接近：级联顺序错、节增益漏乘、状态初值脏，全部在冲激的前几十个样本里现形。
\item \textbf{正弦看稳态。} 通带与阻带各选一个频点跑满建立时间，比幅度与相位：浮点目标容差可设 $10^{-5}$ 档，定点目标的容差就是量化步长乘抽头数，算不出这个数就说明增益账没做完。
\item \textbf{阶跃看暂态。} 过冲量与振铃周期是 Bessel 与椭圆 ``谁是谁''的指纹——阶跃响应错配常意味着两侧用了不同的实现形式（第 13.2 节 note 的第二条）。
\item \textbf{坏消息要能定位。} 对拍脚本输出逐样本差值序列而不是一个 ``max err''，让发散是``从第一拍起''还是``从第 N 拍起''可见——前者是系数问题，后者几乎总是状态问题。
\end{itemize}

\begin{lstlisting}[style=cppstyle,label={lst:ch13-butter},caption={biquad 级联实现四阶 Butterworth 低通：系数以数组给出}]
// 四阶 Butterworth 低通 = 两个二阶节的级联。
// 系数来源：fs = 48 kHz、fc = fs/24（2 kHz），双线性变换
// 含预畸，增益在两节间对半分配；工程上请用
// scipy.signal.butter(4, 2000, fs=48000, output="sos")
// 重新生成并核对（注意其增益分配与节顺序约定）。
#include <cstdint>

struct Biquad {
  float b0, b1, b2, a1, a2;   // 分母首项已归一为 1
  float z1 = 0.f, z2 = 0.f;   // DF2T：全部状态只有两个延迟
  float step(float x) {       // 每样本一次的硬核：无分支
    float y = b0 * x + z1;
    z1 = b1 * x - a1 * y + z2;
    z2 = b2 * x - a2 * y;     // 回环即状态更新，全在寄存器
    return y;
  }
};

// 行格式：b0 b1 b2 a1 a2（每节增益已并入 b 侧，约 0.0146）
constexpr float kSos[2][5] = {
  {0.0146013f, 0.0292025f, 0.0146013f, -1.757732f, 0.819750f},
  {0.0146013f, 0.0292025f, 0.0146013f, -1.559056f, 0.614054f},
};

Biquad stages[2] = {
  {kSos[0][0], kSos[0][1], kSos[0][2], kSos[0][3], kSos[0][4]},
  {kSos[1][0], kSos[1][1], kSos[1][2], kSos[1][3], kSos[1][4]},
};

float lowpass4(float x) {          // 级联顺序：高 Q 节在前
  return stages[1].step(stages[0].step(x));  //（浮点噪声整形惯例）
}

void reset_filter() {              // 换采样率/换链路：状态清零
  stages[0] = Biquad{kSos[0][0], kSos[0][1], kSos[0][2],
                     kSos[0][3], kSos[0][4]};
  stages[1] = Biquad{kSos[1][0], kSos[1][1], kSos[1][2],
                     kSos[1][3], kSos[1][4]};
}
\end{lstlisting}

\begin{warning}[title={状态、增益与格式：搬系数最容易断的三处}]
\begin{itemize}
\item \textbf{忘清状态。} 滤波器对象可以拷贝、可以数组化，但状态是它的私有成员：换采样率、换预设、跳曲库时只重配系数不清状态，旧信号的能量会从上一段``漏''进这一段；在级联里这比单节更难诊断，因为中间节的输入没有物理意义。重配与清状态必须是一个原子动作（见清单~\ref{lst:ch13-butter} 的 \Fn{reset\_filter}）。
\item \textbf{增益只在对拍里现身。} 级联的总增益是各节之积：工具链一边的整型缩放被实现层拆到别处，输出的不是``形状不对''而是``响度差一档''，归一化后的幅度谱测试会过、电平表会炸。对拍必须含一条直流与一条满幅正弦的增益核对，不能只比形状。
\item \textbf{格式错配是双向静音风险。} 把 float 系数数组按 \Tp{int32\_t} 读、或把 Q15 状态喂给 \Tp{q31} 函数，得到的不是噪声就是全零，两种都比崩溃更难定位。序列化格式（类型、字长、Q 档、字节序）写进系数文件头部并在加载时断言，让错配在 load 时而不是运行时爆炸。
\end{itemize}
\end{warning}

\begin{bestpractice}[title={本章的落地清单}]
\begin{itemize}
\item 高阶 IIR 一律走 SOS 级联交付，\Tp{ba} 直传只允许出现在一阶与二阶；交付文件自带 $f_s$、$f_c$、实现形式三样元数据。
\item 实现形式（DF1、DF2T）与级联顺序是数值决策，写进注释与测试，别让``移植后变差''沦为玄学。
\item 执行层的形状对齐清单~\ref{lst:ch13-cmsisfir}：init 一次、热循环零分配零取模，换参数永远走重配加清状态。
\item 每块预算写成一张表贴在代码旁：功能、周期每样本、占预算百分比；总和低于七成才算上线。
\item 设计侧与执行侧用同一激励对拍（冲激、正弦、阶跃各一条），差值序列入版本库，容差写进验收而不是口头。
\end{itemize}
\end{bestpractice}

这一域为什么必须落在 \CL{}/\CPP{}，本章给出的理由是三条别的领域少见的组合：环境上，大量目标是裸机或 RTOS，没有堆纪律、没有异常表、链接脚本管段布局，``运行时''是你自己写的启动代码；成本上，预算以周期计、内存以千字节计，任何一次隐式分配与隐式转换都是负债；工具链上，嵌入式编译器唯一承诺稳定的接口就是 \CL{} ABI——C++ 的名称修饰、异常与标准库支持度随厂商浮动，所以库的对外形态退回 C 函数加实例结构体。于是分工自然浮现：\PY{} 做迭代式的设计（它拥有画图、优化器与交互），\CL{}/\CPP{} 做每拍必跑的执行（它拥有 ABI、布局与周期账），两边靠一次带元数据的序列化对接，靠同一激励的逐点对拍封口。这正是全书主线``C 作 ABI 汇合点、上层语言做门面''在滤波器这一域的最干净样本。
% ====================================================================

% ====================================================================
% ====================================================================
\part{信号处理库：电磁波}

\chapter{电磁仿真：方法选择与 C/C++ 内核}

声音那一部分的痛点是截止期，几何那一部分的痛点是谓词的符号正确性；电磁波这一域的痛点是\emph{规模}。一次三维全波仿真动辄几百万到几亿个未知量、几十 GB 的工作集、跑几十分钟到几天。这个规模把本册的主线推到一个很干净的例子上：方法学必须先按电尺寸选对，选错了任何语言都救不回来；而选对之后的实现几乎全部落在 \CPP{} 加 \CL{} 的内核上，因为只有这一层能同时把 stencil 更新向量化、把工作集分片给 MPI 与 GPU、把色散与各向异性带来的额外状态变量安置在调用方可控的布局里。

这一域还提供一个额外的观察：\emph{门面为什么可以是 Scheme 与 Lua}。几何与材料是一次仿真的输入，改动极频繁、表达极复杂（盒体的布尔堆叠、旋转的各向异性介质、周期性单元），用模板参数描述等于每换一条几何就重编译一次。于是求解器普遍把``描述''交给脚本或声明式文件，解析后送进编译好的内核：\Lib{Meep} 用 Scheme（经由 \Lib{libctl}）、\Lib{openEMS} 用 Lua 再加一份 XML 中间表示、\Lib{gprMax} 用 \PY{} 生成 \Tp{.in} 文本、\Lib{S4} 用 Lua 与 MATLAB。这与本册``接口层：C ABI、Fortran 血统与四种绑定路线''一章的结论同构，只是把 \PY{} 换成了更冷门的脚本语言。

边界声明：本章只讲方法选择、内核形状、数据约定与验证。接收机链路、设备抽象与比特流在``射频链路与软件无线电''一章；时域响应到频谱的那一步交给本册讲 FFT 与谱分析的那一章；射线与几何求交的实现细节参见``网格、点云与光线求交''一章，这里只把它当作一类方法对待。

\section{尺度决定方法：先问电尺寸，再问要什么信息}

问题的第一维度是\textbf{电尺寸} $kL = 2\pi L / \lambda$（$L$ 为结构最大线度，$\lambda$ 取工作频带内最短波长）。第二维度是\textbf{要什么信息}：只要方向图与 RCS，还是要端口上的宽带 S 参数，还是要腔体内的本征频率与品质因数。把这两维摆开，主流方法只剩六条路。

\textbf{时域 FDTD 与 Yee 元胞。} 电场与磁场落在同一网格的错开位置（切向 $E$ 在棱上、切向 $H$ 在面上），时间上再错开半步，构成 \Fn{leapfrog} 推进；二阶中心差分让格式条件稳定，稳定条件是 Courant 数不超过 $1/\sqrt{D}$，$D$ 为有效维数。它一次脉冲激励就能拿到全带宽响应，天然处理色散、损耗、非线性与任意材料分布；代价是网格必须同时跟住最短波长\emph{与}最细特征尺寸，以及共振类问题要跑到场衰减完毕（步数正比于品质因数 $Q$）。

\textbf{频域 FEM。} 变分形式配棱边元（节点元在波导问题里会产生伪模），得到稀疏带状矩阵；带状宽度随极化数与波数增长。适合封闭或半封闭区域的本征问题（本征求解直接给 $\omega$ 与 $Q$）、精细边界、复各向异性张量材料。代价集中在矩阵因子：直接稀疏分解的内存标度在二维接近 $O(N^{1.5})$、三维接近 $O(N^{2})$，这是填充元决定的，不是网格决定的。

\textbf{表面积分法：MoM 与 SIE。} 未知量只放在导体表面或细线上（等效电流、磁流），核是 Green 函数，于是离散对象从体积降到表面——电大尺寸的金属辐射与散射问题因此可行。代价是矩阵\emph{稠密}：$N$ 个表面基函数要 $O(N^{2})$ 个元素，直接解是 $O(N^{3})$，所以这条路线的规模化全靠快速多极子（FMM）或多层 Green 函数加预条件。介质体问题要写成 SIE（PMCHW 型），条件数很差，预处理是必选项而不是优化项。

\textbf{射线类高频近似。} 物理光学、几何绕射理论，以及 SBR（shooting bounce stars：发射一束平行射线，按弹跳次序追踪到多个反射点后相干叠加）用于电大目标（$kL$ 上千）的 RCS、舱内与城市尺度传播。它把随 $\lambda^{-3}$ 增长的未知量换成数量近乎不随频率增长的射线管，代价是绕射、耦合与近场信息基本丢失，需要 PTD、UTT 一类补项。这一类的实现形状（BVH 遍历、分支密集）更接近光线求交那一章，而不是 stencil 更新。

\textbf{RCWA（傅里叶模法）。} 周期性结构沿周期方向做 Fourier 展开，每层解析、层间用散射矩阵递推连接（累积 T 矩阵在消逝波上数值不稳定，这是 S 矩阵写法存在的原因）。未知量是截断谐波数：一维周期按 $2N+1$、二维按 $(2N+1)^{2}$ 增长，每谐波一次本征问题。要点在\emph{分解规则}：介电常数不连续时朴素 Laurent 分解收敛极慢甚至给出错值，正确做法是把切向场与法向通量分别按调和平均与算术平均处理（Li 的自适应傅里叶分解规则）——这正是 \Lib{S4} 的核心卖点。

\textbf{传输线与集总近似。} 横向尺寸远小于波长且截面场型近似已知（同轴、微带、带状线）时，三维问题降为每单位长度的 $RLCG$ 加一段二端口，不连续处用集总元件或模式展开的等效电路表示。成立条件是特征尺寸与 $\lambda$ 之比足够小且沿纵向的参数变化缓慢；一旦不连续的物理尺寸与 $\lambda$ 可比、或激励带宽覆盖高次模的通过频带，就必须回到全波。这一档的价值在系统级：链路预算与匹配网络不需要解 Maxwell 方程，``射频链路与软件无线电''一章走的就是这条路。

\vspace{2pt}
\noindent
表~\ref{tab:ch14-methods} 把六条路按``离散对象''与``成本律''对齐，这两列就是选型判据本身。

\begin{table}[H]
\centering
\caption{六类方法：离散对象、规模来源与成本律}
\label{tab:ch14-methods}
\small
\begin{tabular}{@{}>{\raggedright\arraybackslash}p{0.15\textwidth}
  >{\raggedright\arraybackslash}p{0.16\textwidth}
  >{\raggedright\arraybackslash}p{0.16\textwidth}
  >{\raggedright\arraybackslash}p{0.32\textwidth}@{}}
\toprule
方法 & 离散对象 & 规模由什么决定 & 成本律与失效方式 \\
\midrule
FDTD & 体元胞上的 $E$、$H$ 分量 &
  最短波长、最细特征、仿真时长 &
  内存 $\propto$ 元胞数 $\times$ 分量数，时间 $\propto$ 元胞数 $\times$ 步数；
  高 $Q$ 与大时长时爆炸，另有数值色散 \\
FEM & 单元上的基函数（棱边元） &
  区域体积、多项式阶数、波数 &
  稀疏直接分解，二维近 $O(N^{1.5})$、三维近 $O(N^{2})$；
  内存峰值出现在因子而不是矩阵 \\
MoM、SIE & 表面基函数、细线未知量 &
  表面积与 $\lambda^{2}$ 之比 &
  矩阵稠密，存储 $O(N^{2})$、直接解 $O(N^{3})$；
  靠 FMM 与多层 Green 函数展开，介质体条件数差 \\
射线、SBR & 射线管与弹跳次序 &
  目标电尺寸、所需弹跳数 &
  与频率近乎无关；丢失绕射、耦合与近场，
  电小结构上完全不能用 \\
RCWA & 谐波与层 &
  周期维数、截断谐波数、层数 &
  谐波数增长快、每层一次本征问题；
  非周期、强吸收或陡峭剖面时收敛很慢 \\
传输线、集总 & 二端口级联 &
  不连续尺寸与 $\lambda$ 之比 &
  近似法，几乎不耗算力；
  尺寸比超过约 $0.1$ 就漏掉高次模与辐射 \\
\bottomrule
\end{tabular}
\end{table}

把成本律写成可核对的乘积更实用。FDTD 的元胞数近似 $\bigl(L/\lambda_{\min}\bigr)^{3} \times r^{3}$，$r$ 是每个最短波长的网格数，工程上常取 $10$ 到 $20$；步数是物理时长除以 $\Delta t$，而 $\Delta t$ 被 $\Delta$ 与 Courant 数锁死，于是总浮点量大致随 $r^{4}$ 增长（内存再乘一次 $r^{3}$）。一个二十厘米、工作在 $30$ GHz 的封装结构，若里面有零点二毫米的缝隙，网格就得同时满足波长与特征尺寸两个约束，元胞数轻易过亿——这就是``电尺寸 $\times$ 网格密度''这句话的具体含义。MoM 的存储随频率上升而下降（面积除以 $\lambda^{2}$），因此在电小问题上反而不利；FEM 每波长未知量少但矩阵带宽宽、装配贵。选型不是审美，是这三条曲线的交叉点。

\section{版图：内核语言、门面语言与并行的现实}

\Lib{Meep} 是本册要重点解剖的样本：内核 \Tp{libmeep} 是 \CPP{}（Yee 网格、CPML、色散材料、自适应步长、探针与诊断），描述层是 Scheme（\Lib{libctl} 提供几何与材料的声明式对象，由 Guile 解释），\PY{} 门面是同一内核的第二层绑定，与 numpy 之间靠场数组直接交换（\Tp{get\_array} \Bk{} 一类取场接口返回 \PY{} 侧数组）。PML 是拉伸坐标加记忆变量的 CPML 形式；材料侧用 Drude 与 Lorentz 振子描述色散，内核为每个极化机制维护一组辅助场并与主方程一起推进（辅助微分方程路线）；各向异性介质走张量更新。并行是 MPI 域分解（按元胞切块，交换 $E$、$H$ 的 halo）叠加 OpenMP 线程，较新版本还把同一批更新例程通过 CUDA 后端卸载到 GPU——注意``同一批''这个措辞的含义：内核写成可移植的函数形式，设备版靠编译期分派，而不是重写一份算法。姊妹程序 \Lib{MPB} 用平面波展开求周期结构的本征模与能带（频率以 $a/\lambda$ 这类无量纲量报告），走同一套 Scheme 与 \PY{} 门面；\Lib{Meep} 在这个内核之上加了基于伴随场（adjoint）的反向设计，其梯度成本与设计参量个数几乎无关（一次正向加一次伴随求解），这是它做拓扑优化的算法卖点。

\Lib{openEMS} 是 \CPP{} 的三维 FDTD 引擎，门面是 Lua 脚本：脚本生成一份 XML 的网格与材料描述，引擎读入后执行。它的特色是把一批\emph{子网格}电磁结构做成一等公民——细线、金属片、过孔、微带与带状线段（手册里把这一族建模能力记作 TSVM，缩写展开含义以官方文档为准；Lua 侧能直接看到 \Fn{define\Bk\_Via}、\Fn{define\Bk\_MicrostripLine}、\Fn{AddMeshBox} 这类函数），于是封装与天线馈电不必把导体剖到亚微米网格。时间步由 \Fn{SetTStep} 显式给出（等于把 Courant 数变成用户参数），局部加密用 \Fn{AddMeshBox}，后处理侧提供近场到远场的观察面与端口激励，S 参数与方向图由随附的 Octave 脚本算出。

\Lib{gprMax} 是``\PY{} 门面但内核编译化''的标准案例：\PY{} 程序生成 \Tp{.in} 文本（\Tp{\#PrmLst} 一行给网格密度与 PML 单元数，随后是 \Tp{Box}、\Tp{RxAnt}、\Tp{Tx} 一类对象行），真正的 FDTD 推进内核用 Cython 写并编译成扩展模块，CPU 侧用 OpenMP 分线程、GPU 侧走 OpenCL 内核；色散介质用 Z 变换离散化而不是 ADE——同一格式下两种离散化路线给出的辅助变量布局不同，内存占用也不同。输出是 HDF5 场文件，交给 numpy、h5py 或自带的查看脚本。它的典型用途是探地雷达与生物电磁（比吸收率），恰好说明 FDTD 的可移植性来自网格形状而不是来自物理。

\Lib{S4} 与其重写版 \Lib{S4v2} 是 \CL/\CPP{} 的 RCWA 内核，门面给 Lua 与 MATLAB：几何按分层网格描述，每层的介电常数做二维 Fourier 分解（附自适应分解规则），层内解本征问题、层间用散射矩阵递推，输出反射率、透射率与内部场。它没有时域内核，因此也不受 Courant 约束——代价是频域逐频点求解以及谐波数带来的稠密矩阵。NEC 系（\Lib{NEC2} 及其 \CL{} 翻译血统，加上 \Lib{4nec2} 一类图形前端）代表另一条极窄但极有用的路：细线 MoM，未知量从上千到几万，一次运行几秒到几分钟，专门解决线天线的输入阻抗、方向图与近场耦合。它的内核是稠密复矩阵加复数高斯消元，规模天花板由 $O(N^{2})$ 存储决定，而不是由算法精度决定。

商用全波求解器（本章不比实现，只比方法口径）大致对应上面几条路：频域 FEM、时域积分格式（与 FDTD 同族的有限积分法 FIT）、MoM 加 FMM 与 SBR 的混合，以及通用多物理 FEM 平台。共同形态是专有图形前端加参数化建模，内核实现语言不公开，但从接口形状（脚本化 API、参数扫描、分布式计算）与性能特征看，基本可以判断为 \CPP{} 与 \FTN{} 数值核，许可证均为商业闭源。选型时的真实差别往往不在``精度''而在工作流：能否脚本化、能否把一次扫频拆到集群、能否把网格与材料交给版本控制——这三条恰好是文本门面的主场。

\section{内核的形状：四笔必须付的账}

\textbf{第一笔：算术密度与向量化。} Yee 更新是一条极短的依赖链：两个邻居的差分、两个系数、一次乘加。每元胞的工作量是个位数浮点操作，却要搬五到十个浮点数（场分量加材料系数），所以它是\emph{访存受限}的：上限由带宽而不是峰值浮点决定。这也解释了为什么实现必须把场量按分量拆成连续数组（SoA：每个分量一整块），而不是每元胞一个结构体——后者让 SIMD 通道装的是同一元胞的不同分量，向量化立刻失败。清单~\ref{lst:ch14-fdtd} 的写法就是把唯一连续的内层循环留给编译器。

\textbf{第二笔：工作集。} 三维 $512^{3}$ 网格、六个场分量、单精度约是 $512^{3} \times 6 \times 4 \approx 3$ GB；若色散材料需要四个辅助极化场、CPML 需要六个记忆变量，再翻一倍。凡是``顺手多存一份''的设计决策都以千兆字节计价，而带缓存的结构化对象布局在这一域根本不可能。因此这一域对布局、分配与标量类型（\Tp{float} 还是 \Tp{double}）的敏感度超过本册前面任何一章：门面里一次 \Tp{dtype=float32} \Bk{} 的选择，落到内核就是工作集与 MPI 通信量减半。

\textbf{第三笔：时间步数。} 显式格式每步都要扫全网格，总步数由物理时长除以 $\Delta t$ 决定。谐振与高 $Q$ 腔要把场跑到能拟合出线宽，窄带或连续波激励要跑到稳态；一维六万格跑百万步已是十亿量级的元胞更新，三维再乘几个数量级。这正是频域方法存在的理由：一次稀疏分解能处理多个右端，且耗时不随 $Q$ 变长。

\textbf{第四笔：并行与长期驻留的状态。} MPI 域分解把网格切成块，每块只交换 halo 一层；因为 stencil 深度是 $1$，分解可以切得很细而仍让计算压过通信——这一点与稀疏矩阵求解很不一样。GPU 路线是把同一个更新循环换成线程网格：一线程负责若干元胞，寄存器装 $E$、共享内存装边界行。此外，自适应分辨率、伴随优化与多物理耦合都要求场量与极化状态\emph{长期驻留}：一次运行内反复访问、不能被换出，也不能交给带垃圾回收的运行时去管理。

\begin{lstlisting}[style=cppstyle,label={lst:ch14-fdtd},caption={二维 FDTD 的 leapfrog 更新：Yee 交错、衰减层与向量化提示}]
// 二维切片：z 方向不变时只有两组场各自闭合，
// 要么 (Ez, Hx, Hy)，要么 (Ex, Ey, Hz)；选错会得到零场。
// 下面实现前者（TE_z）。场量按分量分块存放（SoA）。
void step2d(int nx, int ny,
            float* __restrict__ ez, float* __restrict__ hx,
            float* __restrict__ hy,
            const float* __restrict__ ca,
            const float* __restrict__ cb,
            float dt, float dx, float mu0, float eps0) {
  const float kh = dt / (mu0 * dx);      // H 的差分系数
  const float ke = dt / (eps0 * dx);     // E 的差分系数
  // 1) H 推进：hx 位于 j+1/2，沿 y 对 ez 差分
  //    内层 j 连续，可展开成 SIMD；外层 i 交给 OpenMP
#pragma omp parallel for schedule(static)
  for (int i = 0; i < nx; ++i) {
    const int r = i * ny;
#pragma omp simd
    for (int j = 0; j < ny - 1; ++j)
      hx[r + j] -= kh * (ez[r + j + 1] - ez[r + j]);
    // hy 位于 i+1/2，沿 x 对 ez 差分：跨距是 ny，
    // 这一维天然不连续，要靠分块把行段先搬进缓存
    for (int j = 0; j < ny; ++j)
      hy[r + j] += kh * (ez[r + ny + j] - ez[r + j]);
  }
  // 2) E 推进：ca、cb 是按到边界距离算好的衰减层掩膜，
  //    真空处取 ca=1、cb=0，于是退化成普通更新
#pragma omp parallel for collapse(2) schedule(static, 16)
  for (int i = 1; i < nx - 1; ++i)
    for (int j = 1; j < ny - 1; ++j) {
      const int r = i * ny + j;
      const float curl = (hy[r] - hy[r - ny]) - (hx[r] - hx[r - 1]);
      ez[r] = ca[r] * ez[r] + cb[r] * ke * curl;  // 乘加，无分支
    }
}
// 稳定条件：二维取 dt <= dx / (c * sqrt(2))，三维取 sqrt(3)。
// 衰减层不是 PML：PML 要给每个分量配一条拉伸坐标的记忆变量，
// 内存与浮点量各多一份，但掠入射与色散介质下的性能好得多。
\end{lstlisting}

清单~\ref{lst:ch14-fdtd} 里最值得盯的是两个差分方向的\emph{不对称}：$H_x$ 沿连续方向差分（步长 $1$，可向量化），$H_y$ 沿跨行方向差分（步长 \Tp{ny}）。一份二维实现真正被向量化的只有一个方向，另一个方向要靠分块（blocking）或把循环交换成按列遍历；三维情形下三个方向各有两个分量受益，比例才勉强好看。把这两个方向的邻居写反，是本域最经典的静默错误：结果不发散，只是\emph{很稳定地}给出歪斜的散射方向图。

\Lib{Meep} 与 \Lib{openEMS} 的内核形状与这段一致：分量分离的连续数组、材料系数按元胞预乘成两张掩膜（对应 \Tp{ca}、\Tp{cb}），边界与 PML 用额外的记忆变量数组。差别都在工程细节：各向异性介质把标量系数换成 $3\times 3$ 张量乘积，每元胞浮点量上升；色散材料把更新变成 $(E,\ \text{辅助极化})$ 的联立推进；自适应网格加密让 halo 交换变成非对齐的。这些全是内存与分支的账，不是算法的账。

\section{约定与数据：端口、参考阻抗、单位与稳定性}

\textbf{S 参数离开参考阻抗就没有意义。} 一个``端口''在全波求解器里有两种身份：波导端口（给定截面模式与传播常数，模式幅值归一到单位功率）与集总端口（一个间隙电压源串一个阻抗）。端口报出的 $S$ 总是相对参考阻抗 $Z_0$ 的，而电压波归一化与功率波归一化之间相差一个因子 $\sqrt{Z_{0i}/Z_{0j}}$（谁是分子取决于 $a$、$b$ 的定义，跨工具比对前必须问清）。于是同一段结构在 $50$ 欧与 $75$ 欧参考下的 $S_{11}$ 数值本就不同——这不是误差，是定义。多端口结果是一个 $n \times n \times n_f$ 的复数组，其下标顺序（哪个最快变）与复数排布（交错还是分面）要按本册``数值接口的共同约定''一章的口径核对，否则读入端会把 $S_{12}$ 当成 $S_{21}$。

\textbf{近场到远场。} 全波解只在计算域内部可靠，域外靠等效原理：在靠近源处取一个封闭的惠更斯面 $S$，把面上切向 $E$、$H$ 当作等效磁流 $\mathbf{M} = -\hat n \times \mathbf{E}$ 与等效电流 $\mathbf{J} = \hat n \times \mathbf{H}$，再对自由空间 Green 函数积分。远场是该积分在 $r$ 很大时取驻相位的直接结果：

\begin{equation}
\mathbf{E}_{\text{far}}(\hat r) \ \propto\ \oint_{S}
\Bigl[\hat r \times \mathbf{M}
- \eta_0\, \hat r \times (\hat r \times \mathbf{J})\Bigr]\,
e^{\,j k\, \hat r \cdot \mathbf{r}'}\, \mathrm{d}S'
\label{eq:ch14-n2f}
\end{equation}

式~\eqref{eq:ch14-n2f} 的全部工程内容都在采样密度上：相位项 $k\,\hat r \cdot \mathbf{r}'$ 要求观察面上每 $\lambda/2$ 至少一个采样点，否则方向图出现栅瓣；$1/r$ 幅度与归一化方向图分别由式前的系数与 $|\mathbf{E}_{\text{far}}|^{2}$ 给出；观察面尺度为 $D$ 时，远场判据通常取 $2D^{2}/\lambda$（电大口径下这一距离可达数百米，所以测量端也要做一次近远场变换）。各家求解器把这段做成同一个动作：先声明一个包围盒作为惠更斯面，再在后处理里对一组立体角点做式~\eqref{eq:ch14-n2f} 的求和。

\textbf{激励类型与波形。} 硬源每步把某处 $E$ 直接写成给定值，简单但破坏网格的散度相容性并且会在源点反射；软源把激励作为电流项加进更新方程，与 Yee 网格的无散性共存，是主流做法。总场—散射场分解用于平面波入射：用一个闭合面上的软源把区域分成总场区与散射场区，边界上按分量做加减。宽带脉冲（调制高斯脉冲及其导数）一次给全带；窄带或连续波要跑到稳态，成本随 $Q$ 上升。脉冲越窄高频能量越多，也越容易因为网格欠分辨把噪声推进结果——``谱覆盖''与``分辨率''必须在同一张表里权衡。

\textbf{单位与``谁规定 $c=1$''。} \Lib{Meep} 走自然单位：令 $c=1$，长度、时间、频率互为倒数，用户唯一要保证的是\emph{自洽}（所有尺寸同一长度单位、分辨率按此定义、频率取 $1/\text{长度}$）。好处是归一化频率 $a/\lambda$ 与无量纲材料参数直接可比，代价是与现实器件的换算必须由用户做一次，而这一次换算就是错误产地。\Lib{openEMS} 与 \Lib{gprMax} 走 SI 口径（毫米、吉赫、皮秒），Courant 数由 $\Delta$ 与 $c$ 显式给出。两种口径都要求：网格按\emph{介质中}的最短波长取——真空 $30$ GHz 对应十毫米，在折射率 $4.0$ 的半导体里只剩二点五毫米——否则数值色散与导波失真会同时出现。

\textbf{稳定性、数值色散与吸收边界。} 相容的离散 $E$、$H$ 保证 $\nabla \cdot \mathbf{B} = 0$ 在机器精度内成立。离散色散关系给出每格波长的相速误差，它随 $(\Delta/\lambda)^{2}$ 下降且随传播方向变化：于是各向同性的散射体在粗网格下会给出歪斜的方向图，偶极子辐射会出现非物理的方向瓣。PML 的反射系数由层数、拉伸剖面与离散化三者共同决定，掠入射与色散介质最难吸收——很多``谐振频率对但线宽偏''的结果就出在这里。

\begin{warning}[title={稳定但物理错的两条最快路径}]
\begin{itemize}
\item \textbf{把时间步取在稳定性边界上。} Courant 数不超过 $1/\sqrt{D}$ 只保证线性均匀介质的格式不发散。一旦叠加色散材料（辅助方程自身有稳定限制）、集总元件（等效极小的 $\Delta t$）、CPML 或局部加密网格，真实边界会整体下移。轻则几千步后从角点炸出高频噪声，重则不发散、只是缓慢漂移——后者更难发现。对策：$\Delta t$ 由\emph{实际}最小特征尺寸与\emph{实际}最高频率决定并留 $10\%$ 到 $20\%$ 余量，另外跑一次``关掉源''的空运行，检查残差是否单调衰减。
\item \textbf{单位与自然单位之间漏了一次换算。} 折射率 $4.0$ 的介质里仍按真空波长取网格，每波长格数骤减，相位精度被数值色散吞掉，但能量仍然守恒、曲线仍然光滑。同型错误还有把微米当毫米、把 GHz 当 THz：$S_{21}$ 看起来仍然像谐振，只是位置差了一个数量级。对策：脚本开头写死单位声明，并用同一套单位把解析验证式重算一遍。
\item \textbf{端口参考阻抗没记录。} 把两个工具的 $S_{11}$ 叠在同一张图上而参考阻抗不同，看到的``差异''完全由归一化产生。任何导出 $S$ 参数的通路都应把 $Z_0$ 一起带走，并在读入端显式重归一，而不是假设双方一致。
\end{itemize}
\end{warning}

\section{验证：解析解对拍、守恒检查与网格收敛}

电磁仿真没有``看一眼波形''这种验收手段，因为它太容易自我说服。可信的验证只有三类：解析对拍、守恒与对称性检查、网格收敛研究。

\textbf{解析对拍。} 四个必须过的基准：其一，矩形波导主模 $TE_{10}$ 的截止频率 $f_c = 1/(2a\sqrt{\mu\epsilon})$（$a$ 为宽边）以及截止附近的相位常数与群速；其二，微带线的特征阻抗，Hammerstad 与 Jensen 的闭式给出宽厚比、有效介电常数与 $Z_0$（精度约百分之几，够用作对拍）；其三，自由空间损耗 $L = (4\pi d/\lambda)^{2}$ 与两副已知增益天线的链路预算；其四，介质平板的 Fabry--Pérot 谐振间隔 $\Delta f = c/(2nd)$。清单~\ref{lst:ch14-check} 给出前三个的 \CPP{} 实现与三分辨率外推。

\begin{lstlisting}[style=cppstyle,label={lst:ch14-check},caption={三个解析基准的实现与三分辨率 Richardson 外推}]
#include <cmath>
static const double C0 = 299792458.0;   // 真空光速，单位 m/s
static const double PI = 3.141592653589793;

// 1) 矩形波导 TE10 的截止频率：a 是宽边（单位米）
double wg_te10_fc(double a, double er = 1.0) {
  return C0 / (2.0 * a * std::sqrt(er));
}

// 2) 自由空间损耗，返回 dB：d 米、f 赫兹
double fspl_db(double d, double f) {
  return 20.0 * std::log10(4.0 * PI * d * f / C0);
}

// 3) 微带线特征阻抗，Hammerstad-Jensen 近似
double microstrip_z0(double w, double h, double er) {
  double ee = (er + 1.0) / 2.0 + (er - 1.0) / 2.0
            * std::pow(1.0 + 12.0 * h / w, -0.5);  // 有效介电常数
  double u = w / h;
  if (u <= 1.0)                          // 窄带分支
    return 60.0 / std::sqrt(er) * std::log(8.0 / u + u / 4.0);
  return 120.0 * PI / (std::sqrt(ee)      // 宽带分支
         * (u + 1.393 + 0.667 * std::log(u + 1.444)));
}

// 二阶格式的误差按 1/r^2 收缩，因此比值 2 的三分辨率
// 可以做一次 Richardson 外推，给出比最细网格更好的估计。
double richardson(double mid, double fine) {
  return fine + (fine - mid) / 3.0;
}
\end{lstlisting}

清单~\ref{lst:ch14-check} 的要点不在公式（任何一本微波教材都有），而在\emph{对拍这件事必须由你自己实现一遍}：抄工具的输出永远发现不了工具的错。\PY{} 门面侧的对照见清单~\ref{lst:ch14-meep}——同一个二维设置，比较的仍然是同一条量：把归一化频率换算成 Hz 后与 \Fn{wg\_te10\_fc} 对表。

\begin{lstlisting}[style=pythonstyle,label={lst:ch14-meep},caption={Meep 的 Python 门面：分辨率、线源、PML 与通量监视}]
import meep as mp
import numpy as np

res = 30                        # 每单位长度的网格数（Meep 用 c=1）
cell = mp.Vector3(24, 12, 0)    # 第三个分量为 0 即二维
pml = [mp.PML(1.5)]             # 拉伸坐标加记忆变量
fcen, df = 1.0 / 0.8, 0.4       # 中心频率与带宽，均为归一化量

src = [mp.Source(mp.GaussianSource(fcen, fwidth=df),
                 component=mp.Ez,        # 与清单 ch14-fdtd 同一组场
                 center=mp.Vector3(-8, 0),
                 size=mp.Vector3(0, 8))]  # 线源沿 y 展开

sim = mp.Simulation(cell_size=cell, boundary_layers=pml,
                    resolution=res, sources=src,
                    default_material=mp.Medium(epsilon=1))
mon = sim.add_flux(0.5 * fcen, fcen, 1.5 * fcen,
                   mp.FluxRegion(mp.Vector3(6),
                                 size=mp.Vector3(0, 8)))
sim.run(until_after_sources=300)   # 源关后再跑 300 个时间单位

print("通量谱:", mp.get_fluxes(mon))          # 每个频点一个数
line = sim.get_array(direction=mp.Y, center=mp.Vector3(0),
                     components=[mp.Ez])
print("峰值 |Ez|:", np.max(np.abs(line)))
\end{lstlisting}

这段门面代码里有三处值得注意。\Tp{get\_array} \Bk{} 返回的是拷贝还是视图、并行时是否只有根进程拿到完整数据，这类细节决定后处理能否与下一步仿真重叠；\Tp{add\_flux}（新版本可能写作相邻的名字，以所装版本文档为准）把时域场与源做相关从而得到频域通量，这就是``时域一次运行得宽带结果''的具体实现；\Tp{until\_after\_sources} 用``源关闭之后再跑多少''而不是绝对时长来设定运行长度，因为它衡量的正是场衰减完毕所需的时间，而这个量与 $Q$ 直接相关。

\textbf{守恒与对称性检查。} 三条几乎不要钱的检查：功率平衡（入射通量 $=$ 反射 $+$ 透射 $+$ 体积损耗，误差应在百分之一以内）；互易性（无源、线性、非磁化结构上 $S_{ij} = S_{ji}$，端口激励互换后重跑必须复现）；散度残差（$\nabla \cdot \mathbf{D}$ 与自由源密度的偏差，以及 $\nabla \cdot \mathbf{B}$ 在 Yee 网格上应保持在小数点后很多位）。任何一条不过，先怀疑设置而不是接受``数值方法本来就这样''。

\textbf{网格收敛研究。} 至少三个分辨率（相邻比值 $2$），跟踪一个\emph{可测标量}（谐振频率、$S_{11}$ 的零点位置、RCS 峰值），检查误差是否按二阶格式应有的 $1/r^{2}$ 收缩。收缩就外推，不收缩就找不连续：材料界面未对齐网格、PML 与源太近、端口离不连续结构过近、色散参数超出模型适用带宽。这一步是本域唯一可信的单元测试，因为它同时回答``分辨率够不够''与``实现有没有错''——一个稳定但物理错的实现通常表现为\emph{不收敛到解析值}。

\begin{bestpractice}[title={一次全波仿真的四道验收关}]
\begin{itemize}
\item \textbf{先跑空域。} 去掉所有结构，只留源与吸收边界：波形应当干净、边界反射应当低于预期、关掉源之后能量应当单调衰减。这一步抓的是设置错，不是物理错。
\item \textbf{解析对拍先行。} 波导截止频率、微带阻抗、自由空间损耗、Fabry--Pérot 间隔，四选二跑通再谈你的器件；清单~\ref{lst:ch14-check} 就是为这一步准备的。
\item \textbf{守恒与互易做成脚本。} 功率平衡、$S_{ij} = S_{ji}$、散度残差三条自动检查，每次改网格或材料都跑一遍，把它当回归测试而不是临时验证。
\item \textbf{报一个数就附三个分辨率。} 给出漂移量与外推值；没有收敛数据的电磁结果，在本册的口径里等同于未验证。
\end{itemize}
\end{bestpractice}

\begin{longtable}{@{}>{\raggedright\arraybackslash}p{0.13\textwidth}
  >{\raggedright\arraybackslash}p{0.11\textwidth}
  >{\raggedright\arraybackslash}p{0.12\textwidth}
  >{\raggedright\arraybackslash}p{0.12\textwidth}
  >{\raggedright\arraybackslash}p{0.15\textwidth}
  >{\raggedright\arraybackslash}p{0.16\textwidth}@{}}
\caption{六个电磁求解层次：方法、接口语言、许可证基调、并行与典型用途}\label{tab:ch14-solvers}\\
\toprule
库 & 方法 & 接口语言 & 许可证基调 & 并行能力 & 典型用途 \\
\midrule
\endfirsthead
\toprule
库 & 方法 & 接口语言 & 许可证基调 & 并行能力 & 典型用途 \\
\midrule
\endhead
\Lib{Meep} & 时域 FDTD，Yee 加 CPML &
  \CPP{} 内核，Scheme 门面与 \PY{} 绑定 &
  GPL 系 copyleft，以仓库 LICENSE 为准 &
  MPI 域分解、OpenMP 线程，较新版本可卸载 CUDA &
  光子学、色散与非线性材料、伴随式反向设计 \\
\Lib{openEMS} & 时域 FDTD 加子网格模型 &
  \CPP{} 内核，Lua 门面，XML 中间表示 &
  GPL 系 copyleft，以仓库 LICENSE 为准 &
  单机多核为主，分布式支持以版本说明为准 &
  封装、天线馈电、生物电磁与电磁兼容 \\
\Lib{gprMax} & 时域 FDTD，Z 变换色散 &
  \PY{} 门面，Cython 与 OpenCL 内核 &
  GPL 系 copyleft，以仓库 LICENSE 为准 &
  OpenMP 线程、OpenCL 设备、MPI 多进程 &
  探地雷达、无损检测、比吸收率评估 \\
\Lib{S4}、\Lib{S4v2} & RCWA，即傅里叶模法 &
  \CL{} 与 \CPP{} 内核，Lua、MATLAB 门面 &
  宽许可（MIT 基调），以仓库声明为准 &
  单机稠密线性代数，线程来自底层数学库 &
  光栅、超表面与光子晶体薄膜的衍射效率 \\
\Lib{NEC2}、\Lib{4nec2} & 细线 MoM，频域 &
  \FTN{} 原文经 \CL{} 翻译，GUI 前端 &
  公有领域或宽松条款，随分支不同 &
  稠密复矩阵，基本单机运行 &
  线天线输入阻抗、方向图与近场耦合 \\
商用全波平台 & FEM、FIT、MoM 加 FMM、SBR &
  专有 GUI，内核语言不公开 &
  商业闭源，按站点或核数计费 &
  集群与多频点并行，取决于许可与调度 &
  芯片封装、天线平台、雷达截面与合规评估 \\
\end{longtable}

\begin{guideline}[title={本章的选型口径}]
\begin{itemize}
\item 先算 $kL$ 与每波长格数，再选方法：电小且宽带用 FDTD；封闭谐振与本征模用 FEM；开放辐射的电大金属结构用 MoM 加多极子；严格周期结构用 RCWA；$kL$ 上千且只要方向图与 RCS 用射线类。
\item 再看要什么信息：要端口上的宽带 $S$ 参数就必须把端口类型、参考阻抗与模式归一化写进设置；要 $Q$ 就必须把时域窗口或频域扫描范围拉到线宽的三倍以上；要近场分布则优先时域，因为它给的是同一个时刻的全部场。
\item 语言层的判断只有一条：几何与材料交给门面（改动频繁、表达复杂），场推进与线性代数留在 \CPP{}/\CL{} 内核（性能与内存布局敏感），后处理的谱与统计交给 \PY{}。哪个库违反了这条，就能预判它在哪个环节上出问题。
\end{itemize}
\end{guideline}
% ====================================================================

% ====================================================================
% ====================================================================
\chapter{射频链路与软件无线电}

上一张图景是``用 Maxwell 方程算出结构''，这一章是``用样点把信号搬回来''。两者共享同一套底层纪律（复数、采样、单位、参考阻抗），但约束的方向相反：全波求解器跟规模与带宽赛跑，接收链路跟\emph{截止期}与\emph{字节率}赛跑。这一域的语言层形态与本册主线高度一致：设备驱动只给 \CL{} ABI，样点级循环必须由 \CPP{} 承担，厂商差异被抽进一个 C 或 C++ 的抽象层，而\emph{所有}看起来像``软件无线电''的东西——流程图、扫频、标定脚本、频谱图——都在 \PY{} 那一侧。这里有个少见的干净分工：速率以兆样点每秒计的归内核，每小时才动一次的归门面。

边界声明：本章讲数据通路、设备抽象、调度与建模工具。滤波器结构与抽取算法本身交给本册讲 FIR 与 IIR 的那一章，谱估计与窗函数交给讲 FFT 与谱分析的那一章，天线与匹配网络的物理设计参见上一章``电磁仿真：方法选择与 C/C++ 内核''里的端口与参考阻抗一节。

\section{从天线到比特流：六段通路与其表示层}

\textbf{第一段：低噪声放大与线性度预算。} 链路的第一级决定整条链的噪声系数：Friis 公式 $F = F_1 + (F_2 - 1)/G_1 + (F_3 - 1)/(G_1 G_2) + \cdots$，因此第一级的增益要足够大，大到把后级的噪声折算掉，又不能大到让后级或自身的 $P_{1\text{dB}}$ 先撞限。三阶交调点与压缩点之间的余量（常见的经验值是 $10\,\text{dB}$ 上下，器件相关）就是``能同时收下多少个强干扰''的预算，它是硬件指标，软件层只能被动接受。

\textbf{第二段：下变频与中频采样。} 三条路：超外差（一次或多次变频到一个固定中频再数字化，镜频由模拟滤除，优点是选择性好、缺点是频率规划复杂）；直接变频（zero-IF，本振就是载频，把 I、Q 两路各自数字化，省掉中频与镜像滤波器，代价是直流偏置、本振泄漏自混频、以及器件闪烁噪声正好落在基带）；中频欠采样（把信号放在第 $k$ 个 Nyquist 区内，直接以低于射频的速率数字化，频谱按该区间的奇偶发生翻转，抽取与翻折要在数字域显式处理）。三条路的共同结论是：\emph{模拟域省下来的器件，都要在数字域用算力还回去}。

\textbf{第三段：正交解调与复基带表示。} 解析信号 $x(t) = I(t) + jQ(t)$ 把双边带变成单边带，于是采样率的下界从 $2 f_{\max}$ 降到信号带宽本身——这是复基带存在的全部理由，也是为什么后面所有 block 的 \Tp{sizeof(gr\_complex)} \Bk{} 都是 $8$。正交解调的实现误差有三种可测形态：直流偏置（I、Q 各自的中值不为零，表现为载波处一根假谱线）、增益不平衡（两路幅度差）、正交误差（两路相位偏离 $90^\circ$ 若干度）。后两者共同决定镜像抑制比，量级上几十 dB 已算未校正的好值，而它\emph{不会}随 ADC 位数上升——所以设备抽象层都把这两项做成显式接口（\Lib{SoapySDR} 的 \Fn{setIQBalance} 与 \Fn{setDCOffset}，\Lib{UHD} 侧的 \Fn{set\_rx\_dc\_offset} 与 \Fn{set\_rx\_iq\_balance}）。

\textbf{第四段：ADC 位数、噪声底与 SFDR。} 理想 $N$ 位量化给出 $\text{SNR} \approx 6.02 N + 1.76\ \text{dB}$；实际器件的上界由另外两件事决定。其一是采样抖动：$\text{SNR}_{\text{jitter}} = -20\log_{10}(2\pi f_{\text{in}} t_j)$，输入频率越高越苛刻，因此在射频直接采样的方案里时钟质量比位数更值钱。其二是无杂散动态范围 SFDR：谐波、时钟馈通与开关杂散构成的\emph{确定性}谱线，它们不随平均次数消失，也不受位数支配；链路上能同时看住的强信号与弱信号之比基本由 SFDR 而不是 SNR 决定。有效位数 $\text{ENOB} = (\text{SNR} - 1.76)/6.02$ 是这两件事合并后的实测口径。

\textbf{第五段：过采样、处理增益与数字下变频。} 把带宽 $B$ 的信号用 $f_s \gg 2B$ 数字化，量化噪声摊在 $f_s/2$ 的整个区间上，抽取到 $B$ 时信号无损、噪声按比例收缩——每倍过采样约 $3.01\ \text{dB}$ 处理增益（再叠加 dither 时效果更干净）。数字下变频的标准三件套是 NCO（数值振荡器，把载频搬到零）、CIC（级联积分梳状，整数倍抽取的第一级，便宜但通带跌落大）、以及补偿与成形用的 FIR。这三件套的系数与延迟链就是本章后面所有性能讨论的对象。

\textbf{第六段：采样率与批量。} 设备交给软件的永远是一段一段的缓冲：一次调用请求 $n$ 个样点、返回实际交付数、并附带一个时间戳。字节率由三段乘积给出，而这道乘法常常就是方案的死刑或放行：$61.44$ Msps、两路、\Tp{CS16}（每复样点 $4$ 字节）是 $491.5$ MB/s，已经逼近万兆以太网的有效吞吐，因此前传与多通道系统要往更省的格式或更宽的口走；同一速率换成 \Tp{CF32} 是 $983$ MB/s，直接不可行。反过来，USB 2.0 的实际可用吞吐只有几十 MB/s，这就是为什么某些电视芯片方案的接收机上限被定在二十多 Msps——不是射频前端不行，是总线不行。

表~\ref{tab:ch15-format} 把常见的复基带存储格式排开。这里的核心张力很朴素：\emph{总线要窄，内核要直}。

\begin{table}[H]
\centering
\caption{复基带存储的四种形态：字节数、内核代价与产地}
\label{tab:ch15-format}
\small
\begin{tabular}{@{}>{\raggedright\arraybackslash}p{0.14\textwidth}
  >{\raggedright\arraybackslash}p{0.11\textwidth}
  >{\raggedright\arraybackslash}p{0.22\textwidth}
  >{\raggedright\arraybackslash}p{0.35\textwidth}@{}}
\toprule
格式 & 每复样点 & 内核侧代价 & 常见产地与用途 \\
\midrule
\Tp{CS16}（交错 int16） & $4$ B &
  符号扩展、交织转分面（一次 shuffle 或 strided copy）、\Fn{int} 到 \Fn{float} &
  设备与总线的缺省运输格式；\Lib{UHD}、\Lib{SoapySDR} 都支持，短脉冲采集够用 \\
\Tp{CF32}（交错 float32） & $8$ B &
  零转换：驱动里就地复数运算，缩放由定点转浮点那一步付清 &
  通用中间格式，也是所有 \CPP{} 内核的偏好布局；带宽是 \Tp{CS16} 的两倍 \\
分面（I 与 Q 各一块） & $8$ B &
  单路滤波与统计连续可向量化；复数乘加要先把两路重新交织或分别算 &
  文件、GPU 缓冲与图像式管线更常见，设备上少 \\
\Tp{CS16B}（12 位加噪声整形） & $4$ B &
  低 $4$ 位是整形噪声，必须移位掩码；线性度只在标称 $12$ 位内成立 &
  数模转换器把位数省成字节数的产物，见 \Lib{bladeRF} 与部分集成收发芯片 \\
\bottomrule
\end{tabular}
\end{table}

交错与分面的 CPU 代价差在哪，值得说得再具体一点。复数乘加 $y \mathrel{+}= c \cdot x$ 在交错 \Tp{CF32} 上要用一次 shuffle 把实部与虚部分开、反对称乘（$c_r x_r - c_i x_i$ 与 $c_r x_i + c_i x_r$）再拼回去；在分面布局上两条流水线各自连续，代价变成``一次乘法两组累加''，向量化器更容易吃下。但分面把缓存行也拆成两份，于是每样点的内存事务数上升。这就是为什么这一域既没有``最好布局''也没有``最省格式''：总线侧要窄，内核侧要直，抽象层的价值就是让这两件事在同一份代码里各得其所。清单~\ref{lst:ch15-py} 是门面侧的最小对照：同一台设备，\PY{} 拿到的是复数数组，格式协商与缩放都发生在内核。

\begin{lstlisting}[style=pythonstyle,label={lst:ch15-py},caption={Python 门面抓一段复基带并做谱：格式协商与 dBFS 口径}]
import SoapySDR
from SoapySDR import SOAPY_SDR_RX        # 常量也从同一绑定导出
import numpy as np

dev = SoapySDR.Device({"driver": "remote"})   # 键值对是唯一的寻址方式
print(dev.getStreamFormat(SOAPY_SDR_RX, 0))   # 问一句比假设便宜
dev.setSampleRate(SOAPY_SDR_RX, 0, 8e6)
dev.setFrequency(SOAPY_SDR_RX, 0, 2.44e9)
print("增益名称:", dev.listGains(SOAPY_SDR_RX, 0))  # 单位是档不是 dB

n = 4096
buf = np.empty(n, dtype=np.complex64)    # CF32：与内核零转换
stream = dev.setupStream(SOAPY_SDR_RX, "CF32", [0])
dev.activateStream(stream)
# 返回形态随绑定版本不同：新版给结果对象，旧版给三元组
got = dev.readStream(stream, buf, n, 0, 1000000)
buf = buf[:int(got.read())]              # 实际交付数才是长度
dev.closeStream(stream)

w = buf * np.hanning(buf.size)           # 窗：抑制泄漏，不改变 dBFS
sp = np.fft.fftshift(np.fft.fft(w, 4096))
db = 20 * np.log10(np.abs(sp) + 1e-9)    # 参考是满刻度，即 dBFS
print("峰值:", db.max(), "噪声底:", np.median(db))
# 噪声底与峰值之差读作 SFDR；单 bin 与等效噪声带宽要分清
\end{lstlisting}

\section{设备抽象层：\Lib{SoapySDR} 的核心与 \Lib{UHD} 的板级控制}

\Lib{SoapySDR} 想解决的问题只有一个：把 \Lib{Ettus}、\Lib{bladeRF}、\Lib{LimeSDR}、\Lib{RTL-SDR}、\Lib{PlutoSDR} 这些互不相容的 SDK 收成一份代码。它的做法有两层值得学。

第一层是\textbf{键值对贯穿全部接口}。枚举设备、打开设备、建流、查能力，全都收 \Tp{SoapySDR::Kwargs}（本质是一张字符串到字符串的表，写作 \Tp{driver=remote} 这类形式）：\Fn{find} 返回匹配的设备参数列表，\Fn{make} 用其中一条把设备实例化，\Fn{setupStream} 再收一组流参数。好处不只是省参数类：厂商特有的能力\emph{不需要}进 API——它落成一个新键，老代码看不见它也不会坏。第二层是\textbf{能力先问再调}：\Fn{getSettingInfo} 返回范围对象（步长、单位、默认值），于是采样率、带宽、频率、时延都是``先查界、再设定、然后读回实际值''三步；\Fn{setSampleRate} 之后仍要调 \Fn{getSampleRate}，因为设备的实际速率几乎总被时钟树取整过。这一条纪律后面还会再出现一次：\emph{请求值不是达成值}。

流层是同一套形状的连续剧：\Fn{listStreamArgs} 与 \Fn{getStreamFormat} 问``这个方向能给我什么字节形状''，\Fn{setupStream} 按 \Tp{CS16} 或 \Tp{CF32} 建流，\Fn{acquireReadBuffer} 给出设备侧缓冲（避免多一次拷贝），\Fn{readStream} 返回实际样点数并把状态写进一个 flags 引用参数（标志名都带 \Tp{SOAPY\_} \Bk{} 前缀：\Tp{HAS\_TIME}、\Tp{TIME\_ERROR}、\Tp{END\_OF\_BURST}、\Tp{HAS\_ERROR}），收尾用 \Fn{deactivateStream} 与 \Fn{closeStream}。负返回值是错误码而不是样点数——把它当长度用是这一层最经典的事故。时钟与定时通过 ``clocking'' 那组接口暴露：\Fn{setClockRate} 管主时钟、\Fn{setTimeGen0} 与 \Fn{setTimeUnknownPps} 管时间基准、\Fn{setTimeAdjust} 给一个恒定偏移去贴住外部时基；GPIO、寄存器与仪表分别有各自的枚举接口（\Fn{writeGpio}、\Fn{writeSetting}、\Fn{listSensors}），传感器是``只读、可能很慢、可能异步''的那一类状态源。

最关键的结构决定是：\textbf{核心是 \CPP{}，插件 ABI 是 \CL{}}。宿主看到的是一个不透明句柄 \Tp{SoapySDRDevice*} 与一组以 \Tp{SoapySDR\_} \Bk{} 为前缀的 C 函数符号，厂商模块用 C 接口注册自己。于是模块可以由不同编译器、不同标准库、甚至不同 \CPP{} 标准的代码实现，宿主升级不必重编全部后端；\PY{} 绑定（SWIG 生成）与命令行工具 \Tp{SoapySDRUtil} \Bk{} 也都压在这层 C 上。清单~\ref{lst:ch15-soapy} 给的就是这条链：收流、解交织、抽取加 FIR，全在 \CPP{} 侧。

\begin{lstlisting}[style=cppstyle,label={lst:ch15-soapy},caption={SoapySDR 收流：协商速率、解交织、整数倍抽取加 FIR}]
#include <SoapySDR/Device.hpp>
#include <SoapySDR/Types.h>
#include <cstdio>
#include <vector>

// 输入按 I、Q 分面存放，因此抽头循环沿连续方向走。
// 输出样点数是输入的 1/keep；尾部不足一个抽取比就丢掉。
static void decim_fir(const float* xr, const float* xi, int n,
                      const float* taps, int nt, int keep,
                      std::vector<float>& yr,
                      std::vector<float>& yi) {
  const int m = (n - nt) / keep;              // 保证窗口完整
  yr.assign(m > 0 ? m : 0, 0.f);
  yi.assign(m > 0 ? m : 0, 0.f);
  for (int k = 0; k < m; ++k) {
    float ar = 0.f, ai = 0.f;
    const int b = k * keep;
    // 抽头是实系数：反对称复乘退化成两次乘加，易向量化
    // 提示：把 taps 反序存放可让内层变成纯 forward FMA
    for (int t = 0; t < nt; ++t) {
      ar += taps[t] * xr[b + t];
      ai += taps[t] * xi[b + t];
    }
    yr[k] = ar; yi[k] = ai;
  }
}

int main() {
  auto found = SoapySDR::Device::find(SoapySDR::Kwargs());
  if (found.empty()) { std::puts("没有可见设备"); return 1; }
  SoapySDR::Device* dev = SoapySDR::Device::make(found[0]);

  dev->setSampleRate(SOAPY_SDR_RX, 0, 20e6);  // 请求值常被取整
  const double fs = dev->getSampleRate(SOAPY_SDR_RX, 0);
  std::printf("实际速率 %.0f\n", fs);          // 用它算抽取比
  dev->setFrequency(SOAPY_SDR_RX, 0, 2.44e9, SoapySDR::Kwargs());
  dev->setGain(SOAPY_SDR_RX, 0, 40.0);        // 增益单位由设备定义

  std::vector<float> taps(32, 1.f / 32.f);    // 演示用低通
  auto stream = dev->setupStream(SOAPY_SDR_RX, "CF32", {0});
  dev->activateStream(stream);
  const int req = 4096, keep = 4;
  std::vector<float> buf((size_t)req * 2);
  std::vector<float> yr, yi;
  long long ts = 0; int flags = 0, drops = 0;
  for (int round = 0; round < 64; ++round) {
    void* buffs[1] = { buf.data() };
    int got = dev->readStream(stream, buffs, req, flags, ts,
                              1000000);      // 超时，单位微秒
    if (got < 0) { ++drops; continue; }       // 负值是错误码
    if (flags & SOAPY_SDR_TIME_ERROR) ++drops; // 时间戳不连续也是丢
    std::vector<float> xr(got), xi(got);
    for (int i = 0; i < got; ++i) {           // 交织 -> 分面
      xr[i] = buf[2 * i]; xi[i] = buf[2 * i + 1];
    }
    const float* xp = xr.data();             // 拆成短名避免超长行
    const float* ip = xi.data();
    decim_fir(xp, ip, got, taps.data(), 32, keep, yr, yi);
    // 这里可以消费 yr、yi：谱、解调、或交给下一级缓冲
  }
  std::printf("丢弃或错误块数: %d\n", drops);
  dev->deactivateStream(stream);
  dev->closeStream(stream);
  SoapySDR::Device::unmake(dev);
  return 0;
}
\end{lstlisting}

\Lib{UHD} 是 Ettus 设备的原生层，形态上是同一套思想的深度版：\CPP{} API 为主体，附一份 \CL{} 绑定与一份 \PY{} 绑定，三者共享实现。它的抽象比 SoapySDR 多几格：\Tp{multi\_usrp} \Bk{} 管多通道与板间，\Tp{radio\Bk controller} 管收发芯片，\Tp{dboard} 管子板，\Tp{mboard} 管主控与温度、频率、电源等\emph{仪表}（\Fn{get\_sensor} 与 \Fn{get\_mtd}）。数据面是 \Fn{get\_stream\_cmd} 加 \Fn{issue\_stream\_cmd}：命令里带模式（\Tp{START}、\Tp{STOP}）、样点数与起始时间，时间用 \Tp{time\_spec\_t} 表达（板载 tick 与秒之间由 \Fn{get\_ticks\_per\_sec} 换算，这就是``请求值不是达成值''的另一种表现），接收侧从元数据里读 \Tp{timeSpec}、\Tp{out\_of\_burst} 与错误码。时钟同步的动词在这里最直白：\Fn{set\_clock\_source} 与 \Fn{set\_time\_source} 取 \Tp{internal}、\Tp{external} 或 \Tp{gpsdo}，\Fn{set\_time\_next\_pps} 在下一个秒脉冲把时间归零，多设备共用同一参考与同一秒脉冲就能保持相位关系。流格式在 UHD 侧写作 \Tp{fc32}、\Tp{sc16} 一类小写串，另有``on the wire''格式（板间或主机间实际传输的字节形状）与内存格式之分——这个区分正是本章第六段那道乘法的来源。

两层的取舍没有悬念：跨厂商工具链选 \Lib{SoapySDR}（\Lib{GNU Radio} 的 osmosdr 与 Soapy 模块就是用它把五种设备变成一种），追单厂商的能力上限（最高采样率、时间戳精度、突发模式、固件与寄存器通路）选 \Lib{UHD}。代价也很清楚：前者把差异压成键值对，因而有些能力只能靠额外键访问；后者把细节摊开，于是换设备要重写。

\section{图、调度与内核：\Lib{GNU Radio} 及其替代者}

\Lib{GNU Radio} 的分层是本册主线最典型的一例：\emph{内核是 \CPP{} 的类，图是 \PY{} 或 GRC 的对象}。一个 block 的实质是三个函数签名加一段系数。其一，\Fn{work} 或 \Fn{general\_work}：拿到本轮可用的输入指针与输出空间，算完报告产出数。其二，\Fn{forecast}：告诉调度器``要产出 $n$ 个输出我需要多少输入''，写少了会读到未初始化历史，写大了会永远等不到而卡住。其三，\Fn{set\_history} 与 \Fn{consume}：前者声明每次产出所需的输入长度（调度器据此在缓冲尾部多留一段），后者精确报告吃了多少；两者不一致时症状是``输出里混进重复样点''。标签（\Tp{gr::tag\_t}）带相对样点偏移随数据流走，是频率切换、突发边界与时延补偿的载体；控制面另有一条 message 通路（\Fn{message\_port\_register\_in} 与 \Fn{set\_msg\_handler}，参数用 PMT），慢变参数走它而不是走每个样点都查的共享变量。

调度粒度是这一层的性能核心。调度器每次交给 \Fn{work} 的输出配额以数千样点为上限（常见默认在 $4\text{K}$ 量级，随版本的缓冲参数可变），于是每次调用的固定开销要摊在几千个样点上——这正是把 block 写成 \CPP{} 的收益来源，也是 \PY{} block 只适合低速率与控制流的原因：每次跨语言调用都要穿过 pybind11 的转换层并持有解释器全局锁，而这笔开销按\emph{调用次数}而不是按样点数计。缓冲本身是环形且历史上用两块重叠的虚拟地址映射（让回绕不需要拷贝，\Lib{GNU Radio} 的实现沿用了这一族技术），配 \Fn{set\_min\_input\_buffer} 与 \Fn{set\_max\_output\_buffer} 调节驻留量；\Tp{TPB} 调度器给每个 block 一个线程，块数一多，线程切换与缓存颠簸就成了新的预算项，于是``用几层容器把小 block 合成大 block''在高采样率下是实际的优化手段。

out-of-tree 模块的工程路径也值得一提：\Tp{gr\_modtool} 生成骨架，\CPP{} 类加一层 pybind11 注册（\Lib{GNU Radio} 3.10 起用它替掉早年的 SWIG 路线），GRC 用 YAML 描述方块的外观与端口。代价是 ABI 耦合：模块与核心版本、Boost 或 ASIO 之类的基础库版本绑在一起分发。性能不足的部分普遍下沉到 \Lib{liquid-dsp} 这样的 \CL{} 内核库（滤波、重采样、快速傅里叶、同步算法都在这里），让 block 只做数据搬运。

替代品有两类。同一作者的 \Lib{Pothos} 走更纯粹的 ``C++ 图加插件''路线，取消 Python 参与运行时；\Lib{Orchestra} 一类的桌面框架把重点放在应用装配与可视化。两者的社区规模与模块生态都远小于 \Lib{GNU Radio}，因此本册把它们当作``形态参照''而不是默认选择——具体能力以各自文档为准，因为这条赛道的版本变化很快。

清单~\ref{lst:ch15-gr} 是一段可直接扩展的 \CPP{} block 骨架：抽取加 FIR，同时演示 \Fn{forecast}、history、\Fn{consume} 与标签透传。

\begin{lstlisting}[style=cppstyle,label={lst:ch15-gr},caption={GNU Radio 的 C++ block 骨架：forecast、history 与精确消耗}]
#include <gnuradio/io_signature.h>
#include <gnuradio/block.h>
#include <gnuradio/gr_complex.h>

class fir_decim_cc : public gr::block {
 public:
  fir_decim_cc(int decim, std::vector<float> taps)
      : gr::block("fir_decim_cc",
          gr::io_signature::make(1, 1, sizeof(gr_complex)),
          gr::io_signature::make(1, 1, sizeof(gr_complex))),
        d_dec(decim), d_taps(std::move(taps)) {
    // history：每产出一个样点要读多长的历史。
    // 调度器据此在输入缓冲末端多留 d_taps.size()-1 个样点。
    set_history((int)d_taps.size());
    set_relative_rate(1.0 / (double)decim);   // 供 tag 换算与下游推算
    set_output_multiple(1);                   // 按整块产出可再调
  }

  // forecast：产出 noutput_items 需要多少输入。
  // 写小了会算错，写大了本轮永远凑不齐，表现为吞吐塌陷。
  void forecast(int noutput_items,
                std::vector<int64_t>& need) override {
    need[0] = (int64_t)noutput_items * d_dec
            + (int64_t)d_taps.size() - 1;
  }

  int general_work(int noutput_items, std::vector<int64_t>& avail,
                   const gr::const_buffer_vector& in,
                   gr::buffer_vector& out) override {
    const gr_complex* x = in[0].cast<const gr_complex*>();
    gr_complex* y = out[0].cast<gr_complex*>();
    int made = 0;
    const int64_t ht = (int64_t)d_taps.size() - 1;  // 历史长度
    // 内层沿输入连续方向走：抽头反序存放以便纯 forward FMA
    while (made < noutput_items &&
           avail[0] >= (int64_t)made * d_dec + ht + 1) {
      const int b = made * d_dec;
      float ar = 0.f, ai = 0.f;
      for (size_t t = 0; t < d_taps.size(); ++t) {
        ar += d_taps[t] * x[b + t].real();
        ai += d_taps[t] * x[b + t].imag();
      }
      y[made] = gr_complex(ar, ai);           // 实抽头：无复数交叉项
      ++made;
    }
    consume(0, (int64_t)made * d_dec);        // 必须精确声明
    // 标签默认透传（策略 TPP_ALL_TO_ALL），突发边界要在此处改写
    return made;
  }

 private:
  int d_dec;
  std::vector<float> d_taps;
};
\end{lstlisting}

这段骨架里唯一会让人反复调试的是 \Fn{consume} 与 \Fn{forecast} 的一致性：本轮``吃多少''必须与``产出多少''按抽样率比严格对应，多吃会跳样点、少吃会重样点，而两种症状在频谱上都只是``多了几根不该有的线''。新版本把 \Tp{std::vector<int64\_t>} 一类签名换成了专用缓冲类型（\Tp{gr::const\Bk\_buffer\_vector} 与 \Tp{gr::buffer\Bk\_vector}，跨版本编译时以所装头文件为准），并把``速率''从 tag 迁到属性与 message；骨架的三个概念不变。

\section{前端与无源链路的建模：网络级、解析级与电路级}

链路层与电磁层之间的\textbf{通用货币是 S 参数}，因此这一节的主角没有样点，只有复矩阵。

\Lib{scikit-rf} 是这条路上最完整的 \PY{} 实现：\Tp{Network} \Bk{} 装一个 $n \times n \times n_f$ 的复数 $S$ 张量加频率轴与参考阻抗，Touchstone 文件（\Tp{.s2p} 一类）的头里就带参考阻抗，读入时不能丢。级联用 \Fn{connect}（按端口号互连，内部做一次信号流图求解）或运算符写法；重归一化用 \Fn{renormalize}；时延与损耗的去嵌入、差分与混合模式 $S$ 参数、插值到均匀频率栅（这是做时域变换的前提，DC 外推的方式直接决定 TDR 台阶的质量）都在同一层。\Lib{scikit-rf} 还提供传输线介质对象（同轴、微带、共面波导），把``几何尺寸加介电常数''换算成 $Z_0$ 与有效介电常数——上一章``电磁仿真：方法选择与 C/C++ 内核''里讲端口与解析基准的那两节，说的正是同一组公式的另一面。这一层是 \PY{} 主场的理由很干净：没有样点级循环，全部成本都在 numpy 与 scipy 已经付过钱的批量线性代数上，而交互式试错的价值远高于每调用一次的开销。

\Lib{libmobi} 代表另一形态：\CL{} 库加一份 \Tp{XML} 描述文件，把射频前端写成器件树——滤波器（原型综合与耦合矩阵）、匹配网络、传输线段、集总元件、放大器与混频器的行为模型——然后按扫频点算出整棵树的 $S$ 参数、噪声温度与非线性指标。它的定位是``链路计算器里的解析内核''：不需要网格、不需要时域推进，几微秒给一个频点，因此适合在优化循环里被调用上百万次。选择 \CL{} 的原因也正是本章反复出现的那条：要能被任何宿主（MATLAB、\PY{}、LabVIEW、命令行、FPGA 侧的软核）以最小代价嵌进去，且没有运行时依赖。

电路级还是 \Lib{SPICE} 家族的地盘：\Lib{ngspice} 是经典的 \CL{} 移植实现，\Lib{Xyce} 是 \CPP{} 重写并把稀疏直接解算与 MPI 并行放进同一架构（大规模电路与参数扫描是它的靶心），两者的 AC 分析做的就是复数节点方程——与本节的复基带是同一套代数，只是未知量从样点变成相量。场与路的协同有两条现实路径：其一是把全波结果降成端口参数交给电路核，通常需要一次有理拟合（vector fitting）把 $S(\omega)$ 变成可用的宽带宏模型，否则时域仿真里频率相关的数据无法进入网表；其二是把非线性或集总负载当作场域的边界条件，在场求解器的时间步里与电路状态联立推进（上一章提到的 \Lib{openEMS} 与 SPICE 内核协同的通道属于这一类）。两条路都贵，区别在于贵的东西是内存还是同步。

\section{实时通路的工程约束，以及为什么仍是 C/C++ 主场}

\textbf{缓冲与驻留。} 驱动线程到消费者线程之间只有一个被允许的结构：单生产单消费的环形缓冲，预分配、永不扩容。三个细节决定成败：其一，内存要\emph{预先触碰}（分配后写一遍），否则首次访问的页错误发生在数据通路上；其二，进程侧要锁定（\Tp{mlockall} \Bk{} 加 \Tp{MCL\_CURRENT}，Windows 上是 \Fn{VirtualLock}；交给 GPU 的宿主缓冲要用 pinned 分配），否则换页或 DMA 重映射会带来无界延迟；其三，生产者的写指针与消费者的读指针要各占一条缓存行，不然伪共享把吞吐打回标量水平。

\textbf{线程放置与优先级。} 数据通路线程要设实时策略（\Tp{SCHED\_FIFO} 配 rtprio），并用亲和性绑到隔离出来的核（\Tp{isolcpus} 与 \Tp{nohz\_full} 一类内核参数）；解释器、垃圾回收、日志与网络栈都留在其余核上。这不是``性能调优''而是\emph{噪声隔离}：一次几十微秒的调度抖动就能吃穿一个块的预算。软实时系统的通用形态是：内核路径绑核、门面线程不绑、二者通过带计数器的无锁队列对话。

\textbf{突发、丢样与它的静默性。} 总线路径上的时间粒度是硬的：USB 的微帧是 $125\ \mu s$，以太网的一包是一整批样点加一个时间戳。驱动排队、网卡抖动或消费者慢了，后果是\emph{整批}消失。检测手段只有两种：读时间戳的连续性（返回的 \Tp{HAS\_TIME} \Bk{} 标志与时间戳是否按预期推进），以及对比请求交付数与实际交付数。危险在于丢样在数据里看起来像``信号变了''：谱线仍在，只是时间轴上多了一个洞，任何依赖帧计数的算法（相干积累、时延估计、锁相）会静默错位。因此每个流都要有丢弃计数与告警，而不是在回调里试图补齐。

\textbf{时钟与相位一致性。} 多台设备要协同，共享 $10\ \text{MHz}$ 参考是下界，再加一个秒脉冲做时间对齐（\Fn{set\_time\_next\_pps}），必要时用 GPS 驯 oscillator 做保持（holdover）或用 PTP（IEEE 1588 家族）在以太网上分发时间与相位。剩下的残余是采样钟的相对偏差，量级在几十 ppm，工程处理是周期性重采样或一条恒定时间调整量（\Lib{SoapySDR} 的 \Fn{setTimeAdjust}）；分布式阵列里``共参考但时间戳各自打''是最省事的可用形态。

\textbf{增益与饱和。} 数字域的 $0\ \text{dBFS}$ 与射频域的 $0\ \text{dBm}$ 之间隔着一次阻抗与满刻度换算：$50\ \Omega$ 上 $0\ \text{dBm}$ 对应约 $0.224\ \text{V}$ 有效值，即 $632\ \text{mV}$ 峰峰值——把这条写进脚本，AGC 与衰减器的位置才有意义。AGC 让噪声底跟着信号跳，因此标定、频谱比对与相干处理都倾向固定增益加显式衰减，把动态范围交给过采样与位深而不是交给环路。

\begin{note}[title={短返回、参考基准与带宽口径}]
\begin{itemize}
\item \textbf{请求数与交付数的错配。} 一次 \Fn{readStream} 请求 $4096$、实际交付 $2048$ 是正常语义：短返回不是错误，把请求数当长度就读到未初始化数据；而``必须至少交付这么多''的假设会在设备降速时整块丢弃。抽取器与 FIR 还要处理第二个错配：本轮凑不满一个抽头窗口的那几个尾部样点必须\emph{留到下一轮}（靠 history 或显式的残料缓冲），直接丢弃会让输出抽样时刻漂移，表现为谱上出现边带。
\item \textbf{三个参考基准不可互换。} $\text{dBm}$ 以 $1\ \text{mW}$ 为参考，是功率单位，换算需要阻抗；$\text{dBFS}$ 以转换器满刻度为参考，是无量纲的电平口径，所有正常样本都是负值；$\text{dBc}$ 以载波为参考，是比值，杂散与相位噪声用它。把仪器读数（$\text{dBm}$）与文件的谱（$\text{dBFS}$）直接相减得到的``噪声系数''里混着增益、衰减与满刻度设定；把 $\text{dBc}$ 当 $\text{dBm}$ 用则会在链路预算里丢掉一整级。
\item \textbf{采样率与被测带宽的换算。} 复基带的可用带宽是实际采样率本身（正负频率都在用），不是它的一半——这与实信号的习惯相反。报告``带宽 $8\ \text{MHz}$''时先确认指的是 \Tp{CS16} 的两路实采样还是复采样，两种口径在同一份脚本里混用会凭空少掉一半带宽。
\end{itemize}
\end{note}

\textbf{为什么这一域仍是 \CL/\CPP{} 的主场。} 第一，每样点的预算算得出来：$61.44$ Msps 给每个复样点约 $16\ \text{ns}$，按 SIMD 一次处理八个样点也只有约 $130\ \text{ns}$，而一次解释器级的每样点操作就是几十纳秒量级、一次装箱分配更不用谈——所以``每样点都要跑的代码''只能在编译层。第二，驱动只给 \CL{} ABI：\Lib{libusb}、\Lib{libiio}、内核态的 \Tp{uio} 与 \Tp{vfio} 用户态映射、高速网卡的 verbs 与 DPDK 路径，形状全是 C；跨厂商抽象层要么把这些缝进一个 C 接口（这正是 \Lib{SoapySDR} 的插件 ABI 是 C 的原因），要么直接绑一家（\Lib{UHD}）。第三，硬件侧的工具链本身以 C 为汇合点：FPGA 与 SoC 的 HLS 路线把内核编译成 C 可调用组件、嵌入式 Linux 侧的 PetaLinux 一类的发行版给的是 C 库（这条线的细节不在本册范围，此处只需要认清：从板卡到主机之间\emph{没有}任何一层是脚本语言的接口）。第四，\Lib{GNU Radio} 这类框架的性能边界已经被社区反复测出：高采样率的 signal path 属于 \CPP{} block，\PY{} 属于控制流、低速率与图形装配——这不是风格偏好，而是上面三道预算共同压出来的结论。

\begin{longtable}{@{}>{\raggedright\arraybackslash}p{0.15\textwidth}
  >{\raggedright\arraybackslash}p{0.15\textwidth}
  >{\raggedright\arraybackslash}p{0.17\textwidth}
  >{\raggedright\arraybackslash}p{0.16\textwidth}
  >{\raggedright\arraybackslash}p{0.17\textwidth}@{}}
\caption{五个射频与软件无线电库：所在层次、语言层、许可证基调与实时能力}\label{tab:ch15-sdr}\\
\toprule
库 & 所在层次 & 语言层 & 许可证基调 & 实时能力 \\
\midrule
\endfirsthead
\toprule
库 & 所在层次 & 语言层 & 许可证基调 & 实时能力 \\
\midrule
\endhead
\Lib{SoapySDR} & 跨厂商设备抽象与流协商 &
  \CPP{} 核心、\CL{} 插件 ABI、\PY{} 绑定 &
  MIT 基调，以仓库 LICENSE 为准 &
  提供时间戳与缓冲接口，本身不含调度器，实时性取决于宿主线程 \\
\Lib{UHD} & 单一厂商的设备层与板级控制 &
  \CPP{} API，另有 \CL{} 与 \PY{} 绑定 &
  GPL 系 copyleft，以仓库 LICENSE 为准 &
  突发命令、板载时间戳与秒脉冲对齐，可稳定持续采流 \\
\Lib{GNU Radio} & 图、调度与 DSP block 集合 &
  \CPP{} 内核，\PY{} 与 GRC 装配图 &
  GPL 系 copyleft，以仓库 LICENSE 为准 &
  每 block 一线程的软实时，截止期不保证，靠缓冲量兜底 \\
\Lib{scikit-rf} & 频域网络运算与后处理 &
  \PY{} 门面，批量线性代数交给 numpy &
  宽许可（BSD 或 MIT 基调），以仓库为准 &
  无实时通路：离线标定、去嵌入与 TDR \\
\Lib{libmobi} & 前端与滤波器的解析建模 &
  \CL{} 库，\Tp{XML} 描述，可绑其他语言 &
  条款随版本，以仓库声明为准 &
  扫频或参数扫描，单点耗时短但非流式实时 \\
\end{longtable}

\begin{guideline}[title={本章要点}]
\begin{itemize}
\item 把``请求值不等于达成值''写进代码：采样率、频率、增益、缓冲时长都在设定之后读回实际值，后续所有换算（抽取比、字节率、时间戳差）都基于实际值。
\item 格式只在边界转换一次：总线用窄格式（\Tp{CS16} 一类），内核用 \Tp{CF32} 与分面或交错的固定一种，交织与分面的置换只出现在一个函数里；解交织顺手做缩放，避免两次全数组遍历。
\item 计数一切异常：短返回、负返回、时间戳跳变、丢弃块数都进同一个计数器并周期上报。丢样不是崩溃而是静默的数据缺陷，只有计数能让它变成可诊断事件。
\item 分层按速率切：每样点的进 \CPP{} block 或 \CL{} 内核，每批的进调度器与环形缓冲，每小时一次的进 \PY{} 脚本与 GRC 图；跨层调用走 message 或属性，不共享可变状态。
\item 建模沿 ``$S$ 参数'' 这条缝走：全波工具导出端口参数，\Lib{scikit-rf} 做级联与去嵌入，\Lib{libmobi} 或 SPICE 做系统级，最后与实测的噪声系数和 SFDR 对表；任何一环没写参考阻抗，链路预算就没有意义。
\end{itemize}
\end{guideline}
% ====================================================================

% ====================================================================
\part{封装、门面与实战}

\chapter{Python 科学栈其实是 C/C++ 的门面}

前文各章回答的是``内核怎么写''，从本章起回答``内核怎么被用''。本章的论断很硬：\PY{} 科学计算栈不是一个用 \PY{} 写成的科学计算系统，而是\emph{一批 \CL/\CPP{} 内核之上的一层门面}——用户看到的是 \Tp{ndarray}、关键字参数与异常，跑掉真正时间的是 \Tp{gemm}、\Tp{fftw} 与 \Tp{.pyx} 编译出来的本地码。这不是修辞，是可以逐包核对的事实，而且核对成本极低：列出每个包的顶层模块里\emph{没有} \Tp{for} 循环的那些调用，再问一句``它落到哪个符号''。

把这层关系说清楚，是为了后面两件事：一是判断一个现成门面能不能用于长时间运行的流程（判据是它能不能说出自己底下的库名与符号名），二是决定自己发布时要不要写门面（本章末尾给出反向判据）。本章只谈选型与生态映射；\Tp{pybind11} 的语法细节、GIL 的作用域规则、\Tp{py::array\_t} 的 flag、类型桩文件与 \Tp{cibuildwheel} 矩阵都在并行算法分册里讲过，此处一律引用而不重复。

\section{谁在底下：一张生态映射表}

表~\ref{tab:ch16-under} 是本册涉及到的门面层与其内核的对应关系。``交汇点''一列填的是\emph{接口形态}，不是语言名：平面 C 函数、句柄加返回码、编译式绑定、运行期 JIT，就这四种，正是讲 C ABI 的那一章给出的三类接口形态加一条例外。

\begin{longtable}{@{}>{\raggedright\arraybackslash}p{0.16\textwidth}
  >{\raggedright\arraybackslash}p{0.22\textwidth}
  >{\raggedright\arraybackslash}p{0.25\textwidth}
  >{\raggedright\arraybackslash}p{0.17\textwidth}@{}}
\caption{\PY{} 门面层与它底下的实现}\label{tab:ch16-under} \\
\toprule
包 & 门面自己做的事 & 热路径实际落点 & 交汇点形态 \\
\midrule
\endfirsthead
\toprule
包 & 门面自己做的事 & 热路径实际落点 & 交汇点形态 \\
\midrule
\endhead
\midrule
\multicolumn{4}{@{}r}{（接下页）}\\
\endfoot
\bottomrule
\endlastfoot
\Lib{numpy} & \Tp{ndarray} 的形状与 dtype 调度、\Fn{dot} 分派 &
  逐元素核在 \CL{} 循环里；\Fn{dot}、\Fn{matmul} 交给 BLAS 的
  \Tp{dgemm}、\Tp{sgemm}；\Fn{linalg.svd} 交给 LAPACK 的
  \Tp{gesdd} 一族 &
  平面函数加缓冲协议 \\
\Lib{scipy} & 参数解析（\Fn{get\_window}）、后端选择 &
  \Fn{scipy.linalg} 走 LAPACKE；\Fn{signal}、\Fn{interpolate}
  的原生模块是 \CL{}；\Fn{sparse.linalg\Bk.eigsh} 走 \Lib{ARPACK}
  的 \Tp{dsaupd}；近年版本的 \Fn{scipy.fft} 默认后端是自编入的
  pocketfft（\CPP{}） &
  句柄（工作区）加平面函数 \\
\Lib{scikit-learn} & 估计器协议、交叉验证与并行编排 &
  树、距离与损失函数是 \Tp{.pyx} 编译出的本地码；矩阵运算交给
  BLAS；并行交给 OpenMP；线程数靠 \Fn{threadpoolctl} 统管 &
  编译式绑定加 BLAS \\
\Lib{librosa} & \Fn{stft}、\Fn{feature.mfcc} 的帧与单位约定 &
  数值靠 numpy 与 scipy；起波与分离一类是 \Tp{numba} JIT
  （\Tp{guvectorize}）；\Fn{load} 委托 \Lib{soundfile} &
  JIT 本地码加文件层 \CL \\
\Lib{soundfile} & \Tp{SF\_INFO} 映射、dtype 与缩放 &
  容器与编码全部在 \Lib{libsndfile}；门面不含 DSP &
  句柄加返回码 \\
\Lib{matplotlib} & 布局、坐标轴与数据到 RGBA 的映射 &
  光栅化在 \CPP{} 后端 \Tp{Agg}；凸包与三角化走
  \Lib{Qhull}（\CL{}） &
  内部 \CPP{}，不对外 \\
\Lib{PyTorch} & autograd 图、设备分派与算子名 &
  \Tp{ATen}（\CPP{}）逐元素核；\Fn{torch.mm} 交给 BLAS 后端
  （\Lib{MKL} 或 \Lib{oneDNN}）、设备侧交给 \Lib{cuBLAS} &
  \CPP{} 内部加 \CL 数值 ABI \\
\Lib{astropy} & 坐标框架、表格与 \Tp{units} &
  时间尺度与岁差等标准算法在 \Lib{erfa}（\CL{} 移植）；球面坐标
  变换在 \Lib{wcslib}（\CL{}）；其余多为纯 \PY{} &
  \Tp{ctypes} 或自写扩展 \\
\Lib{xarray} & 标签、维度对齐与选择 &
  没有自己的数值内核：计算全部委派给 numpy 与 dask &
  无（例外） \\
\end{longtable}

最后一行是本章唯一需要展开的反例。\Lib{xarray} 证明的不是``\PY{} 也能做科学计算''，而是\emph{编排层可以独立存在，前提是下面已经有一批别人写好的内核}：它自己不做一次浮点循环，代价是把 numpy 的每条约定（形状、dtype、连续性、视图还是拷贝）都升级成了带标签的版本。同理，\Lib{astropy} 的 \Tp{units} 与 \Tp{table} 是纯 \PY{}，而一旦涉及每样本都要算的球面变换与时间尺度换算，它立刻下沉到 \CL{} 的 \Lib{erfa}——一个包内部``哪部分留在 \PY{}、哪部分下沉''的分工，比任何宣传语都能说明这一域的成本模型。

为什么汇合点一定是 \CL{} 而不是别的语言，本章不重复展开，只引用讲接口层那一章的结论并把它落到表上：\CL{} 的调用约定与结构体布局是平台 ABI 的子集，且边界上没有第二套运行时要谈条件。\CL{} 的表达力不是理由，模板头文件库（\Lib{Eigen}）表达力最强，恰恰最没法跨语言；\Lib{numpy} 与 \Lib{scipy} 愿意在 wheel 里带一份 BLAS，也是因为那一层符号是稳定的。

\section{门面层的四个职责，与它不该承担的第五件事}

把表~\ref{tab:ch16-under} 里``门面自己做的事''那一列归并，得到四件事。它们的共同特征是\emph{成本与数据规模无关}：做一次和做一百万次的单价一样。

\textbf{职责一：参数校验与形状推导。}\Fn{sklearn.utils.check\_array} 与 \Fn{check\_X\_y} 是这一类最标准的实现：它把 ``array-like'' 变成 ``二维、\Tp{float64}、C 连续、有限值''，并给出可诊断的异常；\Fn{scipy.signal.get\_window} 把 \Tp{hann} 这个字符串解析成实际系数；\Fn{np.broadcast\_arrays} 在调用之前先把输出形状算出来。机制价值是错误出现在\emph{调用点}而不是内核深处；代价是一次微秒级的检查。真正要当心的是它顺手做的转换：\Tp{dtype=float64} 遇上 \Tp{float32} 输入就是一次全量拷贝，形状``看着对''却非连续的输入也一样，这两笔钱远比内核本身贵。

\textbf{职责二：类型与单位换算。}包括 \Fn{np.result\_type} 的提升规则、整型样本到 $[-1,1]$ 的缩放（讲音频的那一章算过 \Tp{/32768.0} 与 \Tp{/32767.0} 的不对称）、采样率与时间基的换算、以及 \Tp{units} 的量纲折算。换算放在门面的\emph{入口与出口}是正确的，放在循环里就是第五件事。

\textbf{职责三：生命周期与错误翻译。}\Lib{LAPACK} 的 \Tp{info} 出参到了门面手里要变成异常：\Tp{info} 为负表示第几个参数非法，\Lib{numpy} 据此抛 \Tp{ValueError}；为正表示迭代不收敛，据此抛 \Tp{LinAlgError}。\Lib{soundfile} 把 \Fn{sf\_error\_str} 的字符串包成 \Tp{SoundFileError}；\Tp{fftw\_plan} 与 \Tp{SNDFILE*} 这类句柄要绑一个析构回调，否则``谁创建谁销毁''在 \PY{} 的作用域里根本无从谈起。这一职责只在边界做一次，不该在每层重复。

\textbf{职责四：绘图与文件路径。}\Lib{matplotlib} 只消费数组，路径、编码与元数据归门面：\Tp{pathlib} 与 \Fn{os.fspath} 负责把平台的宽字符路径换成 \CL{} 侧能接受的字节串，Windows 上这一步做错的典型症状是``中文目录打不开文件''；谱图的横轴单位、CSV 表头、sidecar 元数据也都在这一层。这一层没有性能问题，但它是用户看到的\emph{全部}契约。

\textbf{第五件事：逐元素计算——门面不该承担。}判据回到第~\ref{lst:ch16-three}清单的每元素成本。反面样本很清楚：把 \Tp{Quantity} 对象乘进热循环、用 \Fn{df.apply} 逐行算、或者把一次融合过的表达式拆成十条 \Fn{np.where} 与 \Fn{np.clip}（每条都产生一个全尺寸临时数组，内存带宽付十遍，而一个 C 内核一次遍历就够）。门面的四个职责里没有一项要求你自己写循环。

\begin{warning}[title={门面的四种越权写法}]
\begin{itemize}
\item \textbf{顺手拷贝}：校验里写了 \Tp{dtype=\allowbreak float64}，用户传 \Tp{float32}，每次调用多付一份全量拷贝，症状是``结果对但比预期慢五倍''。
\item \textbf{把单位换算写进内层}：每样本一次量纲检查或多标量乘法，热路径成本乘以采样数。
\item \textbf{重复校验}：门面校一遍、绑定层再校一遍、内核里再校一遍；同一笔开销付三次，而真正需要的是一次边界校验加一次断言。
\item \textbf{把句柄藏进闭包}：\Tp{SNDFILE*} 或 \Tp{fftw\_plan} 被 \PY{} 闭包捕获，销毁顺序由引用计数决定，于是退出阶段出现``plan 先于缓冲被释放''。生命周期要显式写成 \Tp{with} 或 finalizer。
\end{itemize}
\end{warning}

\section{零拷贝的三种协议与 C 侧的读法}

\PY{} 与内核之间每一次数据交接，都是下面三种协议之一。它们的差别不在``能不能零拷贝''，而在\emph{谁需要链接什么}。

\textbf{协议一：\Tp{\_\_array\_interface\Bk\_\_}}。它是一个 \PY{} 字典（version 为 3），关键字段是 \Tp{shape}、\Tp{strides}、\Tp{itemsize}、\Tp{typestr}（如 \Tp{<f8}，三个字符依次给出字节序、类型与位宽）与 \Tp{data}——后者是一个二元组，第一项是\emph{整数形式的设备无关主机地址}，第二项是只读标志。它最大的优点是 C 侧什么都不用链接：\PY{} 侧取整数、\Tp{ctypes} 转成指针、把整数交给自己 \Tp{.so} 里的裸函数。代价是每个数组要构造一次字典，且只描述主机内存（设备张量要走 DLPack，见并行算法分册）。

\textbf{协议二：缓冲协议（buffer protocol）}。\PY{} 侧叫 \Fn{memoryview}，C 侧是 \Fn{PyObject\_GetBuffer} 填一个 \Tp{Py\_buffer\_view}，字段与协议一几乎一一对应，另多一个 \Tp{format} 串。\Lib{numba} 的数组参数、\Tp{cffi} 的 \Fn{ffi.from\_buffer}、\Tp{hashlib} 与 socket 的写入接口都读它。它的约束有两条：只有实现者暴露了相应回调的对象才可取（\Tp{str} 可以，非连续的 \Tp{ndarray} 视图在多维步长下常不可用）；以及\emph{必须} \Fn{PyBuffer\_Release}，且取到的 \Tp{buf} 指针有效期只到释放为止。它要求链接解释器，因此只适合 Python 专用扩展。

\textbf{协议三：strides、dtype、ptr 三件套摊平成参数}。签名长这样：\Tp{(const char* ptr, int dtype, int ndim, const int64\_t* shape, const int64\_t* strides)}。这是把前两种协议的\emph{结论}写进 \CL{} 平面函数：不需要字典、不需要解释器、可以在纯 C ABI 上跨语言，且能同时接受主机与设备指针。DLPack 交付的就是这个形状，\Lib{Eigen} 的 \Tp{Map} 是它在 \CPP{} 侧的写法，本册实战章的 \Tp{bind/} 层用它。\PY{} 侧一行就够：\Tp{x.ctypes.data}、\Tp{x.dtype}、\Tp{x.strides}。

清单~\ref{lst:ch16-buffer} 是 C 侧两种读法的最小可运行对照，注意三条纪律：\Tp{strides} 的单位是\emph{字节}而不是元素；任何提前出口都要释放缓冲；内核\emph{不要}顺手 \Fn{asarray}——把连续性检查写成拒绝，让调用方决定拷贝。

\begin{lstlisting}[style=cstyle,caption={C 侧读 numpy 数组：缓冲协议与摊平的三件套},label={lst:ch16-buffer}]
// 路线 A：缓冲协议。要链接解释器，只能在 Python 进程里用。
#include <Python.h>

int energy_buf(PyObject* obj, double* out) {
  Py_buffer v;
  if (PyObject_GetBuffer(obj, &v, PyBUF_ND | PyBUF_STRIDES) < 0)
    return -1;                 /* 原因已作为 Python 异常挂起 */
  if (v.ndim != 1 || v.itemsize != (Py_ssize_t)sizeof(float)) {
    PyBuffer_Release(&v);      /* 每个出口都必须释放 */
    return -2;
  }
  double acc = 0.0;
  for (Py_ssize_t i = 0; i < v.shape[0]; ++i) {
    const char* q = (const char*)v.buf + i * v.strides[0];
    const float* p = (const float*)q;          /* 按字节算偏移 */
    acc += (double)(*p) * (double)(*p);
  }
  PyBuffer_Release(&v);        /* 此前 v.buf 才有效 */
  *out = acc;
  return 0;
}

// 路线 B：三件套摊平。不碰解释器，可单测、可换语言。
//   int spec_block(const char* ptr, int64_t n, int64_t stride,
//                  float* out, int64_t m);
// Python 侧传 x.ctypes.data, x.shape[0], x.strides[0]。
// 代价：解释器不再替你保证这块内存活着。调用期间必须有人在
// 拿着 x 的引用；把它存进 C 结构体跨过本次调用就是悬垂指针。
\end{lstlisting}

\begin{note}[title={零拷贝的四个失效瞬间}]
\begin{itemize}
\item \Tp{strides[0]} 为 $0$：那是广播视图。读没问题，写会一次改掉一整行，``输出正确率诡异''的常见根源。
\item 只读缓冲上当场索取可写视图：\Tp{info.flags.writeable} 为假、或内存来自 \Tp{mmap} 的只读映射。
\item 交出指针后又把 \PY{} 变量重绑定：原数组被回收，C 侧手里是一个语法正确的整数。
\item dtype 不符：\Tp{float32} 交给按 \Tp{double} 解释的步长，结果不是崩溃而是``平滑的错值''。
\end{itemize}
\end{note}

\section{同一个算法的三种写法，与它们的成本}

把三种写法放在同一个算法上（单边功率谱）最容易看清分工：写法一调现成库（\Lib{scipy.signal}），写法二用 numpy 向量化自己拼，写法三用 \Tp{ctypes} 调自家 \Tp{.so}。它们的关系不是``谁更快''，而是``谁在什么规模上付什么钱''。

\begin{lstlisting}[style=pythonstyle,caption={单边谱的三种写法：现成库、向量化、ctypes 调自家 .so},label={lst:ch16-three}]
# 输入 x：float32、长度 16384 的单通道块，fs = 48000
import ctypes
import numpy as np
import scipy.signal as ss

# 写法一：调现成库。分段、加窗、平均、归一化都在 C 内核里。
f, pxx = ss.welch(x, fs, nperseg=1024, noverlap=512,
                  window="hann", detrend=False)

# 写法二：numpy 向量化。用 sliding_window_view 避免拷贝，
# 但 seg * w 与 **2 各产生一个全尺寸临时数组。
w = np.hanning(1024).astype(np.float32)
seg = np.lib.stride_tricks.sliding_window_view(x, 1024)[::512]
psd = (np.fft.rfft(seg * w, axis=1) ** 2).mean(axis=0)
hz = np.fft.rfftfreq(1024, 1.0 / fs)
# 与写法一对不上的话先查口径：welch 给的是密度（除 fs 与窗
# 功率），此处是幅度平方；detrend 与 noverlap 也要一致。

# 写法三：ctypes 调自家 .so。只交指针与步长，零拷贝。
lib = ctypes.CDLL("./libspeck.so")        # 无签名校验：抄错即崩
lib.spec_block.argtypes = [ctypes.c_void_p, ctypes.c_int,
                          ctypes.c_long, ctypes.c_void_p]
lib.spec_block.restype = ctypes.c_int     # 0 表示成功
rc = lib.spec_block(x.ctypes.data, x.shape[0],
                    x.strides[0], out.ctypes.data)
if rc != 0:                               # 返回码翻译成异常
    raise RuntimeError("native call failed, rc=%d" % rc)
\end{lstlisting}

三条路的成本结构：写法一付一次 ``参数解析加后端分派''，内核归平均摊薄到每帧，是绝大多数生产代码的正确起点；写法二多付若干个临时数组的带宽，换来的是可组合性（要一个库没有的口径时就自己拼）；写法三只付一次边界穿越，但它把校验、错误与内存生命周期都推回给你自己，适合``内核已经定形、每帧预算以微秒计''的场景。三者之间没有优劣，只有分工。

\begin{warning}[title={七个常见性能陷阱，按出现频率排序}]
\begin{itemize}
\item \textbf{\PY{} 层循环}：每元素一次解释器往返、一次装箱、一次引用计数。判据是每秒元素数，不是行数。
\item \textbf{GIL 让并行库退化}：绑定层调用前没放锁，\Lib{OpenBLAS} 起的线程全卡在解释器上；表现是``CPU 只用一个核''。放锁规则见并行算法分册。
\item \textbf{隐式 \Fn{asarray} 拷贝}：dtype 或连续性不符时门面静默拷贝；单次不明显，每帧一次就是十倍的内存带宽。
\item \textbf{\Tp{float32} 与 \Tp{float64} 走不同 BLAS 路径}：\Tp{sgemm} 与 \Tp{dgemm} 是两套内核，分块尺寸与寄存器占用不同，单精度\emph{未必}快一倍；反过来 \Tp{float32} 交给只认双精度的入口，先付一次升精度拷贝。
\item \textbf{小数组上调用开销压过计算}：$n$ 为几十时，\Fn{np.dot} 的分派、校验与后端查找占大头；要合并调用（batch 一次进）而不是每帧一次进。
\item \textbf{线程超订}：外层线程池乘以内层 OpenMP 乘再加一份 BLAS 线程池。诊断顺序是先数线程，再怀疑算法。
\item \textbf{把单位换算写进内层循环}：每样本一次量纲检查、一次 ``先除后乘'' 的 dB 折算、逐 bin 的 \Tp{rfftfreq} 重算。这些都在循环外做一次。
\end{itemize}
\end{warning}

\section{分发、绑定选型与``不写门面''的判据}

发布这一域的门面，真正的约束不在绑定语法而在\emph{链接}。

\textbf{wheel 里 vendored 的 BLAS。}官方 \Lib{numpy} wheel 通常把一份 \Lib{OpenBLAS} 装进包目录，\Lib{scipy} 的 wheel 也带自己的数值后端，各包之间互不知情。后果有三条：同一进程可能出现两份各自初始化线程池的 BLAS（于是超订成倍恶化）、内存按份数翻倍、包体积明显变大（一份后端可达几十 \Tp{MB} 量级，具体数字随发行而变）。核对手段是实测而不是回忆：\Fn{numpy.show\_config}（老写法读 \Tp{np.\_\_config\_\_}）报告编译进去的后端，\Fn{threadpoolctl} 能列出运行时\emph{实际}加载了哪几份并行库，Linux 与 macOS 用 \Tp{ldd} 与 \Tp{otool -L} 看真实依赖，Windows 用 \Tp{dumpbin /dependents}。conda 与发行版打包走的是另一条路（一份共享库大家用），这也是同一段代码在 venv 里 ``线程失控'' 而在 conda 里正常的常见原因。

\textbf{与 numpy 的 ABI 耦合。}扩展模块是按编译期的 \PY{} 次版本与 numpy 头编译的：\Tp{cp312} 的模块在 \Tp{cp313} 上直接 \Tp{ImportError}；numpy 2.x 把 C-API 分层之后，按 2.x 头编译的轮子可以跑在 2.x 区间，但\emph{不}回头兼容 1.x，早先 \Tp{oldest-supported\Bk-numpy} 的做法正是为这一段历史包袱服务的。落到操作层面只有两条：构建依赖里显式钉住 numpy 策略，发布前在目标区间两端各跑一次 \Tp{import} 冒烟测试（细节见并行算法分册的构建分发一章）。

\textbf{绑定路线的适配度。}表~\ref{tab:ch16-bind} 按本册这几类库的接口形态打分，判断依据只有一条：句柄式平面函数谁都能包，模板头文件库只能走编译式绑定。

\begin{table}[H]
\centering
\caption{五条绑定路线在本册几类库上的适配度}
\label{tab:ch16-bind}
\small
\begin{tabular}{@{}>{\raggedright\arraybackslash}p{0.11\textwidth}
  >{\raggedright\arraybackslash}p{0.10\textwidth}
  >{\raggedright\arraybackslash}p{0.16\textwidth}
  >{\raggedright\arraybackslash}p{0.15\textwidth}
  >{\raggedright\arraybackslash}p{0.14\textwidth}
  >{\raggedright\arraybackslash}p{0.16\textwidth}@{}}
\toprule
路线 & 需要编译器 & 句柄式 \CL{} 接口 & 模板头文件库 & 实时回调进 \PY{} & 本册里的位置 \\
\midrule
\Tp{ctypes} & 否 & 好，签名手抄 & 不行 & 不该（取锁无界） &
  \Lib{soundfile} \\
\Tp{cffi} & ABI 否、API 是 & 好，API 模式能校验布局 & 不行 & 不该 &
  系统库门面一类 \\
\Tp{Cython} & 是 & 好 & 勉强，要手写声明 & 不可预算 &
  \Lib{scikit-learn}、\Lib{h5py} \\
\Tp{pybind11} & 是，且要 \CPP{} 头 & 好，先包一层 & 天然支持 & 不可预算 &
  \Lib{Open3D}、\Lib{CGAL} \\
\Tp{numba} & 否（运行期 JIT） & 勉强，指针参数很脆 & 无 & \Tp{cfunc} 裸调仍超支 &
  \Lib{librosa} 的局部核 \\
自写 \CL{} 扩展 & 是 & 天然 & 无 & 无（本册实战走这条） &
  \Tp{bind/} 层 \\
\bottomrule
\end{tabular}
\end{table}

\textbf{什么时候干脆不写门面。}四种情形：探索期（形状、单位与边界条件还在变，两份实现会互相拖住）；只有你自己调用（直接在 \CPP{} 单测或命令行工具里调内核，验证快得多）；内核还没定形（门面的四个职责会随签名天天改）；实时域（回调绝不能进解释器，此时``门面''只能是块与块之间的统计与配置）。反向判据同样清楚：一旦消费方不是写内核的人，或者同一内核要交给三种语言，门面的成本就摊薄到可忽略——因为它省下的是每个使用者各自抄一遍签名的钱。

\begin{guideline}[title={全书主线，收束成一则判据}]
\begin{itemize}
\item 内核用 \CL/\CPP{} 不是因为``它们快''，而是因为它们是唯一能把\emph{数据布局、所有权与线程语义写死、并且跨语言核对}的层。快只是副产品，可核对才是制度性优势。
\item 门面的四个职责（校验与形状、类型与单位、生命周期与错误翻译、绘图与路径）成本与规模无关；第五件事（逐元素计算）成本与规模成正比，永远不属于门面。
\item 数据交接只有三种形状：\Tp{\_\_array\_interface\Bk\_\_}、缓冲协议、三件套摊平。跨语言发布选第三种，Python 专用扩展选第二种，零链接调试选第一种。
\item 分发要实测不要回忆：谁带了哪份 BLAS、带了几份、线程池开了几个，\Fn{threadpoolctl} 与 \Fn{ldd} 一次就能看清。
\item ``要不要写门面''的答案由消费方决定，不由语言偏好决定。下一章把这一条落到一个具体工程上。
\end{itemize}
\end{guideline}
% ====================================================================

% ====================================================================
\chapter{实战：一条采集、滤波、谱分析与导出的处理链}

这一章把前面各章的组件装成一个可以构建、可以对拍、可以分发的工程。工程名 \Tp{spechain}，链路只有一条：设备回调取到音频块 $\to$ 环形缓冲 $\to$ 双二阶级联滤波 $\to$ 加窗 FFT 出单边谱 $\to$ 写回 WAV 与 CSV，\PY{} 侧负责设计系数、驱动这条链并与 \Lib{scipy} 对拍。它的全部价值不在算法新颖，而在把\emph{职责边界}演示一遍：哪一层该写 \CL、哪一层该写 \CPP、哪一层只允许是门面的四个职责，以及``同一功能在两侧各实现一遍''怎样才能变成一条可执行的测试而不是两份互相漂移的代码。

先给三条贯穿全章的纪律。其一，实时线程与离线线程分工，块之内在 \CL、块之间在 \PY{}；其二，滤波器\emph{设计}放在 \PY{} 侧的 \Fn{scipy.signal} 里，C 侧只读系数；其三，任何跨边界交付的数据都走三件套（指针、步长、类型），任何跨版本交付的判据都写成``两层''（逐位一致只在同一二进制内要求，跨实现用容差）。

\section{分层与目录}

目录结构就是分层结构，一个目录对应一种语言层与一组禁令：

\begin{lstlisting}[style=shellstyle,caption={工程目录：目录边界即语言与依赖边界},label={lst:ch17-tree}]
spechain/
  include/spechain.h     唯一的对外头文件：句柄、配置、错误码
  core/                  C 内核：环形缓冲、biquad、能量归约
  io/                    C：包 libsndfile，容器、编码、交织置换
  dsp/                   C++：FFTW plan、窗与重叠、分段谱
  bind/                  C 平面函数门面（extern "C"），给 ctypes
  py/                    Python：校验、系数设计、编排、绘图
  tests/                 C 单测加 pytest 对拍加 golden
  CMakeLists.txt
\end{lstlisting}

表~\ref{tab:ch17-layers} 把每层的职责与\emph{禁令}写在一起——禁令比职责重要，因为跨层污染是这一域唯一会让工程变慢的东西：一旦 \Tp{core/} 能看见 \Lib{libsndfile}，内核就没法脱离文件系统单测；一旦 \Tp{py/} 里出现逐样本循环，对拍就变成了两份实现各自出错。

\begin{longtable}{@{}>{\raggedright\arraybackslash}p{0.09\textwidth}
  >{\raggedright\arraybackslash}p{0.09\textwidth}
  >{\raggedright\arraybackslash}p{0.32\textwidth}
  >{\raggedright\arraybackslash}p{0.32\textwidth}@{}}
\caption{分层职责与禁令}\label{tab:ch17-layers} \\
\toprule
目录 & 语言 & 职责 & 禁止事项 \\
\midrule
\endfirsthead
\toprule
目录 & 语言 & 职责 & 禁止事项 \\
\midrule
\endhead
\midrule
\multicolumn{4}{@{}r}{（接下页）}\\
\endfoot
\bottomrule
\endlastfoot
\Tp{core/} & \CL &
  环形缓冲、biquad 状态更新、能量归约、错误码 &
  禁止依赖 \Lib{libsndfile} 与 BLAS；禁止堆分配与阻塞 IO；禁止全局可变状态 \\
\Tp{io/} & \CL &
  打开与关闭、容器与子格式、交织与平面置换、短读处理 &
  禁止做 DSP；禁止替调用方猜采样率与通道数（一律回读 \Tp{SF\_INFO}） \\
\Tp{dsp/} & \CPP &
  plan 生命周期、窗与重叠、分段谱与 dB 折算 &
  禁止 \Tp{std::} 类型进 \Tp{bind/}；禁止在 \Fn{process} 里建 plan；禁止 \Tp{-ffast-math} \\
\Tp{bind/} & \CL 门面 &
  句柄配对、参数一次校验、错误码与消息 &
  禁止模板、禁止异常、禁止把 \PY{} 指针存过本次调用 \\
\Tp{py/} & \PY &
  形状校验、单位与 dtype 换算、系数设计与序列化、绘图与路径 &
  禁止逐元素循环；禁止热路径里的单位换算；禁止重做滤波器设计 \\
\Tp{tests/} & \PY 加 C &
  两侧对拍、能量与阶跃判据、golden 比对、环境基线 &
  禁止只比平均值；禁止把库版本与环境变量留在基线之外 \\
\end{longtable}

\section{C 句柄接口的最终形态与两条线程}

句柄式接口在上一章已经露过一次，这里是它定形之后的样子（清单~\ref{lst:ch17-handle}）。四条约定全部\emph{写在头文件注释里}，因为它们都不是编译器能检查的东西：生命周期成对、错误走返回码、无全局状态、线程约束。

\begin{lstlisting}[style=cstyle,caption={对外的 C 句柄接口：成对生命周期、返回码、无全局状态},label={lst:ch17-handle}]
/* include/spechain.h —— 对外只有这一层：无模板、无异常。
 * 单位约定：样本为 [-1,1] 浮点，采样率与频率列都是 Hz。 */
typedef struct ScCtx ScCtx;               /* 不透明：实现藏在里面 */

typedef enum { SC_OK = 0, SC_E_ARG = -1, SC_E_STATE = -2,
               SC_E_IO = -3, SC_E_FULL = -4,
               SC_E_NOMEM = -5 } ScErr;

typedef struct {                          /* 纯数据，调用方可自建 */
  int samplerate;                         /* 必须与系数文件一致 */
  int nsec;                               /* 二阶节个数，1..8 */
  float gain;                             /* 整体增益，系数已归一 */
  float sos[8][6];                        /* 每行 b0 b1 b2 1 a1 a2 */
} ScConfig;

/* 线程约束：一个句柄同一时刻只能被一个线程调用。create 与
 * destroy 可以在别的线程上做，但必须配对，且不与其它调用重叠。
 * 库内部不起线程；并行由调用方按通道或按块分片。 */
ScCtx* sc_create(const ScConfig* cfg, char* err, int errlen);
int sc_process(ScCtx* h, const float* in, float* out, int frames);
int sc_spectrum(ScCtx* h, float* magdb, int* nbins);
int sc_drain(ScCtx* h);                   /* 冲掉重叠缓冲的尾块 */
void sc_destroy(ScCtx* h);                /* 与 create 严格成对 */
const char* sc_strerror(int code);        /* 静态串：不要 free */
\end{lstlisting}

\textbf{实时线程只做两件事}：把设备给的块拷进环、更新两个索引。其余一切（滤波、谱、写文件、换算 dB、统计）都在离线工作线程上。划分依据不是``工作量''而是\emph{截止期}：回调的截止期由声卡时钟决定，块周期减去安全检查就是全部预算。清单~\ref{lst:ch17-ring} 是这一层的完整实现要点：两个原子索引、无锁、容量取 $2$ 的幂以便用位与取模、容量对齐到缓存行以避开伪共享。

\begin{lstlisting}[style=cstyle,caption={SPSC 环形缓冲：生产者只写 head，消费者只写 tail},label={lst:ch17-ring}]
/* 容量是 2 的幂：取模换成位与，回卷分两段 memcpy。 */
#define SC_CAP 8192                  /* 48k 下单通道约 170 ms */

typedef struct {
  alignas(64) _Atomic size_t head;  /* 各自独占一条缓存行 */
  alignas(64) _Atomic size_t tail;
  float buf[SC_CAP];
} ScRing;

/* 采集回调里调用：不分配、不阻塞；满则丢并计数。 */
int sc_ring_write(ScRing* r, const float* x, size_t n) {
  size_t h = atomic_load_explicit(&r->head, memory_order_relaxed);
  size_t t = atomic_load_explicit(&r->tail, memory_order_acquire);
  if (SC_CAP - ((h - t) & (SC_CAP - 1)) < n)
    return SC_E_FULL;                /* 丢块，由计数上报为 XRUN */
  size_t w = h & (SC_CAP - 1);
  size_t first = SC_CAP - w;         /* 到环尾还剩多少 */
  if (first > n) first = n;
  memcpy(r->buf + w, x, first * sizeof(float));
  memcpy(r->buf, x + first, (n - first) * sizeof(float));
  atomic_store_explicit(&r->head, h + n, memory_order_release);
  return SC_OK;
}
\end{lstlisting}

\begin{warning}[title={环形缓冲的四笔账}]
\begin{itemize}
\item \textbf{延迟预算}：$48\,\text{kHz}$ 下 $256$ 帧一块是 $5.33$\,ms，加工作线程的一整块等待就是 $10.7$\,ms。要更低就把块减小，代价是每块的固定开销摊不薄。
\item \textbf{容量预算}：环的作用是把抖动与突发吸收掉，容量小于 ``磁盘写入最坏耗时'' 等于没装环。
\item \textbf{丢块语义}：环满时\emph{必须}丢而不是等，否则生产者的等待时间进入回调，XRUN 立刻从环转移到设备上。
\item \textbf{索引宽度}：\Tp{size\_t} 单调递增而不是回绕成环内下标，``满还是空''靠差值判断；一旦两边都回绕，SPSC 就分不清零与满。
\end{itemize}
\end{warning}

\section{内核：双二阶级联与谱分析}

滤波内核用转置直接 II 型（TDF-II），与 \Lib{SciPy} 的 \Fn{sosfilt} 同一形态，这不是巧合：\emph{同一形态才谈得上逐位对拍}。每级两个状态、每样本一次输出读取，级与级串联（清单~\ref{lst:ch17-biquad}）。

\begin{lstlisting}[style=cstyle,caption={core/biquad.c：TDF-II 级联的单声道状态更新},label={lst:ch17-biquad}]
/* 系数来自 JSON（由 scipy.signal.butter 设计），本文件不设计。
 * 每行 sos 是 b0 b1 b2 1 a1 a2，即 a0 已归一，a1 a2 取负号约定
 * 与 scipy 一致：y = b0*x + z0, z0 = b1*x - a1*y + z1。 */
void chain_run(Chain* c, const float* x, float* y, size_t n) {
  const int k = c->nsec;
  for (size_t i = 0; i < n; ++i) {
    float v = c->gain * x[i];          /* 增益在首级之前，一次乘 */
    for (int s = 0; s < k; ++s) {
      const Biquad* b = &c->sec[s];
      float* z = c->z[s];              /* 每级两个状态 */
      float o = b->b0 * v + z[0];
      z[0] = b->b1 * v - b->a1 * o + z[1];
      z[1] = b->b2 * v - b->a2 * o;
      v = o;                           /* 上一级的输出即下一级输入 */
    }
    y[i] = v;                          /* 单态内层：无分支无判空 */
  }
}
/* 立体声：两个 Chain 实例，状态彼此独立；不要共享 z 数组。
 * 顺序是硬约束：级联顺序换掉，极点先衰减还是先放大就变了，
 * 中间量的动态范围与最终舍入都不同，逐位一致性立刻失效。 */
\end{lstlisting}

三个细节值得单独点出。第一，\textbf{状态归谁}：\Tp{z} 存在句柄里而不是栈上，因为滤波是有记忆的；跨块调用之间 \Tp{Chain} 必须活着，这是 ``无全局状态'' 与 ``有状态实例'' 的分界。第二，\textbf{溢出防护}：把整体 \Tp{gain} 提到首级之前而不是摊进每级的 \Tp{b}，可以让中间量的峰值可查；定点化只在写文件那一步做（见讲音频格式那一章的饱和与抖动）。第三，\textbf{向量化}：TDF-II 每样本依赖上一样本，级内不可向量化，可向量化的维度是\emph{通道}与\emph{级}——所以多通道要按通道分片，而不是指望编译器把 \Tp{for s} 拉宽。

谱分析一侧有三条纪律，都体现在清单~\ref{lst:ch17-fft} 里。缓冲要交给库自己的分配器，因为只有 \Fn{fftwf\_malloc} 给出 SIMD 对齐；plan 只在 \Fn{create} 里建一次；单边输出的长度是 $N/2+1$，不是 $N$（见讲复数与半谱那一章）。

\begin{lstlisting}[style=cstyle,caption={dsp/spec.cpp：plan 复用、对齐缓冲与半谱长度},label={lst:ch17-fft}]
/* 长度 1024 的实数到复数变换：输出 513 个 bin，bin k 对应
 * k * fs / 1024 Hz。plan 是句柄的成员，绝不每次调用重建。 */
int spec_init(Spec* s, int n) {
  s->in = (float*)fftwf_malloc(sizeof(float) * n);   /* 对齐 */
  s->out = (fftwf_complex*)fftwf_alloc_complex(n / 2 + 1);
  if (!s->in || !s->out) return SC_E_NOMEM;
  fftwf_make_planner_thread_safe();     /* threads 版才有的锁 */
  s->pl = fftwf_plan_dft_r2c_1d(n, s->in, s->out, FFTW_ESTIMATE);
  return s->pl ? SC_OK : SC_E_STATE;    /* ESTIMATE：不碰 wisdom */
}
\end{lstlisting}

用 \Tp{FFTW\_ESTIMATE} 而不是 \Tp{FFTW\_PATIENT} 是\emph{启动时间}与\emph{全局状态}的折中：更聪明的规划要搜索、耗秒级，并读写 wisdom 文件——那是跨线程与跨容器的共享状态。若确实需要离线调优，就在构建阶段跑一次 \Fn{fftwf\_export\_wisdom\_to\_file}，把产物当数据文件交付，而不是让运行期自己去写。

\section{构建：CMake 与依赖探测}

构建脚本是这份工程里最容易被低估的一层，因为它同时决定三件事：内核能不能单独测试、发布产物带几份库、以及浮点语义会不会变。清单~\ref{lst:ch17-cmake} 给出要点。

\begin{lstlisting}[style=cmakestyle,caption={CMake 要点：两条探测路径、目标分层、显式关掉宽松浮点},label={lst:ch17-cmake}]
cmake_minimum_required(VERSION 3.21)
project(spechain LANGUAGES C CXX)

# 依赖有两条路：包自带的 CMake config 优先，退到 pkg-config。
find_package(Sndfile CONFIG QUIET)      # imported 名随安装版本
if(NOT TARGET sndfile::sndfile)
  find_package(PkgConfig REQUIRED)
  pkg_check_modules(SNDFILE REQUIRED sndfile)
endif()

# FFTW3：较新 CMake 有官方模块；单精度常需自己补一条。
find_package(FFTW3 CONFIG QUIET)
if(NOT TARGET FFTW3::FFTW3)
  find_package(PkgConfig REQUIRED)
  pkg_check_modules(FFTW3F REQUIRED fftw3f)
endif()

add_library(core STATIC core/ring.c core/biquad.c)
add_library(dsp   STATIC dsp/chain.cpp dsp/spec.cpp)
target_link_libraries(dsp PUBLIC core)
target_include_directories(core PUBLIC include)

# 浮点语义是接口的一部分：禁用宽松优化，关掉 FMA 合约。
target_compile_options(core PRIVATE
  $<$<CXX_COMPILER_ID:GNU,Clang>:-fno-fast-math -ffp-contract=off
     -fvisibility=hidden>)

add_library(scbind SHARED bind/chain_capi.cpp)
target_link_libraries(scbind PRIVATE dsp core)
set_target_properties(scbind PROPERTIES
  POSITION_INDEPENDENT_CODE ON
  C_VISIBILITY_PRESET hidden
  VISIBILITY_INLINES_HIDDEN ON)
\end{lstlisting}

\textbf{静态与动态的选择。}\Tp{core/} 与 \Tp{dsp/} 编成静态库，好处是单测可以直接链它们、符号不外泄、发布时可以把它们并进一个模块；\Tp{scbind} 编成动态库，因为 \PY{} 侧要 \Tp{ctypes} 装载的只有这一个文件。\Lib{libsndfile} 与 \Lib{FFTW} 优先用动态：省体积、省许可证纠缠、也省得自己维护编码后端。只有做嵌入式或``一个 \Tp{.so} 带走一切''的分发时才考虑静态链它们，此时必须显式检查符号可见性，否则会把两份不同的 \Lib{OpenBLAS} 或两份 \Tp{zlib} 带进同一进程——门面章里那条 ``vendored BLAS 冲突'' 在这里会以链接错误的形式出现。

\textbf{可见性与导出宏。}\Tp{-fvisibility=hidden} 只在 ELF 与 Mach-O 上有意义；\Tp{bind/} 层要显式标注导出的符号：\Tp{\_\_attribute\_\_} 里给 \Tp{visibility(\textquotedbl default\textquotedbl)}，再加一个跨平台的 \Tp{SC\_API} 宏（Windows 上展开为 \Tp{dllexport}，用 \Fn{generate\_export\_macro} 生成也可以）。反过来，\Tp{core/} 与 \Tp{dsp/} 的任何符号都不该出现在 \Fn{nm -D} 的输出里——这条检查值得进 CI。

\textbf{浮点开关。}\Tp{-ffast-math} 不在这里出现是有原因的（讲热路径的那一章已经给过结论：它会破坏收敛判据与对拍基准）。\Tp{-ffp-contract=off} 则是更细的一刀：默认 \Tp{fast} 会把 \Tp{a*b+c} 合成一条 FMA，少一次舍入，于是同一份源码在不同向量宽度下\emph{不逐位一致}。GCC 较新版本可写 \Tp{-ffp-contract=off}，Clang 的默认值与开关名以所用版本文档为准，MSVC 侧对应 \Tp{/fp:precise} 并配合 \Tp{/arch:} 系列；把这三条写进构建说明，比事后调 golden 文件便宜。

\section{系数、编排、对拍与优化位点}

\subsection{系数为什么在 \PY{} 侧设计}

设计一次、执行百万次，这是\emph{最该分层}的地方。\Fn{scipy.signal.butter} 给出的是归一化 \Tp{sos} 矩阵；C 侧只把它读进 \Tp{ScConfig}（清单~\ref{lst:ch17-handle}），于是极点求解、双精度到单精度的转换、以及 \Tp{sos} 的配节策略这三件 ``\PY{} 与 C 一定对不齐'' 的事\emph{只存在于} \PY{} 侧，对拍失败时能立刻归因到执行而不是设计。

反过来，什么时候应该在 C 里做设计？当设计本身发生在产品运行期——用户拖动截止频率滑块的连续滤波器。那一类需求存在，而且讲 IIR 的那一章给出的库就自带设计；本工程的滤波器只在启动时装配一次，所以选 JSON。

序列化格式取 JSON 而不是二进制：系数只有 48 个数，解析成本与运行无关，而可读的系数文件等于免费的调试工具（\Tp{tests/} 里可以直接断言某个 \Tp{b1} 的符号）。字段就是 \Tp{ScConfig} 的字段，一行注释写清 ``samplerate 必须与实际数据一致''——采样率不匹配是这一域最隐蔽的错，它不改变形状，只改变频率轴。

\subsection{门面与对拍}

\begin{lstlisting}[style=pythonstyle,caption={py/spechain.py：设计系数、驱动句柄、与 scipy 对拍},label={lst:ch17-py}]
import ctypes
import json
import numpy as np
import soundfile as sf
from scipy import signal as sg

lib = ctypes.CDLL("./libscbind.so")     # 签名自己抄：错就是崩
lib.sc_create.argtypes = [ctypes.c_void_p, ctypes.c_char_p,
                          ctypes.c_int]
lib.sc_create.restype = ctypes.c_void_p
lib.sc_process.argtypes = [ctypes.c_void_p, ctypes.c_void_p,
                           ctypes.c_void_p, ctypes.c_int]
lib.sc_process.restype = ctypes.c_int

def design(fp, fs, order=8):
    """只设计，不执行：sos 与增益一起进 JSON 契约。"""
    sos = sg.butter(order, fp, btype="low", fs=fs, output="sos")
    return {"samplerate": int(fs), "nsec": int(len(sos)),
            "gain": float(np.prod(sos[:, 0])), "sos": sos.tolist()}

def run(path, cfg, block=4096):
    x, fs = sf.read(path, dtype="float32")      # 解码归门面
    if x.ndim > 1:
        x = x.mean(axis=1)                      # 混单声道：一次
    assert fs == cfg["samplerate"], "采样率与系数不符"
    y = np.empty_like(x)                        # 预分配，不 append
    for i in range(0, x.size - block, block):   # 一次边界穿越一块
        rc = lib.sc_process(CTX, x[i:].ctypes.data,
                            y[i:].ctypes.data, block)
        if rc:                                  # 返回码翻译成异常
            raise RuntimeError("sc_process rc=%d" % rc)
    lib.sc_drain(CTX)
    return y

def compare(y, x, cfg):
    yref = sg.sosfilt(np.array(cfg["sos"]), x) * cfg["gain"]
    np.testing.assert_allclose(y, yref, rtol=2e-4, atol=1e-6)
    # 谱的口径必须逐参数对齐，否则差的是归一化而不是算法
    f, t, S = sg.spectrogram(x, fs=cfg["samplerate"], window="hann",
                             nperseg=1024, noverlap=512,
                             detrend=False, scaling="spectrum",
                             mode="magnitude")
    return f, t, S
\end{lstlisting}

这段代码是门面四个职责的完整演示：\Fn{run} 里有形状与采样率校验、dtype 换算、预分配、返回码翻译；而\emph{没有}一行逐元素计算——那是 \Tp{core/} 的事。谱的横轴、CSV 表头与 PNG 图在 \Tp{py/} 的另一半里，用 \Lib{matplotlib} 消费数组完成。

\subsection{导出、测试判据与基线}

三条出口各有判据。\textbf{WAV 写回}：沿用读进来的那份 \Tp{SF\_INFO}（容器与子格式都不换），写成 \Tp{SF\_FORMAT\_FLOAT}，并且把 \Fn{sf\_writef\_float} 的返回值与请求帧数比对——短写是错误而不是常态。\textbf{CSV 谱}：三列（bin、Hz、dB），用 \Tp{\%.9g} 一类固定精度而不是 locale 敏感的默认转换，表头带上窗、重叠与口径。\textbf{golden}：小端 \Tp{float32} 裸数组加一份 sidecar，sidecar 里记 shape、dtype、\Fn{sf\_version\_string} 与 \Fn{fftwf\_version} 的报告值、\Tp{OMP\_NUM\_THREADS} 等环境变量、编译器与块大小。没有 sidecar 的 golden 会在三个月后变成谁也解释不了的字节串。

判据分四层，从宽到严：\textbf{量级}（总能量在十倍以内、谱峰所在 bin 相同）；\textbf{物理一致性}（低通不该放大带外，正弦输入的稳幅应与理论 $|H(f)|$ 对得上——这就是可用的``能量守恒''测试，直接断言``输入输出总能量相等''是错的，滤波器按定义就是能量不守恒的）；\textbf{统计一致}（\Fn{assert\_allclose}，\Tp{rtol} 按数据尺度定，不写死绝对值）；\textbf{逐位一致}（只允许在同一二进制、同一输入、关闭 FMA 合约与重排的前提下要求）。跨 BLAS 与跨 FFT 后端不要求逐位——那是讲数值接口那一章列出的第三层可复现，成本与收益不成比例。阶跃响应是另一条独立判据：零初始状态下的过冲与调整时间与 \Fn{sosfilt} 对拍，另加一条\emph{非零初始状态}的用例，专门抓 \Tp{z} 数组跨块丢失这类错误。

\begin{lstlisting}[style=shellstyle,caption={三个位点的测法：读数换成 ns 每帧，环境写进 sidecar},label={lst:ch17-meas}]
# 1) 同一输入、同一块大小、多次取最小值（尾部的频率噪声不算）
./sc_bench --in tone48k.wav --block 256 --reps 31 --csv bench.csv

# 2) 唯一有意义的读数：每帧纳秒，与块周期的预算比
awk -F, '{s+=$3} END {print 1e6*s/NR, "ns_per_frame"}' bench.csv

# 3) 基线带环境：换后端、换线程数、换 SIMD 都会移动结果
python -c "import numpy; numpy.show_config()" > env.blas.txt
\end{lstlisting}

\textbf{三个最有收益的位点}，按本工程的实测口径给出（口径可复述，倍数留给现场）：

\begin{enumerate}
\item \textbf{交织到平面的置换与滤波融合}。原来 \Tp{io/} 做一次置换、\Tp{core/} 再遍历一次；融合成一次遍历后，读的是同一份缓存。测法：固定 $48$\,kHz 立体声、块 $256$，比较两版的 \Fn{ns per frame}，并同时看 \Fn{perf stat} 一类的缓存缺失计数——收益来源要能指认，不然只是噪声。
\item \textbf{plan 复用与缓冲对齐}。从``每块建一个 plan'' 改成 ``create 建一次、process 只执行''，同时把 \Tp{malloc} 换成 \Fn{fftwf\_malloc}。测法：看\emph{首块}与\emph{稳态块}的耗时差、以及 P99；planner 带锁，多线程下这条差异会被放大。
\item \textbf{\PY{} 侧的逐帧拼接}。把 \Tp{list.append} 加 \Fn{np.concatenate} 换成清单~\ref{lst:ch17-py} 里的预分配加视图写入。测法：\Tp{time.perf\_counter\_ns} 计时，配 \Tp{tracemalloc} 看分配峰值——这条的收益一半在时间一半在内存抖动。
\end{enumerate}

第四个候选位点是关内层线程（把 BLAS 与 FFTW 的 threads 版换成串行版，让外层按通道并行）。它的收益随核数与通道数变化很大，只能用 ``在同一台机器上跑一条线程数曲线'' 的方式定标，不该写进文档当承诺。

\begin{bestpractice}[title={只改这五处}]
\begin{itemize}
\item 把逐样本的循环从 \PY{} 挪进 \Tp{core/}，边界穿越按块而不是按帧。
\item 关掉 \Tp{-ffast-math} 与 FMA 合约，把``谁逐位一致、谁只讲容差''写进测试判据。
\item plan、工作缓冲、环形缓冲全部在 \Fn{create} 里分配；\Fn{process} 里零分配、零阻塞、零日志。
\item 滤波器设计移到 \PY{} 侧并序列化，C 只读系数；级联顺序与增益位置写成注释。
\item golden 旁边放一份 sidecar：库版本、环境变量、块大小、编译器。缺任何一项，基线都会在换机后失效。
\end{itemize}
\end{bestpractice}

\begin{guideline}[title={分层与禁令：本章的收束}]
\begin{itemize}
\item 分层的意义不是``代码整洁''，而是让每层可以被单独构建、单独测试、单独替换：禁令比职责更能说明这一点。
\item 同一功能在两侧各写一遍是\emph{代价}，只有当它变成一条可执行的对拍测试时才是\emph{收益}；判据分四层，逐位一致只在自己二进制内要求。
\item 设计、校验、换算、翻译、编排归门面；逐元素计算、状态、缓冲与截止期归 \CL/\CPP{} 内核。这条线画错了，代码会同时变慢和变难测。
\item 全书主线在这个工程里落成一句话：内核用 \CL/\CPP{} 不是因为快，而是因为它们是唯一能把\emph{布局、所有权与线程语义写死}，并让另一门语言逐条核对的层。
\end{itemize}
\end{guideline}
% ====================================================================

% ====================================================================
\appendix

\chapter{库、语言层与绑定矩阵}

本附录把全书出现的实现收进三张速查表：谁用什么语言写成、对外暴露什么接口、以及从别的语言进去要付什么代价。``语言层''一栏只写\emph{内核}与\emph{对外接口}两件事，因为这一域的常见误解正是把两者混为一谈：一个 \CPP{} 写的库若只暴露 \CL{} 句柄，它的可达范围与一个纯 \CL{} 库相同；反之一个 \CL{} 写的库若把 \CPP{} 类放进头文件，它的可绑定性就降到 \CPP{} 一档。许可证一栏一律``以仓库 LICENSE 原文为准''，本表只标注大类。

\section{数学与数值计算}

\begin{longtable}{@{}>{\raggedright\arraybackslash}p{0.16\textwidth}
  >{\raggedright\arraybackslash}p{0.19\textwidth}
  >{\raggedright\arraybackslash}p{0.19\textwidth}
  >{\raggedright\arraybackslash}p{0.28\textwidth}@{}}
\caption{数学与数值计算类库的接口形态}\label{tab:apA-math}\\
\toprule
库 & 内核语言 & 对外接口 & 绑定与备注 \\
\midrule
\endhead
\Lib{BLAS} 参考 & \FTN{} & \FTN{} 符号加 \CL{} 的 \Tp{cblas\_} &
  任何语言可及；同一符号可有多个实现 \\
\Lib{OpenBLAS} & \CL{} 加汇编 & \FTN{} 符号加 \Tp{cblas\_} &
  运行期 ISA 分派；宽许可证 \\
\Lib{MKL} & \CL{} 与 \CPP{} 加汇编 & \CL{} 句柄与函数 &
  专有许可，体积大，线程模型自成一套 \\
\Lib{BLIS} & \CPP{} 内核 & \CL{} 微内核加接口 &
  教学与移植新硬件的首选形态 \\
\Lib{LAPACK} & \FTN{} & \FTN{} 符号 &
  上层通常经 \Lib{LAPACKE} 的 \CL{} 接口使用 \\
\Lib{Eigen} & \CPP{} 模板 & 头文件模板 &
  无 ABI；绑定要实例化，性能靠内联 \\
\Lib{GSL} & \CL{} & \CL{} 句柄加返回码 &
  \Tp{ctypes}/\Tp{cffi} 直取；错误处理器需关闭 \\
\Lib{Boost.Math} & \CPP{} 模板 & 头文件模板 &
  policy 控制精度与错误行为 \\
\Lib{Sundials} & \CL{} & \CL{} 长生命周期对象加回调 &
  仿真器接口形态的典型 \\
\Lib{Ceres} & \CPP{} & \CPP{} 加模板 &
  自动微分靠模板实例化 \\
\Lib{NLopt} & \CL{} 加 \CPP{} & \CL{} 函数指针加结构体 &
  多种局部与全局优化器集合 \\
\Lib{OSQP} & \CL{} & \CL{} 工作空间句柄 &
  QP 内核，跨语言绑定成熟 \\
\end{longtable}

\section{几何与物理}

\begin{longtable}{@{}>{\raggedright\arraybackslash}p{0.16\textwidth}
  >{\raggedright\arraybackslash}p{0.19\textwidth}
  >{\raggedright\arraybackslash}p{0.19\textwidth}
  >{\raggedright\arraybackslash}p{0.28\textwidth}@{}}
\caption{几何与物理类库的接口形态}\label{tab:apA-geo}\\
\toprule
库 & 内核语言 & 对外接口 & 绑定与备注 \\
\midrule
\endhead
\Lib{CGAL} & \CPP{} 模板 & 头文件模板加可选 \CL{} 包装 &
  Kernel 参数化精确性；依赖任意精度库 \\
\Lib{GEOS} & \CL{} 加 \CPP{} 内部 & \CL{} 上下文加返回码 &
  带 \Tp{\_r} 的可重入入口是线程安全做法 \\
\Lib{PROJ} & \CL/\CPP{} & \CL{} 句柄加错误码 &
  上下文式 API；数据库与单位约定在库内 \\
\Lib{libigl} & \CPP{} 模板 & 头文件加 \Lib{Eigen} 类型 &
  数据层即 \Lib{Eigen}，绑定用 pybind11 \\
\Lib{Open3D} & \CPP{} & \CPP{} 类加官方绑定 &
  同一内核两套门面（传统与张量式） \\
\Lib{Embree} & \CPP{} & \CPP{} 类加 \CL{} 实例 API &
  运行期 ISA 分派；SoA 批量求交 \\
\Lib{Bullet} & \CPP{} & \CPP{} 类为主 &
  实时域内禁止分配与阻塞 \\
\Lib{ODE} & \CL{} & \CL{} 空间与对象句柄 &
  接口老但跨语言可达性极好 \\
\Lib{MuJoCo} & \CL{} 与 \CPP{} & \CL{} 结构体加函数 &
  官方门面直接映射 C 状态对象 \\
\Lib{FCL} & \CPP{} & \CPP{} 碰撞对象加 \CL{} 风格层 &
  机器人运动规划里的窄相常用 \\
\end{longtable}

\section{信号处理：声音与电磁波}

\begin{longtable}{@{}>{\raggedright\arraybackslash}p{0.16\textwidth}
  >{\raggedright\arraybackslash}p{0.19\textwidth}
  >{\raggedright\arraybackslash}p{0.19\textwidth}
  >{\raggedright\arraybackslash}p{0.28\textwidth}@{}}
\caption{信号处理类库的接口形态}\label{tab:apA-dsp}\\
\toprule
库 & 内核语言 & 对外接口 & 绑定与备注 \\
\midrule
\endhead
\Lib{FFTW} & \CL{}（\FTN{} 谱系） & \CL{} plan 句柄加数组 &
  交错复数布局；plan 与 wisdom 有状态 \\
\Lib{KissFFT} & \CL{} & \CL{} 配置结构体加裸数组 &
  小、可嵌入、可定点；无 SIMD 依赖 \\
\Lib{pffft} & \CL{} & \CL{} 裸数组 &
  单精度，实数变换用 split 布局 \\
\Lib{PocketFFT} & \CPP{}（\CL{} 风格） & 头文件加函数 &
  常作为打包分发用的折中选择 \\
\Lib{cuFFT} & 专有 & \CL{} 句柄加计划 &
  跨设备指针与流语义是接口的一部分 \\
\Lib{libsndfile} & \CL{} & \CL{} 句柄加格式宏 &
  封装与编码解耦；\PY{} 门面的底层 \\
\Lib{miniaudio} & \CL{} & 单头文件 \CL{} &
  零依赖，适合直接嵌入分发 \\
\Lib{PortAudio} & \CL{} & \CL{} 流加回调 &
  回调线程是实时域 \\
\Lib{Cubeb} & \CPP{} & \CL{} 门面 &
  浏览器音频栈里的兼容层 \\
\Lib{SoX} & \CL{} & 命令行加 \CL{} 库 &
  离线批处理与格式转换的主力 \\
\Lib{iir1} & \CPP{} 模板 & 头文件模板 &
  系数与状态类型分离 \\
\Lib{r8b} & \CPP{} & \CPP{} 类加 \CL{} 风格层 &
  多相重采样与质量档 \\
\Lib{CMSIS-DSP} & \CL{} & \CL{} 实例结构体加函数 &
  目标为 Cortex-M 与定点 \\
\Lib{Meep} & \CPP{} & Scheme 与 \PY{} 门面 &
  FDTD 内核加 MPI 并行 \\
\Lib{openEMS} & \CPP{} & \CPP{} 加 Lua 门面 &
  时域为主，含远场变换 \\
\Lib{gprMax} & 编译化内核 & \PY{} CLI &
  \PY{} 门面之下的 FDTD 循环不在解释器里 \\
\Lib{SoapySDR} & \CPP{} & \CL{} 插件 API &
  跨厂商设备抽象；键值协商格式 \\
\Lib{UHD} & \CPP{} & \CPP{} 类加 \CL{} 绑定 &
  流式传输与时间戳语义 \\
\Lib{GNU Radio} & \CPP{} 内核 & \PY{} 图加 \CPP{} block &
  样点级循环在 block 的 work 里 \\
\Lib{scikit-rf} & \PY{} 加数值内核 & \PY{} &
  S 参数与网络级联的门面层 \\
\Lib{libmobi} & \CL{} & \CL{} 加 XML 描述 &
  射频前端与滤波器综合 \\
\end{longtable}

\begin{guideline}[title={怎么读这三张矩阵}]
\begin{itemize}
\item 内核语言决定你能否读懂与改动它；对外接口决定你能否从别的语言安全进入它。
\item 凡是``\CL{} 句柄加返回码''的库，绑定成本几乎恒定，选型时应看性能与语义而不是可绑定性。
\item 凡是``头文件模板''的库，可获得最高性能，但每次升级都要重编译，跨语言绑定必须自己写 \CPP{} 门面。
\item 表中未列的版本、后端与许可证细节，一律以官方仓库的 LICENSE 与文档为准。
\end{itemize}
\end{guideline}

\chapter{构建、链接与诊断速查}

这一域的缺陷有一半在编译与链接阶段就已注定：符号找不到、找到了但不是你以为的那个实现、两份库同时装进一个进程。本章给的是定位顺序与查询手段，不是某个库的构建教程。

\section{名称映射与发现}

同一个库在 CMake、pkg-config、包管理器与 \PY{} 生态里常常有四个不同的名字。下表的 CMake 列写的是探测名，实际用法是把它作为参数交给 \Tp{find\_package}，而导入目标名（带命名空间的那个）以库自己导出的 config 为准；Module 模式与 Config 模式的差别、以及第三方发行版可能改名，都要在配方里核对。

\begin{longtable}{@{}>{\raggedright\arraybackslash}p{0.16\textwidth}
  >{\raggedright\arraybackslash}p{0.20\textwidth}
  >{\raggedright\arraybackslash}p{0.18\textwidth}
  >{\raggedright\arraybackslash}p{0.26\textwidth}@{}}
\caption{库名的四种写法}\label{tab:apB-names}\\
\toprule
库 & CMake 名 & pkg-config & \PY{} 侧常见对应 \\
\midrule
\endhead
BLAS 与 LAPACK & \Tp{BLAS}、\Tp{LAPACK} &
  \Tp{blas}、\Tp{lapack} &
  \Tp{numpy} 与 \Tp{scipy} 各带一份 \\
\Lib{LAPACKE} & \Tp{LAPACKE} & \Tp{lapacke} &
  \Tp{scipy.linalg} 的底层之一 \\
\Lib{FFTW} & \Tp{FFTW3}（随发行） & \Tp{fftw3}、\Tp{fftw3f} &
  \Tp{pyfftw} \\
\Lib{Eigen} & \Tp{Eigen3}，目标带命名空间 &
  无（头文件） & 无（编译期依赖） \\
\Lib{GSL} & \Tp{GSL} & \Tp{gsl} & \Tp{numpy} 或自写绑定 \\
\Lib{libsndfile} & \Tp{SndFile} & \Tp{sndfile} & \Tp{soundfile} \\
\Lib{PortAudio} & \Tp{PortAudio} & \Tp{portaudio} &
  \Tp{sounddevice}、\Tp{pyaudio} \\
\Lib{PROJ} & \Tp{PROJ} & \Tp{proj} & \Tp{pyproj} \\
\Lib{GEOS} & \Tp{GEOS} & \Tp{geos-c} & \Tp{shapely} \\
\Lib{Open3D} & \Tp{Open3D} & 视配方 & \Tp{open3d} \\
\Lib{Sundials} & \Tp{SUNDIALS} & \Tp{sundials} & \Tp{appsi} 一类 \\
\Lib{Ceres} & \Tp{Ceres} & \Tp{ceres} & 自写绑定 \\
\Lib{SoapySDR} & \Tp{SoapySDR} & \Tp{soapysdr} & 官方绑定包 \\
\end{longtable}

\begin{lstlisting}[style=cmakestyle,caption={把这一域的常见依赖接在一起：显式线程库与显式数学库}]
find_package(BLAS REQUIRED)
find_package(LAPACK REQUIRED)
find_package(Threads REQUIRED)
find_package(SndFile CONFIG QUIET)
find_package(Eigen3 CONFIG QUIET)

add_library(dsp_core STATIC src/filters.c src/spectrum.c src/io.c)
target_link_libraries(dsp_core PUBLIC
  ${BLAS_LIBRARIES} ${LAPACK_LIBRARIES} Threads::Threads)
if(TARGET SndFile::sndfile)
  target_link_libraries(dsp_core PUBLIC SndFile::sndfile)
endif()
# 显式 -lm 在 Linux 上仍然必要（链接顺序：库在 -lm 之前）
target_compile_features(dsp_core PUBLIC c_std_11)
# 禁止宽松浮点：发布构建里这一条要写死
target_compile_options(dsp_core PRIVATE -fno-fast-math)
\end{lstlisting}

\section{符号查不到的四种因}

\textbf{名字修饰}。\FTN{} 加尾下划线、\CPP{} 加参数类型编码。先用工具看真实符号，再改声明。

\textbf{按值与按引用}。\FTN{} 的标量默认按引用，签名声明错一个整型参数，链接能过而结果是错的。

\textbf{实现与头文件版本不一致}。头来自一处、库来自另一处（尤其 Windows 上多个发行版共存），表现为少符号或结构体尺寸不符。

\textbf{静态运行库冲突}。Windows 上把用了静态 CRT 的库链进用了动态 CRT 的可执行文件，跨模块释放内存时崩溃；Linux 上表现为符号被 \Tp{RTLD\_GLOBAL} 抢走。

\begin{lstlisting}[style=shellstyle,caption={定位手段：看符号、看依赖、看运行时线程数}]
# 这个库里到底有没有我要的符号（注意 Fortran 的尾下划线）
nm -D --defined-only libopenblas.so | grep gemm
ldd ./build/dsp_app | grep -Ei 'blas|lapack|gfortran|omp'
# 上面第二行：我的可执行文件实际链到了哪一份 BLAS
# Windows 下等价的两步（符号与依赖）
#   dumpbin /symbols kernel.lib | findstr /i dgemm
#   dumpbin /dependents app.exe
# 运行时到底起了多少线程
python -c "import threadpoolctl as t; print(t.threadpool_info())"
top -H -p $(pgrep -f dsp_app)     # 线程级 CPU，判断是否超订
# FFTW 的 wisdom 落在哪、是否首次调用在做规划
strace -e trace=file ./build/dsp_app 2>&1 | grep -i wisdom
\end{lstlisting}

\begin{warning}[title={一个进程里有两份 BLAS}]
表现是``能用但慢''或``随机崩溃''：\Tp{numpy} 自带一份，你自己又链了一份；或者插件与宿主各带一份 PROJ。诊断顺序是先查依赖闭包（上面的 \Tp{ldd}），再查符号抢占（\Tp{LD\_DEBUG\_BINDINGS}、\Tp{RTLD\_GLOBAL}），最后才怀疑代码。根治办法是统一由宿主提供，或全部走同一发行包。
\end{warning}

\section{构建期环境变量与开关}

\begin{longtable}{@{}>{\raggedright\arraybackslash}p{0.26\textwidth}
  >{\raggedright\arraybackslash}p{0.21\textwidth}
  >{\raggedright\arraybackslash}p{0.33\textwidth}@{}}
\caption{构建与打包阶段常用的环境变量}\label{tab:apB-env}\\
\toprule
变量 & 影响 & 误用后果 \\
\midrule
\endhead
\Tp{CMAKE\_PREFIX\_PATH} & 探测安装的库与配置包 &
  找不到库，或链到系统里的旧版本 \\
\Tp{PKG\_CONFIG\_PATH} & pkg-config 搜索路径 &
  头与库来自两个前缀，符号缺失 \\
\Tp{MKLROOT} & \Lib{MKL} 的路径约定 &
  链接成功但运行时缺 DLL \\
\Tp{OMP\_NUM\_THREADS} 等 & 运行期线程数 &
  把并行域套两层，吞吐塌到谷底 \\
\Tp{MACOSX\_}\Bk\Tp{DEPLOYMENT\_}\Bk\Tp{TARGET} & 目标系统版本 &
  产物在新系统能跑、旧系统缺符号 \\
\Tp{AUDITWHEEL\_POLICY} & \PY{} wheel 的标签策略 &
  上传即失败或装了跑不起来 \\
\end{longtable}

\chapter{术语与缩写表}

\begin{longtable}{@{}>{\raggedright\arraybackslash}p{0.16\textwidth}
  >{\raggedright\arraybackslash}p{0.26\textwidth}
  >{\raggedright\arraybackslash}p{0.38\textwidth}@{}}
\caption{全书术语与缩写}\label{tab:apC-glossary}\\
\toprule
缩写 & 展开 & 一句话定义 \\
\midrule
\endhead
ABI & 应用二进制接口 &
  跨编译器都遵守的调用约定、布局与对齐规则 \\
AoS & 结构数组 &
  按元素把字段排在一起的布局，遍历物体友好 \\
SoA & 数组结构 &
  按字段把元素排在一起的布局，向量化友好 \\
BLAS & 基础线性代数子程序 &
  一层向量与矩阵运算的接口标准，实现可替换 \\
BVH & 层次包围盒 &
  把几何包成树做遍历剪枝，光线与碰撞的常用结构 \\
C2R、R2C & 复到实、实到复 &
  利用共轭对称只存半谱的 FFT 类型 \\
DFT、FFT & 离散傅里叶变换、快速傅里叶变换 &
  同一变换的两种算法复杂度，工程上常混用两词 \\
dBFS & 相对满量程的分贝 &
  数字音频的值域基准，零即满刻度 \\
dBm & 相对一毫瓦的分贝 &
  射频功率基准，与参考阻抗一起才有意义 \\
FDTD & 时域有限差分 &
  在 Yee 元胞上做时间推进的电磁仿真方法 \\
FEM & 有限元法 &
  分片基函数上的弱形式离散，频域与结构常用 \\
FIR & 有限冲激响应 &
  无反馈的卷积滤波器，稳定但阶数高 \\
IIR & 无限冲激响应 &
  有反馈的递归滤波器，阶数低但要管稳定性 \\
FMA & 融合乘加 &
  一次舍入的乘加指令，改变浮点结果末位 \\
GIL & 全局解释器锁 &
  \PY{} 里限制多线程并行的运行时机制 \\
GJK & 相邻体积算法 &
  凸体间最近距离的迭代算法，窄相基础 \\
GEMM & 通用矩阵乘 &
  带 alpha 与 beta 的稠密矩阵乘，多数内核的性能源 \\
ICP & 迭代最近点 &
  点云配准的迭代算法，依赖初值与降采样 \\
KDTree & K 维树 &
  最近邻与半径查询的空间划分结构 \\
LAPACK & 线性代数程序包 &
  在 BLAS 之上的分解、特征值与求解器集合 \\
MoM & 矩量法 &
  表面积分方程的离散，天线与散射常用 \\
MPI & 消息传递接口 &
  跨进程跨节点的并行标准，大规模仿真的骨干 \\
PML & 完全匹配层 &
  电磁仿真里吸收外行波的边界处理 \\
PIMPL & 指向实现的指针 &
  \CPP{} 里用间接隐藏布局以稳定 ABI 的惯用法 \\
PSD & 功率谱密度 &
  每单位带宽的功率，注意与每 bin 幅度区分 \\
Q15、Q31 & 定点格式 &
  小数点固定在低位若干位的整数表示，嵌入式常用 \\
RCWA & 严格耦合波分析 &
  周期性结构（光栅）的频域方法 \\
S 参数 & 散射参数 &
  端口网络在参考阻抗下的入射与反射关系 \\
SDR & 软件无线电 &
  把混频、滤波与解调交给通用处理器完成 \\
SIMD & 单指令多数据 &
  一条指令作用于多个通道的执行方式 \\
SPICE & 带偏置的综合仿真程序 &
  电路时域仿真方法与一族同名工具 \\
STFT & 短时傅里叶变换 &
  分帧加窗做 FFT，得到时频谱 \\
TLP & 线程级并行 &
  以线程为并行单元的执行方式 \\
TSDF & 截断符号距离场 &
  体素化的表面表示，用于多视角重建 \\
ULP & 末位单位 &
  浮点数间隔，衡量实现与参考值的误差 \\
VNA & 矢量网络分析仪 &
  测 S 参数的仪器，其数据是仿真对拍对象 \\
XRUN & 过欠载 &
  实时音频缓冲被耗尽或写满，表现为爆音 \\
Yee & 元胞 &
  场分量在网格上交错排布，FDTD 的离散模板 \\
\end{longtable}

\section{全书主线回顾}

把本册的论证压成四句话。第一，这一域的性能来自可核对的物理量：运算密度、访存模式、分支可预测性与能不能停，而不是语言名字（第 1 章）。第二，跨语言的汇合点只有 \CL{} ABI，因此``以 C/C++ 为主''的准确含义是把 \CL{} 当作对外契约，把 \FTN{} 与 \CPP{} 留在契约之后（第 2 章）。第三，主序、步长、复数形态、单位、错误传递与线程模型是六类必须写死的约定，它们出错时不崩溃也不报错（第 3 章）。第四，从 \Lib{BLAS} 到 \Lib{Meep}，各章反复出现同一个结构：内核用 \CL{} 或 \CPP{} 写死布局与算术，句柄式接口负责可达，薄门面负责编排，而门面上的每一层抽象都在收税。

\begin{bestpractice}[title={接手一个陌生库的六个问题}]
\begin{itemize}
\item 内核语言是什么，对外接口是什么？二者不一致时以接口为准。
\item 它要求什么数据布局：主序、步长、复数形态、对齐？
\item 它的状态活在哪里：句柄、全局、plan 还是编译期？
\item 谁起线程，线程数受哪个环境变量控制，能否与我的并行域共存？
\item 错误如何传出：返回码、错误字符串还是异常，是否线程局部？
\item 许可证是什么大类，静态链接进发布物时是否可接受？
\end{itemize}
\end{bestpractice}
% ====================================================================

\end{document}
